% !TeX program = xelatex
% Standalone English volume. Use XeLaTeX until the references stabilize.
\documentclass[a4paper,11pt]{article}
\usepackage[T1]{fontenc}
\usepackage{lmodern}
\usepackage[english]{babel}
\usepackage[left=26mm,right=26mm,top=25mm,bottom=25mm,headheight=15pt,headsep=7mm,footskip=11mm]{geometry}
\usepackage{amsmath,amssymb,mathtools,amsthm}
\newtheorem{proposition}{Proposition}[section]
\newtheorem{lemma}[proposition]{Lemma}
\newtheorem{theorem}[proposition]{Theorem}
\newtheorem{corollary}[proposition]{Corollary}
\newtheorem{maintheorem}[proposition]{Main Theorem}
\theoremstyle{definition}
\newtheorem{definition}[proposition]{Definition}
\theoremstyle{remark}
\newtheorem{remark}[proposition]{Remark}
\makeatletter
\renewcommand*\l@section{\@dottedtocline{1}{0em}{2.4em}}
\renewcommand*\l@subsection{\@dottedtocline{2}{1.45em}{3.45em}}
\renewcommand*\l@subsubsection{\@dottedtocline{3}{3.8em}{4.55em}}
\makeatother
\usepackage{xcolor,fancyhdr,titlesec,needspace,xurl}
\usepackage{longtable,multicol}
\usepackage{hyperref}
\definecolor{ink}{HTML}{1A3542}
\definecolor{muted}{HTML}{4F606A}
\hypersetup{unicode=true,colorlinks=true,linkcolor=ink,filecolor=ink,urlcolor=ink,citecolor=ink,bookmarksnumbered=true,
 pdfdisplaydoctitle=true,pdflang={en},
 pdftitle={The generalized Riemann hypothesis for Dirichlet L-functions. Edition 930, Volume XIX},
 pdfsubject={A natural-language reconstruction of the proposed proof},
 pdfauthor={A mathematical reconstruction}}
\titleformat{\section}{\Large\bfseries\color{ink}}{\thesection}{.65em}{}
\titlespacing*{\section}{0pt}{.8em}{.65em}
\setlength{\parindent}{1.2em}
\linespread{1.045}
\setlength{\parskip}{.22em plus .05em minus .03em}
\setlength{\emergencystretch}{2em}
\setcounter{topnumber}{2}
\setcounter{tocdepth}{3}\setcounter{secnumdepth}{3}
\widowpenalty=4500
\clubpenalty=4500
\pagestyle{fancy}\fancyhf{}
\fancyhead[L]{\small\color{muted}Dirichlet GRH \textperiodcentered\ Volume I}
\fancyhead[R]{\small\color{muted}Dirichlet zeros and power intervals}
\fancyfoot[C]{\small\thepage}
\renewcommand{\headrulewidth}{.3pt}
\newcommand{\lead}[1]{\par\smallskip\noindent{\bfseries\color{ink}#1.}\enspace}
\newcommand{\sour}[2]{\unskip\penalty150\hspace{.35em}{\small\hyperlink{ref#1}{[#2]}}}
\titleformat{\subsection}{\large\bfseries}{\thesubsection}{.6em}{}
\titlespacing*{\subsection}{0pt}{1.15em plus .15em}{.4em}
\setlength{\abovedisplayskip}{9pt plus 2pt minus 3pt}
\setlength{\belowdisplayskip}{9pt plus 2pt minus 3pt}
\newcommand{\RR}{\mathbb R}\newcommand{\CC}{\mathbb C}\newcommand{\QQ}{\mathbb Q}
\newcommand{\ZZ}{\mathbb Z}\newcommand{\ord}{\operatorname{ord}}
\newcommand{\Frob}{\operatorname{Frob}}\newcommand{\Gal}{\operatorname{Gal}}
\newcommand{\len}{\operatorname{length}}\newcommand{\Res}{\operatorname*{Res}}
\newcommand{\wmean}{\operatorname{av}}\newcommand{\Ker}{\operatorname{ker}}
\newcommand{\Lh}{\mathcal L}\newcommand{\Gg}{\mathcal G}
\newcommand{\real}{\operatorname{Re}}\newcommand{\imag}{\operatorname{Im}}
\newcommand{\abs}[1]{\left|#1\right|}\newcommand{\eps}{\varepsilon}
\newcommand{\anchoritem}[2]{\par\Needspace{2.8\baselineskip}\hypertarget{ref#1}{}\noindent{\bfseries\color{ink}#2.}\enspace}
\newcommand{\srcloc}[2]{{\small\nolinkurl{#1.lean}}, line~#2}
\numberwithin{equation}{section}
\raggedbottom

\newcommand{\Sel}{\mathcal S}
\newcommand{\Cl}{\operatorname{Cl}}
\newcommand{\Aut}{\operatorname{Aut}}
\newcommand{\GL}{\operatorname{GL}}
\newcommand{\OO}{\mathcal O}
\hypersetup{bookmarksdepth=3}


\addto\captionsenglish{\renewcommand{\partname}{Volume}}


\makeatletter
\expandafter\def\csname r@cor:amplitude204\endcsname{{110.17}{52}{}{proposition.110.17}{Dirichlet-GRH-v930-Volume-XXI.pdf}}
\expandafter\def\csname r@cor:content125\endcsname{{92.19}{15}{}{proposition.92.19}{Dirichlet-GRH-v930-Volume-XVIII.pdf}}
\expandafter\def\csname r@cor:match209\endcsname{{93.12}{30}{}{proposition.93.12}{Dirichlet-GRH-v930-Volume-XVIII.pdf}}
\expandafter\def\csname r@cor:match210\endcsname{{93.19}{35}{}{proposition.93.19}{Dirichlet-GRH-v930-Volume-XVIII.pdf}}
\expandafter\def\csname r@cor:original200\endcsname{{83.7}{57}{}{proposition.83.7}{Dirichlet-GRH-v930-Volume-XVI.pdf}}
\expandafter\def\csname r@cor:rank111\endcsname{{9.14}{119}{}{proposition.9.14}{Dirichlet-GRH-v930-Volume-I.pdf}}
\expandafter\def\csname r@cor:trace235\endcsname{{110.7}{43}{}{proposition.110.7}{Dirichlet-GRH-v930-Volume-XXI.pdf}}
\expandafter\def\csname r@eq:direct76\endcsname{{66.16}{63}{}{equation.66.16}{Dirichlet-GRH-v930-Volume-XIII.pdf}}
\expandafter\def\csname r@eq:label209\endcsname{{93.6}{28}{}{equation.93.6}{Dirichlet-GRH-v930-Volume-XVIII.pdf}}
\expandafter\def\csname r@eq:lastdegree181\endcsname{{104.2}{6}{}{equation.104.2}{Dirichlet-GRH-v930-Volume-XX.pdf}}
\expandafter\def\csname r@eq:letters111\endcsname{{9.1}{117}{}{equation.9.1}{Dirichlet-GRH-v930-Volume-I.pdf}}
\expandafter\def\csname r@eq:main127\endcsname{{92.17}{20}{}{equation.92.17}{Dirichlet-GRH-v930-Volume-XVIII.pdf}}
\expandafter\def\csname r@eq:match209\endcsname{{93.7}{30}{}{equation.93.7}{Dirichlet-GRH-v930-Volume-XVIII.pdf}}
\expandafter\def\csname r@eq:real127\endcsname{{92.15}{19}{}{equation.92.15}{Dirichlet-GRH-v930-Volume-XVIII.pdf}}
\expandafter\def\csname r@eq:region198\endcsname{{20.12}{33}{}{equation.20.12}{Dirichlet-GRH-v930-Volume-IV.pdf}}
\expandafter\def\csname r@eq:seed127\endcsname{{92.16}{20}{}{equation.92.16}{Dirichlet-GRH-v930-Volume-XVIII.pdf}}
\expandafter\def\csname r@lem:basis63\endcsname{{65.14}{34}{}{proposition.65.14}{Dirichlet-GRH-v930-Volume-XIII.pdf}}
\expandafter\def\csname r@lem:classes304\endcsname{{119.2}{4}{}{proposition.119.2}{Dirichlet-GRH-v930-Volume-XXIII.pdf}}
\expandafter\def\csname r@lem:coefficient199\endcsname{{83.1}{52}{}{proposition.83.1}{Dirichlet-GRH-v930-Volume-XVI.pdf}}
\expandafter\def\csname r@lem:conjugation292\endcsname{{104.30}{26}{}{proposition.104.30}{Dirichlet-GRH-v930-Volume-XX.pdf}}
\expandafter\def\csname r@lem:cosets163\endcsname{{86.21}{18}{}{proposition.86.21}{Dirichlet-GRH-v930-Volume-XVII.pdf}}
\expandafter\def\csname r@lem:covol127\endcsname{{92.23}{19}{}{proposition.92.23}{Dirichlet-GRH-v930-Volume-XVIII.pdf}}
\expandafter\def\csname r@lem:degree181\endcsname{{104.6}{6}{}{proposition.104.6}{Dirichlet-GRH-v930-Volume-XX.pdf}}
\expandafter\def\csname r@lem:exchangeroot336\endcsname{{86.33}{26}{}{proposition.86.33}{Dirichlet-GRH-v930-Volume-XVII.pdf}}
\expandafter\def\csname r@lem:forms203\endcsname{{110.10}{46}{}{proposition.110.10}{Dirichlet-GRH-v930-Volume-XXI.pdf}}
\expandafter\def\csname r@lem:gauss89\endcsname{{92.8}{8}{}{proposition.92.8}{Dirichlet-GRH-v930-Volume-XVIII.pdf}}
\expandafter\def\csname r@lem:match209\endcsname{{93.10}{28}{}{proposition.93.10}{Dirichlet-GRH-v930-Volume-XVIII.pdf}}
\expandafter\def\csname r@lem:openflags353\endcsname{{120.7}{18}{}{proposition.120.7}{Dirichlet-GRH-v930-Volume-XXIII.pdf}}
\expandafter\def\csname r@lem:orbvol263\endcsname{{111.18}{73}{}{proposition.111.18}{Dirichlet-GRH-v930-Volume-XXI.pdf}}
\expandafter\def\csname r@lem:pivots298\endcsname{{107.5}{78}{}{proposition.107.5}{Dirichlet-GRH-v930-Volume-XX.pdf}}
\expandafter\def\csname r@lem:power291\endcsname{{107.11}{82}{}{proposition.107.11}{Dirichlet-GRH-v930-Volume-XX.pdf}}
\expandafter\def\csname r@lem:size97\endcsname{{87.1}{33}{}{proposition.87.1}{Dirichlet-GRH-v930-Volume-XVII.pdf}}
\expandafter\def\csname r@lem:tensor293\endcsname{{104.31}{27}{}{proposition.104.31}{Dirichlet-GRH-v930-Volume-XX.pdf}}
\expandafter\def\csname r@lem:termparent221\endcsname{{109.10}{11}{}{proposition.109.10}{Dirichlet-GRH-v930-Volume-XXI.pdf}}
\expandafter\def\csname r@lem:translate163\endcsname{{86.19}{16}{}{proposition.86.19}{Dirichlet-GRH-v930-Volume-XVII.pdf}}
\expandafter\def\csname r@nav:VI-15\endcsname{{92.8}{18}{}{subsection.92.8}{Dirichlet-GRH-v930-Volume-XVIII.pdf}}
\expandafter\def\csname r@nav:VI-26\endcsname{{9.9}{117}{}{subsection.9.9}{Dirichlet-GRH-v930-Volume-I.pdf}}
\expandafter\def\csname r@nav:VII-00\endcsname{{86.1}{3}{}{subsection.86.1}{Dirichlet-GRH-v930-Volume-XVII.pdf}}
\expandafter\def\csname r@prop:basechange258\endcsname{{71.17}{34}{}{proposition.71.17}{Dirichlet-GRH-v930-Volume-XIV.pdf}}
\expandafter\def\csname r@prop:cartan290\endcsname{{105.8}{39}{}{proposition.105.8}{Dirichlet-GRH-v930-Volume-XX.pdf}}
\expandafter\def\csname r@prop:class236\endcsname{{109.40}{34}{}{proposition.109.40}{Dirichlet-GRH-v930-Volume-XXI.pdf}}
\expandafter\def\csname r@prop:column110\endcsname{{86.29}{23}{}{proposition.86.29}{Dirichlet-GRH-v930-Volume-XVII.pdf}}
\expandafter\def\csname r@prop:complement210\endcsname{{93.20}{36}{}{proposition.93.20}{Dirichlet-GRH-v930-Volume-XVIII.pdf}}
\expandafter\def\csname r@prop:depth245\endcsname{{95.3}{67}{}{proposition.95.3}{Dirichlet-GRH-v930-Volume-XVIII.pdf}}
\expandafter\def\csname r@prop:finite237\endcsname{{94.20}{53}{}{proposition.94.20}{Dirichlet-GRH-v930-Volume-XVIII.pdf}}
\expandafter\def\csname r@prop:graphexchange336\endcsname{{71.20}{37}{}{proposition.71.20}{Dirichlet-GRH-v930-Volume-XIV.pdf}}
\expandafter\def\csname r@prop:hyp209\endcsname{{93.13}{30}{}{proposition.93.13}{Dirichlet-GRH-v930-Volume-XVIII.pdf}}
\expandafter\def\csname r@prop:index90\endcsname{{94.6}{41}{}{proposition.94.6}{Dirichlet-GRH-v930-Volume-XVIII.pdf}}
\expandafter\def\csname r@prop:loops302\endcsname{{119.1}{3}{}{proposition.119.1}{Dirichlet-GRH-v930-Volume-XXIII.pdf}}
\expandafter\def\csname r@prop:measure283\endcsname{{105.19}{50}{}{proposition.105.19}{Dirichlet-GRH-v930-Volume-XX.pdf}}
\expandafter\def\csname r@prop:normalframe349\endcsname{{104.34}{29}{}{proposition.104.34}{Dirichlet-GRH-v930-Volume-XX.pdf}}
\expandafter\def\csname r@prop:phase301\endcsname{{126.5}{14}{}{proposition.126.5}{Dirichlet-GRH-v930-Volume-XXIV.pdf}}
\expandafter\def\csname r@prop:plane281\endcsname{{111.15}{71}{}{proposition.111.15}{Dirichlet-GRH-v930-Volume-XXI.pdf}}
\expandafter\def\csname r@prop:positions264\endcsname{{121.4}{33}{}{proposition.121.4}{Dirichlet-GRH-v930-Volume-XXIII.pdf}}
\expandafter\def\csname r@prop:representative294\endcsname{{111.16}{71}{}{proposition.111.16}{Dirichlet-GRH-v930-Volume-XXI.pdf}}
\expandafter\def\csname r@prop:rowfinite262\endcsname{{88.15}{70}{}{proposition.88.15}{Dirichlet-GRH-v930-Volume-XVII.pdf}}
\expandafter\def\csname r@prop:scale284\endcsname{{105.14}{45}{}{proposition.105.14}{Dirichlet-GRH-v930-Volume-XX.pdf}}
\expandafter\def\csname r@prop:shift203\endcsname{{110.11}{47}{}{proposition.110.11}{Dirichlet-GRH-v930-Volume-XXI.pdf}}
\expandafter\def\csname r@prop:short107\endcsname{{9.9}{114}{}{proposition.9.9}{Dirichlet-GRH-v930-Volume-I.pdf}}
\expandafter\def\csname r@prop:split173\endcsname{{72.15}{58}{}{proposition.72.15}{Dirichlet-GRH-v930-Volume-XIV.pdf}}
\expandafter\def\csname r@prop:tracelat253\endcsname{{122.7}{63}{}{proposition.122.7}{Dirichlet-GRH-v930-Volume-XXIII.pdf}}
\expandafter\def\csname r@prop:tracemass301\endcsname{{112.4}{77}{}{proposition.112.4}{Dirichlet-GRH-v930-Volume-XXI.pdf}}
\expandafter\def\csname r@prop:vol127\endcsname{{92.25}{20}{}{proposition.92.25}{Dirichlet-GRH-v930-Volume-XVIII.pdf}}
\expandafter\def\csname r@prop:weight163\endcsname{{86.20}{17}{}{proposition.86.20}{Dirichlet-GRH-v930-Volume-XVII.pdf}}
\expandafter\def\csname r@sec:absolute163\endcsname{{86.7}{15}{}{subsection.86.7}{Dirichlet-GRH-v930-Volume-XVII.pdf}}
\expandafter\def\csname r@sec:char94\endcsname{{121.16}{51}{}{subsection.121.16}{Dirichlet-GRH-v930-Volume-XXIII.pdf}}
\expandafter\def\csname r@sec:exchange336\endcsname{{86.11}{26}{}{subsection.86.11}{Dirichlet-GRH-v930-Volume-XVII.pdf}}
\expandafter\def\csname r@sec:global110\endcsname{{86.10}{23}{}{subsection.86.10}{Dirichlet-GRH-v930-Volume-XVII.pdf}}
\expandafter\def\csname r@sec:identity127\endcsname{{92.8}{18}{}{subsection.92.8}{Dirichlet-GRH-v930-Volume-XVIII.pdf}}
\expandafter\def\csname r@sec:match209\endcsname{{93.4}{28}{}{subsection.93.4}{Dirichlet-GRH-v930-Volume-XVIII.pdf}}
\expandafter\def\csname r@sec:missing\endcsname{{6.6}{32}{}{subsection.6.6}{Dirichlet-GRH-v930-Volume-I.pdf}}
\expandafter\def\csname r@sec:orbits203\endcsname{{110.6}{46}{}{subsection.110.6}{Dirichlet-GRH-v930-Volume-XXI.pdf}}
\expandafter\def\csname r@sec:pathnorm352\endcsname{{119.7}{11}{}{subsection.119.7}{Dirichlet-GRH-v930-Volume-XXIII.pdf}}
\expandafter\def\csname r@sec:vectorchain353\endcsname{{73.2}{74}{}{subsection.73.2}{Dirichlet-GRH-v930-Volume-XIV.pdf}}
\expandafter\def\csname r@seriesguide\endcsname{{1.1}{4}{}{subsection.1.1}{Dirichlet-GRH-v930-Volume-I.pdf}}
\expandafter\def\csname r@thm:columns235\endcsname{{109.38}{32}{}{proposition.109.38}{Dirichlet-GRH-v930-Volume-XXI.pdf}}
\expandafter\def\csname r@thm:common111\endcsname{{9.13}{118}{}{proposition.9.13}{Dirichlet-GRH-v930-Volume-I.pdf}}
\expandafter\def\csname r@thm:complement209\endcsname{{93.14}{32}{}{proposition.93.14}{Dirichlet-GRH-v930-Volume-XVIII.pdf}}
\expandafter\def\csname r@thm:cross132\endcsname{{86.3}{4}{}{proposition.86.3}{Dirichlet-GRH-v930-Volume-XVII.pdf}}
\expandafter\def\csname r@thm:direct76\endcsname{{66.24}{63}{}{proposition.66.24}{Dirichlet-GRH-v930-Volume-XIII.pdf}}
\expandafter\def\csname r@thm:exchangefourier336\endcsname{{86.34}{26}{}{proposition.86.34}{Dirichlet-GRH-v930-Volume-XVII.pdf}}
\expandafter\def\csname r@thm:first250\endcsname{{71.10}{25}{}{proposition.71.10}{Dirichlet-GRH-v930-Volume-XIV.pdf}}
\expandafter\def\csname r@thm:firstsquare255\endcsname{{71.3}{16}{}{proposition.71.3}{Dirichlet-GRH-v930-Volume-XIV.pdf}}
\expandafter\def\csname r@thm:gap132\endcsname{{86.4}{5}{}{proposition.86.4}{Dirichlet-GRH-v930-Volume-XVII.pdf}}
\expandafter\def\csname r@thm:identity127\endcsname{{92.24}{20}{}{proposition.92.24}{Dirichlet-GRH-v930-Volume-XVIII.pdf}}
\expandafter\def\csname r@thm:last244\endcsname{{67.11}{84}{}{proposition.67.11}{Dirichlet-GRH-v930-Volume-XIII.pdf}}
\expandafter\def\csname r@thm:lastcolumn321\endcsname{{123.6}{83}{}{proposition.123.6}{Dirichlet-GRH-v930-Volume-XXIII.pdf}}
\expandafter\def\csname r@thm:levels88\endcsname{{7.11}{97}{}{proposition.7.11}{Dirichlet-GRH-v930-Volume-I.pdf}}
\expandafter\def\csname r@thm:levelsum203\endcsname{{110.13}{48}{}{proposition.110.13}{Dirichlet-GRH-v930-Volume-XXI.pdf}}
\expandafter\def\csname r@thm:markedglobal339\endcsname{{71.25}{41}{}{proposition.71.25}{Dirichlet-GRH-v930-Volume-XIV.pdf}}
\expandafter\def\csname r@thm:matrixfiber295\endcsname{{107.4}{77}{}{proposition.107.4}{Dirichlet-GRH-v930-Volume-XX.pdf}}
\expandafter\def\csname r@thm:midlevel131\endcsname{{86.1}{3}{}{proposition.86.1}{Dirichlet-GRH-v930-Volume-XVII.pdf}}
\expandafter\def\csname r@thm:mix277\endcsname{{90.8}{91}{}{proposition.90.8}{Dirichlet-GRH-v930-Volume-XVII.pdf}}
\expandafter\def\csname r@thm:moment286\endcsname{{105.3}{34}{}{proposition.105.3}{Dirichlet-GRH-v930-Volume-XX.pdf}}
\expandafter\def\csname r@thm:norm91\endcsname{{94.9}{44}{}{proposition.94.9}{Dirichlet-GRH-v930-Volume-XVIII.pdf}}
\expandafter\def\csname r@thm:one61\endcsname{{65.13}{33}{}{proposition.65.13}{Dirichlet-GRH-v930-Volume-XIII.pdf}}
\expandafter\def\csname r@thm:orbits204\endcsname{{110.16}{51}{}{proposition.110.16}{Dirichlet-GRH-v930-Volume-XXI.pdf}}
\expandafter\def\csname r@thm:ordinary163\endcsname{{86.22}{18}{}{proposition.86.22}{Dirichlet-GRH-v930-Volume-XVII.pdf}}
\expandafter\def\csname r@thm:ordinary164\endcsname{{110.2}{39}{}{proposition.110.2}{Dirichlet-GRH-v930-Volume-XXI.pdf}}
\expandafter\def\csname r@thm:pathnorm352\endcsname{{119.11}{11}{}{proposition.119.11}{Dirichlet-GRH-v930-Volume-XXIII.pdf}}
\expandafter\def\csname r@thm:phase282\endcsname{{111.5}{61}{}{proposition.111.5}{Dirichlet-GRH-v930-Volume-XXI.pdf}}
\expandafter\def\csname r@thm:prime210\endcsname{{93.17}{34}{}{proposition.93.17}{Dirichlet-GRH-v930-Volume-XVIII.pdf}}
\expandafter\def\csname r@thm:product281\endcsname{{111.4}{60}{}{proposition.111.4}{Dirichlet-GRH-v930-Volume-XXI.pdf}}
\expandafter\def\csname r@thm:ramifiedcoset324\endcsname{{123.13}{94}{}{proposition.123.13}{Dirichlet-GRH-v930-Volume-XXIII.pdf}}
\expandafter\def\csname r@thm:ramifiedflag324\endcsname{{123.10}{89}{}{proposition.123.10}{Dirichlet-GRH-v930-Volume-XXIII.pdf}}
\expandafter\def\csname r@thm:rangenew50\endcsname{{6.29}{62}{}{proposition.6.29}{Dirichlet-GRH-v930-Volume-I.pdf}}
\expandafter\def\csname r@thm:ranktrace302\endcsname{{112.7}{80}{}{proposition.112.7}{Dirichlet-GRH-v930-Volume-XXI.pdf}}
\expandafter\def\csname r@thm:real282\endcsname{{105.15}{46}{}{proposition.105.15}{Dirichlet-GRH-v930-Volume-XX.pdf}}
\expandafter\def\csname r@thm:region198\endcsname{{20.13}{33}{}{proposition.20.13}{Dirichlet-GRH-v930-Volume-IV.pdf}}
\expandafter\def\csname r@thm:rel155\endcsname{{106.20}{72}{}{proposition.106.20}{Dirichlet-GRH-v930-Volume-XX.pdf}}
\expandafter\def\csname r@thm:saturation220\endcsname{{109.32}{28}{}{proposition.109.32}{Dirichlet-GRH-v930-Volume-XXI.pdf}}
\expandafter\def\csname r@thm:skew289\endcsname{{105.6}{37}{}{proposition.105.6}{Dirichlet-GRH-v930-Volume-XX.pdf}}
\expandafter\def\csname r@thm:splitcompact323\endcsname{{123.7}{84}{}{proposition.123.7}{Dirichlet-GRH-v930-Volume-XXIII.pdf}}
\expandafter\def\csname r@thm:splitmass323\endcsname{{126.4}{13}{}{proposition.126.4}{Dirichlet-GRH-v930-Volume-XXIV.pdf}}
\expandafter\def\csname r@thm:splitpath322\endcsname{{123.9}{87}{}{proposition.123.9}{Dirichlet-GRH-v930-Volume-XXIII.pdf}}
\expandafter\def\csname r@thm:suffix237\endcsname{{110.9}{45}{}{proposition.110.9}{Dirichlet-GRH-v930-Volume-XXI.pdf}}
\expandafter\def\csname r@thm:sum287\endcsname{{111.8}{64}{}{proposition.111.8}{Dirichlet-GRH-v930-Volume-XXI.pdf}}
\expandafter\def\csname r@thm:tail283\endcsname{{111.6}{62}{}{proposition.111.6}{Dirichlet-GRH-v930-Volume-XXI.pdf}}
\expandafter\def\csname r@thm:thetacells340\endcsname{{71.27}{42}{}{proposition.71.27}{Dirichlet-GRH-v930-Volume-XIV.pdf}}
\expandafter\def\csname r@thm:unique221\endcsname{{109.15}{13}{}{proposition.109.15}{Dirichlet-GRH-v930-Volume-XXI.pdf}}
\expandafter\def\csname r@thm:vectorchain353\endcsname{{73.2}{74}{}{proposition.73.2}{Dirichlet-GRH-v930-Volume-XIV.pdf}}
\makeatother
\begin{document}
\setcounter{part}{18}
\fancyhead[L]{\small\color{muted}Dirichlet GRH \textperiodcentered\ Volume XIX}
\fancyhead[R]{\small\color{muted}Local orbital integrals}
\thispagestyle{plain}\begin{center}
{\LARGE\bfseries The generalized Riemann hypothesis\\[2pt]for Dirichlet \(L\)-functions}\\[6pt]
{\small 27 November 2026 \quad\textperiodcentered\quad Edition 930}\end{center}
\part{Local orbital integrals and flag lattices}\label{part:XIX}
\begin{abstract}\noindent We compute orbital integrals at finite places on Hermitian frames. At split auxiliary primes we compute the signed integral on each admissible core position and the elementary divisors for primitive terminal frequencies. We match the auxiliary compact row characters to the flag and determine the largest possible overlap of the original and transported auxiliary theta tests along the chain.\end{abstract}
\setcounter{section}{97}
\noindent The \hyperref[seriesguide]{guide in Volume~I} lists all twenty-seven volumes.
\begingroup\small\setlength{\parskip}{0pt}\tableofcontents\endgroup
\clearpage

\section{Local measures on Hermitian frames}\label{nav:g-IX-0}\label{nav:volume-XIX}

\subsection{The symmetry of an orbital integral}\label{nav:IX-00}
\label{sec:orbital133}

Use the seeds of \eqref{eq:real127} and \eqref{eq:seed127} at
the level of Theorem~\ref{thm:midlevel131}, with \(\beta=\varpi\).
Put \(F=F_+\). In unfolding the product of two periodizations,
fix a rational double coset represented by \(\gamma\in G(F)\).
For a right finite coset \(t_fK^0\) common to \(B_fK^0\) and
\(\gamma B_fK^0\), put
\[
 \Gamma=B(F)\cap\gamma B(F)\gamma^{-1}\cap t_fK^0t_f^{-1},
 \qquad D=N\cap\gamma N\gamma^{-1}.
\]
Intersection with a finite adelic subgroup in this notation
means imposing the indicated condition on the finite component.
Let \(W,W'\) be column-valued functions of types \(\rho,\rho'\)
on the real group, with the transformation in \eqref{eq:real127}.
Write \(vw^*\) for the map \(u\mapsto v\langle w,u\rangle\),
with the inner product conjugate-linear in its first variable.
The real map before evaluation on the finite column covectors is
\begin{equation}
   J_{\gamma,t}=\int_{\Gamma\backslash G_\infty}
                 W(g)\,W'(\gamma^{-1}g)^*\,dg .
   \label{eq:orb133}
\end{equation}
The finite factors in \eqref{eq:seed127} contract this tensor
and give scalar phases and the volume of the finite coset.
Indeed, after unfolding the first rational summation, the
second splits into orbits of its left rational period group.
The stabilizer of the ordered pair of flags is
\(B(F)\cap\gamma B(F)\gamma^{-1}\). On decomposing the finite
variable further into right level cosets its stabilizer is
precisely \(\Gamma\).

\begin{lemma}[The common unipotent character]\label{lem:zero133}
The group \(\Gamma\) is a uniform lattice in \(D_\infty\).
The integral \eqref{eq:orb133} converges absolutely for smooth
\(W,W'\) of compact support modulo \(B_\infty\).
Whenever it converges absolutely, it vanishes unless
\begin{equation}
     \chi_{\varpi,\infty}(d)=
        \chi_{\varpi,\infty}(\gamma^{-1}d\gamma)
                   \qquad(d\in D_\infty).
     \label{eq:relevant133}
\end{equation}
If the equality holds, its value is the covolume of \(\Gamma\)
in \(D_\infty\) times the integral of the same tensor over
\(D_\infty\backslash G_\infty\), with compatible measures.
\end{lemma}
\begin{proof}
Write \(t_f=bk\) in its first finite support. For every
\(\alpha\in B(F)\) stabilizing the finite coset, the projection
to \(H(F)\) stabilizes the column class of the \(H_f\)-component
of \(b\), by Theorem~\ref{thm:midlevel131}. That class has
trivial rational stabilizer by Lemma~\ref{lem:covol127}.
Hence \(\alpha\in N(F)\). Apply the same reasoning to
\(\gamma^{-1}t_fK^0\) for the second flag. It follows that
\(\Gamma=D(F)\cap t_fK^0t_f^{-1}\).

In either upper unipotent group, exponential and logarithm are
mutually inverse finite polynomials and identify the group with
its nilpotent Lie algebra in characteristic zero. On the
intersection they therefore identify \(D\) with the intersection
of the two Lie algebras, an \(F\)-vector space. Thus \(D(F)\) is
dense in \(D_f\) by additive approximation. The compactness of
the adelic quotient of a unipotent group follows by induction
along a central series with vector-group quotients, starting
with the compact additive quotient \(\mathbb A_F/F\). Compact
sets of representatives lift across these quotients in polynomial
exponential coordinates. Reducing its finite coordinate by
density, as in Lemma~\ref{lem:covol127}, identifies the double
quotient by \(D_f\cap t_fK^0t_f^{-1}\) with
\(\Gamma\backslash D_\infty\). This proves compactness.

For \(\alpha\in\Gamma\), both finite unipotent characters are
trivial. To see this for the first, \(b^{-1}\alpha_f b\)
belongs to \(N_f\cap K^0\), and conjugation by \(B_f\) preserves
the character. For the second use \(\gamma^{-1}t_f\).
The adelic character of a rational unipotent element is one, so
the corresponding real characters are also one on \(\Gamma\).
This proves descent in \eqref{eq:orb133}.
On the compact fibre \(\Gamma\backslash D_\infty\), the
integrand varies by the ratio of the two characters. If it is
nontrivial, right translation by a point where it differs from
one shows its integral is zero. Otherwise the integrand is
constant along the fibre. The invariant quotient measures exist
because both \(D_\infty\) and \(G_\infty\) are unimodular.
Disintegration gives the two assertions. Equality of the real
characters is also equivalent to equality of their real
differentials on the common Lie algebra, by applying the
characters along each one-parameter exponential subgroup. In
the frame coordinates these differentials are the real traces
of the successive simple coordinates with the marked terminal
vector. This gives the condition at
\srcloc{Formalism/SharpRestriction/CuspSpanNormBounds/BlockEnds}{79}.

For the asserted convergence, use as scalar envelopes the
functions \(\|W(g)\|\) and \(\|W'(g)\|\) times the indicators of
the respective finite supports before translation. Each descends
to \(B(F)\backslash G(\mathbb A_F)\) because the real
transformations are unitary. Each has compact support on this
quotient. Indeed, there are finitely many finite classes, and
over a compact part of \(B_\infty\backslash G_\infty\) their
period fibres are compact by Lemma~\ref{lem:covol127} and
compactness of \(H_\infty\). The ordinary rational periodization
of each envelope is bounded with compact support on the
automorphic quotient. Locally one can lift each compact support
to a compact set in the adelic group. Only finitely many rational
elements can move this set to meet a given compact set, by
discreteness and closedness of \(G(F)\). Their two periodizations
therefore have integrable product. Unfolding the nonnegative
product by Tonelli bounds each individual double coset and each
of its contributing finite cells, whose finite measure is
positive. The norm of a rank-one tensor is the product of the
two vector norms. This proves absolute convergence of
\eqref{eq:orb133}.
\end{proof}

\begin{theorem}[Restriction to the common compact group]
\leavevmode\label{thm:common133}
Take \(C\) from Theorem~\ref{thm:cross132} for the frames
\(y,\gamma y\). In any absolutely convergent cell, \(J_{\gamma,t}\)
intertwines the two restrictions to \(U(C)_\infty\). If \(C\)
has dimension \(k\), this map vanishes for different words. For
identical words it is a scalar times the unitary identification
of the two irreducible actions, unique up to phase.

If \(\dim C=k-1\), let \(\mathcal M\) be the set of highest weights
of \(U(C)_\infty\) occurring in both restrictions, and let
\(d_\mu\) be the dimension for such a weight. Then the rank
is at most \(\sum_{\mu\in\mathcal M}d_\mu\).
If \(\mathcal M\) is nonempty, write \(d_*=\min_{\mu\in\mathcal M}d_\mu\)
and put
\[
   A_{\gamma,t}=\int_{\Gamma\backslash G_\infty}
                       \|W(g)\|\,\|W'(\gamma^{-1}g)\|\,dg .
\]
The operator and Hilbert--Schmidt norms of \(J_{\gamma,t}\) are at
most \(A_{\gamma,t}/d_*\) and \(A_{\gamma,t}/\sqrt{d_*}\),
respectively.
\end{theorem}
\begin{proof}
The case of unequal unipotent characters follows from
Lemma~\ref{lem:zero133}. For equal characters use its integral
over \(D_\infty\backslash G_\infty\). A member \(h\in U(C)_\infty\)
fixes both frames, normalizes both unipotent radicals, and
normalizes \(D_\infty\). Conjugation on \(D_\infty\) by this
compact group has modulus one. Left translation by \(h\) hence
preserves the quotient measure, whether or not it normalizes
the discrete lattice. In either decomposition in the two period
groups its unipotent part has trivial character. Indeed it fixes
the \(u_i\) and the terminal vector, and so does its projected
column action in the corresponding Witt splitting. Equivariance
of \(W,W'\) now shows that \(J_{\gamma,t}\) intertwines the two
actions. Each identification of \(C\) with its image in a
column space is complex linear and isometric at each real place,
so these actions are unitarily equivalent to the usual
inclusions in the two definite groups. For full dimension
Schur's lemma gives the assertion. Two different words have
different highest weights.

For dimension \(k-1\), ordinary compact branching at one place
gives precisely the weights \(\mu\) for which
\(\lambda_i\ge\mu_i\ge\lambda_{i+1}\) for \(1\le i<k\), each
with multiplicity one \cite[Theorem 2.1]{XIX-XIV-IV-molev106}. Taking the
product over the real places gives the same multiplicity statement
there. An intertwiner vanishes between inequivalent irreducibles
and has image in the sum over the common weights. This proves
the rank bound.

More explicitly, average \(vw^*\) over the compact group with
Haar probability. Between two copies of a common irreducible the
average is a scalar times an isometric identification by Schur's
lemma. Composing with its inverse and taking the trace bounds
the modulus of that scalar by
\(\|v_\mu\|\|w_\mu\|/d_\mu\), where the subscripts denote
orthogonal projections. Distinct blocks have orthogonal domains
and images. Taking the largest scalar bounds the operator norm
by \(\|v\|\|w\|/d_*\). For the Hilbert--Schmidt norm square the
bound on each scalar, multiply it by \(d_\mu\), and sum. The
result is at most \(\|v\|^2\|w\|^2/d_*\).
Averaging each pointwise tensor in the quotient integral leaves
the integral unchanged. Fubini and the triangle inequality give
both estimates, including the fibre covolume.
\end{proof}

For example, for the weights of
Theorem \ref{thm:common111},
write \(\lambda_i=u_i-(k+1-2i)/2\) at one place. The gaps in
\(u\) are at least \(3\Delta\). Interlacing gives, for \(j-i\ge2\),
\[
 \frac{\mu_i-\mu_j+j-i}{j-i}
  \ \ge\ \frac{u_{i+1}-u_j+1}{j-i}
  \ \ge\ \frac{3\Delta+1}{2}.
\]
The adjacent factors are at least one. Weyl's dimension
formula gives \(d_*\ge ((3\Delta+1)/2)^{f\binom{k-2}{2}}\).
Applying the norm estimates to a full period pairing still
requires control of the scalar envelopes and their sum. The
choice of a compact cutoff alone gives no uniform bound for
that sum as it is removed.

\subsection{The rank of an orbital integral between two different words}\label{nav:IX-01}
Use the fields and primes of Theorem~\ref{thm:common111} and its
set \(\mathcal A\) of words. Write \(d(\rho_c)=\dim_\CC V_{\rho_c}\)
for the real-place representation dimension, to distinguish it
from the mass-weighted column number \(d_c\) there. Use the level
with terminal vector \(z=\varpi e\) from
Section \ref{nav:VII-00}.
For any rational pair of frames and a fixed finite cell in the
unfolded pairing, let \(J_{c,c'}\) be the integral of the
rank-one maps from \(V_{\rho_{c'}}\) to \(V_{\rho_c}\) in \eqref{eq:orb133}, before evaluation on finite class
covectors. The assertions below hold for compactly supported
smooth seeds modulo the real period group and, more generally,
whenever the individual tensor integral converges absolutely.

We first examine compact branching at one real place. Write
\(\lambda_i=u_i-(k+1-2i)/2\) for the highest weight of the
word \(c\). On restriction to \(U(k-1)\) the summands are
multiplicity-free. Their weights satisfy
\(\lambda_i\ge\mu_i\ge\lambda_{i+1}\). With
\(y_i=\mu_i+(k-2i)/2\), the choices of \(y_i\) fill the
following product of intervals, with step size one. Explicitly,
\[
    I_i(c)=[u_{i+1}+1/2,u_i-1/2],\qquad 1\le i<k .
\]
Successive intervals are separated by at least one, so each
choice is a valid dominant weight. Give each summand the
probability \(\dim V_\mu/\dim V_{\lambda}\). These probabilities
sum to one by the branching decomposition. By the Weyl dimension
formula they are proportional, on this product, to
\(\prod_{i<j}(y_i-y_j)\). The normalization depends only on \(c\)
and the place.

\begin{lemma}[Conditional mass near either endpoint]\label{lem:end134}
There is an absolute \(\delta>0\) with the following property
at every real place. For \(k\) sufficiently large and
\(k/64\le i\le 63k/64\), fix the other \(k-2\) coordinates at
any allowed values. Each of the two end intervals of \(\Delta\)
consecutive points of \(I_i(c)\) has conditional probability at
least \(\delta\), uniformly in the place, word and \(\Delta\).
\end{lemma}
\begin{proof}
For \(2\le j\le k-2\) the endpoints from \eqref{eq:letters111}
give
\[
 I_j(c)=[4\Delta(k-j-1)-\Delta c_{j+1},\
               4\Delta(k-j)-1-\Delta c_j].
\]
It contains \(w_j=\Delta(4-c_j+c_{j+1})\) points, between
\(3\Delta\) and \(5\Delta\). Its real coordinates differ from
\(4\Delta(k-j)\) by a number in \([-5\Delta,0]\). For real
\(y\in I_i(c)\), separate from the conditional density the two
factors involving the immediate neighbors. The other factors
form \(R(y)=\prod_{|j-i|\ge2}|y-y_j|\). For an interior \(j\)
with \(d=|j-i|\ge2\), the difference \(y-y_j\) differs from
\(4\Delta(j-i)\) by at most \(5\Delta\). Taking reciprocals gives
\[
 \left|\frac1{y-y_j}-\frac1{4\Delta(j-i)}\right|
 \le \frac{5}{4\Delta d(4d-5)} .
\]
The total of these errors is \(O(1/\Delta)\). The signed main
sum on \(2\le j\le k-2\), omitting the two neighbors and \(i\),
is \((H_{k-2-i}-H_{i-2})/(4\Delta)\) with
\(H_n=\sum_{d=1}^n 1/d\). Both indices are bounded below by
\(k/128\) for large \(k\), so their harmonic sums differ by an
absolute constant at most. For \(j=1\), the constraint
\(y_1\ge u_2+1/2\) puts it at distance \(\gg k\Delta\) from
\(I_i(c)\). The other endpoint satisfies
\(y_{k-1}\le u_{k-1}-1/2\) and has the same distance bound.
These one-sided inequalities apply even at the high places.
It follows that \(|(\log R)'|\le C/\Delta\) on the interval for
an absolute \(C\). The ratio of its maximum to its minimum there
is at most \(\exp(5C)\).

Write \(L,U\) for the endpoints of \(I_i(c)\). Of the lowest
\(\Delta\) points take those from \(L+\lfloor\Delta/2\rfloor\)
to \(L+\Delta-1\). The neighbor below is at most \(L-1\).
On these points its distance is at least
\((\lfloor\Delta/2\rfloor+1)/w_i\) times its largest distance
on the whole interval, as is seen by taking the closest allowed
neighbor \(L-1\). This ratio is at least \(1/10\).
The neighbor above is at least \(U+1\), giving the ratio at
least \((w_i-\Delta+1)/w_i\ge 2/3\) for that factor. Reflection
gives the same estimates on at least half the points of the top
slice. In each case the chosen points constitute at least a
tenth of the full interval. Comparing their weights with the
maximum weight proves the lemma with
\(\delta=\exp(-5C)/1000\).
\end{proof}

\begin{theorem}[Small relative rank for every mixed cell]
\label{thm:wordrank134}
There is an absolute \(\kappa>0\) such that, for all sufficiently
large \(x\) on the common intervals and every \(c\ne c'\) in
\(\mathcal A\),
\begin{equation}
       \operatorname{rank} J_{c,c'}
       \ \le\ e^{-\kappa fk}
                \min\{d(\rho_c),d(\rho_{c'})\}.
       \label{eq:wordrank134}
\end{equation}
The bound holds for each fixed rational pair of frames and finite
cell, with the same \(\kappa\) for every pair.
\end{theorem}
\begin{proof}
Theorems \ref{thm:cross132} and \ref{thm:common133}, give either full column rank
for the common compact group, in which case a mixed map is zero,
or a copy of \(U(k-1)\) at each real place. In the latter case,
fix one place and compare its two branching boxes. If \(c_t=0\)
and \(c'_t=1\), the top \(\Delta\) points of \(I_t(c)\) have
been excluded by the upper bound for \(I_t(c')\). If \(c_t=1\)
and \(c'_t=0\), the bottom \(\Delta\) points of \(I_{t-1}(c)\)
have been excluded by the lower bound for \(I_{t-1}(c')\).
The definition of the boxes keeps the same lattice for the two
words. Discarding mismatches whose assigned coordinate is outside
\([k/64,63k/64]\) costs at most \(2k/64+O(1)\) bits. At least
\((k-2)/8\) bits differ by Theorem~\ref{thm:common111}. Each
coordinate can be assigned to at most two bits. For sufficiently
large \(k\) this leaves at least \(k/32\) different coordinates,
each with an excluded end slice.

Under the dimension probability for \(c\), the event of lying in
the box of \(c'\) requires avoidance of all these slices. For
each coordinate, even conditional on all the other coordinates,
its avoidance probability is at most \(1-\delta\) by
Lemma~\ref{lem:end134}. Multiply the indicators of the avoidance
events and condition on all coordinates except one of those
still present. Its indicator can be removed at cost at most
\(1-\delta\), since the other indicators are measurable under
that conditioning. Iteration gives at most
\((1-\delta)^{k/32}\). This bounds the fraction of dimension
in summands also present for \(c'\). Interchanging the two words
proves the same bound relative to the dimension for \(c'\).

The real column space is the tensor product over the \(f\) real
places. Its dimension probabilities on tuples of summands are
products, and the tuple is common precisely when it is common
at each place. Thus \(\kappa=-\log(1-\delta)/32\) works.
The image of the intertwining integral lies in the sum of the
common summands, and its kernel contains the entire orthogonal
complement of those summands in the domain, by Theorem \ref{thm:common133}. This proves \eqref{eq:wordrank134} in both
dimensions.
\end{proof}

\begin{corollary}[Simultaneous kernels for a collection of cells]
\label{cor:many134}
For a fixed word \(c\), take any \(M\ge1\) ordered frame pairs
with finite cells at this level and identify their input column
spaces with \(V_{\rho_c}\) by isometries. At each cell take every
other selected word as an output. There is a subspace of
\(V_{\rho_c}\) annihilated by all these maps whose codimension
is at most
\[
             M\,|\mathcal A|e^{-\kappa fk}\,d(\rho_c).
\]
For \(M\le\exp(C_0 f m^5 x^\eps)\) with fixed \(C_0\), its
relative codimension is at most \(e^{-\kappa fk/2}\) for all
sufficiently large \(x\) on these same intervals.
\end{corollary}
\begin{proof}
Intersect the indicated kernels. The codimension of an
intersection is at most the sum of the individual codimensions
by mapping to the direct sum of the corresponding quotient
spaces. Apply \eqref{eq:wordrank134} to each map. Finally
\(\log|\mathcal A|\le k\log 2\) and \(f\to\infty\) by the field
degree in Proposition~\ref{prop:short107}. The lengths in the
proof of Theorem~\ref{thm:common111} give \(m^5x^\eps=o(k)\).
The claimed exponential bound follows.
\end{proof}

These subspaces may depend on the cells. The corollary estimates
how many dimensions can couple different words in a prescribed
collection. Estimating the full period still requires the
unrestricted sum over cells and the terms pairing a word with
itself. The arithmetic normalization on the integral theta
lattice is a further requirement even when the complex pairing
has been bounded.

\subsection{The finite measure of a nondegenerate Bessel orbit}\label{nav:IX-02}
\label{sec:cellmass205}

Fix \(\gamma\in\mathcal R\) as in Section~\ref{sec:orbits203}.
The rational label \(\gamma\) specifies the two decorated flags.
The untruncated period also sums over finite right level cosets
supporting both seeds. We compute the measure of this second sum.
Let \(W\) be the nondegenerate span of the two frames, let
\(C=W^\perp\), and write \(H_C=U(C)\). For
\(J=B\cap\gamma B\gamma^{-1}\), \(D=N\cap\gamma N\gamma^{-1}\)
and a finite place \(v\) of \(F\), put
\[
 X_v=\{tK_v:t\in B_vK_v,\ \gamma^{-1}t\in B_vK_v\},\qquad
          K_v=K_v^0,\qquad K^0=\prod_{v<\infty}K_v .
\]
Here and below a local algebra \(E_v=E\otimes_F F_v\) may be a
product of two fields. For \(t_fK^0\) in the restricted product
\(X_f\) of the \(X_v\), set \(\Gamma_t=J(F)\cap t_fK^0t_f^{-1}\).
The finite cells over \(\gamma\) are \(J(F)\backslash X_f\).

\begin{lemma}[The stabilizer of the two decorated flags]
\label{lem:stab205}
Let \(h_1,\ldots,h_l\) be the block lengths in
Lemma~\ref{lem:forms203}. For \(U_h\) the upper unitriangular group
of degree \(h\), there are \(F\)-group isomorphisms
\[
 D\simeq \prod_{\nu=1}^l\operatorname{Res}_{E/F}U_{h_\nu},
 \qquad J=D\times H_C .
\]
The number of independent \(E\)-coordinates of \(D\) is
\(a=\sum_\nu h_\nu(h_\nu-1)/2\).
\end{lemma}
\begin{proof}
Choose rational representatives of the two flags in monomial normal
form \(M\). An element of \(J\) acts on them by upper unitriangular
matrices \(U,V\). It preserves \(C\) and \(U^* M V=M\), so
\(V=M^{-1}(U^*)^{-1}M\). In positions \(i<j\), an entry of
\(U^{-1}\) can be nonzero only when \(w(i)>w(j)\), since otherwise
its conjugate would lie strictly below the diagonal of \(V\).
These inversions occur precisely within the individual reversed
blocks. Thus \(U^{-1}\), and hence \(U\), is block upper unitriangular
with no entries between the blocks. Conversely every such \(U\)
gives an upper unitriangular \(V\). These two matrices preserve the
self pairings and their cross pairing. They therefore define an
isometry on \(W\). Extending by the identity on \(C\) gives a
unipotent subgroup \(D'\), parametrized over \(F\) by the stated
product. Any element of \(J\) is uniquely an element of \(D'\)
times an element acting only on \(C\). The two actions commute.

At a real place the image of \(D'\) in either column group is both
unipotent and unitary for a positive definite form, so is trivial.
This gives the same identity of algebraic maps over \(F\), since
the coordinate entries are polynomials and the real points of the
displayed unipotent product are Zariski dense. The action of \(H_C\)
in either column group is faithful by Theorem~\ref{thm:cross132}.
Taking the kernels of the column actions gives \(D=D'\).
\end{proof}

Define a finite set \(S_0\) of places for the level and field. Include
the places over \(2TQ\), those ramified in \(E/F\), and those where the
chosen integral row lattice is not self-dual or \(K_v\) is smaller
than its full stabilizer. Include any remaining places where the
initial frame is not integral. These conditions hold at only finitely
many places, since the level is a product of the prescribed local
modifications of the integral lattice stabilizers. Enlarge \(S_0\)
to \(S(\gamma)\) by adding places where \(\mathcal G\) or a nonzero
endpoint minor \(D_j\) is a nonunit in at least one factor.

\begin{proposition}[The local orbits of finite level cosets]
\label{prop:good205}
The set \(J(F_v)\backslash X_v\) is finite at every finite place.
Outside \(S(\gamma)\) it consists of one orbit. The stabilizer there
is conjugate in \(J(F_v)\) to
\[
       \prod_\nu U_{h_\nu}(\mathcal O_{E_v})\times U(\Lambda_C)
\]
for a self-dual lattice \(\Lambda_C\) in \(C_v\).
\end{proposition}
\begin{proof}
First consider an arbitrary finite place. Write \(tK_v=bK_v
=\gamma b'K_v\) with \(b,b'\in B_v\). The assignment
\(tK_v\mapsto b^{-1}\gamma b'\) induces a bijection
\[
 J(F_v)\backslash X_v \ \longrightarrow\
  (B_v\cap K_v)\backslash\bigl((B_v\gamma B_v)\cap K_v\bigr)
                                      /(B_v\cap K_v).
\]
Indeed the two choices of representatives produce the left and right
compact factors on the right. Equal matrices from two pairs of
representatives give an element of \(J(F_v)\) carrying one \(tK_v\)
to the other. Each matrix in the intersection gives such a pair.

The set \(B_v\gamma B_v\) is closed in \(G(F_v)\). To see this,
take a convergent sequence in the orbit. The cross matrices converge.
At each step of Gaussian elimination the entries before the prescribed
pivot are zero and the pivot is the same fixed nonzero scalar for the
whole sequence. All eliminations are continuous, since they divide
only by that scalar. The limit has the same monomial normal form.
Its concatenated Gram determinant is still \((-1)^d\mathcal G\ne0\).
The two pairs with this normal form are isometric on their spans.
Hermitian Witt extension gives a local isometry of the whole space
preserving both decorated flags in order. In the split case this
extension amounts to extending a basis in the first factor and using
the dual transformation in the second. Thus the limit belongs to
the original double coset.

Each double coset of \(B_v\cap K_v\) is open relative to that orbit.
Here are details at a point \(g_0\) of the orbit. For \(g\) in the
orbit close to \(g_0\), the elimination above chooses representatives
for \((Y,gY)\) continuously. Their concatenated Gram matrix is fixed.
The map between the two normalized lists for \(g_0\) and \(g\) is
close to the identity on their spans. It admits a unitary extension
tending to the identity. In fact, let \(P\) be orthogonal projection
from the old complement to the new one. It tends to the identity by
the projection formula using the inverse of the fixed Gram matrix.
Thus \(P^*P\) on the old complement tends to \(1\). For sufficiently
close \(g\), the binomial series for \((P^*P)^{-1/2}\) converges,
including at places over \(2\), since the \(n\)-th coefficient is
\((-1)^n4^{-n}\binom{2n}{n}\). The map
\(P(P^*P)^{-1/2}\) is the required isometry on the complement.
The resulting extension \(b_g\) preserves the first decorated flag
and carries \(g_0Y\) to the decoration of \(gY\). Both \(b_g\) and
\(g_0^{-1}b_g^{-1}g\) lie in \(B_v\) and tend to the identity.
They lie in \(B_v\cap K_v\) for all sufficiently close \(g\).
This proves openness. The intersection of the closed orbit with
\(K_v\) is compact. Its disjoint open double cosets give the
finiteness assertion.

Now suppose \(v\notin S(\gamma)\) and let \(\Lambda_t=t\Lambda_v\).
The decompositions \(tK_v=bK_v=\gamma b'K_v\) supply integral
representatives \(bY,\gamma b'Y\) for both flags in \(\Lambda_t\).
Their cross matrix is integral. All its Gaussian pivots are units,
because the endpoint ratios are \(h_\nu\)-th powers of the initial
block entries and the endpoints are units. The elimination can
therefore be performed integrally, in both components at a split
place. The resulting \(2d\) vectors have a common, fixed Gram
matrix with unit determinant. Their span over \(\mathcal O_{E_v}\)
is a self-dual lattice in \(W_v\). Pairing with this basis and
inverting its Gram matrix gives the integral orthogonal decomposition
of \(\Lambda_t\) into that lattice and a self-dual lattice in \(C_v\).

For two choices of \(tK_v\) match the two normalized bases on \(W_v\).
This isometry preserves the two decorations. On \(C_v\) their
self-dual lattices are isometric. For an unramified quadratic
extension of odd residue characteristic, the reduction of a
unimodular Hermitian form has a vector of nonzero norm
by polarization. Lift the vector, split off its line, and continue.
Each resulting unit coefficient is a norm of a unit. Indeed the
residue-field norm is surjective and a principal unit in \(F_v\)
has a square root by Hensel's lemma. This proves integral equivalence
to the unit diagonal form. In the split case a lattice in one
factor determines its dual in the other, and any two bases are
related by an invertible linear map. The isometries on \(W_v\)
and \(C_v\) together lie in \(J(F_v)\) and carry one row lattice
to the other. A row lattice determines its right \(K_v\)-coset at
this level, proving transitivity. Nonemptiness follows from the
assumed finite support of \(\gamma\).

In the normalized integral bases, an element on \(W_v\) preserves
the lattice precisely when \(U\) and \(V\) in Lemma~\ref{lem:stab205}
are integral. All block scalars are units, so integrality of \(U\)
already suffices. Two normalized representatives for the same pair
differ on \(W_v\) by an element of \(D(F_v)\), which conjugates the
stabilizer. On \(C_v\) the stabilizer is the full lattice unitary
group, proving the stated description.
\end{proof}

Use the product of additive Haar measures in the entries of
\(D=\prod_\nu \operatorname{Res}_{E/F}U_{h_\nu}\), normalized by
\(\operatorname{vol}(\mathcal O_{E_v})=1\) at finite places and
ordinary complex area at infinity. Write \(c_E=2^{-f}\sqrt{|\Delta_E|}\),
where \(\Delta_E\) is the absolute field discriminant. Choose
product Haar measures on \(G_f,H_C(\mathbb A_F)\), with the
self-dual lattice stabilizers of volume one at the places outside
\(S(\gamma)\). Put
\[
 m_C=\frac{\operatorname{vol}(H_C(F)\backslash H_C(\mathbb A_F))}
                 {\operatorname{vol}(H_C(F_\infty))},\qquad
 L_{\gamma,v}
  =\sum_{tK_v\in J(F_v)\backslash X_v}
    \frac{\operatorname{vol}(K_v)}
         {\operatorname{vol}_{J_v}(J(F_v)\cap tK_vt^{-1})} .
\]
Here \(F_\infty=F\otimes_{\QQ}\RR\) and the Haar measure on \(J_v\)
is the product for \(D_v,H_{C,v}\). Other fixed choices of measure
can be used by retaining the corresponding local volume ratios.

\begin{theorem}[Total finite cell measure of one rational orbit]
\label{thm:mass205}
With the actual level \(K^0\) and its rational column stabilizers,
\begin{equation}
 \sum_{tK^0\in J(F)\backslash X_f}
    \operatorname{vol}(K^0)\operatorname{vol}
                (\Gamma_t\backslash D_\infty)
       =c_E^a\,m_C\prod_{v<\infty} L_{\gamma,v}.       \label{eq:mass205}
\end{equation}
Every factor \(L_{\gamma,v}\) outside \(S(\gamma)\) equals one.
All sums and the product in this formula are finite after removing
these unit factors.
\end{theorem}
\begin{proof}
For a fixed \(t_fK^0\) write
\(K_J=J_f\cap t_fK^0t_f^{-1}\),
\(K_D=K_J\cap D_f\), and \(K_C=\operatorname{pr}_{H_C}(K_J)\).
These are compact open subgroups in the indicated groups.
Lemma~\ref{lem:zero133} gives
\(\Gamma_t=D(F)\cap K_D\) because a rational element stabilizing
the finite coset projects trivially into both column groups at
the marked level.

The rational points of \(D\) are dense in \(D_f\) by additive
approximation in its matrix entries. Additive product measure is
Haar since translation is triangular in the successive superdiagonals,
with diagonal one. For the standard compact open of integral
unitriangular matrices, the real lattice quotient has volume \(c_E^a\).
Indeed the integer lattice of \(E\) has Euclidean covolume \(c_E\)
in the chosen embedding into \(\CC^f\), by the discriminant
determinant formula. Reducing entries in successive superdiagonals
by \(\mathcal O_E\) gives one additive fundamental region per
entry. Earlier superdiagonals stay fixed at each reduction.
For any compact open \(K_D\), intersect it with the standard one.
Density identifies the two finite group indices with the indices of
their rational intersections. The comparison of covolumes gives
\[
        \operatorname{vol}(\Gamma_t\backslash D_\infty)
                     =c_E^a/\operatorname{vol}(K_D).
\]

The rational cells above one \(J_f\)-orbit are
\(J(F)\backslash J_f/K_J\). Projection to \(H_C\) identifies
this set with \(H_C(F)\backslash H_{C,f}/K_C\). To prove injectivity,
lift a right \(K_C\) factor to \(K_J\) and a left rational column
factor to \(H_C(F)\subset J(F)\). Once the column components agree,
density of \(D(F)\) lets one align the remaining components modulo
\(K_D\). Surjectivity follows by inclusion of \(H_{C,f}\).
Each of the column class stabilizers is trivial. Indeed for a
representative \(j_f\) and a rational column stabilizer, the possible
unipotent components of its lifts in \(j_fK_Jj_f^{-1}\) form a
nonempty open coset. Choose a rational component by density.
The resulting rational element of \(J\) stabilizes \(j_ft_fK^0\)
and belongs to \(D(F)\) by Lemma~\ref{lem:zero133}, so its column
component was the identity.

The column quotient is compact at infinity and has finite class
number, by finiteness of the genus of a positive definite Hermitian
lattice and commensurability of compact open subgroups. Its class
number at \(K_C\) is therefore \(m_C/\operatorname{vol}(K_C)\).
Indeed each class contributes the volume of
\(H_C(F_\infty)\times K_C\) to the adelic quotient by the just
proved triviality of stabilizers. Across these representatives
\(K_D\) changes only by conjugation in the unimodular group \(D_f\).
Finally \(\operatorname{vol}_{J_f}(K_J)=
\operatorname{vol}_{D_f}(K_D)\operatorname{vol}_{H_{C,f}}(K_C)\)
by the product group formula, since each nonempty fibre is a coset
of \(K_D\). These identities give the contribution of one
\(J_f\)-orbit. Its choices separate over finite places. Almost
everywhere \(\gamma_v\in K_v\), and Proposition~\ref{prop:good205}
gives a unique local orbit with stabilizer of volume one. Thus the
sum over adelic orbits factors into the asserted finite product.
\end{proof}

For the actual seeds \eqref{eq:seed127}, write \(B_\ell\) and
\(B_{\ell'}\) for the maximal norms of the two covectors on the
finite column class sets. Let \(J_{\gamma,t}\) now denote the
scalar contribution including its finite measure and both covectors.
The common real quotient gives the explicit upper bound
\begin{align}
 \sum_{t\in J(F)\backslash X_f}|J_{\gamma,t}|
  &\le B_\ell B_{\ell'}\,c_E^a m_C
          \Bigl(\prod_{v\in S(\gamma)}L_{\gamma,v}\Bigr)
          I_\infty(\gamma),                              \label{eq:amp205}\\
 I_\infty(\gamma)
  &=\int_{D_\infty\backslash G_\infty}
       \|W_1(g)\|\,\|W_2(\gamma^{-1}g)\|\,d\dot g .       \notag
\end{align}
Indeed the unipotent characters have absolute value one, the real
product of norms is \(D_\infty\)-invariant, and each contraction of
the rank-one operator is bounded by this product times the two
covector norms. Disintegration and \eqref{eq:mass205} give the
inequality. The untruncated inputs and their convergence before
finite contraction are treated in Theorem~\ref{thm:ordinary163}.
All quotient measures in \eqref{eq:amp205} use the stated Haar on
\(D_\infty\).

\begin{theorem}[The primes of a nonunit local factor]\label{thm:primecost205}
Put \(P_\gamma=\prod_{v\in S(\gamma)\setminus S_0}\mathrm N v\),
where \(\mathrm N v\) is the residue-field size. Then
\(P_\gamma\le w_\gamma=w(D(\gamma))\). For \(0\le\delta<1/2\)
and all sufficiently large \(Q\) in terms of \(\delta\),
\begin{equation}
   \sum_{\gamma\in\mathcal R} P_\gamma^\delta w_\gamma^{-1}
      \le\left(\frac{128\log(2Q)}
                         {Q^{4(1-\delta)}}\right)^{fd}.
                                                      \label{eq:primecost205}
\end{equation}
\end{theorem}
\begin{proof}
Every extra place divides the norm to \(F\) of some endpoint
minor or divides \(\mathcal G\) in \(F\). All these quantities
are nonzero algebraic integers. Thus
\[
 P_\gamma\le
 |N_{E/\QQ}(\textstyle\prod_{j\in I}D_j)|\,N_{F/\QQ}(\mathcal G).
\]
The first factor is at most \(w_\gamma^{1/2}\) by the embedding
formula. At each real place \(0<\tau\mathcal G\le|\tau D_d|^2\),
which bounds the second factor by \(w_\gamma^{1/2}\).
It suffices to sum \(w_\gamma^{-1+\delta}\).

In the proof of Theorem~\ref{thm:levelsum203} retain all the boxes,
separations \(P_j\), and multiplicities. For \(b_1=2(1-\delta)>1\)
replace the one-coordinate density by
\(2b_1 r(1+r^2)^{-b_1-1}\,dr\), whose integral from \(u\) to
infinity is \((1+u^2)^{-b_1}\). For \(0\le t<2b_1-2\) the new
moment is \(M=\int_0^\infty 2b_1 r^{3+t}(1+r^2)^{-b_1-1}\,dr\).
The same volume and binomial bounds give
\((9M+2)^{fd}Q^{-(2+t)fd}\). Take
\(t=2b_1-2-\eta\) with \(\eta=1/\log(2Q)\), which is nonnegative
for all sufficiently large \(Q\). Integrating over \(r\le1\) and
\(r>1\) gives \(M\le1+4/\eta\), since \(b_1\le2\).
The inequality \(Q^\eta\le e\) gives \eqref{eq:primecost205}.
\end{proof}

For example \(\delta=1/4\) allows a factor \(P_\gamma^{1/4}\)
with \(Q^{-3fd}\) up to the displayed logarithm. This estimate
controls which primes can contribute. The sizes of the integrals
\(L_{\gamma,v}\) at those primes, the mass \(m_C\), and the real
integral in \eqref{eq:amp205} also enter the mixed period.

\begin{proposition}[Supported orbits outside the regular minor locus]
\label{prop:oneblock205}
For each period field, translate and level in Section~\ref{sec:orbits203} are infinitely many distinct relevant rational double cosets
with \(\gamma_f\in K^0\), \(\mathcal G\ne0\), and \(D_j=0\)
for every \(1\le j<d\). Their common unipotent group is
\(\operatorname{Res}_{E/F} U_d\).
\end{proposition}
\begin{proof}
Use the Witt basis in Section~\ref{sec:char94}, with the last
isotropic basis \(v_1,\ldots,v_s\) ordered dually to the first.
Let \(W_{\rm mid}=W_0\perp Ee\). Choose a positive coordinate
vector \(w\in W_0\) of norm one in the signed row space. Take
\(a_0\in E^\times\) with \(\bar a_0=(-1)^d a_0\), choosing
\(a_0=1\) for even \(d\) and any nonzero element of trace zero
for odd \(d\). Set \(\alpha=\bar a_0/\bar\varpi\). Define
\(A:E^s\longrightarrow W_{\rm mid}\) by sending the first basis
vector to \(\alpha(w+e)\) and the others to zero. With \(A^\dagger\)
given by the Hermitian pairing of its columns with a middle vector,
we have \(A^\dagger A=0\), since \(w+e\) is isotropic.
Let \(L\) be the \(s\)-square matrix with entries
\[
       L_{i,d+1-i}=(-1)^{i-1}a_0\quad(2\le i\le d-1)
\]
and all other entries zero. The choice of \(\bar a_0\) gives
\(L^*=-L\). In the ordered three summands define
\[
   g_n=\begin{pmatrix}
     1 &0&0\\ n A&1&0\\ nL&-n A^\dagger&1
   \end{pmatrix}\qquad(n\in\QQ).
\]
The Witt matrix pairs the first and third summands by the identity.
Direct multiplication, using \(L+L^*=A^\dagger A=0\), gives
\(g_n\in G(F)\) and \(g_n g_m=g_{n+m}\).
The induced map on finite adeles of \(\QQ\) is continuous since its
entries are polynomials with fixed algebraic coefficients integral
at almost all places. The preimage of \(K^0\) is an open neighborhood
of zero. It contains \(N\widehat{\ZZ}\) for some \(N\ge1\).
Take positive integral multiples \(n\) of \(N\).

For these translates \(Z=h(Y,g_nY)=M_n-TE_{dd}\), where \(M_n\)
is the full reversed block with \(i\)-th row entry
\((-1)^{i-1}n a_0\). Indeed the middle antidiagonal entries come
from \(nL\). The top right entry is \(n\bar\alpha\varpi=na_0\)
and the bottom left is \(-n\bar\varpi\alpha=-n\bar a_0\).
The terminal vector keeps its middle component. Adding \(T/(na_0)\)
times the first row to the last row eliminates the extra diagonal
entry. Thus \(D_d=(na_0)^d\), all proper leading minors vanish,
and \(\mathcal G=D_d\overline{D_d}\) is totally positive.
The distinct positive \(n\) give distinct final determinants.

Finally on this full reversed block any upper unitriangular \(U\)
gives an upper \(V=M_n^{-1}(U^*)^{-1}M_n\). Its simple entries
in reverse order are the conjugates of those of \(U\), since the
ratio of consecutive antidiagonal entries is \(-1\). Their real
traces agree. Lemma~\ref{lem:stab205} describes the entire common
unipotent group, so this checks relevance on that group.
\end{proof}

The formal-degree orbital identity in
\srcloc{Development/AbelianVarietyTheory/Proofs/FlagMeasures/BesselOrbitalIntegralEqInvDegSmulStatement}{7}
takes its group element in the regular set defined in
\srcloc{Development/AbelianVarietyTheory/Definitions/SemilocalFrames/MatrixOfBesselTwoSidedAction}{77}.
That definition requires each proper leading minor and the Gram
determinant to be nonzero. Proposition~\ref{prop:oneblock205}
exhibits supported, character-compatible data outside its domain.
In the convention using the
imaginary parts of the simple coordinates, one can realize constant
antidiagonal entries directly. Choose a nonzero \(a_0\in E\) of
trace zero, and set each displayed antidiagonal entry of \(L\) equal
to \(a_0\). Now \(L^*=-L\) and the two extreme entries are both
\(na_0\). The simple entries of \(V\) are the negatives of the
conjugates of those of \(U\) in reverse order. Their imaginary
parts agree. All proper leading minors still vanish. The definition
of the regular set uses only the minors and the Gram determinant.

\begin{corollary}[Absolute integrability on the whole real group]
\label{cor:divergence205}
Let \(W_1,W_2\) be two nonzero smooth Bessel sections belonging
to the prescribed real discrete series and finite under the right
maximal compact group. For every translate in
Proposition~\ref{prop:oneblock205},
\[
     \int_{G_\infty}\|W_1(g)\|\,\|W_2(\gamma^{-1}g)\|\,dg
                         =+\infty .
\]
\end{corollary}
\begin{proof}
Each section is real analytic. To verify this for scalar components,
decompose the vector into its finitely many right compact types.
On the Lie algebra of the real group use the invariant real trace
form, negative on the compact summand \(\mathfrak k\) and positive
on the complementary summand \(\mathfrak p\). In orthonormal bases
for these signed forms the real group Casimir is
\(\Omega=\sum_{\mathfrak p}X_i^2-\sum_{\mathfrak k}Y_i^2\).
It acts by the infinitesimal character of the prescribed representation.
The compact Casimir \(\sum_{\mathfrak k}Y_i^2\) acts by a scalar
on each compact type. Hence each scalar component in one type is
an eigenfunction of the analytic differential operator
\(R(\Omega)+2\sum_{\mathfrak k}R(Y_i)^2\). Its principal part is
the sum of squares of a full invariant tangent frame and is elliptic.
Analytic elliptic regularity gives the assertion
(C.~B. Morrey and L.~Nirenberg, \emph{On the analyticity of the
solutions of linear elliptic systems of partial differential
equations}, Ann.\ Scuola Norm.\ Sup.\ Pisa \textbf{11} (1957),
271--290).

The real group is connected. Indeed every factor has a Cartan
decomposition with connected maximal compact group and connected
exponential factor. A nonzero analytic scalar component therefore
has open dense nonvanishing set. The two sections have positive
product of norms on a nonempty open set even after translation.
This product is invariant under left \(D_\infty\), since the
two corresponding unipotent groups act through unitary characters.
Thus it is positive on an open set of positive quotient measure.
Here \(D_\infty\) is a noncompact unipotent group of positive
dimension by Proposition~\ref{prop:oneblock205}, so its Haar
measure is infinite. The quotient integral formula and Tonelli's
theorem give the divergence.
\end{proof}

The integral in \eqref{eq:amp205} is over
\(D_\infty\backslash G_\infty\). The rational lattice contributes
the finite covolume in \eqref{eq:mass205}. The value of each signed
projected matrix entry still depends on its sections and covectors.

\subsection{The Bessel integral on a nonregular orbit}\label{nav:IX-03}
\label{sec:quot207}

Fix a real place of a field in Section~\ref{sec:absolute163}. Write
\(d=s+1\), \(r=k+2d-1\), and let \(B=N\rtimes U(k)\) be the
stabilizer of the decorated frame \(Y=(u_1,\ldots,u_{d-1},z)\).
The negative index is \(d\). We use the same Haar measures on
\(G\) and \(B\) as in Theorem~\ref{thm:rel155}. For two words,
write \(E_i\) for their finite-dimensional Hilbert spaces at this
place and \(\tau_i(nh)=\chi_i(n)\rho_i(h)\) for their unitary
actions of \(B\). The row representation \(\sigma\) is the same.
Let \(D=N\cap\gamma N\gamma^{-1}\), and write
\(d_0^\gamma=\gamma^{-1}d_0\gamma\) for \(d_0\in D\).
We assume that the two frames span a nondegenerate space with
positive definite orthogonal complement and that their characters
agree on \(D\). These hypotheses hold for all the relevant orbits
of nonzero Gram determinant in Section~\ref{sec:orbits203}.
The real versions of Lemmas~\ref{lem:forms203} and~\ref{lem:stab205}
follow by the same matrix operations at this place.

\begin{lemma}[Volume on a two-sided orbit]\label{lem:orbitvol207}
Give \(D\) any Haar measure, and give
\(\mathcal X_\gamma=(B\times B)/\{(d_0,d_0^\gamma):d_0\in D\}\)
the corresponding right quotient measure. With \(w(g)=\|g\|_{\rm op}\)
in signed orthonormal coordinates, for \(R_0\ge1\),
\[
 \operatorname{vol}\{[b_1,b_2]\in\mathcal X_\gamma:
                   w(b_1\gamma b_2^{-1})\le R_0\}
       \ \ll_\gamma\ R_0^{\,8r^2(d^2+2)} .
\]
The constant can depend on the actual real pair and the Haar measures.
\end{lemma}
\begin{proof}
Choose representatives \(Y_1,Y_2\) of the two copies of the
standard decorated frame so that \(h(Y_1,\gamma Y_2)=M\) is
monomial. These are fixed changes by upper unitriangular matrices.
Put \(x=b_1\gamma b_2^{-1}\) and write
\(b_iY_i=Y_i U_i\). Then
\[
             h(Y_1,xY_2)= (U_1^{-1})^* M U_2^{-1}.
\]
All entries on the left are \(O_\gamma(R_0)\).
We first obtain a decomposition of this matrix into a lower
unitriangular matrix, \(M\), and an upper unitriangular matrix
such that the two triangular factors and their inverses have
norm \(O_\gamma(R_0^{d^2})\).

Use the elimination in Lemma~\ref{lem:forms203}. If the pivot
columns in successive rows are \(w_1,\ldots,w_d\), order the
columns temporarily in this order to calculate the residual matrix
after each elimination. Its entry at remaining row \(i\), column
\(j\), after \(t\) pivots is the determinant with ordered rows
\(1,\ldots,t,i\) and ordered columns \(w_1,\ldots,w_t,j\),
divided by the determinant with the first \(t\) rows and pivot
columns. This follows by subtracting the linear combination of
the first \(t\) rows that clears the selected entries, and expanding
the bordered determinant in its last row. The denominator is the
product of the first \(t\) pivots. Each pivot equals the fixed
entry in \(M\) by uniqueness of the monomial form. Thus each
eliminating multiplier is \(O_\gamma(R_0^d)\). In the original
column order only columns to the right of the current pivot require
clearing. At one step each of the two eliminating matrices is the
identity plus a matrix supported on one row or one column with
square zero. Its inverse has the same bound. At most \(d\) such
matrices are multiplied on either side, proving the assertion.

We can arrange that the \(U_i\) of an actual pair \(b_i\) realize
this decomposition. Indeed two triangular decompositions with the
same \(M\) differ by simultaneous right multiplication of the frame
matrices \(U_i\) by \(A_1,A_2\) with \(A_1^* M A_2=M\).
The isometry on the two spans acts on \(Y_1,\gamma Y_2\) by these
matrices. Extend it by the identity on the orthogonal complement.
The proof of Lemma~\ref{lem:stab205} puts this element in \(D\).
Thus the change comes from right multiplication by
\((d_0,d_0^\gamma)\) in the required quotient.

For the resulting representative both \(b_1Y_1=Y_1 U_1\) and
\(b_1\gamma Y_2=xY_2 U_2\) are \(O_\gamma(R_0^{d^2+1})\).
Their original spans contain a fixed maximal negative subspace.
A bound for the images of a basis of such a subspace bounds the
operator norm of an isometry up to a fixed constant. To see this,
first take the negative coordinate space, and write an isometry
in positive and negative blocks as
\(\left(\begin{smallmatrix}A&B_0\\ C&E\end{smallmatrix}\right)\).
The relations \(AA^*=1+B_0B_0^*\) and \(CC^*=EE^*-1\)
bound the other two blocks by the negative columns. A fixed
isometry taking that coordinate space to the chosen negative
subspace accounts for the change of basis.
Hence \(w(b_1)\ll_\gamma R_0^{d^2+1}\) and
\(w(b_2)=w(x^{-1}b_1\gamma)\ll_\gamma R_0^{d^2+2}\).

Take a compact identity neighborhood of positive measure in \(D\).
Right multiplication above by this neighborhood multiplies the
two bounds by at most a fixed constant. Along each fibre in the
stated set the enlarged product of balls therefore has measure
bounded below by the measure of this neighborhood. The quotient
integration formula bounds its quotient volume by the product of
the two Haar ball volumes divided by that positive number.
Each ball of radius \(S_0\ge1\) in \(B\) has measure
\(O(S_0^{4r^2})\) by Lemma~\ref{lem:translate163}.
\end{proof}

For column vectors \(e_i\), \(e_1 e_2^*\) denotes the map
\(u\mapsto\langle u,e_2\rangle e_1\), with inner products linear
in the first variable. Given two left equivariant sections put
\[
 K_\gamma(W_1,W_2)=
 \int_{D\backslash G} W_1(g)W_2(\gamma^{-1}g)^*\,d\dot g .
\]
Equality of the two unipotent characters makes the integrand
independent of its representative. For an operator-valued
coefficient \(F\) of \(\sigma\), put
\[
 \mathcal O_\gamma(F)=
  \int_{\mathcal X_\gamma}
   \tau_1(b_1)^{-1} F(b_1\gamma b_2^{-1})\tau_2(b_2)
                                                     \,d[b_1,b_2].
\]
Here both \(\tau_i\) act by matching scalars on the two images
of \(D\). These quotients also occur in
\srcloc{Formalism/SharpRestriction/CuspSpanNormBounds/RootMeasBig}{27}
and line~53 there, with the two integration variables inverted.

\begin{theorem}[Schur orthogonality on the common unipotent quotient]
\label{thm:quot207}
Let \(J_i v(g)=\sqrt{d_\sigma}\Lambda_i(\sigma(g)v)\) be the
isometric sections of Theorem~\ref{thm:rel155}. Choose compact-finite
\(u_i\) and \(\xi_i\in E_i\) such that
\(\ell_i=\langle\Lambda_i(u_i),\xi_i\rangle\ne0\).
Such choices exist for each word. Set
\[
          F(g)=\langle\sigma(g)u_2,u_1\rangle\,\xi_1\xi_2^* .
\]
For compact-finite \(v_1,v_2\), on every sufficiently large field
with the spacings of Section~\ref{sec:absolute163}, one has
\begin{equation}
 K_\gamma(J_1 v_1,J_2 v_2)
     =\frac{\langle v_1,v_2\rangle}
                 {\overline{\ell_1}\ell_2}\,\mathcal O_\gamma(F).
                                                        \label{eq:quot207}
\end{equation}
Both integrals are absolutely convergent, including for a single
reversed block of length \(d\).
\end{theorem}
\begin{proof}
For \(a_i(g)=\langle\sigma(g)v_i,u_i\rangle\), start with
\[
        W_i(g)=\int_B a_i(bg)\tau_i(b)^{-1}\xi_i\,db.
\]
Lemma~\ref{lem:translate163} and the proof of
Proposition~\ref{prop:weight163} give absolute convergence.
Testing against \(z\in E_i^\vee\), the relative character in the
proof of Theorem~\ref{thm:rel155} gives
\(z(W_i(g))=\overline{\ell_i}\,z(\Lambda_i(\sigma(g)v_i))\).
Indeed the compact factor of that scalar coefficient pairs \(z\)
with \(\rho_i(h)^{-1}\xi_i\), and the unipotent factor is
\(\overline{\chi_i(n)}\). Nonvanishing of the functional and density
of compact-finite smooth vectors give the choices of \(u_i,\xi_i\).

There are two ways to integrate on the space with coordinates
\((b_1,b_2,g)\) and identification
\((b_1d_0,b_2d_0^\gamma,g)\sim(b_1,b_2,d_0g)\).
One first integrates in \(B^2\) and then in \(D\backslash G\).
The other integrates in \(G\) and then in \(\mathcal X_\gamma\).
For clarity, first disintegrate \(B^2\) by the right diagonal
action in the first integral. In the fibre move \(d_0\) to
the left of \(g\) by the identification. The fibre integral
followed by the \(D\backslash G\) integral is precisely Haar
integration on \(G\). Tonelli gives this identity for nonnegative
functions. The groups involved are unimodular, since \(N,D\)
are unipotent and the compact groups preserve their Haar measures.

Apply the nonnegative identity to
\(|a_1(b_1g)a_2(b_2\gamma^{-1}g)|\). In the second integration
order put \(y=b_2\gamma^{-1}g\), \(x=b_1\gamma b_2^{-1}\).
Write \(m_*=8r^2(d^2+2)+1\). For the family at hand,
\(M-m_*>4r^2\) for sufficiently large \(x\) in the interval
parameter. Indeed \(M=8(f+1)(r+1)^3+8r^2+4\) and the ranks
in Section~\ref{sec:absolute163} satisfy \(r/d^2\to\infty\).
Here \(2d=m_0\asymp x^{1/2-\eps/16}\) and \(k\asymp x^{9/2+10\eps}\) for the selected family. These give \(r/d^2\to\infty\).
Since \(w(x)\le w(xy)w(y)\), the product of the two coefficient
majorants is bounded by
\[
                C\,w(x)^{-m_*}w(y)^{-(M-m_*)}.
\]
The \(y\)-integral is finite by Lemma~\ref{lem:translate163}.
Lemma~\ref{lem:orbitvol207} and summation over dyadic balls give
finiteness of the remaining integral. The same ball estimate
applied to \(F\) proves absolute convergence of its orbital
integral. Thus Fubini applies coordinate by coordinate with the
unitary column factors inserted.

Schur orthogonality now gives
\[
 \int_G a_1(xy)\overline{a_2(y)}\,dy
   =d_\sigma^{-1}\langle v_1,v_2\rangle
                    \langle\sigma(x)u_2,u_1\rangle .
\]
This follows by moving \(\sigma(x)^{-1}\) onto \(u_1\) in the
first coefficient. Multiplication by the column factors and
integration on \(\mathcal X_\gamma\) evaluates \(K_\gamma(W_1,W_2)\).
Substituting \(W_i=\overline{\ell_i}\,J_i v_i/\sqrt{d_\sigma}\)
gives \eqref{eq:quot207}.
\end{proof}

\begin{proposition}[A weighted bound with its field fixed]
\label{prop:tail207}
On the original level \(K^0\), write \(J_{\gamma,c}\) for the scalar
contribution of a finite cell as in \eqref{eq:amp205}. Let
\(H_1(\gamma)=\prod_{\tau}\max(1,|\tau\zeta_1|)\), where \(\zeta_1\)
is the scalar in the first reversed block. For \(p,t\ge0\) with
\[
            h=2p d(d+1)+t<L-4fr^2
\]
one has, for all relevant supported nondegenerate labels,
\begin{equation}
 \sum_{\gamma,c} w_\gamma^p H_1(\gamma)^t |J_{\gamma,c}|<\infty,
 \qquad
 \sum_c |J_{\gamma,c}|
       \le C_{x,p,t}\,Q^{-ft} w_\gamma^{-p}.             \label{eq:tail207}
\end{equation}
Here one constant works for every label and for unit input
vectors in the prescribed finite-dimensional spaces at this \(x\).
\end{proposition}
\begin{proof}
The first \(j\) frame vectors have exterior product \(Y^{(j)}\)
fixed by \(B\). In signed coordinates the norm of the Hermitian
pairing on exterior powers is one for the Euclidean exterior norm.
Evaluate \(D_j\) as the pairing between
\(g^{-1}Y^{(j)}\) and \(g^{-1}\gamma Y^{(j)}\). At every place,
\[
 |\tau D_j|
 \le \|Y^{(j)}\|^2
       \{\nu_\tau(g)\nu_\tau(\gamma^{-1}g)\}^{j}.
\]
Indeed one can vary the representative on each side independently
and \(w(a)=w(a^{-1})\). With
\(C_Y=\prod_{j=1}^d(1+\|Y^{(j)}\|^4)^2\) this bounds \(w_\gamma\)
by \(C_Y^f\{V(g)V(\gamma^{-1}g)\}^{2d(d+1)}\).
The Euclidean squared norms here are \(2^j\) for \(j<d\) and
\(2^{d-1}T\) for \(j=d\).

In a monomial matrix with first block length \(h_1\), every
entry with \(i+j<h_1+1\) is zero. Consequently the entries on
the first antidiagonal \(i+j=h_1+1\) in that block stay fixed
under both unitriangular changes. At least one such entry pairs
two isotropic frame vectors, of Euclidean norm \(\sqrt2\) each.
Take \(i=1,j=h_1\) for \(h_1<d\), and \(i=2,j=d-1\)
for \(h_1=d\), using \(d\ge3\) on this family. The same two
infima as above give \(|\tau\zeta_1|\le
2\nu_\tau(g)\nu_\tau(\gamma^{-1}g)\).
Moreover \((\zeta_1/Q)^{h_1}\) is integral, by
Proposition~\ref{prop:shift203}. When \(h_1=d\) the preceding
leading minor vanishes. Integrality of \(\zeta_1/Q\) follows
because the ring of integers is integrally closed. Taking its
nonzero absolute norm gives \(H_1(\gamma)\ge Q^f\).

For each section use its positive seed
\({\bf1}_{B_fK^0}\|W_i\| V^h\) on the adelic group.
Its periodization over \(B(F)\backslash G(F)\) is bounded because
Lemma~\ref{lem:cosets163} gives \(O_x(R_0^{4fr^2})\) indices in a
height ball, while Proposition~\ref{prop:weight163} gives decay
\(O_x(R_0^{h-L})\) on a dyadic shell. It is integrable by unfolding
to the finite list of seed classes and using the weighted integral
in that proposition. It therefore belongs to \(L^2\). Multiply
the two positive periodizations, integrate, and unfold by Tonelli
as in Theorem~\ref{thm:ordinary163}. The weight inequalities bound
the first sum in \eqref{eq:tail207} by the product of the two
\(L^2\) norms times \(2^{ft} C_Y^{fp}B_\ell B_{\ell'}\), with the
finite covector norms of \eqref{eq:amp205}. Retaining one label
and using \(H_1\ge Q^f\) gives the second inequality. At fixed
\(x\), expansion in orthonormal bases of the finitely many input
spaces gives a common constant for the unit vectors. This proof
does not bound its growth with \(x\).
\end{proof}

\subsection{Change of frame and the arithmetic cell measure}\label{nav:IX-04}
\label{sec:frames207}

Keep the two sections in \eqref{eq:quot207} and the Haar measure
on \(G\). At one real place allow terminal norms \(T_1,T_2>0\)
on the two sides. The application to \eqref{eq:amp205} has
\(T_1=T_2=T\). Let \(j_0=0<j_1<\cdots<j_l=d\) be the nonzero
minor endpoints, \(h_\nu=j_\nu-j_{\nu-1}\), and \(\zeta_\nu\)
the first antidiagonal entry in block \(\nu\). Put
\[
 \begin{split}
 S_j&=j(2d-j-1),\qquad S=S_d=d(d-1),\\
 e_\nu&=2d-j_{\nu-1}-j_\nu-1,\qquad
 a=\sum_\nu h_\nu(h_\nu-1)/2,\qquad
 H=\sum_\nu h_\nu(h_\nu^2-1)/6 .
 \end{split}
\]
Write \(\mathcal G=|D_d|^2-T_1T_2|D_{d-1}|^2\).
The local density \(P\) is the product of
\(|\zeta_\nu|^{2h_\nu(k-1+e_\nu)}\) over blocks, with the factor
for a last block of length one replaced by
\((|\zeta_l|^2-T_1T_2)^{k-1}\). This is the positive density at
the realized points of nonzero Gram determinant in
\srcloc{Formalism/SharpRestriction/CuspSpanNormBounds/RelevantZeroGramTerms}{240}.

\begin{lemma}[The local density in endpoint minors]\label{lem:density207}
Every such orbit satisfies
\begin{equation}
 P=\mathcal G^{\,k-1}
       \prod_{\nu<l}|D_{j_\nu}|^{2(h_\nu+h_{\nu+1})}
                         |D_d|^{2(h_l-1)}.              \label{eq:density207}
\end{equation}
For a relevant rational orbit supported at the prescribed finite
level, let \(P_\infty=\prod_{\tau}P_\tau\) in the arithmetic
frames. Then \(P_\infty^{1/2}\ge Q^{fd(k+d-2)}\).
\end{lemma}
\begin{proof}
For each block \(|\zeta_\nu|^{h_\nu}=|D_{j_\nu}/D_{j_{\nu-1}}|\),
where \(D_0=1\). The factors involving \(k-1\) multiply to
\(\mathcal G^{k-1}\). Indeed, if \(h_l>1\), then \(D_{d-1}=0\).
If \(h_l=1\), then \(\mathcal G=|D_{d-1}|^2
(|\zeta_l|^2-T_1T_2)\) and \(e_l=0\).
For the remaining factors telescope the endpoint ratios using
\(e_\nu-e_{\nu+1}=h_\nu+h_{\nu+1}\) and \(e_l=h_l-1\).
This proves \eqref{eq:density207}.

By Theorem~\ref{thm:gap132}, \(N_{F/\QQ}\mathcal G\ge Q^{2df}\).
For each proper endpoint the product of the \(f\) absolute values
of \(D_j\) is at least \(Q^{jf}\), since \(D_j/Q^j\) is a
nonzero integer of \(E\) by Proposition~\ref{prop:shift203}.
The same applies to \(D_d\) when \(h_l>1\), since its shifted
congruence then has \(D_{d-1}=0\). Its exponent is zero in the
other case. The sum of the endpoint exponents in the square root,
weighted by their indices, is
\[
 \sum_{\nu<l}j_\nu(h_\nu+h_{\nu+1})+d(h_l-1)
       =\sum_\nu h_\nu e_\nu=\sum_{i=1}^d 2(d-i)=S .
\]
The middle equality uses the fact that \(e_\nu\) is the mean of
\(2(d-i)\) on the \(\nu\)-th block. Combining the two norm
bounds gives the assertion.
\end{proof}

Use the character \(\exp(i\Re(c\sum_{i<d}U_{i,i+1}))\) in the
frame \(Y\), where \(c\ne0\). Give \(D\) the measure
\[
       dd_c=(|c|/(2\pi))^a\prod_{i<j\ {\rm in\ one\ block}}d^2U_{ij},
\]
where \(d^2u\) is Euclidean area in \(\CC\).
These are the upper entries of the first frame action on \(D\)
in Lemma~\ref{lem:stab205}. Translation is triangular in order
of increasing superdiagonal, with all diagonal Jacobian entries
one. The formula therefore defines Haar measure. Passing back
from the monomial frame by an upper unitriangular matrix conjugates
each principal block by its own upper unitriangular principal
block. This follows by block upper triangular multiplication and
preserves the displayed area. Thus one may also read the within-block
entries in the original frame, as in
\srcloc{Formalism/SharpRestriction/CuspSpanNormBounds/RootMeasBig}{48}.
We denote the two choices of quotient kernel by \(K^{\rm root}\)
for \(dd_c\) and \(K^{\rm area}\) for unweighted area.

\begin{theorem}[Transport of the normalized orbital kernel]
\label{thm:scale207}
For \(t,\kappa>0\) replace the terminal vectors by \(t z_i\).
Let \(A\) multiply \(u_i\) by \(\lambda_i=t\kappa^{d-i}\)
for \(i<d\), multiply its isotropic dual by \(\lambda_i^{-1}\),
and fix the middle space. Put \(\gamma'=A^{-1}\gamma A\).
Write \(j_B=|\det(\operatorname{Ad}_A|\mathfrak b)|\).
The new coefficient of the character is \(c'=\kappa c\).
For the transported sections \(W'_i(g)=j_B^{-1/2}W_i(Ag)\),
with the two \(B\)-Haar measures the original ones, one has
\begin{equation}
  (P')^{1/2}K_{\gamma'}^{\rm root}(W'_1,W'_2)
     =t^{\,2(k-d)}\kappa^{-S-a}
                         P^{1/2}K_{\gamma}^{\rm root}(W_1,W_2).
                                                        \label{eq:scale207}
\end{equation}
The transport preserves the \(L^2(B\backslash G)\) norm of every
section. The formula applies to every block decomposition.
\end{theorem}
\begin{proof}
Set \(S_0=\operatorname{diag}(t\kappa^{d-i})_{i=1}^d\).
If \(Y_i'\) is the new standard frame, then \(A Y'_i=Y_i S_0\).
Thus the new cross matrix is \(S_0^* Z S_0\). The action of
\(A^{-1}bA\) in the new frame is \(S_0^{-1} U S_0\) when
\(b\) acts by \(U\) in the old frame. Each simple coordinate
is divided by \(\kappa\), giving the asserted \(c'\). The
column group action is preserved. These facts also apply after
unitriangular changes to a monomial frame, by conjugating those
changes by \(S_0\). In particular
\[
 D'_j=t^{2j}\kappa^{S_j}D_j,\qquad
 \zeta'_\nu=t^2\kappa^{e_\nu}\zeta_\nu,\qquad
 \mathcal G'=t^{4d}\kappa^{2S}\mathcal G .
\]
For the last equality \(S_{d-1}=S_d\) and the product of
the new terminal norms is \(t^4T_1T_2\).

Let \(E_*=\sum_{i=1}^d 4(d-i)^2\). The sum of squared deviations
of the numbers \(2(d-i)\) from their mean on a block of length
\(h\) is \(h(h^2-1)/3\), by the sums of the first \(h\) integers
and their squares. Hence
\(\sum h_\nu e_\nu^2=E_*-2H\) and \(\sum h_\nu e_\nu=S\).
Substitution in the product defining \(P\) gives
\begin{equation}
       P'/P=t^{4d(k+d-2)}
               \kappa^{2((k-1)S+E_*-2H)}.               \label{eq:pscale207}
\end{equation}
A terminal singleton obeys the same scaling since \(e_l=0\)
and both terms of its difference acquire the factor \(t^4\).

To compute \(j_B\), take the three independent sets of coordinates
of the unipotent Lie algebra in the Witt basis. In its isotropic
upper triangular block \(i<j<d\), the complex entry has factor
\(\lambda_i/\lambda_j\). Each of the \(k+1\) complex middle entries
in row \(i<d\) has factor \(\lambda_i\). In the skew-Hermitian
block from the dual isotropic space to the isotropic space, a real
diagonal coordinate has factor \(\lambda_i^2\) and the complex
off-diagonal coordinate \(i<j\) has factor \(\lambda_i\lambda_j\).
These are all free coordinates by the equation \(X^*J+JX=0\)
for the Lie algebra. The compact column algebra is fixed.
Taking real determinants yields
\[
 j_B=\prod_{i<d}\lambda_i^{2(k+2(d-i))}
      =t^{\,2(d-1)(k+d)}\kappa^{kS+E_*}.
\]
In the disintegration of \(dg\) into \(db\,d\dot g\), left
multiplication by \(A\) preserves \(dg\) and carries each left
fibre by \(b\mapsto AbA^{-1}\). Consequently
\(\int_{B\backslash G} f(Ag)d\dot g=j_B\int f(g)d\dot g\).
This proves the isometry assertion and gives the stated
normalization for the transported sections.

The map \(d_0\mapsto A^{-1}d_0 A\) carries \(D\) onto \(D'\).
It multiplies the \((i,j)\) coordinate within a block by
\(\kappa^{i-j}\). The sum of \(j-i\) over these pairs is
\(H\), since on a block it is
\(\sum_{m=1}^{h-1}m(h-m)=h(h^2-1)/6\).
The pullback of \(dd'_{c'}\) is therefore
\(\kappa^{a-2H}dd_c\). Disintegrating group Haar measure on
both sides of \(g\mapsto Ag\) shows that the pushforward
of the new quotient measure is \(\kappa^{2H-a}\) times the old
one. Including both section scalars gives
\(K_{\gamma'}^{\rm root}(W'_1,W'_2)
 =j_B^{-1}\kappa^{2H-a}K_\gamma^{\rm root}(W_1,W_2)\).
The integrals exist by Theorem~\ref{thm:quot207}, or by the
change of variables from that theorem. Multiplying by the square
root of \eqref{eq:pscale207} gives \eqref{eq:scale207}.
\end{proof}

\begin{corollary}[The regular squared norm]\label{cor:regular207}
Let \(\Omega=\{q\in\CC^d:q_1\cdots q_{d-1}\ne0,\
|q_d|^2>T_1T_2\}\). On this set let \(\gamma(q)\) have diagonal
cross matrix with entries \(q_1,\ldots,q_d\), and put
\[
  \mathfrak J(U)=\int_U P(q)
      \|K_{\gamma(q)}(W_1,W_2)\|_{\rm HS}^2\,d^{2d}q
               \qquad(U\subset\Omega\ \text{Borel}).
\]
For \(q'_i=t^2\kappa^{2(d-i)}q_i\) and the image \(U'\) of \(U\),
\begin{equation}
                   \mathfrak J'(U')=t^{4k}\mathfrak J(U).
                                                       \label{eq:regular207}
\end{equation}
This holds on every truncation and on all of \(\Omega\) as an
identity of nonnegative extended integrals. Whenever the full
norm is positive and finite, the pointwise quantity
\(P\|K_\gamma^{\rm root}\|_{\rm HS}^2/\mathfrak J(\Omega)\)
is multiplied by \(t^{-4d}\kappa^{-2(S+a)}\) at every realized
nondegenerate orbit.
\end{corollary}
\begin{proof}
Diagonal cross data are realized by taking the first \(d-1\)
vectors of the second frame to be \(q_iv_i\), where \(v_i\)
are the dual isotropic vectors. For its last vector take
\(-q_d z_1/T_1+\sqrt{|q_d|^2/T_1-T_2}\,w_0\), with \(w_0\)
a positive unit middle vector perpendicular to \(z_1\). Its norm
is \(-T_2\) and its pairing with \(z_1\) is \(q_d\).
These choices lift locally smoothly to isometries taking the
second standard frame to the displayed frame. On the hyperbolic
planes interchange the isotropic basis and its dual with the
reciprocal conjugate scalings. In the middle plane extend the
negative unit vector to a signed orthonormal basis by projecting
a positive vector onto its perpendicular and normalizing, and
extend by the identity on the positive complement. A fixed change
of the original middle basis can first put \(z_2\) on its
negative axis. This also gives measurable representatives.
Any two representatives with these diagonal data are in the
same \(B\)-double coset by Lemma~\ref{lem:forms203} over \(\CC/\RR\).
For singleton blocks \(D\) is trivial. Changing the representative
by \(b_1\gamma b_2^{-1}\) multiplies the kernel on the two sides
by \(\tau_1(b_1)\) and \(\tau_2(b_2)^*\), so its Hilbert--Schmidt
norm is well defined. The integrands are measurable by integration
of their measurable components over \(G\) and absolute convergence
on every such orbit.

On \(\Omega\), equation \eqref{eq:scale207} squared gives the
factor \(t^{4(k-d)}\kappa^{-2S}\). The real Jacobian of the map
\(q\mapsto q'\) is \(t^{4d}\kappa^{2S}\).
The change of variables for a nonnegative integral gives
\eqref{eq:regular207}. Dividing the squared identity on an arbitrary
realized orbit by the full regular norm, under the stated finiteness
and positivity conditions, gives the last assertion. The regular
density used here is the one in
\srcloc{Formalism/FormalDegreeBessel/StatementBF}{146}.
\end{proof}

\begin{theorem}[The discriminant in the arithmetic quotient]
\label{thm:rd207}
Return to \(E/F\) and \(T_1=T_2=T\). Set
\(R=|\Delta_E|^{1/(2f)}\). Choose any \(A_0>0\) and apply
Theorem~\ref{thm:scale207} at all real places with \(t=1\)
and \(\kappa=A_0^{-1}\). Indicate transported data by primes.
For the product quotient and the actual sections one has
\begin{equation}
        c_E^a P_\infty^{1/2}K_{\infty,\gamma}^{\rm area}
       = A_0^{-fS}(R/A_0)^{fa}
                 (P'_\infty)^{1/2}K_{\infty,\gamma'}^{\rm root}.
                                                        \label{eq:rd207}
\end{equation}
In particular \(A_0=R\) removes all dependence on the block
stabilizer dimension from the scalar on the right. The finite
cell contributions at the original level then satisfy
\begin{align}
 \sum_c |J_{\gamma,c}|
 &\le B_\ell B_{\ell'}\,m_C
       \Bigl(\prod_{v\in S(\gamma)}L_{\gamma,v}\Bigr)
          R^{-fS}Q^{-fd(k+d-2)}
          \big\|(P'_\infty)^{1/2}
                       K_{\infty,\gamma'}^{\rm root}\big\|_{\rm op}.
                                                        \label{eq:levelscale207}
\end{align}
\end{theorem}
\begin{proof}
In Section~\ref{sec:char94} the real character has \(|c|=4\pi\),
since the trace from \(\CC\) to \(\RR\) is twice the real part.
Thus the product of the root Haar measures is \(2^{fa}\) times
area and \(K_\infty^{\rm area}=2^{fa}K_\infty^{\rm root}\).
The discriminant covolume in \eqref{eq:mass205} is
\(c_E=2^{-f}|\Delta_E|^{1/2}=(R/2)^f\). Multiply the inverses
of \eqref{eq:scale207} over the \(f\) places. The two area and
lattice factors combine to \(2^{fa}c_E^a=R^{fa}\), giving
\eqref{eq:rd207}. The argument works first on tensors by Fubini,
and then on the actual compact-finite vectors by their finite
tensor expansions.

At a fixed finite cell the two covectors, including their finite
unitary factors, are constant in the real variable. The real
integrand between them descends to \(D_\infty\backslash G_\infty\)
before contraction by the relevance identity. Its integral over
the rational unipotent quotient is the covolume in
\eqref{eq:mass205} times that contraction of
\(K_{\infty,\gamma}^{\rm area}\). Bound the latter in modulus
by its operator norm times the two covector norms and sum the
covolumes by Theorem~\ref{thm:mass205}. Now use \eqref{eq:rd207}
at \(A_0=R\) and Lemma~\ref{lem:density207}.
\end{proof}

The two scales in the source can be compared explicitly.
Its real block multiplier is \(R^{-e_\nu}\), whereas its parameter
for the derivative of the additive character is \(2\pi/(RQ)\).
For the physical trace character the latter magnitude, after the
conversion \(s\mapsto-2is\), corresponds to \(A_0=RQ\).
In \eqref{eq:rd207} this choice gives the scalar
\((RQ)^{-fS}Q^{-fa}\) and block multipliers \((RQ)^{-e_\nu}\).
At \(A_0=R\) the block multipliers are \(R^{-e_\nu}\) and the
absolute coefficient is \(4\pi/R\). Thus a use of the source's
pointwise estimate at its joint data still requires a comparison
at its specified character and block coordinates.
See \srcloc{Formalism/SharpRestriction/CuspSpanNormBounds/LocalDataRealPlace}{16},
\srcloc{Formalism/SharpRestriction/CuspSpanNormBounds/LeadMinorK}{168}
and \srcloc{Development/AbelianVarietyTheory/Definitions/GlobalTwisting/CcKPrimitive}{109}.

\subsection{The contribution of the two auxiliary levels}\label{nav:IX-05}
\label{sec:refine211}

The bound in Corollary~\ref{cor:match209} uses the largest norm of a
column covector over all its finite classes. We now sum its squared
norm over the refinements of a single class. This evaluates the
effect of both auxiliary levels on the nondegenerate part of the
original Gram matrix.

Retain \(K=K^0\), the terminal vector \(\varpi e\), the two split
primes \(g_1,g_2\), and \(Q=g_1g_2\) from
Section~\ref{sec:match209}. Let \(K^*\) enlarge \(K\) to the full
lattice stabilizer at \(v\mid Q\), with \(K^*_v=K_v\) elsewhere.
Write \(K_H^*=H_f\cap K^*\) and \(K_H=H_f\cap K\) for the
levels in the definite column group. Set
\[
 I_\sigma=[K^*:K],\qquad I_\pi=[K_H^*:K_H],\qquad
 c_\rho^*=
 \frac{\operatorname{vol}(K^*)v_N\operatorname{vol}(H_\infty)}
      {\operatorname{vol}(N_f\cap K^*)\,\dim\rho}.
\]
Here all Haar measures are those of \eqref{eq:main127}. The symbol
\(c_\rho^*\) denotes this scalar even when a column class at the
larger level has a nontrivial rational stabilizer. For a fixed
relevant label of nonzero Gram determinant put
\(J=B\cap\gamma B\gamma^{-1}\) and \(D=N\cap\gamma N\gamma^{-1}\).
For either local level \(K'\) let
\[
 L_v(K')=\sum_{tK'\in J(F_v)\backslash
            ((B_vK'\cap\gamma B_vK')/K')}
      \frac{\operatorname{vol}(K')}
           {\operatorname{vol}_{J_v}(J(F_v)\cap tK't^{-1})}.
\]
Use the complement measures of Theorem~\ref{thm:complement209}.

\begin{lemma}[The fraction of refinements supporting two flags]
\label{lem:flags211}
At \(v\mid g\) for \(g\in\{g_1,g_2\}\), a full-level coset
supporting both flags has either zero or
\(i_{\pi,v}=|\mathrm{GL}_k(\mathbb F_g)/U_k(\mathbb F_g)|\)
right \(K_v\)-cosets supporting both. Write \(L_v^{=}\) for the
part of \(L_v(K_v^*)\) on which these refinements exist. Then
\begin{equation}
 L_v(K_v)=\frac{i_{\pi,v}}{i_{\sigma,v}}L_v^{=},\qquad
 \frac{i_{\sigma,v}}{i_{\pi,v}}=
            g^k\prod_{j=k+1}^r(g^j-1),\quad
 i_{\sigma,v}=[K_v^*:K_v].
 \label{eq:fraction211}
\end{equation}
In particular \(0\le L_v^{=}\le L_v(K_v^*)\).
\end{lemma}
\begin{proof}
Use the first split factor and the Witt basis in
Theorem~\ref{thm:gap132}. Put \(s=d-1\). The blocks have sizes
\(1,\ldots,1,k+1,1,\ldots,1\), with \(s\) singletons on either
side. Let \(\bar N\) be the upper block unipotent radical over
\(\mathbb F_g\). On the middle block embed \(\mathrm{GL}_k\)
by \(h\mapsto\operatorname{diag}(h,1)\). The reductions of
\(B_v\cap K_v^*\) and \(K_v\) are respectively
\[
 \bar B=\bar N\rtimes\mathrm{GL}_k(\mathbb F_g),
 \qquad \bar K=\bar N\rtimes U_k(\mathbb F_g).
\]
Indeed the singleton actions are one and the middle action of
the period group fixes the terminal vector and its dual on the
middle quotient. Its arbitrary invertible \(k\)-block and all
the unipotent blocks lift entry by entry over the integer ring.
The row level has exactly the additional upper unitriangular
condition on that \(k\)-block. This follows as well from the
zero entries in the terminal column of the specified row reduction.
The principal congruence subgroup is contained in \(K_v\).
It follows that \((B_v\cap K_v^*)K_v\) is the entire inverse
image of the subgroup \(\bar B\).

Choose \(b_1,b_2\in B_v\) and \(h_1,h_2\in K_v^*\) with
\(t=b_1h_1\) and \(\gamma^{-1}t=b_2h_2\). A refinement represented
by \(tu\), \(u\in K_v^*\), supports the first flag precisely when
\(\bar u\in\bar h_1^{-1}\bar B\), and the second precisely when
\(\bar u\in\bar h_2^{-1}\bar B\). These two left translates are
equal or disjoint. When equal, their right \(\bar K\)-cosets are
a translate of \(\bar B/\bar K\), of cardinal \(i_{\pi,v}\).
The equality test is independent of the decompositions just chosen.

In the full-level orbit of \(tK_v^*\), its stabilizer
\(J_t=J(F_v)\cap tK_v^*t^{-1}\) permutes the supporting refinements.
For a \(J_t\)-orbit of length \(u_0\), its stabilizer has volume
\(\operatorname{vol}(J_t)/u_0\). Its contribution at the finer
level is therefore \(u_0\operatorname{vol}(K_v)/\operatorname{vol}(J_t)\).
Summing proves the first equality. The second follows from
\(|\mathrm{GL}_n(\mathbb F_g)|=g^{n(n-1)/2}\prod_{j=1}^n(g^j-1)\)
and \(|\bar N|=g^{(r(r-1)-k(k+1))/2}\). The local sums at
both levels are finite by the compact-orbit argument in
Proposition~\ref{prop:good205}.
\end{proof}

At the other primes set \(L_v^{=}=L_v(K_v)\). Define
\(\mathcal M_\gamma^{=}=\mathfrak m_{k-1}\prod_v L_v^{=}\) and
\(\mathcal M_\gamma^*=\mathfrak m_{k-1}\prod_v L_v(K^*_v)\).
These are finite products after the factors one are removed.
Lemma~\ref{lem:flags211} and Theorem~\ref{thm:complement209}
give the exact equality \(M_\gamma=(I_\pi/I_\sigma)\mathcal
M_\gamma^{=}\). At the two split primes we can recover the other
factor \(I_\pi\) in the estimate for unit columns as follows.

\begin{proposition}[The column norms above one rational cell]
\label{prop:fibres211}
Consider a rational cell at level \(K^*\) admitting a supported
refinement at level \(K\). Fix a representative \(t\) of such a
refinement. Put
\[
 \Gamma^*=J(F)\cap tK^*t^{-1},\quad
 \Gamma_D=D(F)\cap tK^*t^{-1},\quad a_t=[\Gamma^*:\Gamma_D],
 \qquad
 m_t^*=\frac{\operatorname{vol}(K^*)}{a_t}
               \operatorname{vol}(\Gamma_D\backslash D_\infty).
\]
The two flags give column classes \(y_1,y_2\) at level \(K_H^*\).
Choose finite adelic representatives for these classes and set
\[
 A_i=H(F)\cap y_i K_H^* y_i^{-1},\qquad
 U_i(y_i)^2=\sum_{y\mapsto y_i}\|\ell_{i,y}\|^2\quad(i=1,2).
\]
The sum runs over the column classes of the original level
\(K_H\) mapping to \(y_i\), and \(\ell_i\) is an arbitrary
tuple of covectors for the \(i\)-th word. For the original signed
period summands above this cell,
\begin{equation}
 \sum_{c\mapsto t}|J_{\gamma,c}|
 \le \|K_{\infty,\gamma}^{\rm area}\|_{\rm op}\,
       \frac{m_t^*}{I_\sigma}
       \sqrt{|A_1||A_2|}\,U_1(y_1)U_2(y_2).
 \label{eq:fibres211}
\end{equation}
Moreover the sum over all coarse rational cells admitting such
a refinement has \(\sum_t m_t^*=c_E^a\mathcal M_\gamma^{=}\).
\end{proposition}
\begin{proof}
Use \(t=b_1h_1,\ \gamma^{-1}t=b_2h_2\) with \(b_i\in B_f\)
and \(h_i\in K\). We may take \(y_i=p_H(b_i)\) as the two
representatives. Before dividing by \(\Gamma^*\), the common
refinements are parametrized by
\[
             \mathcal C=\prod_{v\mid Q}
                 \mathrm{GL}_k(\mathbb F_g)/U_k(\mathbb F_g).
\]
This follows from the proof of Lemma~\ref{lem:flags211}, since
both \(\bar h_i\) lie in \(\bar K\). For a parameter \(z\in\mathcal C\)
lift it by an integral block \(h(z)\) in \(H\) at the marked
places, extended by one elsewhere. The lift specifies the finite
row coset \(th(z)K\). Its column class on side \(i\) is represented
by \(y_i\) times a lift of \(u_i z\), where \(u_i\) is the fixed
tuple of \(k\)-blocks of the reductions of \(h_i\). Indeed the
reduction of \(h_i h(z)\) has that middle block. Its unipotent
part lifts to \(N_v\cap K_v^*\), and all errors reducing to the
identity lie in \(K_v\). Thus each of the two maps to
\(K_H^*/K_H\) runs through all \(I_\pi\) cosets once.

On passing to rational column classes the fibres are the
\(A_i\)-orbits in these cosets, with the action conjugated by
\(y_i\). The groups \(A_i\) are finite because they preserve a
positive definite integral trace lattice. Each action is free,
since at the original level the column stabilizers are trivial
by Proposition~\ref{prop:column110} and
Theorem~\ref{thm:midlevel131}. Consequently if \(y_i(z)\) is
the fine class on side \(i\), then
\[
              \sum_{z\in\mathcal C}\|\ell_{i,y_i(z)}\|^2
                         =|A_i|\,U_i(y_i)^2 .
\]
The norms are independent of representatives because the
archimedean column representation is unitary.

The projection of \(\Gamma^*\) into the first rational column
group has kernel \(\Gamma_D\), by Lemma~\ref{lem:stab205}, and
its image lies in \(A_1\). To see the latter inclusion conjugate
by \(b_1\) at the finite places and take the column component.
This proves that \(a_t\) is finite. An element of \(\Gamma_D\)
acts trivially on \(\mathcal C\). In fact its conjugate by
\(b_1^{-1}\) belongs to \(N_f\), and at the marked places
it is integral because \(t=b_1h_1\) with \(h_1\in K\).
It reduces into \(\bar N\), which is normal in \(\bar B\)
and fixes all of \(\bar B/\bar K\). Conjugating also by \(h_1\)
does not change this property. At the other places the level
has a single refinement. Conversely if an element of \(\Gamma^*\)
fixes a fine coset, its rational column image fixes a fine
column class in the sense of belonging to its compact stabilizer.
That rational stabilizer is trivial. Thus
\(\Gamma^*/\Gamma_D\) acts freely on \(\mathcal C\) and every
fine row stabilizer is \(\Gamma_D\).

For each fine rational cell, the orbital integral contracts
\(K_{\infty,\gamma}^{\rm area}\) with its two column covectors
and a scalar of modulus one. Its measure is
\(\operatorname{vol}(K)\operatorname{vol}(\Gamma_D\backslash D_\infty)\)
by the disintegration giving \eqref{eq:amp205}. The product of
the two covector norms is constant on each rational orbit of
refinements. We may sum it over all \(z\in\mathcal C\) and
divide by \(a_t\). Cauchy--Schwarz and the two marginal identities
give \eqref{eq:fibres211}, since
\(\operatorname{vol}(K)=\operatorname{vol}(K^*)/I_\sigma\).

For the assertion about \(m_t^*\), count just the measures of
these same cells. There are \(I_\pi/a_t\) above \(t\). Their
total is \(m_t^* I_\pi/I_\sigma\). Summing and using
\eqref{eq:mass205} at the original level gives
\((I_\pi/I_\sigma)\sum_t m_t^*=c_E^a M_\gamma\).
Substitution of the formula following Lemma~\ref{lem:flags211}
proves the asserted identity, retaining possible torsion at
the larger level.
\end{proof}

\begin{theorem}[Auxiliary level indices in the normalized period]
\label{thm:norm211}
Let \(A_*\) be the largest order of a rational column stabilizer
at level \(K_H^*\). Suppose the two column tuples have unit norm
for the original identity terms in \eqref{eq:main127} and the real
row inputs are normalized as in Section~\ref{sec:match209}. For
each relevant supported rational label of nonzero Gram determinant,
\begin{align}
 \sum_c|J_{\gamma,c}|
 &\le \frac{A_*\,\mathcal M_\gamma^{=}}
            {\sqrt{c_{\rho_1}^*c_{\rho_2}^*}}\,
    R^{-fS}Q^{-f(dk+a)}
       (P^\#_\infty)^{1/2}\|\widehat K_{\gamma^\#}\|_{\rm op},
                                                      \label{eq:norm211}\\
 \log A_*&=O\bigl(fk\log^2(2fk)\bigr)=o(fk\log b).
 \notag
\end{align}
This holds on the common cohomology intervals, with \(S=d(d-1)\)
and \(a=\sum_\nu\binom{h_\nu}{2}\). One may use
\(\mathcal M_\gamma^*\) in place of \(\mathcal M_\gamma^{=}\) on
the right.
\end{theorem}
\begin{proof}
At \(v\mid Q\) the integral unipotent blocks all reduce into
\(\bar K\). Hence \(N_f\cap K^*=N_f\cap K\). Formula
\eqref{eq:main127} gives \(c_{\rho,K}=c_\rho^*/I_\sigma\).
A unit tuple therefore has
\(\sum_y\|\ell_{i,y}\|^2=I_\sigma/c_{\rho_i}^*\). This bounds
each partial squared norm \(U_i(y_i)^2\). The two square roots
cancel the entire \(I_\sigma^{-1}\) in \eqref{eq:fibres211}.
Summing over \(t\) now bounds the period by
\[
   \frac{A_*\,\mathcal M_\gamma^{=}}
        {\sqrt{c_{\rho_1}^*c_{\rho_2}^*}}\,
                           c_E^a\|K_{\infty,\gamma}^{\rm area}\|_{\rm op}.
\]
Use \eqref{eq:match209} and \(P_\infty^{1/2}\ge Q^{fd(k+d-2)}\)
from Theorem~\ref{thm:rd207}. The resulting power of \(Q\) is
\(f(d(d-2)-a)-fd(k+d-2)=-f(dk+a)\), giving \eqref{eq:norm211}.
The inequality \(\mathcal M_\gamma^{=}\le\mathcal M_\gamma^*\)
follows locally from Lemma~\ref{lem:flags211}.

Every coarse column stabilizer acts faithfully on the underlying
rational space of dimension \(n=2fk\). Minkowski's theorem for
finite subgroups of \(\mathrm{GL}_n(\mathbb Q)\) gives
\[
 v_p(|A_i|)\le \sum_{j\ge0}
                \left\lfloor\frac{n}{p^j(p-1)}\right\rfloor
 \qquad\text{for every prime }p
\]
\cite[Lecture I, Theorem 1]{XIX-XIV-IV-serre211}. Only \(p\le n+1\) occur
and the sum is at most \(np/(p-1)^2\). Summing its product with
\(\log p\) and bounding the sum by that over all integers
\(2\le j\le n+1\) gives \(O(n\log^2(2n))\), by the integral
estimate for \(\sum_{j\le n+1}(\log j)/j\). On the selected
intervals \(f,k\) are bounded by fixed powers of \(x\) and
\(\log b=x^\epsilon\). This proves the claimed uniform estimate.
\end{proof}

For example, in the Cauchy--Schwarz estimate following
Proposition~\ref{prop:hyp209} use \eqref{eq:norm211} before
summing over the rational labels. For any finite set
\(\mathcal A\) of such labels the resulting upper bound is
\begin{align*}
 &\frac{A_*R^{-fS}Q^{-fdk+2f}}
           {\sqrt{c_{\rho_1}^*c_{\rho_2}^*}}
 \left(\sum_{\gamma\in\mathcal A}
      \frac{(\mathcal M_\gamma^{=})^2}{w_\gamma Q^{2fa}}\right)^{1/2}\\
 &\hspace{24mm}\cdot
 \left(\sum_{\gamma\in\mathcal A}
 w(\beta_\gamma)P_\infty^\#
             \|\widehat K_{\gamma^\#}\|_{\rm HS}^2\right)^{1/2}.
\end{align*}
Indeed the comparison of the two weights in \eqref{eq:label209}
and the integral injection there apply to the same original
labels. Both sums are nonnegative so finite truncations may be
increased. The remaining finite local integrals include those
for the two full lattice stabilizers subject to the equality
test in Lemma~\ref{lem:flags211}. Their dependence on the
nonunit determinants, the archimedean sum, and the lower estimate
for the specified integral pairing still enter the period comparison.

\subsection{Local factors at divisors of the Gram determinant}\label{nav:IX-06}
\label{sec:localmass213}

The part \(\mathcal M_\gamma^{=}\) in Theorem~\ref{thm:norm211}
contains local factors at the divisors of \(\mathcal G\). We
evaluate them at full lattice level when the preceding nonzero
endpoint minors are units and the last block is a singleton.
This includes every place outside the fixed set \(S_0\) where
the sole exceptional condition of Section~\ref{sec:cellmass205}
is divisibility of the Gram determinant.

Fix a finite place \(v\notin S_0\). Write \(F_0=F_v\), \(E_0=E_v\),
\(O=\mathcal O_{F_0}\), \(\mathcal O=\mathcal O_{E_0}\),
\(\pi\) for a uniformizer of \(F_0\), and \(q=|O/\pi O|\).
The residue characteristic is odd, \(T\) is a unit, and \(E_0\)
is either the split algebra or an unramified quadratic field.
Give the full row lattice stabilizer volume one. Use the Haar
measure of Theorem~\ref{thm:complement209} on \(H_C=U(C)\)
and give \(\mathcal O\) volume one in each common unipotent
entry. Unit conditions in the split algebra refer to both factors.

\begin{lemma}[Reduction to one integral vector]
\label{lem:vector213}
Let the last block have length one. Suppose every nonzero proper
endpoint minor \(D_j\), \(j<d\), is a unit at \(v\). Then the local
mass \(L_v(K^*_v)\) of Section~\ref{sec:refine211} is the sum,
over \(H_C(F_0)\)-orbits of self-dual lattices \(\Lambda\) in a
rank \(k\) unramified Hermitian space \(V\) which contain a fixed
vector \(w\), of
\(\operatorname{vol}_{H_C}(\operatorname{Stab}(\Lambda))^{-1}\).
Here \(\operatorname{Stab}(\Lambda)\) denotes its stabilizer in
\(H_C(F_0)\). The vector has nonzero norm
\[
                 a_0=\frac{\mathcal G}
                              {T D_{d-1}\overline{D_{d-1}}}.
\]
In particular \(e=v(a_0)=v(\mathcal G)\ge0\).
\end{lemma}
\begin{proof}
For a local coset supporting both flags, take representatives of
both decorated frames in its row lattice. The changes to monomial
cross matrix of Lemma~\ref{lem:forms203} are integral here.
Every earlier block pivot is a unit since a power of it, up to
a sign, is a ratio of unit endpoint minors. Elimination of the
off-diagonal entries in the last row and column uses only those
earlier pivots. The first \(d-1\) vectors of the two frames together
now span a unimodular lattice in the sum of their hyperbolic planes.
Its Gram matrix has zero diagonal blocks and invertible off-diagonal
blocks. The inverse matrix gives an integral orthogonal projection
onto this lattice, so it splits off the whole row lattice.

Fix one rationally normalized pair of local frames. All pairs
in the same decorations with the same monomial matrix differ by
the block unitriangular action of \(D(F_0)\), as in
Lemma~\ref{lem:stab205}. Thus by this action the preceding
hyperbolic lattice can be taken to be the lattice in the fixed
normalized basis. The two remaining vectors \(z,z'\) are in its
orthogonal complement. They have norm \(-T\) and cross pairing
\(q_d=D_d/D_{d-1}\), since the last block is a singleton.
The second vector has integral projection on the first because
\(q_d\in\mathcal O\) and \(T\) is a unit. Split off \(\mathcal O z\)
and let \(w\) be the perpendicular projection of \(z'\) onto the
remaining space \(V\). Its norm is \(-T+q_d\bar q_d/T=a_0\), and
\(C=w^\perp\) inside \(V\).

Conversely every self-dual lattice in \(V\) containing \(w\)
gives an allowed row lattice by adjoining these fixed summands.
Indeed the lattice in the remaining middle space contains both
\(z,z'\), each of unit norm, so each individually splits off
an orthogonal unimodular line. For either flag, use its first
\(d-1\) vectors with dual partners from the other flag obtained
by inverting the unit cross matrix. The resulting hyperbolic
bases, the chosen terminal vector, and a basis of its unimodular
perpendicular give an integral isometry of the corresponding
original marked lattice to the proposed lattice. In the unramified
field case unimodular lattices of a given rank are integrally
isometric. Over the residue field a nonsingular Hermitian form has
a vector of nonzero norm, which lifts to split off a unit line.
Induction and surjectivity of the norm on units give a basis of
unit diagonal ones. The norm on units is surjective by the
finite-field norm and Hensel lifting of principal units at odd
residue characteristic. For the split algebra use a basis in one
factor and the dual basis in the other. These isometries preserve
the individual ordered decorations, which is precisely the two
support conditions in \(X_v\).

The subgroup \(D(F_0)\) acts on the preceding hyperbolic summand
and fixes the entire middle space because its final block has
length one. The stabilizer of the chosen hyperbolic lattice in
\(D(F_0)\) consists of the integral block unitriangular matrices.
Indeed the induced matrices on the other frame are their adjoint inverses
conjugated by monomial blocks of units and are also integral.
This subgroup has volume one. The other direct factor of the
common stabilizer acts on \(V\) by \(H_C(F_0)\), fixing \(w\).
This gives the asserted orbit masses, including for nonsingleton
blocks earlier in the frame.
\end{proof}

\begin{proposition}[The split local factor]\leavevmode\label{prop:split213}
Under the assumptions of Lemma~\ref{lem:vector213}, if
\(E_0=F_0\times F_0\) then
\[
                      L_v(K_v^*)=\sum_{i=0}^{e}q^{(k-1)i}.
\]
\end{proposition}
\begin{proof}
A self-dual lattice in \(V\) consists of an \(O\)-lattice \(L\)
in \(F_0^k\) and its dual. Write \(w=(u,\eta)\), with \(u\)
a vector and \(\eta\) a covector, so that \(\eta(u)=a_0\).
The two integrality conditions are \(u\in L\) and
\(\eta(L)\subset O\). Put \(C_+=\ker\eta\). The stabilizer of
the pair \(u,\eta\) acts as the full \(\mathrm{GL}(C_+)\).
It moves \(L\cap C_+\) to a fixed lattice \(L_0\) in \(C_+\).
Projection of \(L\) onto \(F_0 u\) along \(C_+\) has image
\(\pi^{-i}Ou\) for some \(0\le i\le e\), by the two conditions.
Over this image \(L\) is generated by \(L_0\) and a vector
\(\pi^{-i}u+y\). Exactly the classes
\(y\in\pi^{-i}L_0/L_0\) work, since \(\pi^i y\in L_0\)
is equivalent to \(u\in L\). The compact group \(\mathrm{GL}(L_0)\)
has volume one and acts on this finite set. Each stabilizer has
inverse volume equal to the length of its orbit. Summing therefore
gives \(q^{(k-1)i}\) at \(i\), as claimed.
\end{proof}

\begin{lemma}[The nonsplit masses at valuations zero and one]
\leavevmode\label{lem:base213}
Let \(E_0/F_0\) be the unramified quadratic field and let \(V\)
admit a self-dual lattice. For a vector of unit norm the mass
in Lemma~\ref{lem:vector213} is one. For a vector \(u\) with
\(v(h(u,u))=1\) it is \(q^{k-1}+1\).
\end{lemma}
\begin{proof}
At unit norm the line splits from any containing self-dual lattice,
and the perpendicular is self-dual. All such lattices form one
orbit with stabilizer of volume one. For valuation one put
\(a=h(u,u)\). Every containing lattice has \(u\) primitive.
In it choose a vector \(v_0\) with \(h(u,v_0)=1\), by self-duality.
There is a \(b\in O\) such that \(v_0+bu\) is isotropic. Its norm
is \(h(v_0,v_0)+2b+ab^2\), which has a simple root modulo \(\pi\).
The new pairing with \(u\) is a unit. Divide by it to get an
isotropic \(v_1\) with \(h(u,v_1)=1\). The plane on \(u,v_1\)
has unit Gram determinant, so the lattice splits as this plane
and a unimodular lattice \(\Lambda_0\) of rank \(k-2\).
All such marked lattices form one orbit by integral isometry
of the remaining lattices. Such a lattice exists for any norm
of valuation one. A unit diagonal plane contains a hyperbolic
integral basis since \(-1\) is a unit norm and \(2\) is invertible.
In that basis \((1,a/2)\) has norm \(a\). Matching this vector
with the given one by Hermitian Witt extension gives a lattice
in the required marked space.

Put \(y=u-av_1\), of norm \(-a\), and \(L=\mathcal O y\perp\Lambda_0\).
This is the intersection with \(C=u^\perp\). On \(C\) use the
scaled form \(h'=\pi^{-1}h\). It gives the hyperspecial volume
convention of Theorem~\ref{thm:complement209}. In fact the
valuation of \(\det h|_C\) is odd, since the determinant of \(V\)
has even valuation and \(u\) has norm of valuation one. The two
scalings therefore differ by a scalar of even valuation. This
scalar is a norm, and scalar multiplication transports self-dual
lattices while preserving their stabilizers. At unit norm the
same argument applies with \(h'=h\).

For the form \(h'\) one has \(L^\vee=\mathcal O y\perp\pi\Lambda_0\).
Put \(P=\operatorname{Stab}_{U(C)}(L)\). The stabilizer of the
original lattice in \(U(C)\) is the subgroup satisfying
\(g y\equiv y\bmod aL\), since \(v_1=(u-y)/a\) and
the original lattice is generated by \(L,v_1\). Any \(g\in P\)
preserves \(L^\vee\), so \(gy\equiv \xi y\bmod\pi L\) for a scalar
\(\xi\) whose norm is one in the residue field. This follows by
evaluating its unit norm for \(h'\), since the \(\pi\Lambda_0\)
part has norm divisible by \(\pi\). Every such scalar lifts to a
norm-one unit acting on the \(y\)-line. Indeed lift the scalar
arbitrarily, and correct its norm in \(1+\pi O\) by a principal
unit. The index of the stabilizer in \(P\) is \(q+1\).

It remains to determine \(\operatorname{vol}(P)\). Write
\(t=(k-2)/2\), an integer since the column rank is even.
Self-dual lattices \(M\) between \(L^\vee\) and \(L\) correspond
to maximal isotropic \(t\)-planes in the nonsingular
\(2t\)-dimensional residue space \(L/L^\vee\) with form
\(\pi h'\bmod\pi\). The condition \(M^\vee=M\) means that the
plane equals its perpendicular in this space. The group \(P\)
is transitive on the choices, since the integral unitary group
on \(\Lambda_0\) induces its full residue unitary group. For a
fixed \(M\), the conjugate lattices \(L'\) containing \(M\)
in its \(\operatorname{Stab}(M)\)-orbit correspond to the isotropic
\(t\)-planes \(\pi L'/\pi M\) in the \((2t+1)\)-dimensional
nonsingular space \(M/\pi M\) with form \(h'\bmod\pi\).
Indeed \(M\subset L'\subset\pi^{-1}M\) is determined by this
plane, and \(\pi h'(L',L')\subset\mathcal O\) gives isotropy.
These planes all occur by transitivity. In both uses finite-field
transitivity follows by extending an isotropic basis to a Witt
basis. Reduction of the integral unitary group of a unimodular
lattice is onto. If \(Y\) lifts a residue isometry, then
\(Y(Y^*Y)^{-1/2}\) is a unitary lift. The binomial series
converges because \(Y^*Y-1\) vanishes modulo \(\pi\) and the
denominators of its coefficients are powers of two.

For clarity, let \(N_{n,t}\) count isotropic \(t\)-planes in a
nonsingular residue Hermitian space of dimension \(n\). The
number of nonzero isotropic vectors is
\[
            z_n=(q^n-(-1)^n)(q^{n-1}-(-1)^{n-1}).
\]
To count, sum a nontrivial additive character of the residue
field over the norm form in a diagonal unit basis. The sum in
one coordinate is \(-q\), since a nonzero norm has \(q+1\)
preimages. Averaging over all additive characters gives
\(q^{2n-1}+(-1)^n(q-1)q^{n-1}\) vectors of norm zero,
including zero. Counting an isotropic line in an isotropic
\(t\)-plane now gives
\(N_{n,t}=z_n N_{n-2,t-1}/(q^{2t}-1)\) for \(t\ge1\).
Consequently
\[
 \frac{N_{2t,t}}{N_{2t+1,t}}=
 \frac{\prod_{j=1}^t(q^{2j-1}+1)}
      {\prod_{j=1}^t(q^{2j+1}+1)}
 =\frac{q+1}{q^{2t+1}+1}.
\]
The ratio also holds at \(t=0\). The two indices over
\(P\cap\operatorname{Stab}(M)\) calculated above show that
this ratio is \(\operatorname{vol}(P)\), because
\(\operatorname{vol}(\operatorname{Stab}(M))=1\).
Dividing by \(q+1\) gives the volume of the desired stabilizer.
Its inverse is \(q^{2t+1}+1=q^{k-1}+1\).
\end{proof}

\begin{theorem}[The factor at every terminal Gram valuation]
\label{thm:local213}
For both splitting types under Lemma~\ref{lem:vector213},
\begin{equation}
                     L_v(K_v^*)=1+q^{k-1}+\cdots+q^{(k-1)e}.
                                                      \label{eq:local213}
\end{equation}
\end{theorem}
\begin{proof}
Proposition~\ref{prop:split213} proves the split case. In the
field case fix a reference self-dual lattice \(\Lambda\) of \(V\)
and give its stabilizer in \(U(V)\) volume one. On the surface
\(\Sigma_a=\{u:h(u,u)=a\}\) for \(a\ne0\), the group is transitive
by Witt extension. The mass in the lemma for the vector \(u\) is
the quotient measure of \(\{g:u\in g\Lambda\}\) in
\(U(u^\perp)\backslash U(V)\). Indeed each orbit of lattice
cosets has measure equal to the inverse volume of its stabilizer.
The map \(g\mapsto g^{-1}u\) therefore gives this mass as
\(\mu_u(\Sigma_a\cap\Lambda)\) for the corresponding invariant
measure on \(\Sigma_a\).

There is a second invariant measure \(\sigma_a\) on the same
surface. Give \(\Lambda\) additive volume one, and normalize
additive measure on \(F_0\) by \(\operatorname{vol}(O)=1\).
For a compact open set \(D'\subset V\), its surface measure on
\(D'\) is
\[
 \lim_{n\to\infty} q^n\operatorname{vol}
 \{z\in D':h(z,z)-a\in\pi^nO\}.
\]
The derivative in direction \(y\) is \(\operatorname{Tr}h(z,y)\),
surjective onto \(F_0\) at every point of \(\Sigma_a\).
By the local inverse function theorem the norm is one coordinate
on a sufficiently small box around each such point, with constant
absolute Jacobian there. Finitely many disjoint small boxes cover
the compact intersection in \(D'\). The compact remainder avoids
the surface, hence makes zero contribution for large \(n\).
This proves the formula. The measure is positive and invariant
since unitary linear maps preserve additive measure on \(V\).
Under multiplication by \(\pi\) one has
\(\sigma_{\pi^2 a}(\pi D')=q^{-2(k-1)}\sigma_a(D')\), by scaling
the ambient \(2k\)-dimensional measure and the norm coordinate.
In this formula intersection with the relevant surface is understood.
By uniqueness of invariant quotient measure, write
\(\mu_u=c_u\sigma_a\) with \(c_u>0\). The stabilizers and their
Haar measures for \(u,\pi u\) coincide, so \(\mu_{\pi u}\)
is the image of \(\mu_u\) under multiplication by \(\pi\).
It follows that \(c_{\pi u}=q^{2(k-1)}c_u\).

For a unit target and for a target divisible by \(\pi\), the
primitive part of \(\sigma_a(\Lambda)\) equals respectively
\[
 H_1=1-q^{-k},\qquad
 H_0=1+(q-1)q^{-k}-q^{1-2k}
                  =(1-q^{-k})(1+q^{1-k}).
\]
Indeed for each nonzero residual vector in the specified fibre
the trace derivative is onto the residue field. Each lift at
level \(n\) therefore has \(q^{2k-1}\) lifts at level \(n+1\).
The limiting density is its residue count divided by \(q^{2k-1}\).
The zero-norm count is \(z_k\) in Lemma~\ref{lem:base213}. All
nonzero target norms occur equally often by multiplying by units,
so each has \(q^{2k-1}-q^{k-1}\) preimages. Here \(k\) is even.

Set \(A=q^{k-1}\) and \(\delta=e\bmod2\) in \(\{0,1\}\). The
vector \(u_0=\pi^{-(e-\delta)/2}w\) has norm valuation \(\delta\).
At valuations zero and one integral vectors are primitive.
Lemma~\ref{lem:base213} therefore gives \(c_{u_0}=1/H_1\)
if \(\delta=0\) and \(c_{u_0}=(A+1)/H_0\) if \(\delta=1\).
By the scaling rule \(c_w=A^{e-\delta}c_{u_0}\).
For the additive density at \(w\), decompose its integral
vectors according to their exact divisibility by \(\pi^j\).
The range is \(0\le2j\le e\) and the contribution is
\(A^{-2j}H_0\) when \(2j<e\), or \(A^{-2j}H_1\) at \(2j=e\).
For \(e=2h\), multiplication by \(c_w\) gives
\(1+(1+A^{-1})\sum_{j=1}^h A^{2j}\).
For \(e=2h+1\) it gives \((A+1)\sum_{j=0}^h A^{2j}\).
Both expressions equal \(\sum_{i=0}^e A^i\), proving the formula.
\end{proof}

\begin{corollary}[The product at primes dividing only the Gram determinant]
\label{cor:product213}
For any supported relevant label with \(\mathcal G\ne0\), let
\(\mathcal U_\gamma\) be the set of \(v\notin S_0\) where every
nonzero endpoint minor, including the final one, is a unit.
Define the integral ideal
\(\mathfrak b_\gamma=\prod_{v\in\mathcal U_\gamma}
 \mathfrak p_v^{v(\mathcal G)}\). For the marked-period local
factors of Section~\ref{sec:refine211},
\[
 \prod_{v\in\mathcal U_\gamma}L_v^{=}
 =(\mathrm N\mathfrak b_\gamma)^{k-1}
       \sum_{\mathfrak a\mid\mathfrak b_\gamma}
                         (\mathrm N\mathfrak a)^{1-k}.
\]
The divisor sum is between \(1\) and \(1+O(f2^{-k})\),
uniformly in the rational label on the chosen intervals.
\end{corollary}
\begin{proof}
At these unmarked places \(L_v^{=}=L_v(K_v^*)\). When the last
block has length greater than one, \(D_{d-1}=0\). Then
\(\mathcal G=D_d\overline{D_d}\) is a unit on this set, so the
factor is one by Proposition~\ref{prop:good205}. When the last
block is a singleton, use Theorem~\ref{thm:local213}.
Factor \(q^{(k-1)e}\) from each local polynomial and multiply
over \(v\), using unique factorization of ideals. For \(k>2\)
the divisor sum is bounded by \(\zeta_F(k-1)\le\zeta(k-1)^f\)
by the Euler products, since above a rational prime \(p\) there are
at most \(f\) places, each of norm at least \(p\).
Comparison with the integral beyond the first term gives
\(\zeta(k-1)-1=O(2^{-k})\) for large \(k\). Since \(f\) is
polynomial in \(x\) and \(k\asymp x^{9/2+10\epsilon}\),
\(f2^{-k}\to0\). This proves the uniform conclusion.
\end{proof}

In particular the local product records the full valuations of
the Gram ideal. Even a single factor with \(e=1\) contributes
\(1+q^{k-1}\) to the finite mass. The analytic comparison must
account for this factor together with the real density in
\eqref{eq:match209}. Primes of nonunit endpoint minors and the
fixed finite levels can carry further factors.

\subsection{The local mass of an ordered Hermitian frame}\label{nav:IX-07}
\label{sec:frames215}

Let \(K\) be a finite extension of \(\QQ_p\), with integer ring
\(O\), uniformizer \(\pi\), normalized valuation \(v\), and
residue field of order \(q\). Let \(A\) be the unramified
quadratic field over \(K\) or \(K\times K\). Write
\(\mathcal O\) for its integer ring, \(x\mapsto\bar x\) for its
involution, and \(\eta=-1\) or \(1\), respectively. Give the
integer lattices in additive coordinates volume one. Let
\(\Lambda\) be a self-dual Hermitian lattice of rank \(r\), with
form \(H\). We use \(n\) columns, put \(h=r-n\), and assume
\(n\ge1\), \(h\) odd, and \(m=h-n+1\ge2\). In the application
\(n=2d\) and \(h=k-1\). A tuple is \emph{primitive} if its columns
give a split injection \(\mathcal O^n\to\Lambda\), or equivalently
are independent after reduction in each residue field of \(A\).

For a nonsingular Hermitian matrix \(S\), let
\(\Sigma_S=\{X:X^*HX=S\}\). We use it when nonempty.
On this space let \(\mu_S\) be quotient measure for the ambient
unitary group by the pointwise stabilizer of a frame. Give the
stabilizer of \(\Lambda\) volume one. In the perpendicular space
of dimension \(h\), scale the form by the inverse of its
determinant and give a self-dual lattice stabilizer volume one,
as in Theorem~\ref{thm:complement209}. A change of basis multiplies
this scaling by a norm, so it gives the same measure. Define
\[
 P_0=\prod_{j=h+1}^r(1-\eta^j q^{-j}),\qquad
 p(h,u)=\prod_{j=h-u+1}^h(1-\eta^j q^{-j}).
\]
Empty products are one. For a compact open test in the ambient
matrix space, define \(\sigma_S\) on its intersection with
\(\Sigma_S\) by the limit of the additive volumes of
\(X^*HX\equiv S\bmod\pi^j\), multiplied by \(q^{jn^2}\).
Here the target lattice consists of integral Hermitian matrices.
The derivative \(Y\mapsto X^*HY+Y^*HX\) is onto over \(K\)
at every point of \(\Sigma_S\). Indeed \(X^*H\) is onto and
the trace \(A\to K\) is onto. The local inverse function theorem
gives coordinates consisting of the Gram entries and complementary
coordinates. On small compact open coordinate balls the absolute
Jacobian is constant. A finite covering of the test near its
compact intersection with the fibre proves that the limit exists
and gives an invariant measure there.

\begin{lemma}[Primitive frames]\label{lem:primitive215}
If \(S\) is integral and its reduction has radical dimension \(u\),
the \(\sigma_S\)-measure of its primitive frames is \(P_0p(h,u)\).
In the split algebra the two component matrices have the same rank.
\end{lemma}
\begin{proof}
First count over the residue field and divide by \(q^{2rn-n^2}\).
In the split case choose an injection of the first \(n\)-space
into the \(r\)-space. The number is
\(q^{rn}\prod_{j=h+1}^r(1-q^{-j})\). The second map, written
as a map from the \(r\)-space onto an \(n\)-space, is prescribed
on the first image. In the \(u\)-space quotient by its prescribed
image, the remaining \(h\)-space must map onto. Of its
\(q^{hn}\) unrestricted extensions the fraction doing so is
\(\prod_{j=h-u+1}^h(1-q^{-j})\), by counting independent rows.
This proves the formula since \(rn+hn=2rn-n^2\).

In the field case diagonalize the residue form on the source
into \(n-u\) unit lines and a zero block of size \(u\). A
nonsingular residue Hermitian space has an orthonormal basis.
Indeed, vanishing of all norms would imply vanishing of all
pairings by polarization with arbitrary residue scalars and
nondegeneracy of the trace pairing. Split off a vector of
nonzero norm and use surjectivity of the residue norm on units.
In dimension \(j\) the number of vectors of norm one is
\(q^{j-1}(q^j-(-1)^j)\), and the number of nonzero isotropic
vectors is
\(z_j=(q^j-(-1)^j)(q^{j-1}-(-1)^{j-1})\).
Both follow by additive-character orthogonality, since the sum
of a nontrivial additive character of the norm in one coordinate
is \(-q\). Choose the unit vectors in successive perpendiculars.
In the remaining space of dimension \(h+u\), after choosing
\(i\) independent isotropic vectors with pairings zero, the next
one has \(q^{2i}z_{h+u-2i}\) choices. It lifts a nonzero
isotropic vector in the perpendicular quotient by their span.
The unit-vector factors give normalized factors indexed from
\(h+u+1\) through \(r\). The isotropic-vector factors give the
adjacent pairs of indices \(h+u-2i,h+u-2i-1\) for \(0\le i<u\).
Their union is the interval from \(h-u+1\) through \(r\).
The leading power of \(q\) is \(q^{2rn-n^2}\), obtained by
multiplying the leading powers \(q^{2(r-i)-1}\) at steps
\(0\le i<n\). This proves the residue count.

For a primitive residue tuple the map \(Y\mapsto X^*HY\)
is onto all \(n\)-square matrices, by a left inverse to the
columns and self-duality. The map \(B\mapsto B+B^*\) is onto
Hermitian matrices, using the residue trace on the diagonal.
Each solution therefore has exactly \(q^{2rn-n^2}\) lifts at
each successive digit for the prescribed target. This argument
also applies in residue characteristic two. Taking the defining
limit proves the assertion. The residue counts correspond to
\srcloc{Formalism/LocalToGlobal/SiegelLatticeTheory/SiegelMassRatios/DensExact}{22}
and \srcloc{Formalism/LocalToGlobal/SiegelLatticeTheory/SiegelMassRatios/DensExact}{32}.
\end{proof}

\begin{proposition}[Quotient measure on the Gram fibre]\label{prop:quot215}
For every represented nonsingular \(S\),
\[
                 d\mu_S=P_0^{-1}q^{h v(\det S)}\,d\sigma_S .
\]
\end{proposition}
\begin{proof}
Witt extension is transitive on \(\Sigma_S\), so the two measures
are proportional. For \(P\in\mathrm{GL}_n(A)\) right
multiplication by \(P\) transports the quotient measure unchanged,
since it preserves the span and its pointwise stabilizer. On
the additive matrix space its absolute Jacobian is
\(|N_{A/K}\det P|_K^r\). On the Hermitian target the map
\(B\mapsto P^*BP\) has Jacobian \(|N_{A/K}\det P|_K^n\).
For diagonal \(P\) multiply the diagonal and off-diagonal
Jacobians. Integral invertible factors preserve the target
lattice, so elementary divisors give the result for every \(P\),
separately in the two factors in the split case. The fibre measure
of a transported test is therefore multiplied by
\(|N_{A/K}\det P|_K^h\). This cancels the change of the factor
\(q^{h v(\det S)}\).

It remains to calibrate one target in each congruence class.
The split algebra has a single nonsingular class. Over the
unramified field there are two, represented by \(I_n\) and
\(\operatorname{diag}(I_{n-1},\pi)\). To see this, diagonalize
over the field. Every unit is a norm by residue norm surjectivity
and successive lifts using the surjective residue trace. Reduce
each valuation to zero or one. Two diagonal entries \(\pi,\pi\)
make a hyperbolic plane, as do \(1,1\), since \(-1\) is a norm.
At \(I_n\) every integral frame is primitive and splits off a
unimodular summand. Unimodular Hermitian lattices of the same
rank are isometric, by unit-vector diagonalization and unit norms
in the field case and dual bases in the split case. The compact
orbit of integral frames thus has quotient volume one and
additive fibre measure \(P_0\).

At the other target every integral frame is again primitive.
Otherwise its saturation would decrease the determinant valuation
by at least two. Splitting off the \(n-1\) unit columns leaves
the choice of a vector of norm \(\pi\) in rank \(h+1\).
Lemma~\ref{lem:base213} gives this marked lattice orbit mass
as \(q^h+1\). The stabilizer of the unit columns has volume one
in that intermediate group, by the same unimodular splitting.
Here are the details extending the lemma to residue characteristic
two. A primitive vector \(u_0\) of norm \(\pi\) has an integral
partner \(v_0\) pairing to one. Choose \(\lambda\in\mathcal O\)
of unit trace. The norm of \(v_0+b\lambda u_0\), for \(b\in O\),
is
\[
 H(v_0,v_0)+b\operatorname{Tr}(\lambda)+\pi b^2N_{A/K}(\lambda).
\]
It has a simple root modulo \(\pi\), which lifts. Rescale the
new partner by a unit to restore pairing one. Existence follows
in an integral hyperbolic plane by taking \((1,z)\) with
\(\operatorname{Tr}(z)=\pi\). Such a plane splits from every
unimodular lattice of rank at least two by the unit norm of
\(-1\), self-duality and the surjective integral trace. A residue
unitary isometry lifts a digit at a time using surjectivity of
\(B\mapsto B+B^*\) on residue Hermitian matrices. These facts
supply the partner and reduction steps in Lemma~\ref{lem:base213}.
Its two finite isotropic-plane indices give the volume
\((q+1)/(q^h+1)\) of the perpendicular lattice stabilizer used
there. Fixing the marked generator has index \(q+1\), giving
the asserted orbit mass at every \(q\). The perpendicular
scaling by \(\pi^{-1}\) in that calculation differs from the
determinant scaling by a norm since its determinant has odd
valuation. By Lemma~\ref{lem:primitive215} the additive fibre
measure at this target is \(P_0(1+q^{-h})\), since \(h\) is odd.
The ratio is \(q^h/P_0\), as required.
\end{proof}

\begin{theorem}[Integral frames in the stable range]\label{thm:saturate215}
For integral nonsingular \(S\), let \(M_0=\mathcal O^n\) with
form \(S\), and let \(M_0^\vee\) denote its Hermitian dual. Then
\[
 q^{-h v(\det S)}\mu_S(\Sigma_S\cap\Lambda^n)
   =\sum_{\substack{M_0\subset M'\subset M_0^\vee\\ S(M',M')\subset\mathcal O}}
       [M':M_0]^{-h}p\bigl(h,u(M')\bigr).
\]
Here \([M':M_0]\) is cardinality of the quotient and \(u(M')\)
is the residue radical dimension in a basis of \(M'\). For any
measurable subset \(\mathcal D\) of the integral frames containing
all primitive frames and an absolute constant \(C\),
\begin{equation}
 e^{-Cq^{-m}}\le q^{-h v(\det S)}\mu_S(\mathcal D)
                                    \le e^{Cq^{-m}} . \label{eq:fibre215}
\end{equation}
\end{theorem}
\begin{proof}
For each integral tuple \(X\), the lattice of \(w\in A^n\)
with \(Xw\in\Lambda\) is the unique \(M'\) on which the induced
tuple is primitive. The form is integral on \(M'\), hence
\(M'\subset M_0^\vee\). There are only finitely many such
lattices. Write \(M'=PM_0\), so \(|N_{A/K}\det P|_K=[M':M_0]\).
The change \(X\mapsto XP\) identifies this stratum with all
primitive frames for \(P^*SP\). The quotient measure is
unchanged, and the determinant valuation factor is multiplied by
\([M':M_0]^{-h}\). The proposition and the lemma give the sum.

The term \(M'=M_0\) gives the lower bound, since
\(\lvert\log p(h,u)\rvert\le C_1 q^{-m}\) by summing the
geometric tail beginning at \(h-n+1\). For the upper bound sum
\([M':M_0]^{-h}\) over all finite over-lattices. Over a discrete
valuation ring of residue order \(q_0\) the sum is
\(\prod_{i=0}^{n-1}(1-q_0^{\,i-h})^{-1}\). Indeed duality
turns it into the sum over sublattices. A triangular normal
form with diagonal \(\pi^{b_1},\ldots,\pi^{b_n}\) has
\(q_0^{\sum_{i=1}^n(n-i)b_i}\) possibilities, by reducing
each entry to the right of its diagonal entry modulo that entry.
Independent geometric sums give the product. Use \(q_0=q^2\)
in the field case, and two such sums with \(q_0=q\) in the
split case. Its logarithm is at most \(C_2q^{-m}\), since
\(m\ge2\). This proves both estimates uniformly. An over-lattice
expansion in the nonsplit case also occurs in the formula of
Cho and Yamauchi given by Li and Zhang,
\href{https://arxiv.org/abs/1908.01701}{Theorem 3.5.1}.
\end{proof}

\subsection{The finite orbital masses of the original period}\label{nav:IX-08}
\label{sec:flags215}

Return to the fields, intervals and row levels of
Section~\ref{sec:refine211}. Normalize the local row measure by
\(\operatorname{vol}(K_v^*)=1\), using this same convention in
the identity terms. The row lattice has diagonal entries
\(\pm1\) before its conjugation at \(T\). After conjugation the
two split factors there are mutually dual. Thus at every finite
place \(K_v^*\) is the full stabilizer of a self-dual lattice
\(\Lambda_v\). Our \(E/F\) is unramified at all finite places,
as in Section~\ref{sec:match209}. Put
\[
 h=k-1,\quad n=2d,\quad m=h-2d+1,\quad S_*=d(d-1).
\]
Here \(m\ge2\) on the intervals under consideration for large
\(x\), since \(k\gg d^2\).
For the two frames \(Y_1,Y_2\) take representatives in the
monomial normal form of Lemma~\ref{lem:forms203}. Their two
self pairings are \(H_0=-TE_{dd}\) and their cross pairing is
\(M\), with signed reversed blocks \(M_\nu\) of length
\(h_\nu\), parameter \(c_\nu\), and endpoint \(j_\nu\).
Assume throughout that the rational label \(\gamma\) is relevant
and supported and that \(\mathcal G\ne0\). Its \(2d\) vectors
span a nondegenerate space with perpendicular \(C\).

At \(v\) put \(A=E_v\), \(\mathcal O=\mathcal O_{E_v}\),
\(q=q_v\) and \(\mathcal U=U_d(A)\times U_d(A)\), where \(U_d\)
is upper unitriangular. Normalize \(\mathcal U\) by giving the
integral subgroup volume one. Its subgroup \(D\) preserving \(M\)
has the measure in Lemma~\ref{lem:stab205}, with unit integral
entry volumes in its first frame. In \(D\backslash\mathcal U\)
put \(Z=U^*MV\) at the class of \((U,V)\). Let \(V_v\) be the
quotient volume where \(Z\) is integral. Put \(V_v^{=}=V_v\)
away from \(Q=g_1g_2\). At \(v\mid Q\) define \(V_v^{=}\)
by imposing both congruences
\begin{equation}
   Z-H_0\in\pi M_d(\mathcal O),\qquad
   2H_0-Z-Z^*\in\pi^2 M_d(\mathcal O).
   \label{eq:marked215}
\end{equation}

\begin{lemma}[Integral representatives of one decorated flag]\label{lem:flag215}
A full row lattice admits a decorated flag precisely when some
representative is integral and primitive. Such representatives
form one right orbit under \(U_d(\mathcal O)\), acting freely.
\end{lemma}
\begin{proof}
It suffices to use the reference lattice, by translation.
Theorem~\ref{thm:midlevel131} gives a primitive integral
representative for the reference flag. Conversely take an
integral primitive representative with the same self pairing.
The first \(d-1\) vectors are primitive isotropic columns.
Unimodularity supplies integral dual partners with pairings the
identity. Adding combinations of the isotropic vectors makes the
partners mutually isotropic, since \(B\mapsto B+B^*\) is onto
the integral Hermitian matrices, including at two. Split off
these hyperbolic planes. The terminal vector is orthogonal to
the isotropic columns. Subtract its integral components along
them, an allowed change of representative. Its image in the
middle lattice is primitive by independence of the original
columns after reduction.

At a nonsplit place its norm \(-T\) is a unit, since \(T\)
splits completely. Its line splits off. The remaining self-dual
lattice is determined up to isometry by its rank, by the
unit-vector argument in Proposition~\ref{prop:quot215}. At a
split place the middle vector is a primitive vector and a
primitive covector of fixed pairing \(c=-T\). Make the vector
the first basis vector and write the covector as \((c,t)\).
If \(c\) is a unit, shears fixing the vector remove \(t\).
If \(c\) is a nonunit, \(t\) is primitive, so a basis change
on the other coordinates makes it \((1,0,\ldots,0)\).
These integral normal forms show that the flag belongs to the
reference lattice orbit, also at \(T\). Finally a primitive
integral matrix has an integral left inverse. Applying it to
another integral representative shows that their change of basis
is integral. Every integral unitriangular change is allowed,
and independence of the columns gives freeness.
\end{proof}

\begin{proposition}[Reduction of the finite mass to a matrix volume]
\label{prop:marked215}
Write \(\delta_v=1\) at \(v\mid Q\), and \(0\) elsewhere. With
the measures just chosen, the mass of compatible full-level
cells satisfies
\begin{equation}
 L_v^{=}=q^{h v(\mathcal G)+\delta_v(S_*-2dh)}
             V_v^{=}\exp\bigl(O(q^{-m})\bigr).
 \label{eq:finite215}
\end{equation}
The implied constant is absolute. The same formula with
\(\delta_v=0\) and \(V_v\) gives \(L_v(K_v^*)\) at every place.
\end{proposition}
\begin{proof}
For \(g\) in the local row group set
\((X_1,X_2)=(g^{-1}Y_1U,g^{-1}Y_2V)\).
At full level integrate the indicator that both individual
tuples are primitive and integral, first over \(\mathcal U\).
For each supported pair in \(g\Lambda_v\) the integral is one,
by the single coset of representatives for each flag in
Lemma~\ref{lem:flag215}. For all other pairs it vanishes.
Integration over \(J\backslash G\) now gives \(L_v(K_v^*)\),
since each lattice orbit has measure
\(\operatorname{vol}(K_v^*)/\operatorname{vol}(J\cap gK_v^*g^{-1})\).

We may instead integrate over \(D\backslash\mathcal U\) on the
outside and over \(U(C)\backslash G\) on the inside. Indeed
\(J=D\times U(C)\) by Lemma~\ref{lem:stab205}. If a change of
two representatives is left multiplied by the image of
\(d_0\in D\), replace \(g\) in the integrand by \(d_0^{-1}g\).
The quotient integration formula then absorbs the \(D\)-integral
in the \(G\)-integral. All measures are Haar quotient measures
of unimodular groups, and the integrands are nonnegative.
For fixed \(U,V\) the inner integral is exactly \(\mu_B\) on
the fibre with Gram matrix
\[
                    B=\begin{pmatrix}H_0&Z\\Z^*&H_0\end{pmatrix}.
\]
Its determinant is \((-1)^d\mathcal G\). Its test includes all
jointly primitive integral frames and is contained in the set
of all integral frames. Integrality of \(B\) is equivalent to
integrality of \(Z\). Integrating \eqref{eq:fibre215} over these
\(U,V\) gives the full-level assertion.

At a marked place Lemma~\ref{lem:flags211} tests whether the
two reduced decorated flags agree in the full-level lattice.
The stabilizer of the reduced decoration is precisely its
\(\bar B\) there. Indeed preserving the two split isotropic
strings with fixed graded vectors gives the block unipotent
radical, and fixing the terminal vector and covector on the
middle quotient leaves \(\mathrm{GL}_k\). For equal decorations
the fraction of pairs of integral representatives satisfying
\(X_1\equiv X_2\bmod\pi\) is \(q^{-S_*}\). Once the first
reduction is fixed, exactly one of the
\(|U_d(\mathcal O/\pi)|=q^{d(d-1)}\) changes of the second
reduction matches it, by freeness and surjectivity of reduction.
For different decorations that fraction is zero. Multiply the
test of equal reduced representatives by \(q^{S_*}\) to obtain
the indicator of compatibility in the integral above.

In its inner integral now use the ordered tuple
\(X_1,(X_2-X_1)/\pi\). Its Gram matrix has off-diagonal
block \((Z-H_0)/\pi\) and final diagonal block
\((2H_0-Z-Z^*)/\pi^2\). Its determinant valuation is
\(v(\mathcal G)-2d\). The field change of basis preserves the
quotient measure by Proposition~\ref{prop:quot215}. The new
test is that both tuples be integral and the first primitive.
Indeed the old second tuple then reduces to the old first and
is primitive as well. The new test again lies between all
jointly primitive frames and all integral frames.
Theorem~\ref{thm:saturate215} applies exactly on
\eqref{eq:marked215}. This proves the marked assertion.
\end{proof}

\begin{theorem}[A uniform bound at all finite places]\leavevmode\label{thm:allfinite215}
Set
\[
 a_\nu=\binom{h_\nu}{2},\quad a=\sum_\nu a_\nu,\quad
 b_\nu=2h_\nu(d-j_\nu)+a_\nu,\quad
 W_\gamma=\prod_\nu |N_{E/\QQ}c_\nu|^{-a_\nu}.
\]
Let \(P_\infty\) be the product of real densities in
Lemma~\ref{lem:density207} and \(f=[F:\QQ]\). Then
\begin{align}
 \mathcal M_\gamma^{=}
 &=\mathfrak m_h |N_{F/\QQ}\mathcal G|^h
          Q^{f(S_*-2dh)} \Bigl(\prod_{v<\infty}V_v^{=}\Bigr)
          \exp\bigl(O(f2^{-m})\bigr),                    \label{eq:fprod215}\\
 \frac{\mathcal M_\gamma^{=}}{\mathfrak m_h P_\infty}
 &\le \exp\bigl(O(f2^{-m})\bigr)
            Q^{f(2a-S_*-2dh)} W_\gamma .                \label{eq:fratio215}
\end{align}
The constants are independent of the label and of the selected
field. At the regular points, where every block has size one,
\(W_\gamma=1\).
\end{theorem}
\begin{proof}
By definition \(\mathcal M_\gamma^{=}=\mathfrak m_h\prod_v L_v^{=}\).
Each of \(g_1,g_2\) has \(f\) split places of residue order the
corresponding rational prime. The product of the Gram factors in
\eqref{eq:finite215} is the indicated ideal norm. Also
\(\sum_v q_v^{-m}\le f\sum_{j\ge2}j^{-m}\ll f2^{-m}\)
for \(m\ge2\). This gives \eqref{eq:fprod215}.

We bound the remaining volume in explicit coordinates. Write
\(|c|_A=|N_{A/F_v}c|_v\). In \(Z=LMV\), with \(L=U^*\)
lower unitriangular, choose the representative of
\(D\backslash\mathcal U\) having all diagonal blocks of \(L\)
equal to the identity. This leaves every lower off-block entry
and all upper entries of \(V\) free. These coordinates give
additive product quotient measure. In fact write \(L=L_0 L_b\)
with \(L_b\) its block diagonal part. Right multiplication
by its blocks has determinant one on lower off-diagonal blocks.
The accompanying change in \(V\) is a left translation, and
involution and inversion have modulus one on the unipotent
block coordinates of \(D\).

Expose row and column block \(\nu\) in order, keeping previous
rows and columns fixed. The diagonal block of \(Z\) is a
known matrix plus \(M_\nu V_{\nu\nu}\). Each of its \(a_\nu\)
free entries is multiplied by \(\pm c_\nu\) in a distinct
position, so integrality gives volume at most
\(|c_\nu|_A^{-a_\nu}\). For \(j>\nu\) the new row block has
variable part \(M_\nu V_{\nu j}\). For \(i>\nu\) the new
column block has variable part \(L_{i\nu}M_\nu V_{\nu\nu}\).
In each case the known summand uses only previously exposed
coordinates. The factor \(V_{\nu\nu}\) has determinant one.
The volumes for all these rectangular blocks therefore
contribute \(|c_\nu|_A^{-2h_\nu(d-j_\nu)}\).
At a marked place the first congruence of \eqref{eq:marked215}
prescribes a residue class in each of the same free output
entries. Each such class has volume \(q^{-2}\) in the quadratic
integer algebra. Dropping the second congruence gives the bounds
\[
       V_v^{=}\le q^{-2\delta_v(S_*-a)}
                         \prod_\nu |c_\nu|_A^{-b_\nu},
       \qquad \sum_\nu b_\nu=S_*-a .
\]
At almost all places all \(c_\nu\) are units and integral free
coordinates give volume one. The same bound is then one, so
the product of the \(V_v^{=}\) is finite after omitting its
unit factors.

The finite product formula gives
\(\prod_{v<\infty}|c_\nu|_A^{-1}=|N_{E/\QQ}c_\nu|\).
By Lemma~\ref{lem:density207}, with
\(e_\nu=2(d-j_\nu)+h_\nu-1\),
\[
 P_\infty=|N_{F/\QQ}\mathcal G|^h
                       \prod_\nu |N_{E/\QQ}c_\nu|^{h_\nu e_\nu}.
\]
Since \(b_\nu-h_\nu e_\nu=-a_\nu\), division proves
\eqref{eq:fratio215}. This cancellation includes primes dividing
several nonzero endpoint minors simultaneously.
\end{proof}

\begin{corollary}[Finite measure in the squared real kernel]\label{cor:square215}
Write \(R=|\Delta_E|^{1/(2f)}\), and use the original and
transported kernels of Corollary~\ref{cor:match209}. Then
\begin{align}
 \mathcal M_\gamma^{=}\,c_E^{2a}\|K_{\infty,\gamma}^{\rm area}\|_{\rm HS}^2
 &\le \mathfrak m_h\exp\bigl(O(f2^{-m})\bigr) R^{-2fS_*}
              Q^{f(d(d-3)-2dh)} \notag\\
 &\hspace{15mm}{}\times W_\gamma P_\infty^\#
                           \|\widehat K_{\gamma^\#}\|_{\rm HS}^2 .
 \label{eq:square215}
\end{align}
For the integral relabeling of Lemma~\ref{lem:match209}, put
\(B_\gamma^\#=|N_{F/\QQ}\mathcal G^\#|
 \prod_{\nu}|N_{E/\QQ}D^\#_{j_\nu}|\). Then
\[
               W_\gamma\le Q^{f(d-1)}
                                  (B_\gamma^\#)^{(d-1)/2}.
\]
\end{corollary}
\begin{proof}
Square \eqref{eq:match209} in Hilbert--Schmidt norm and multiply
both sides by \(\mathcal M_\gamma^{=}/P_\infty\).
The exponent of \(Q\) after \eqref{eq:fratio215} is
\(f(2d(d-2)-S_*-2dh)=f(d(d-3)-2dh)\).
For the weight let \(A_j=|N_{E/\QQ}D_j|\), \(A_0=1\).
Since \(|N_{E/\QQ}c_\nu|^{h_\nu}=A_{j_\nu}/A_{j_{\nu-1}}\),
the exponent of each proper endpoint in \(W_\gamma\) is
\((h_{\nu+1}-h_\nu)/2\), and that of the final endpoint is
\(-(h_l-1)/2\). Every nonzero \(A_j\) is an integer at least one.
Drop the negative exponents and bound the others by \((d-1)/2\).
In \eqref{eq:label209} the reindexed endpoint norm is at least
\(A_j\) for \(j\ge2\). At \(j=1\) it is \(Q^{-2f}A_1\).
The reindexed endpoint and Gram norms are positive integers,
proving the bound. This norm weight is defined in
\srcloc{Formalism/SharpRestriction/CuspSpanNormBounds/LeadMinorK}{164}.
\end{proof}

For unit column tuples the proof of Theorem~\ref{thm:norm211}
bounds the absolute period sum for a fixed label by
\(A_*\mathcal M_\gamma^{=}c_E^a\|K_{\infty,\gamma}^{\rm area}\|_{\rm op}
/\sqrt{c_{\rho_1}^*c_{\rho_2}^*}\).
Formula \eqref{eq:fratio215} supplies its finite mass at every
place. Formula \eqref{eq:square215} controls one power of that
mass times a squared kernel. Taking a square of the sum over
cells introduces products from different cells. A bound on the
sum and a lower bound for the designated integral pairing still
enter the comparison in Volume~I.

\subsection{The finite measure above each rational label}\label{nav:IX-19}
In Corollary~\ref{cor:amplitude204} the finite cells over one
\(\gamma\) all have the same real unipotent stabilizer \(D_\infty\),
but can have different rational lattices in that group. On the
same fields \(E_x\) and at the level \(K^0\) these lattices and
their multiplicities admit the following mass formula.

\begin{theorem}[Finite cell measure on the period fields]
\label{thm:mass206}
Fix \(\gamma\in\mathcal R\) of Theorem~\ref{thm:orbits204}. Let
\(W\) be the span of its two frames, \(C=W^\perp\),
\(H_C=U(C)\), and \(a=\sum_\nu\binom{h_\nu}{2}\) for its block
lengths \(h_\nu\). Write \(J=B\cap\gamma B\gamma^{-1}\),
\(D=N\cap\gamma N\gamma^{-1}\), and \(K_v=K_v^0\). Set
\[
 X_v=(B_vK_v\cap\gamma B_vK_v)/K_v,\qquad
 L_{\gamma,v}=
 \sum_{tK_v\in J(F_v)\backslash X_v}
  \frac{\operatorname{vol}(K_v)}
       {\operatorname{vol}_{J_v}(J(F_v)\cap tK_vt^{-1})}.
\]
All these local sums are finite. For \(X_f\) the restricted product
of the \(X_v\) and \(\Gamma_t=J(F)\cap tK^0t^{-1}\), the total
finite measure with its real unipotent covolumes equals
\begin{equation}
 \sum_{tK^0\in J(F)\backslash X_f}
     \operatorname{vol}(K^0)\operatorname{vol}(\Gamma_t\backslash D_\infty)
   =\bigl(2^{-f}\sqrt{|\Delta_{E_x}|}\bigr)^a
       \frac{\operatorname{vol}(H_C(F)\backslash H_C(\mathbb A_F))}
            {\operatorname{vol}(H_C(F_\infty))}
       \prod_{v\in S(\gamma)}L_{\gamma,v}.                \label{eq:mass206}
\end{equation}
Use ordinary complex area and give the finite integer rings volume
one in the independent entries of \(D\). Away from the fixed bad
places of the field and level and the divisors of the nonzero
minors and \(\mathcal G\), use volume one for the self-dual lattice
stabilizers. This exceptional union is \(S(\gamma)\).
\end{theorem}
\begin{proof}
The two nondegenerate frames determine their cross matrix up to
unitriangular changes. In monomial form the common stabilizer on
\(W\) has one upper unitriangular matrix over \(E_x\) in each
reversed block. Its action fixes \(C\) pointwise. The remaining
factor acts only on \(C\), giving \(J=D\times H_C\).
At an ordinary place with unit minors and Gram determinant, integral
Gaussian elimination produces a self-dual lattice on \(W\). The
orthogonal complement is a self-dual lattice in \(C_v\), unique up
to local isometry. There is one local \(J(F_v)\)-orbit and its
stabilizer has volume one. At a remaining place the compact level
meets the closed double coset \(B_v\gamma B_v\) in finitely many
double cosets for the integral period group.
Lemma \ref{lem:stab205} and Proposition \ref{prop:good205} establish these assertions, including the local
extension of isometries at places over \(2\).

For a finite representative let \(K_J=J_f\cap tK^0t^{-1}\),
\(K_D=K_J\cap D_f\) and \(K_C\) be its projection to \(H_{C,f}\).
Density of \(D(F)\) in \(D_f\) and the discriminant covolume of
\(\mathcal O_{E_x}\) give
\(\operatorname{vol}(\Gamma_t\backslash D_\infty)
=(2^{-f}\sqrt{|\Delta_{E_x}|})^a/\operatorname{vol}(K_D)\).
Above this adelic orbit the rational cells are the definite
column classes at level \(K_C\). Their rational stabilizers
are trivial by the two marked column level conditions, so their
number is the definite quotient mass in \eqref{eq:mass206}
divided by \(\operatorname{vol}(K_C)\). Since
\(\operatorname{vol}(K_J)=\operatorname{vol}(K_D)\operatorname{vol}(K_C)\),
summing over the independent local orbits proves the identity.
The covolume and class calculations in full are in
Theorem \ref{thm:mass205}.
\end{proof}

\begin{corollary}[Allowable loss at the extra finite places]
\label{cor:primecost206}
Let \(S_0\) be the places exceptional for the field and level in
Theorem~\ref{thm:mass206}, let \(S'_\gamma=S(\gamma)\setminus S_0\),
and let \(P_\gamma=\prod_{v\in S'_\gamma}\mathrm Nv\).
For any fixed \(0\le\delta<1/2\) and sufficiently large \(V\),
\[
  \sum_{\gamma\in\mathcal R}\frac{P_\gamma^\delta}{w_\gamma}
     \le\left(\frac{128\log(2Q)}{Q^{4(1-\delta)}}\right)^{fd}.
\]
For each fixed \(B\ge1\) and \(\eta>0\), throughout the same
intervals for sufficiently large \(V\),
\[
  \sum_{\gamma\in\mathcal R}\frac{B^{|S'_\gamma|}}{w_\gamma}
            \le \exp\!\bigl(-(64-\eta)fd\,x^\eps\bigr).
\]
\end{corollary}
\begin{proof}
The product of the extra residue sizes divides the product of the
absolute integer norms of the endpoint minors and \(\mathcal G\).
The former product is at most \(w_\gamma^{1/2}\), and
\(0<\tau\mathcal G\le|\tau D_d|^2\) gives the same bound for
the latter norm. Thus \(P_\gamma\le w_\gamma\). The integral of
\(2b_1r(1+r^2)^{-b_1-1}\,dr\) from \(u\) to infinity equals
\((1+u^2)^{-b_1}\). Taking \(b_1=2(1-\delta)>1\) in the box
counts of Theorem~\ref{thm:orbits204} bounds the sum of
\(w_\gamma^{-1+\delta}\) by the displayed expression, as computed in
Theorem \ref{thm:primecost205}.

For the second conclusion choose \(0<\delta<1/2\) with
\(64\delta<\eta\). Each place with \(\mathrm Nv\ge B^{1/\delta}\)
contributes \(B\le(\mathrm Nv)^\delta\). The remaining places lie
above rational primes at most \(B^{1/\delta}\), with at most \(f\)
places above each. Their factors together are at most
\(B^{f\lceil B^{1/\delta}\rceil}\). This has logarithm \(O_{B,\delta}(f)\).
Using \(\log Q=(16+o(1))x^\eps\) in the first estimate gives the
assertion uniformly on the interval.
\end{proof}

For the integral itself, multiply the right side of \eqref{eq:mass206}
by the maximal norms of the two finite covectors and by
\[
   \int_{D_\infty\backslash G_\infty}
       \|W_1(g)\|\,\|W_2(\gamma^{-1}g)\|\,d\dot g .
\]
This bounds the sum of the absolute cell contributions by contraction
of the two real sections and integration along the common unipotent
fibre. A uniform bound on this expression relative to
\(w_\gamma^{-1}\) is still needed. It includes the values of the
exceptional local factors and the mass of the definite complement.
The integral lower bound for the designated theta pairing in
Corollary~\ref{cor:original200} and the total order of the
exceptional cohomology groups in \eqref{eq:region198} remain open.

\begin{proposition}[A family to which the regular orbital identity does not apply]
\label{prop:oneblock206}
On every field and at every translate in Theorem~\ref{thm:mass206} are infinitely many supported, relevant rational double
cosets with \(D_1=\cdots=D_{d-1}=0\) and \(\mathcal G\ne0\).
\end{proposition}
\begin{proof}
In the signed Witt basis choose a positive coordinate vector \(w\)
of norm one orthogonal to the negative terminal vector \(e\).
Their sum is isotropic. Use a shear of the first isotropic basis
vector by a multiple of \(w+e\), and its negative adjoint shear
on the dual isotropic space. A skew-Hermitian antidiagonal matrix
on the intervening indices fills out the remaining cross entries.
Explicitly the resulting cross matrix is the sum of \(-TE_{dd}\)
and a single reversed block with successive nonzero entries
\(na_0,-na_0,\ldots\), where \(\bar a_0=(-1)^d a_0\).
The three block matrices and the verification that they preserve
the Hermitian form are given in
Proposition \ref{prop:oneblock205}.
The shears form a rational one-parameter group. For every sufficiently
divisible integer \(n>0\) its finite component lies in \(K^0\).
Adding a multiple of the first row to the last removes the terminal
diagonal entry. Then \(D_d=(na_0)^d\) distinguishes the rational
cosets and gives \(\mathcal G=|D_d|^2\) at every real place.
The alternating signs make the two sets of simple unipotent
coordinates conjugate in reverse order, so their trace characters
agree.
\end{proof}

For instance the source's formal-degree identity for a regular
orbital integral requires every proper leading minor to be nonzero.
The examples in Proposition~\ref{prop:oneblock206} have a common
unipotent group of positive dimension. For nonzero compact-finite
discrete-series Bessel sections their product of norms has infinite
integral on the whole real group. Indeed the two sections are analytic
by elliptic regularity for the sum of the real and compact Casimir
operators with the elliptic signs. On the connected group their
product of norms is positive on an open set and is constant along
the noncompact common unipotent fibres.
Corollary \ref{cor:divergence205}
gives the full argument. The period formula uses the quotient by
that unipotent group. An estimate on this quotient or a vanishing
result for the particular projected entries is still needed.

\subsection{The contribution of the common unipotent stabilizer}\label{nav:IX-20}
Keep the fields \(E=E_x\), the translates, the prime
\(T=\varpi\bar\varpi\) and the level \(K^0\) of
Theorem~\ref{thm:mass206}. At a real place the two sections
\(W_i=J_i v_i\) for two words have the same discrete-series row
\(\sigma\). Their column actions \(\tau_i\) are unitary. On
a relevant orbit put \(D=N\cap\gamma N\gamma^{-1}\) and fix
Haar measure on \(D\). The operator entering each finite
contraction is
\[
 K_\gamma(W_1,W_2)=
 \int_{D\backslash G}W_1(g)W_2(\gamma^{-1}g)^*\,d\dot g .
\]
Here \(e_1e_2^*\) maps \(u\) to \(\langle u,e_2\rangle e_1\),
with inner products linear in the first variable. Equality of
the two characters on \(D\) makes the integrand well defined.

\begin{theorem}[The orbital identity at the original level]
\label{thm:quot208}
For all sufficiently large \(x\), the quotient kernel on every
relevant orbit with \(\mathcal G\ne0\), including the one-block
orbits in Proposition~\ref{prop:oneblock206}, equals
\begin{equation}
 \frac{\langle v_1,v_2\rangle}{\overline{\ell_1}\ell_2}
 \int_{(B\times B)/\Delta_\gamma D}
   \tau_1(b_1)^{-1}\xi_1\xi_2^*\tau_2(b_2)\,
       \langle\sigma(b_1\gamma b_2^{-1})u_2,u_1\rangle\,d[b_1,b_2].
                                                        \label{eq:quot208}
\end{equation}
Here \(\Delta_\gamma d_0=(d_0,\gamma^{-1}d_0\gamma)\) acts on the
right, and \(\ell_i=\langle\Lambda_i(u_i),\xi_i\rangle\ne0\)
for a compact-finite choice of \(u_i,\xi_i\) in the normalized
Bessel model. Both integrals converge absolutely.
\end{theorem}
\begin{proof}
Average the coefficient \(\langle\sigma(bg)v_i,u_i\rangle\)
against \(\tau_i(b)^{-1}\xi_i\) on \(B\). By the normalized
relative character, this gives
\(\overline{\ell_i}W_i(g)/\sqrt{d_\sigma}\), where \(d_\sigma\)
is the formal degree. Insert the two averages into the quotient
integral and disintegrate \(B\times B\) along the indicated
diagonal group. Its fibre combines with \(D\backslash G\)
to give \(G\). After \(y=b_2\gamma^{-1}g\), the two row
coefficients have arguments \(b_1\gamma b_2^{-1}y\) and \(y\).
Their integral on \(G\) is
\(d_\sigma^{-1}\langle v_1,v_2\rangle
\langle\sigma(b_1\gamma b_2^{-1})u_2,u_1\rangle\)
by Schur orthogonality.

For the interchange, if \(\|b_1\gamma b_2^{-1}\|_{\rm op}\le U\),
elimination in the cross matrix bounds representatives of both
\(B\)-coordinates by \(O_\gamma(U^{d^2+2})\) modulo the common
unipotent group. On eliminating each pivot, the denominators
are products of nonzero fixed pivot entries and the numerators
are minors of order at most \(d\). The two frames span all
negative directions, so bounds on their coordinates bound
the operator norms of the representatives. The quotient ball has
measure \(O_\gamma(U^{8r^2(d^2+2)})\). The row gaps of
Theorem~\ref{thm:ordinary164} give coefficient decay of greater
order even after reserving \(4r^2+1\) powers for integration
in \(y\). This proves absolute convergence in the expansion.
The coordinate and measure computations are given in
Lemma \ref{lem:orbitvol207} and Theorem \ref{thm:quot207}.
\end{proof}

\begin{proposition}[A bound uniform over the labels of one field]
\label{prop:tail208}
For fixed \(x\), put \(L=8(f+1)(r+1)^3\). For \(p,t\ge0\)
with \(2pd(d+1)+t<L-4fr^2\), there is a finite \(C_{x,p,t}\)
such that at the original level
\[
      \sum_c |J_{\gamma,c}|
           \le C_{x,p,t} Q^{-ft} w_\gamma^{-p}
\]
for all relevant nondegenerate labels and unit inputs.
\end{proposition}
\begin{proof}
The determinant \(D_j\) pairs the first \(j\) frame vectors in
the exterior power with the corresponding translated vectors.
The product of the two quotient heights to power \(j\) bounds
its absolute value up to the squared Euclidean norm of that wedge.
In the first reversed block there is an antidiagonal entry pairing
two isotropic frame vectors. Its modulus is bounded by twice the
product of the two heights and equals the first scalar modulus.
The first scalar divided by \(Q\) is a nonzero algebraic integer
by the endpoint congruences. Thus the product over embeddings
of the maxima of one and its unscaled modulus is at least \(Q^f\).
Weight the two positive seeds by the quotient heights to exponent
\(2pd(d+1)+t\), unfold their product, and use Cauchy--Schwarz.
The height-ball count of exponent \(4fr^2\) and the seed decay
of exponent \(L\) put both periodizations in \(L^2\).
The bound follows with their two norms in the constant. The
details, including the uniformity for unit inputs, are in
Proposition \ref{prop:tail207}.
\end{proof}

\begin{theorem}[Arithmetic normalization of all nondegenerate cells]
\label{thm:rd208}
Write \(R=|\Delta_E|^{1/(2f)}\), \(S=d(d-1)\), and use the
block stabilizer dimension \(a\) of Theorem~\ref{thm:mass206}.
At a real embedding let \(P\) be the following positive density
at the block endpoints \(j_\nu=h_1+\cdots+h_\nu\):
\[
        P=\mathcal G^{k-1}
            \prod_{\nu<l}|D_{j_\nu}|^{2(h_\nu+h_{\nu+1})}
                              |D_d|^{2(h_l-1)}.
\]
Put \(P_\infty=\prod_\tau P_\tau\). Scale the \(i\)-th isotropic
vector and its dual by reciprocal factors \(R^{-(d-i)}\) and
\(R^{d-i}\), keeping the terminal vector fixed. Use the conjugated
orbit and the isometrically transported Bessel sections, denoted
by primes. Then
\begin{align}
 c_E^a P_\infty^{1/2}K_{\infty,\gamma}^{\rm area}
       &=R^{-fS}(P'_\infty)^{1/2}K_{\infty,\gamma'}^{\rm root},
                                                        \label{eq:rd208}\\
 \sum_c|J_{\gamma,c}|
       &\le B_\ell B_{\ell'}\,m_C\!
         \Bigl(\prod_{v\in S(\gamma)}L_{\gamma,v}\Bigr)
         R^{-fS}Q^{-fd(k+d-2)}
         \big\|(P'_\infty)^{1/2}K_{\infty,\gamma'}^{\rm root}
                                             \big\|_{\rm op}. \notag
\end{align}
Area is Euclidean measure on the \(a\) complex entries of each
real unipotent stabilizer. Root measure multiplies that area at
each place by \((|c'|/(2\pi))^a\), where \(c'\) is the coefficient
of the new simple-root character.
\end{theorem}
\begin{proof}
For a general scale \(\kappa>0\) in place of \(R^{-1}\), the
block scalar in positions \(j_{\nu-1}+1,\ldots,j_\nu\) acquires
the factor \(\kappa^{e_\nu}\), with
\(e_\nu=2d-j_{\nu-1}-j_\nu-1\).
Let \(H=\sum_\nu h_\nu(h_\nu^2-1)/6\) and
\(E_*=\sum_{i=1}^d4(d-i)^2\).
The character coefficient acquires a factor \(\kappa\).
The root Haar pullback has factor \(\kappa^{a-2H}\), since each
upper entry within a block scales by \(\kappa^{i-j}\).
The two section normalizations together give \(\kappa^{-kS-E_*}\),
the inverse modulus of the frame change on \(B\). Finally \(P^{1/2}\)
acquires \(\kappa^{(k-1)S+E_*-2H}\) by the endpoint formula.
Multiplying these factors, taking the inverse for the change in
quotient measure, gives \(\kappa^{-S-a}\) for the scaled
density times the scaled root kernel.

Originally \(|c|=4\pi\) since the additive character uses
the complex trace. Hence \(K_\infty^{\rm area}=2^{fa}K_\infty^{\rm root}\).
The discriminant covolume is \(c_E=(R/2)^f\). Putting
\(\kappa=R^{-1}\) in the preceding computation gives
\eqref{eq:rd208}. The proper endpoint congruences and the integral
Gram determinant \(\mathcal G/Q^{2d}\) imply
\(P_\infty^{1/2}\ge Q^{fd(k+d-2)}\). If the last block has length
one, \(D_d\) occurs to exponent zero in the product of endpoint
factors. Otherwise \(D_{d-1}=0\), so the full endpoint has the
same homogeneous congruence as the others. Now contract the actual
operator with the two finite covectors, use their maximum norms
\(B_\ell,B_{\ell'}\), and sum the covolumes by
Theorem~\ref{thm:mass206}. The general change of frame and the
corresponding change of the regular squared norm are proved in
Section \ref{nav:IX-04}.
\end{proof}

Proposition~\ref{prop:tail208} gives a bound for all labels of
each \(E_x\). Corollary~\ref{cor:amplitude204} requires a bound
as \(x\) varies on the common intervals. For instance, taking
\(p=1\) in the proposition would suffice for its hypothesis if
\(\log C_{x,1,t}-ft\log Q\le(64-\delta)fd\,x^\eps\) for one
fixed \(\delta>0\) at all sufficiently large \(x\).
Theorem~\ref{thm:rd208} evaluates the discriminant and level
factors entering this comparison. The growth of the remaining
operator norms and finite factors, the lower bound for the
designated integral theta pairing in Corollary~\ref{cor:original200},
and the exceptional cohomology sum in \eqref{eq:region198}
remain to be controlled on those intervals.

\subsection{Summing the columns at the two auxiliary primes}\label{nav:IX-21}
Fix \(x\) on the interval of Theorem~\ref{thm:region198}. At the
marked places \(v\mid Q=g_1g_2\), enlarge \(K=K^0\) to the full
row lattice stabilizer and call the resulting finite level \(K^*\).
At the terminal prime \(T\) we use the prescribed level in both.
Write \(I_\sigma=[K^*:K]\),
\(I_\pi=[H_f\cap K^*:H_f\cap K]\). For a word with definite
column representation \(\rho\), let \(c_{\rho,K}\) be its identity
term constant, so the squared norm of a column tuple \(\ell\)
is \(c_{\rho,K}\sum_y\|\ell_y\|^2\). Put
\[
 c_\rho^*=
 \frac{\operatorname{vol}(K^*)\,v_N\operatorname{vol}(H_\infty)}
      {\operatorname{vol}(N_f\cap K^*)\,\dim\rho},
 \qquad v_N=\operatorname{vol}(N(F)\backslash N(\mathbb A_F)).
\]

\begin{lemma}[Two flags over the same residue field]
\label{lem:refine212}
At \(v\mid g\), \(g\in\{g_1,g_2\}\), each full-level row coset
supporting both flags has either zero or
\([\mathrm{GL}_k(\mathbb F_g):U_k(\mathbb F_g)]\) supporting
refinements. In the latter case each flag projects bijectively
from the common refinements onto the local column cosets.
\end{lemma}
\begin{proof}
Use the first split factor in the ordered Witt basis. Let \(\bar N\)
be the unipotent radical for the blocks of sizes
\(1,\ldots,1,k+1,1,\ldots,1\). On the middle block the period
subgroup acts by \(\operatorname{diag}(h,1)\). Its integral
reduction is \(\bar B=\bar N\rtimes\mathrm{GL}_k(\mathbb F_g)\).
All these matrices lift by lifting their individual unipotent
blocks and the invertible middle block. The row level contains
the principal congruence kernel and reduces to
\(\bar K=\bar N\rtimes U_k(\mathbb F_g)\), by the zero entries
in the terminal column of Proposition~\ref{prop:index90}.
For each of the two flags, the allowed residue refinements of
a full row coset are the right \(\bar K\)-cosets in a single
left translate of \(\bar B\). The two translates agree if they
meet. In that case the refinements are parametrized by
\(\bar B/\bar K\). Projection onto the middle \(k\)-block
identifies them with \(\mathrm{GL}_k(\mathbb F_g)/U_k(\mathbb F_g)\).
The choice of a representative for either flag acts on this
set by left translation, a bijection.
\end{proof}

\begin{theorem}[A bound independent of the two level indices]
\label{thm:norm212}
For every relevant original label \(\gamma\) of nonzero Gram
determinant and any two unit column tuples at \(K^0\),
\[
 \sum_c|J_{\gamma,c}|
 \le
 \frac{A_*\,\mathcal M_\gamma^{=}}
      {\sqrt{c_{\rho_1}^*c_{\rho_2}^*}}\,
 R^{-fS}Q^{-f(dk+a)}(P_\infty^\#)^{1/2}
                        \|\widehat K_{\gamma^\#}\|_{\rm op}.
\]
Here the real inputs are normalized as in Corollary~\ref{cor:match210}.
The factor \(\mathcal M_\gamma^{=}\) is \(\mathfrak m_{k-1}\)
times the product of the local sums of Theorem~\ref{thm:mass206}
at \(K^*\), restricting at \(v\mid Q\) to the row cosets for which
the two translates in Lemma~\ref{lem:refine212} agree. Moreover
\(\log A_*=O(fk\log^2(2fk))=o(fk\log b)\), uniformly on the
cohomology intervals.
\end{theorem}
\begin{proof}
Above a compatible rational coset \(tK^*\) there are \(I_\pi\)
simultaneous refinements before dividing by rational symmetries.
Choose \(t\) from a supporting fine refinement. On each side,
the \(I_\pi\) parameters project once onto the column level
cosets by the lemma. Write \(A_i\) for the finite rational
stabilizer of the underlying coarse column class \(y_i\).
It acts freely on that set, since the fine rational column
stabilizers are trivial. Thus the sum of the squared norms
of the covectors on side \(i\) over the common parameters is
\(|A_i|\sum_{y\mapsto y_i}\|\ell_{i,y}\|^2\).

Let \(\Gamma_D=D(F)\cap tK^*t^{-1}\) and let \(a_t\) be its
index in \(J(F)\cap tK^*t^{-1}\). Its finiteness follows by
projecting into \(A_1\), with kernel \(\Gamma_D\). At each
marked place \(\Gamma_D\) acts through \(\bar N\) on the
refinements and fixes them all. The quotient of order \(a_t\)
acts freely by triviality of the fine rational column stabilizers.
Put \(m_t^*=\operatorname{vol}(K^*)\operatorname{vol}
(\Gamma_D\backslash D_\infty)/a_t\). The fine cells over \(t\)
therefore have total measure \(m_t^*I_\pi/I_\sigma\).
Locally the same ratio is obtained by counting the supporting
refinements in each full-level stabilizer orbit. Formula
\eqref{eq:mass206} now gives
\(\sum_t m_t^*=c_E^a\mathcal M_\gamma^{=}\), with
\(c_E=2^{-f}\sqrt{|\Delta_E|}\).

To insert the actual covectors, sum the product of their norms
over the common refinements and apply Cauchy--Schwarz to the
two squared sums just calculated. The coefficient of this sum
is \(m_t^*/I_\sigma\) times the operator norm of the real kernel.
The unipotent levels agree because all integral unipotent blocks
reduce into \(\bar N\). Hence \(c_{\rho,K}=c_\rho^*/I_\sigma\),
and for a unit input its full sum of squared covector norms is
\(I_\sigma/c_\rho^*\). The two square roots in Cauchy--Schwarz
cancel \(I_\sigma^{-1}\). Bound \(\sqrt{|A_1||A_2|}\) by the
maximum \(A_*\) over the coarse column classes, sum over \(t\),
and apply the density and scaling identity of
Corollary~\ref{cor:match210}. This gives the asserted bound.

Each \(A_i\) acts faithfully on a rational vector space of
dimension \(n=2fk\). Minkowski's finite subgroup theorem gives
\(v_p(|A_i|)\le\sum_{j\ge0}\lfloor n/(p^j(p-1))\rfloor\).
Only \(p\le n+1\) contribute, with valuation at most
\(np/(p-1)^2\). Summing their logarithmic contributions bounds
\(\log A_*\) by \(O(n\log^2(2n))\). The ranks and degrees grow
polynomially in \(x\), whereas \(\log b=x^\epsilon\).
The local mass identity, the action of the rational stabilizers
and the normalization appear in full in
Section \ref{nav:IX-05}.
\end{proof}

Thus a column concentrated on one fine class also satisfies the
estimate. A uniform upper bound for the remaining local masses
and real orbital integrals, the lower comparison for the actual
integral theta pairing in Corollary~\ref{cor:original200}, and the
total exceptional strict order in \eqref{eq:region198} remain
necessary on these same fields and intervals.

\subsection{The primes dividing only the Gram determinant}\label{nav:IX-22}
Fix an original relevant period term at \(K^0\) with nonzero
\(\mathcal G=D_d\overline{D_d}-T^2D_{d-1}\overline{D_{d-1}}\).
The finite product in Theorem~\ref{thm:norm212} has an explicit
factor at the primes of \(\mathcal G\) where all nonzero endpoint
minors are units. The multiplicity with which a prime divides
the determinant enters the answer.

\begin{theorem}[Local masses at the terminal Gram determinant]
\label{thm:local214}
Let \(v\notin S_0\) be an unmarked place of odd residue
characteristic and residue size \(q\), with the full row lattice
stabilizer of volume one and the complement measure of
Proposition~\ref{prop:complement210}. Suppose the last block is
a singleton and every nonzero proper endpoint minor is a unit
at \(v\), in both factors in the split case. With \(e=v(\mathcal G)\)
the local mass in Theorem~\ref{thm:norm212} is
\[
             L_v^{=}=1+q^{k-1}+q^{2(k-1)}+\cdots+q^{e(k-1)}.
\]
Here the sum has just the term \(1\) if \(e=0\).
\end{theorem}
\begin{proof}
In the integral pair of decorated frames, Gaussian elimination
uses unit pivots in the first \(d-1\) positions, in the order
given by the reversed blocks. The corresponding vectors of
both frames together split off a unimodular hyperbolic summand.
The remaining two vectors have norm \(-T\), a unit, and mutual
pairing \(D_d/D_{d-1}\). Split off the first vector and project
the second perpendicularly to it. This gives a vector \(w\) in
rank \(k\) of norm \(a=\mathcal G/(T D_{d-1}\overline{D_{d-1}})\).
The local sum is over self-dual lattices containing \(w\), modulo
its unitary stabilizer, each counted with the inverse volume of
its lattice stabilizer. The common unipotent subgroup acts only
on the split-off hyperbolic summand, where its integral block
entries have volume one. The integral extensions of both
decorations and the orbit identification are given in
Lemma \ref{lem:vector213}.

At a split place, a self-dual lattice is given by a lattice \(L\)
in a \(k\)-dimensional space and its dual. Write \(w=(u,\eta)\),
with \(\eta(u)=a\). We need \(u\in L\) and \(\eta(L)\) integral.
After fixing \(L_0=L\cap\ker\eta\) by a change of basis, the
remaining generator is \(\pi^{-i}u+y\), where \(0\le i\le e\)
and \(y\) runs over \(\pi^{-i}L_0/L_0\). Each orbit under the
volume-one group \(\mathrm{GL}(L_0)\) has size equal to the
inverse volume of its stabilizer. There are \(q^{(k-1)i}\)
possibilities at \(i\).

At an unramified field place the two initial masses, for norm
valuations zero and one, are \(1\) and \(A+1\) with \(A=q^{k-1}\).
At zero the perpendicular lattice is self-dual. At one a primitive
vector \(u\) extends to a unimodular plane with an isotropic
partner \(v_1\) pairing to one. In the perpendicular of \(u\)
the intersection lattice is generated by \(u-h(u,u)v_1\) and
the unimodular rank \(k-2\) summand. For the form scaled by
\(\pi^{-1}\), its lattice stabilizer has volume
\((q+1)/(q^{k-1}+1)\). This is the ratio of the numbers of
maximal isotropic \((k-2)/2\)-planes in the residue Hermitian
spaces of dimensions \(k-2\) and \(k-1\), by counting incident
self-dual lattices. Fixing the distinguished generator modulo
\(\pi\) has a further index \(q+1\).
Lemma \ref{lem:base213}
proves these assertions and computes both incidence counts.

For higher valuations, use additive volume on a self-dual
lattice and form the density of the equation \(h(z,z)=a\)
by counting solutions modulo increasing powers of \(\pi\).
On primitive vectors the norm derivative is surjective modulo
\(\pi\). The density for a unit target is \(H_1=1-q^{-k}\),
and for a nonunit target it is \(H_0=H_1(1+A^{-1})\), by
counting the residue norms in \(k\) coordinates. Multiplication
of a vector by \(\pi\) scales additive surface measure by
\(A^{-2}\) and transports the quotient measure with the same
stabilizer. Decomposition according to the exact divisibility
of \(z\), with the two initial masses just computed, gives
\[
 \begin{cases}
  1+(1+A^{-1})\sum_{j=1}^h A^{2j},&e=2h,\\
  (A+1)\sum_{j=0}^h A^{2j},&e=2h+1 .
 \end{cases}
\]
Both are \(\sum_{i=0}^{e}A^i\).
Theorem \ref{thm:local213}
gives the quotient measure comparison and the density calculation.
\end{proof}

\begin{corollary}[The part of the Gram ideal away from the minors]
\label{cor:product214}
Let \(\mathcal U_\gamma\) be the places \(v\notin S_0\) where all
nonzero endpoint minors are units. Put
\(\mathfrak b_\gamma=\prod_{v\in\mathcal U_\gamma}
\mathfrak p_v^{v(\mathcal G)}\). For all the labels on the same
cohomology intervals, uniformly as \(x\) grows,
\[
 \prod_{v\in\mathcal U_\gamma}L_v^{=}
 =(\mathrm N\mathfrak b_\gamma)^{k-1}
      \sum_{\mathfrak a\mid\mathfrak b_\gamma}(\mathrm N\mathfrak a)^{1-k}
 =(\mathrm N\mathfrak b_\gamma)^{k-1}\bigl(1+O(f2^{-k})\bigr).
\]
\end{corollary}
\begin{proof}
If the last block has length greater than one, \(D_{d-1}=0\)
and \(D_d\) is a unit on this set. Hence \(\mathcal G\) is a unit
and the ordinary factor is one by Theorem~\ref{thm:mass206}.
Otherwise use Theorem~\ref{thm:local214}. Factoring out the
highest power in each local sum and multiplying gives the
divisor sum of the ideal. Its excess over one is at most
\(\zeta_F(k-1)-1\le\zeta(k-1)^f-1=O(f2^{-k})\) since
\(f2^{-k}\to0\) on these intervals.
\end{proof}

The resulting power of \(\mathrm N\mathfrak b_\gamma\) enters
the estimate for the original period. Its comparison with the
archimedean orbital integrals and the remaining finite factors,
the lower bound for the designated integral theta pairing, and
the total cohomology order in \eqref{eq:region198} are open here.

\subsection{The finite masses on the common field}\label{nav:IX-23}
Use the fields and intervals of \eqref{eq:region198}. For a
relevant, finitely supported pair of nonzero Gram determinant,
let \(h=k-1\), \(S=d(d-1)\) and \(m=h-2d+1\). In its monomial
cross matrix let \(c_\nu\) denote the first antidiagonal
entry of the block of length \(h_\nu\). Put
\[
 a_\nu=\binom{h_\nu}{2},\qquad a=\sum_\nu a_\nu,\qquad
 W_\gamma=\prod_\nu |N_{E/\QQ}c_\nu|^{-a_\nu}.
\]
Give each full row lattice stabilizer \(K_v^*\) volume one. Keep
the equality test on the reduced flags at both auxiliary primes
in the mass \(\mathcal M_\gamma^{=}\) of
Theorem~\ref{thm:norm212}.

\begin{theorem}[The finite mass relative to the real density]
\label{thm:finite216}
With \(P_\infty\) the product of the real point densities,
uniformly in the field and the supported label,
\[
 \frac{\mathcal M_\gamma^{=}}{\mathfrak m_{k-1}P_\infty}
   \le \exp\bigl(O(f2^{-m})\bigr)
                    Q^{f(2a-S-2dh)}W_\gamma .
\]
This includes every finite place of the original period.
\end{theorem}
\begin{proof}
Here \(E/F\) is unramified at every finite place, the row
lattices are self-dual, and the individual \(d\)-frames are
saturated in the reference lattice. At \(T\) the two components
are the modified lattice and its dual. The ranks are
\(r=k+2d-1\), \(n=2d\), and \(h=r-n\) is odd with \(m\ge2\)
for large \(x\). Thus the local frame theorem applies at every
place. Its proof counts independent columns over the residue
field, lifts them, and sums over the possible saturations of
an integral tuple. For the full row orbit the resulting main
factor is \(q_v^{h v(\mathcal G)}\), up to
\(\exp(O(q_v^{-m}))\) and the volume of integral cross matrices.
At a marked place replace the second frame by its difference
from the first divided by a uniformizer. This changes the
Gram valuation by \(-2d\). Counting equal representatives of
the two reduced flags contributes \(q_v^S\). Prescribing the
residues in the \(S-a\) free cross-matrix coordinates contributes
at most \(q_v^{-2(S-a)}\) in addition to the usual pivot
Jacobians. The two marked rational primes split completely.
After multiplication the Gram factors cancel those in \(P_\infty\).
The pivot Jacobians leave exponent \(-a_\nu\) on each norm of
\(c_\nu\). The error follows from
\(\sum_{v<\infty}q_v^{-m}\ll f2^{-m}\). The complete measure
calculation is in
Sections \ref{nav:IX-07} and \ref{nav:IX-08}.
\end{proof}

\begin{corollary}[The squared kernel for one rational pair]
\label{cor:square216}
For the transported test and kernel of Corollary~\ref{cor:match210},
\begin{align*}
 \mathcal M_\gamma^{=}c_E^{2a}\|K_{\infty,\gamma}^{\rm area}\|_{\rm HS}^2
 &\le \mathfrak m_{k-1}\exp\bigl(O(f2^{-m})\bigr)
             R^{-2fS}Q^{f(d(d-3)-2dh)}\\
 &\hspace{14mm}{}\times W_\gamma P_\infty^\#
                         \|\widehat K_{\gamma^\#}\|_{\rm HS}^2 .
\end{align*}
\end{corollary}
\begin{proof}
The squared change-of-variables identity contributes
\(R^{-2fS}Q^{2f(d(d-2)-a)}P_\infty^\#/P_\infty\).
Apply Theorem~\ref{thm:finite216} and collect the powers of \(Q\).
\end{proof}

The corollary bounds one squared kernel with its finite orbital
measure. The upper estimate for the period summed over rational
cells, the designated lower estimate in
Corollary~\ref{cor:original200}, and the exceptional strict
cohomology orders in \eqref{eq:region198} are still required.

\section{Integral lattices at the marked primes}\label{nav:g-IX-1}

\subsection{The two Hermitian column lattices}\label{nav:II-61}
Use the inert place, lattices and notation of
Proposition~\ref{prop:inertflag289}. Two self-dual column
lattices are called neighbors if their intersection has
one-dimensional residue quotient in each. Let \(A_v\) sum
the translations to all neighbors, with the column compact
subgroup \(K_H=U(\Lambda)\) of volume one.

\begin{proposition}[Column position and multiplicity at an inert place]\leavevmode
\label{prop:neighbors289}
The two column lattices agree for the \(P^{2h}(P-1)\)
terms with \(x=0\). In every other term they are neighbors.
There are \(d=P M_k/(P^2-1)\) neighbors, forming one
\(K_H\)-orbit, each with unsigned multiplicity
\(P^{2h}(P^2-1)\). On spherical columns the signed local
factor is \(a'I+u'A_v\), with
\[
       |a'|\le P^{2h}(P-1),\qquad
       |u'|\le P^{2h}(P^2-1).
\]
\end{proposition}
\begin{proof}
The first lattice \(n(w,z)L\) projects to \(\Lambda\) in the
column. For the second, project the intersection of \(g^{-1}L\)
with the perpendicular of the isotropic flag and terminal
vector, modulo the flag, as in Proposition~\ref{prop:lattice282}.
For integral input \((t,m,\beta)\), the middle and last
coordinates of \(g^{-1}=\gamma_c^{-1}n(w,z)\) are
\[
 m+X_0\beta/c,\qquad -ct-h_M(X_0,m)+(1-Z_0)\beta .
\]
The last must vanish. The other perpendicularity conditions
put \(m\) in the integral sum of the remaining isotropic
vectors and \(\Lambda\), because \(X_0\) belongs to that sum.
Writing \(x_0\) for its column component, the second lattice is
\[
 \Lambda'=\{y+x_0\beta/c:\ y\in\Lambda,\ \beta\in O,\quad
              h_M(x_0,y)\equiv(1-Z_0)\beta\bmod c\}.
\]
It is self-dual. Indeed \(g\in K B(C)\), so the projected
lattice is the image of \(\Lambda\) under the unitary column
factor of \(g^{-1}\). For \(x=0\) the formula gives exactly
\(\Lambda\). For \(x\ne0\) its intersection \(N\) with
\(\Lambda\) is the inverse image of the hyperplane \(x^\perp\).
Integrality forces \(c\mid\beta\), and adding an integral
multiple of \(x_0\) preserves this hyperplane since \(x\)
is isotropic. Conversely every lift of a vector in that
hyperplane occurs by taking \(\beta=0\).
Modulo \(\Lambda\) the vectors of \(\Lambda'\) fill the
fractional line generated by \(x_0/c\), since the residue
functional \(h_M(x,\cdot)\) is onto. Both quotients by \(N\)
have dimension one.

For any self-dual neighbor, its image after multiplication by
\(\pi\) and reduction in \(\Lambda/\pi\Lambda\) is an isotropic
line. Indeed the pairings of its elements are integral before
multiplication. Its intersection with \(\Lambda\) reduces
into the perpendicular of this line and equals the inverse
image by its index. There are \(M_k/(P^2-1)\) such lines.
Lift a generator of any line to a primitive isotropic
\(e\in\Lambda\). Starting with any lift, one corrects its
norm successively modulo \(\pi^{j+1}\) by adding \(\pi^j y\),
\(j\ge1\). The derivative is the residue trace of
\(h_M(e,y)\), which is onto. Completeness gives the lift.
Choose an integral dual vector \(f_1\), and add a multiple of
\(e\) to give \(f_1\) norm zero by trace surjectivity.

In this hyperbolic plane \(N\) is generated by \(e,\pi f_1\),
with the full perpendicular lattice adjoined. Its dual
\(N^\#\) has plane generators \(e/\pi,f_1\). Self-dual
lattices between \(N\) and \(N^\#\) correspond to isotropic
lines in the plane \(N^\#/N\) with pairing induced by
\(\pi h_M\bmod\pi\). This follows by taking the perpendicular
of a line in that nondegenerate plane. Besides the line
of \(f_1\), which gives \(\Lambda\), they are the lines of
\(e/\pi+t f_1\) with \(t+\bar t=0\) in \(\kappa\).
There are \(P\). The integral shears
\(e\mapsto e+\pi z f_1\) with \(z+\bar z=0\) permute these
lines transitively. Every residue skew scalar lifts to a
skew integer, for \(c/\pi\) is a skew unit whose residue
spans the trace-zero line also in characteristic two.
Isotropic lines of the original residue lattice are conjugate
under \(K_H\) by the preceding hyperbolic completion and the
uniqueness of the unimodular complement of a given rank
in Proposition~\ref{prop:quot215}. Thus all neighbors occur
in the double coset that scales \(e,f_1\) by
\(\pi^{-1},\pi\), respectively.

Integral column isometries act on the successful unipotent
cosets by conjugation and transport their column lattices,
so the multiplicity per neighbor is the second count of
Proposition~\ref{prop:inertflag289} divided by \(d\).
Finally the average of a translation between two
\(K_H\)-fixed vectors is the same throughout a double coset.
On the neighbor coset it is \(A_v/d\), by decomposition
into its \(d\) compact right cosets. Each unipotent term
has a character of modulus one, or its compact average of
modulus at most one, exactly as in
Proposition~\ref{prop:split287}. This gives the two signed
coefficient bounds.
\end{proof}

\subsection{The first-plane factor at every good prime power}\label{nav:II-62}
Use the field and levels of Section~\ref{sec:missing}. Let \(F_v(c)\)
be the signed finite orbital factor for \(\gamma_c\) between
the original columns at a place of full lattice levels
away from \(QT\). Put \(P=Nv\) and normalize by the
identity orbital mass as in Theorem~\ref{thm:finite281}.

\begin{proposition}[A uniform bound at arbitrary good valuation]\leavevmode
\label{prop:powers290}
For every \(e=v(c)\ge1\), at a split or inert place,
\[
       \|F_v(c)\|
          \le \bigl[4(k+1)^2\bigr]^e P^{(k+4s+1)e}.
\]
The constants hold for all the prescribed words, all unit
parts of \(c\), and all residue characteristics.
\end{proposition}
\begin{proof}
Use the variables and truncated valuations \(t,u'\) of
Lemma~\ref{lem:depth290}. For each successful unipotent
coset its contribution is a column translation compressed
between \(K_H\)-fixed vectors, with a character multiplier
of modulus at most one. This follows by taking its two
column coordinates and averaging over the right compact
subgroup, as in Proposition~\ref{prop:split287}. No
cancellation between different cosets is needed here.
Write \(j=e-t-u'\) and \(j_+=\max(0,j)\). On a split
column \(2j_+\le d_-+d_+\le2e\). On an inert column
\(j_+\le d\le e\). Hence the bound per coset from
Proposition~\ref{prop:cartan290} is at most
\[
               (k+1)^{2e}P^{-(k-2)j_+}.
\]

At a split place integrality of \(X,Y,Z\) restricts
the upper blocks to fractional entries of denominator
dividing \(c\). In each right integral coset choose
the lifts of \(X,Y\bmod c\) first. Right multiplication
by an integral unipotent adjusts the two blocks independently
and changes \(Z\) by an integral scalar, since \(Y\) is
integral. After they are fixed, only an integral central
shift is possible, leaving at most \(P^e\) choices of
\(Z\bmod c\). The two middle blocks have length
\(n=r-2=k+2s-1\). Their column and dual column coordinates
therefore allow at most \(P^{k(2e-t-u')}\) choices with
these depths. The remaining coordinates give at most
\(P^{2(n-k)e}\).

At an inert place fix \(X\bmod c\) in the quadratic
integer ring. Then \(Y=h_M(X,\cdot)\) is determined and
\(u'=t\). At this depth the column coordinates give at most
\(P^{2k(e-t)}=P^{k(2e-t-u')}\) choices, and the other
coordinates at most \(P^{2(n-k)e}\). Integral right
translations permit a chosen lift of \(X\), using the
trace-one element as in Proposition~\ref{prop:inertflag289}.
For that lift the admissible integral \(Z\)'s satisfy the
same difference equation as there, hence form at most
one translate of \(O_{F_v}\). The remaining integral
central translations change \(Z\) by \(cO^-\).
Since \(c/\pi^e\) is a skew unit, this is precisely
\(\pi^e O_{F_v}\). There are at most \(P^e\) scalar choices.

Thus on either side the contribution in norm at fixed
\(t,u'\) is bounded by
\[
  (k+1)^{2e}
     P^{(2(n-k)+1)e+k(e+j)-(k-2)j_+}.
\]
For \(j\ge0\) the last two terms of the exponent equal
\(ke+2j\le(k+2)e\). For \(j<0\) they are bounded by
\(ke\). Now \(2(n-k)+1=4s-1\). There are at most
\((e+1)^2\le4^e\) possible depths. Summing proves the
bound, with the exact physical character still present.
\end{proof}

\subsection{The first-plane mass at the three marked primes}\label{nav:II-63}
Keep the physical level \(K^0\) of Section~\ref{sec:missing}, with
\(Q=g_1g_2\) and terminal prime \(T\). Write \(F=E^+\),
\(f=[F:\QQ]\), \(n=r-2\), \(R=n+1=k+2s\).
For \(0\ne\xi\in\mathcal O_E^-\) let \(\mathfrak b\) denote
the ideal of \(F\) whose extension is \(\xi\mathcal O_E\), as
in Theorem~\ref{thm:skew289}. Factor it as
\(\mathfrak b=\mathfrak g\mathfrak m\), where \(\mathfrak m\)
is supported above \(QT\) and \((\mathfrak g,QT)=1\).
Set \(c=Q\xi\). The level specified in Section~\ref{sec:global110}
is the full lattice stabilizer away from \(g_1,g_2\).
The change of row lattice in Theorem~\ref{thm:midlevel131}
has support only over \(T\). Thus every divisor of
\(\mathfrak g\) has the full levels used in
Proposition~\ref{prop:powers290}.

\begin{proposition}[A bound including every marked valuation]\leavevmode
\label{prop:marked291}
At \(v\mid QT\), put \(P=Nv\) and \(a=v(\mathfrak m)\).
The unsigned first-plane mass relative to the diagonal is
at most \(P^{(2n+1)a}\). The signed operator at all finite
places consequently has norm at most
\[
    C_k^{\Omega(\mathfrak g)}
       (N\mathfrak g)^{R+2s+1}(N\mathfrak m)^{2R-1},
       \qquad C_k=4(k+1)^2 .
\]
Here \(\Omega\) counts prime factors with multiplicity. The
bound holds at the original column level also when its
reduction at the specified place is a proper subgroup.
\end{proposition}
\begin{proof}
All three primes split completely. Use the first split
component and the two mutually dual middle lattices there.
At \(T\) take a basis of the actual graph lattice in the
middle and the dual basis on the other side. Together
with the remaining Witt pairs these give integral middle
coordinates of dimension \(n\). The first pair \(u,v_1\)
is integral at every marked place by
Proposition~\ref{prop:plane281}. In block order \(u,M,v_1\),
the first unipotent radical has unrestricted upper blocks
\((l,H,z)\), where \(l\) is a covector and \(H\) a vector.
Its integral subgroup \(U_L\) consists of integral blocks
in these bases. Theorem~\ref{thm:finite281} gives the mass
as the number of right \(U_L\)-cosets with
\(g=n(l,H,z)^{-1}\gamma_c\in K_v^0 B(F_v)\), at the graph
and at the two unitriangular levels alike.

Put \(\delta=1\) above \(Q\) and \(\delta=0\) above \(T\).
Then \(v(c)=a+\delta\). In the coordinates of
Lemma~\ref{lem:depth290}, the middle part of \(gu\) is
\(-cH\), and the last row of \(g^{-1}\) is
\((-c,-cl,1-cz)\). The first vector and first dual covector
are fixed by the period subgroup. At \(T\) both images
are integral. Above \(Q\) they agree with the reference
ones modulo \(\pi\) because the reduction of the row level is
upper unitriangular, fixing the first vector and the last
covector. The integral change of first Witt partner adds
a multiple of \(u\) and preserves these two tests, as in
Theorem~\ref{thm:finite281}. It follows that
\[
       l\in\pi^{-a}O^{1\times n},\qquad
       H\in\pi^{-a}O^{n\times1},\qquad z\in\pi^{-a}O .
\]
For the graph these statements use the bases chosen above.
There are at most \(P^{2na}\) choices of \(l,H\) modulo
their integral lattices. Fix representatives of those choices
by right multiplying by integral unipotents. The new \(z\)
differs by an integral scalar plus the product of \(l\)
with an integral column, so it is still in \(\pi^{-a}O\).
Once \(l,H\) are fixed exactly, right integral equivalence
is translation of \(z\) by \(O\), giving at most \(P^a\)
choices. This proves the local count, including \(a=0\).

Each successful coset pairs two column arguments by unitary
translation, possibly averaged over compact stabilizers.
It has norm at most one, since translations are isometries,
the character has modulus one, and the compact averages are
probability averages. More explicitly, a finite coset class
with stabilizer of relative index \(j\) has weight \(j\)
in the integration of Theorem~\ref{thm:finite281}. Its
pairing is the sum of \(j\) such terms, so that summing
their norm bounds gives precisely the unsigned mass.
Multiplying the marked bounds and
Proposition~\ref{prop:powers290} at the good divisors proves
the displayed estimate. At the other places the minimal
valuation has mass one. The product uses the unfolding of
Theorem~\ref{thm:skew289} and \(2n+1=2R-1\).
\end{proof}

\subsection{Integrating the covector in the terminal graph}\label{nav:II-64}
Fix \(v\mid T\). In the first split factor write \(L=F_v\),
\(O=\mathcal O_L\), \(\pi=T\) and \(P=Nv\). With an integral
lift of the primitive residue covector \(\ell\), the graph
coordinates of Theorem~\ref{thm:midlevel131} are
\((w,z')\), where \(z'=(z-\ell w)/\pi\). In these coordinates
the central part of the row lattice is \(O^k\oplus O\).
The row level \(K\) is the full stabilizer of the resulting
row lattice. Put \(C=\mathrm{GL}_k(O)\) and \(y_0=\ell/\pi\).
In the central block consider the group
\[
 D=\left\{\begin{pmatrix}a&0\\ d&1\end{pmatrix}
        :a\in\mathrm{GL}_k(L),\ d\in L^{1\times k}\right\},
 \qquad p\left(\begin{smallmatrix}a&0\\d&1\end{smallmatrix}\right)=a.
\]
It acts on row covectors on the right by \(y b=ya+d\).
The original column group \(H\) is the stabilizer of \(y_0\)
in \(D\). In fact preserving the physical last coordinate
\(z=\ell w+\pi z'\) requires \(d=y_0-y_0a\).
Let \(D_0\) be the subgroup \(a\in C,\ d\in O^{1\times k}\).
Then \(p(H\cap D_0)=C_\ell\), the prescribed subgroup
fixing the residue covector. The orbit
\(S_0=y_0D_0=\pi^{-1}O^{1\times k}\setminus O^{1\times k}\)
is compact and open. Reduction modulo \(O^{1\times k}\)
gives all nonzero residue covectors by transitivity of
\(C\), and integral translations give each entire fiber.

Write \(N\) for the unipotent radical with the prescribed
singleton blocks. It is normalized by \(D\). The group
\(B'=N\rtimes D\) contains the period group \(B=N\rtimes H\).
Use the set of right \(U_L\)-cosets
\[
 \Omega'=\{n\in U(L):n^{-1}\gamma_c\in K B'\}/U_L
\]
in the first unipotent radical of Theorem~\ref{thm:finite281}.
The two first endpoints still give its integrality tests,
so it is finite by the denominator count in
Proposition~\ref{prop:marked291}. Conjugation by \(D_0\)
preserves \(U_L\) and \(\Omega'\), since it preserves the
row lattice and commutes with \(\gamma_c\).

\begin{proposition}[The contribution from an orbit of the enlarged set]\leavevmode
\label{prop:affine293}
Fix \(nU_L\in\Omega'\). Write \(n^{-1}\gamma_c=k_0b'\)
with \(k_0\in K\), \(b'\in B'\), and let \(\beta\in D\)
be the central block of \(b'\). Put \(A=p(\beta)^{-1}\).
The sum of the actual signed column terms over the
\(D_0\)-orbit of \(nU_L\), divided by the size of this
orbit, has norm at most \(\Xi(A)\). A coset failing the
original support condition contributes zero. The estimate
uses the original finite additive character.
\end{proposition}
\begin{proof}
The image of \(B'\cap K\) on the central block is \(D_0\).
Indeed an integral invertible block triangular matrix in
\(B'\) induces an integral automorphism there, and every
element of \(D_0\) embeds as an integral block. The
remaining unipotent factor is then integral as well.
It follows that \(k_0b'\in KB\) precisely when
\(\beta\in D_0 H\). Also \(|\det A|=1\), since \(n\) and
\(\gamma_c\) have determinant one, \(\det k_0\) is a unit,
and \(\det b'=\det p(\beta)\).

Push probability Haar measure on \(D_0\) forward by
\(s\mapsto sns^{-1}U_L\). This gives uniform counting
measure of total mass one on the finite orbit. The new
central block is \(s\beta s^{-1}\). Put \(y=y_0s\) and
\(y'=y\beta^{-1}\). The test \(s\beta s^{-1}\in D_0H\)
is equivalent to \(y'\in S_0\), since \(H\) fixes
\(y_0\) and \(S_0\) is stable under \(D_0\).
In that case choose \(s'\in D_0\) such that \(y_0s'=y'\).
The identity
\[
     (s\beta s^{-1})^{-1}
       =\bigl(s\beta^{-1}(s')^{-1}\bigr)\bigl(s's^{-1}\bigr)
\]
has first factor in \(H\) and second factor in \(D_0\).
Thus the true relative column translation in the unfolding
of Theorem~\ref{thm:sum287} can be taken to have matrix
\(p(s)A p(s')^{-1}\). Changing \(s'\) changes this matrix
on the right by \(C_\ell\). The unipotent terms and the
two character factors contribute only their actual scalar
phase. This follows by factoring off the central blocks
in \(B'\), leaving \(N\) in the period subgroup, and
using the right invariance of the row seed under \(K\).

Let \(\sigma\) be a local column component. For a
\(C_\ell\)-fixed vector \(u\) define a vector-valued
function \(f_u(y_0s)=\sigma(p(s))^{-1}u\) on \(S_0\)
and extend it by zero to \(Y=L^{1\times k}\).
Two possible \(s\)'s differ on the left by \(H\cap D_0\),
so this is well defined. Use additive Haar measure \(dy\)
with \(\operatorname{vol}(S_0)=1\). It is invariant under
\(D_0\), and its restriction to the orbit \(S_0\) is the
pushforward of Haar probability by \(s\mapsto y_0s\).
In particular \(\|f_u\|_2=\|u\|\). Push the phase-weighted
measure on \(D_0\) forward by the same map. Its density
\(\phi(y)\) with respect to \(dy\) has modulus at most
one, since the unweighted pushforward on \(S_0\) is \(dy\).
Use density zero elsewhere. The normalized orbital pairing
of two original vectors is therefore
\[
       \int_Y \phi(y)
        \langle \sigma(A) f_v(y\beta^{-1}),f_u(y)\rangle\,dy .
\]
Here the extension by zero includes the exact success test.
A change to the adjoint convention conjugates the pairing
and interchanges the two test vectors.

On \(L^2(Y)\) the weighted composition
\(S f(y)=\phi(y) f(y\beta^{-1})\) is a contraction, since
the affine map has determinant of absolute value one.
For \(h\in C\), let \(R_h f(y)=f(yh)\). Then \(f_u\)
is fixed by \(\sigma(h)\otimes R_h\), by using \(s h\)
as the representative of \(yh\). The displayed pairing
is a coefficient of \(\sigma(A)\otimes S\) between these
fixed vectors. Proposition~\ref{prop:tempered292} supplies
temperedness at this split place. Lemma~\ref{lem:tensor293}
gives the asserted bound, with the adjoint of the scalar
multiplier if the inner product is linear in its second
argument.
\end{proof}

\subsection{The first-plane flag at an inert place}\label{nav:IV-47}
Let \(D/C=E_w/F_v\) be an unramified quadratic field extension
at a finite place of full lattice level away from \(QT\).
Write \(O\) for the integers of \(D\), \(\pi\) for a uniformizer
of \(C\), and \(\kappa/\kappa_0\) for the residue extension,
where \(\#\kappa_0=P\). Fix a skew \(c\) with \(cO=\pi O\).
In the first-plane decomposition \(Du\oplus M\oplus Dv_1\)
of Theorem~\ref{thm:finite281}, let
\(p_1,\ldots,p_h\) be the remaining isotropic basis vectors
and \(p_{h+1}\) the terminal vector, with \(h=s-1\).
Its norm \(-T\) is a unit. Complete these vectors by integral
Witt partners and the reference column lattice
\(\Lambda\subset W_0\) of rank \(k\). Write \(P_i\) for the
residue span of \(p_1,\ldots,p_i\), with \(P_0=0\).

\begin{proposition}[The inert cosets at valuation one]\leavevmode
\label{prop:inertflag289}
Let \(M_k\) count the nonzero isotropic vectors of
\(\Lambda/\pi\Lambda\). In the unsigned local mass of
Theorem~\ref{thm:finite281} are \(P^{2h}(P-1)\)
cosets with zero residue column and \(P^{2h}P M_k\)
with nonzero residue column. The counts hold in every
residue characteristic.
\end{proposition}
\begin{proof}
Write the unipotent representative as \(n(w,z)\) as in
Theorem~\ref{thm:finite281}, set \(X_0=cw,\ Z_0=cz\), and
put \(g=n(w,z)^{-1}\gamma_c\). Its first vector is
\[
       gu=(1-\bar Z_0)u-X_0+cv_1,\qquad
       Z_0-\bar Z_0=h_M(X_0,X_0)/c .
\]
Here \(h_M\) is conjugate linear in its first variable.
Thus integrality requires \(X_0,Z_0\) integral and the
residue \(X\) of \(X_0\) isotropic. Conversely, for each such
\(X\) fix a lift \(X_0\). The displayed equation has
integral solutions forming a translate of \(O_C\). Indeed
there exists \(\lambda\in O\) of trace one, by surjectivity
of residue trace followed by division by a unit. For an
integral skew \(y\), the difference
\(\lambda y-\overline{\lambda y}\) equals \(y\).

Each right integral unipotent coset with residue \(X\) has a
representative with this fixed \(X_0\). To change \(w\) by an
integral vector \(w'\), multiply on the right by
\(n(w',-\lambda h_M(w',w'))\). Its cross term in \(cz\) is
integral since \(cw\) is integral. Once \(w\) is fixed,
two solutions for \(Z_0\) are equivalent precisely when their
difference belongs to \(cO^-\), where \(O^-\) is the lattice
of skew integers. As \(c/\pi\) is a skew unit,
\(O^-=(c/\pi)O_C\) and \(cO^-=\pi O_C\). There are therefore
exactly \(P\) scalar classes for each \(X\), before the other
flag conditions.

Suppose \(i\) is the first index with
\(h_M(X,p_i)\ne0\) in \(\kappa\). Choose the lift of \(X\)
with \(h_M(X_0,p_j)=0\) exactly for \(j<i\), using the Witt
coordinates. Then \(gp_j=p_j\) there, while
\(gp_i=p_i-h_M(X_0,p_i)u/c\) has a pole of order one.
If the flag admits an integral primitive representative, the
earlier integral columns \(gu,p_1,\ldots,p_{i-1}\) are already
independent modulo \(\pi\). Their exterior product is unchanged
by adding earlier columns to later ones. Clearing the pole by
such additions would force the residue of \(u\) into their
span. Indeed independence forces every coefficient used in
the additions to have valuation at least \(-1\), after which
one multiplies by \(c\) and reduces. The coefficient of \(gu\)
in this residue relation must be nonzero. Projection to \(M\)
therefore gives \(X\in P_{i-1}\). Every member of \(P_{i-1}\)
is orthogonal to \(p_i\), a contradiction.

It follows that \(X\in P_h\oplus(\Lambda/\pi\Lambda)\).
Write \(x\) for its column component. It is isotropic.
Choose \(X_0\) in the corresponding integral direct sum, with
zero column lift when \(x=0\). Now every \(gp_i=p_i\).
If \(x\ne0\), their reductions are independent together with
\(gu\) for all \(P\) scalar choices. If \(x=0\), the norm of
\(X_0\) is zero and \(Z_0\in O_C\). Independence holds
precisely for \(Z_0\not\equiv1\bmod\pi\), giving \(P-1\)
choices. Dependence prevents primitivity by the exterior
product test. In the independent case these integral columns
have the required pairings because \(g\) is an isometry.
Lemma~\ref{lem:flag215} extends the primitive decoration
to an integral isometry, which is precisely \(g\in K B(C)\).
There are \(P^{2h}\) choices in \(P_h\), proving both counts.
\end{proof}

\subsection{Two flag minors at arbitrary valuation}\label{nav:IV-48}
Fix a place outside \(QT\) with full row and column lattice
levels. Let \(O\) be the integers in the first factor at a
split place or those of \(E_w\) at an inert place. Use a
uniformizer \(\pi\) of \(F_v\) in either case. Suppose the
first-plane parameter \(c\) has valuation \(e\ge1\).
In order \(u,M,v_1\), write the three upper unipotent
blocks as \((l,H,z)\), a row, a column and a scalar.
Put \(Y=cl,\ X=cH,\ Z=cz\). In the inert case
\(l=-h_M(H,\cdot)\) and \(Y=h_M(X,\cdot)\).
For \(g=n^{-1}\gamma_c\), multiplication gives
\[
 g^{-1}=\begin{pmatrix}
 1&Y/c&Z/c\\0&I&X/c\\-c&-Y&1-Z
 \end{pmatrix},\qquad
 gu=(1-Z+YX/c)u-X+cv_1 .
\]
If \(g\in K B(F_v)\), its first vector and the first dual
covector (the last row of \(g^{-1}\)) are integral. Hence
\(X,Y,Z\) are integral. Write \(X_W,Y_W\) for the vector
and covector components on the reference column \(\Lambda\).
For a vector or covector let \(v\) be the minimum valuation
of its coordinates, with \(v(0)=+\infty\). Put
\(t=\min(e,v(X_W))\), \(u'=\min(e,v(Y_W))\) and
\(L_0=\max(0,e-t-u')\).

\begin{lemma}[Elementary divisors forced by the dual flag]\leavevmode
\label{lem:depth290}
Let \(\Lambda'\) be the column lattice of \(g^{-1}L\) for a
successful unipotent coset. Its elementary divisors relative
to \(\Lambda\) have at most one negative and at most one
positive exponent. Write these as \(-d_-\) and \(d_+\),
taking zero for an absent exponent. Both \(d_\pm\) lie between
\(L_0\) and \(e\). At an inert place \(u'=t\) and \(d_-=d_+\).
\end{lemma}
\begin{proof}
Let \(\lambda_1,\ldots,\lambda_s\) be the remaining dual
flag covectors on \(M\), the last being the terminal covector.
At the present place its norm parameter is a unit. Use unit
multiples and an adapted integral basis so that these are
coordinate functionals. Write \(q_1,\ldots,q_s\) for their
coordinate vectors, \(S\) for their span, and \(P_h\) for
the span of the \(h=s-1\) remaining primal isotropic vectors.
The middle lattice is \(P_h\oplus\Lambda\oplus S\).
Let \(X_S=\sum_j q_j\lambda_j X\). The transformed list of
dual rows starts with
\[
   \eta_0=(-c,-Y,1-Z),\qquad
   \eta_j=(0,\lambda_j,\lambda_j X/c),\quad 1\le j\le s .
\]
Every minor of each initial list is integral. Indeed the
rows can be made integral by adding earlier rows to later
ones by Lemma~\ref{lem:flag215}, and such additions preserve
the exterior product. This applies to the contragredient
list at a split place and to the Hermitian dual list at
an inert place.

For any column coordinate indexed by \(i\), take the minor
of \(\eta_0,\ldots,\eta_j\) on that coordinate,
\(q_1,\ldots,q_{j-1}\), and the last coordinate \(v_1\).
Its value up to sign is \(Y_{W,i}\lambda_j X/c\).
On \(q_1,\ldots,q_s,v_1\) the determinant for the whole list
is, up to sign, \(D=1-Z+YX_S/c\). Thus
\[
              \pi^{u'}X_S/c\ \text{is integral},
                     \qquad D\in O .
\]
For the first conclusion when \(u'=e\) use \(X_S\) integral.

An integral input \((a,m,\beta)\) of \(g^{-1}\) has middle
coordinate \(m+X\beta/c\). For its column projection one
requires zero last coordinate and zero \(S\)-coordinates of
the middle vector, then projects modulo \(P_h\). Write
\(m=m_P+w+m_S\). The conditions are
\[
   m_S=-X_S\beta/c\in S,\qquad
      Ym_P+Y_W w\equiv D\beta\pmod{cO},
      \qquad w\in\Lambda .
\]
Indeed the latter congruence is exactly the condition that
the last coordinate vanish for an integral \(a\).
The projected vector is \(w+X_W\beta/c\).
Its image modulo \(\Lambda\) lies in a single cyclic module
of length at most \(e\), giving at most one negative
exponent and \(d_-\le e\). When \(L_0>0\), choose
\(\beta=\pi^{u'}\), \(m_P=0\). Both conditions can be
satisfied since \(Y_W\Lambda=\pi^{u'}O\) and \(D\in O\).
The resulting vector has denominator of exact valuation
\(e-t-u'\), which proves \(d_-\ge L_0\).

At an inert place \(\Lambda'\) is self-dual by the unitary
column projection, so its exponents occur in opposite pairs.
Also \(u'=t\) by unimodularity of the column form. At a split
place apply the same calculation to the dual component with
endpoints in reverse order. The upper blocks of the reversed
\(n^{-\mathsf t}\) are \((-H^{\mathsf t},-l^{\mathsf t},lH-z)\),
and the parameter of the reversed \(\gamma_c^{-\mathsf t}\)
is \(-c\). Thus the two column depths interchange. The
projected dual lattice is \((\Lambda')^\vee\), because
\(g^{-1}L\) is a \(B(F_v)\)-translate and its two Levi
actions are contragredient. The negative exponents on that
side are the negatives of the positive exponents on this
side. The same bounds prove the assertion for \(d_+\).
\end{proof}

\subsection{The two unitriangular primes at every valuation}\label{nav:IV-49}
Take a place \(v\mid Q=g_1g_2\), with residue field
\(\kappa\) of size \(P\). Use the first split factor and
put \(O=\mathcal O_{F_v}\), \(C=\mathrm{GL}_k(O)\). The
column at this place is fixed by the inverse image in \(C\)
of the upper unitriangular subgroup. Let \(F_v(c)\) denote
the signed first-plane factor, with the identity mass
normalization of Theorem~\ref{thm:finite281}.

\begin{proposition}[A local estimate at the prescribed unitriangular level]\leavevmode
\label{prop:markedpowers292}
Take \(c=Q\xi\) and \(a=v(\xi)\ge0\). The factor at \(a=0\)
is the identity. For \(a\ge1\) one has
\[
  \|F_v(c)\| \le [4(k+1)^4]^a
                           P^{(k+4s+1)a+k}.
\]
\end{proposition}
\begin{proof}
Let \(K\) be the original row level. By
Lemma~\ref{lem:flags211} it reduces to
\(\bar N\rtimes U_k(\kappa)\), where \(\bar N\) is the
upper unipotent radical for the blocks with central block
of size \(k+1\). Define \(K_1\) to be the inverse image
of \(\bar N\) in the full row lattice stabilizer. Embed
\(C\) in the column block. It normalizes \(K_1\), the
period subgroup \(B=NH\), and the integral first radical
\(U_L\). It commutes with \(\gamma_c\) and preserves the
Bessel character on \(N\), since this character uses the
flag and the fixed terminal vector. The integral first
radical has only upper off-diagonal entries in the first
row and last column, so \(U_L\subset K_1\). The decomposition
of the reduction, lifting the \(U_k(\kappa)\) component
in \(H\cap K\), gives \(KB=K_1B\). Moreover the image
of \(B\cap K_1\) on the column is the principal congruence
subgroup \(C_1=\ker(C\to\mathrm{GL}_k(\kappa))\).
Indeed the middle block of an integral period element is
\(\operatorname{diag}(h,1)\), and membership in \(K_1\)
requires \(h\equiv1\). The remaining factor lies in the
integral unipotent on which the character is trivial.

The successful right \(U_L\)-cosets in
Theorem~\ref{thm:finite281} now form a finite \(C\)-set
under conjugation. For a representative \(n\) choose
\(n^{-1}\gamma_c=k_1 b_n\) with \(k_1\in K_1\), \(b_n\in B\).
Write \(A_n\) for the column component of \(b_n^{-1}\).
In the common-column unfolding of Theorem~\ref{thm:sum287}
the first period at \(n\) translates the column by the
identity, and the second, at \(\gamma_c^{-1}n=b_n^{-1}k_1^{-1}\),
translates it by \(A_n\). The two unipotent factors
contribute a scalar phase of modulus one. Using the adjoint
convention exchanges \(A_n\) with its inverse and gives
the same bounds below. The measure on cosets is counting
measure by Theorem~\ref{thm:finite281}. A new decomposition
changes \(A_n\) on the right by \(C_1\), which fixes each
original column vector. A change of representative \(n\)
by \(U_L\) can be absorbed into \(k_1\), and its phase is
trivial on \(U_L\).

For \(h\in C\) the conjugate decomposition gives \(hA_nh^{-1}\)
as a representative for the translated coset. The average
over each \(C\)-orbit of cosets is therefore a matrix
coefficient of the average in Lemma~\ref{lem:conjugation292},
with a scalar multiplier of modulus at most one. More
explicitly, push Haar probability forward by
\(h\mapsto hnh^{-1}U_L\). The image measure is uniform
on this orbit. In each fiber the relative translation
agrees on the test vectors with the chosen representative
because the ambiguity lies in \(C_1\). The actual phase
can be used as the multiplier. The argument takes place
in the whole local component before testing on the two
original vectors. Proposition~\ref{prop:tempered292}
applies to every such component.

Suppose \(a\ge1\). Put \(e=a+1\) and use \(X=cH,\ Y=cl,\
Z=cz\) from Lemma~\ref{lem:depth290}. By the endpoint
congruences in Proposition~\ref{prop:marked291}, all three
are in \(\pi O\) in their respective spaces. Let
\(t,u'\) be the valuations of \(X_W/\pi,Y_W/\pi\)
on the column and dual column, truncated at \(a\).
These depths are constant on each compact orbit, and
also under a change of representatives modulo \(c\).
Forgetting the congruence level gives a successful coset
for the full stabilizer of the same row lattice. The
terminal norm is a unit at \(v\), so the flag-minor proof
of Lemma~\ref{lem:depth290} applies there. The column
lattice is \(A_n O^k\), since \(k_1\) stabilizes the row
lattice and projecting \(b_n^{-1}\) to the column gives
its Levi action. The two truncated depths for that lemma
are \(t+1,u'+1\). Its two possible nonzero elementary
divisors satisfy
\[
 j_+\le d_\pm\le a+1,\qquad
 j=a-1-t-u',\quad j_+=\max(0,j).
\]
Lemma~\ref{lem:conjugation292} bounds the norm of an
orbit average by \(k^{2(a+1)}P^{-(k-2)j_+}\), using
the weaker exponent \(k-2\) for convenience.

Choose the lifts of \(X/\pi,Y/\pi\) modulo \(\pi^a\)
first, as in Proposition~\ref{prop:marked291}. At these
depths their column coordinates allow at most
\(P^{k(2a-t-u')}\) choices. The other coordinates and
the central scalar allow at most \(P^{(2(n_0-k)+1)a}\)
choices, where \(n_0=r-2=k+2s-1\). Multiplying each
orbit average by its orbit size accounts for every
coset exactly once. To bound their total use
\[
 k(2a-t-u')-(k-2)j_+
       =k(a+1+j)-(k-2)j_+ \le ka+k+2a .
\]
There are at most \((a+1)^2\le4^a\) depth pairs and
\(k^{2(a+1)}\le k^{4a}\). This proves the assertion for
positive \(a\). At \(a=0\) the unique coset is represented
by \(1\) and \(\gamma_c\in K\). Both period arguments have
the same column and trivial relative phase.
\end{proof}

\subsection{Elementary divisors in the terminal graph basis}\label{nav:IV-50}
At a place \(v\mid T\) use \(\Omega'\), \(\beta\) and
\(A=p(\beta)^{-1}\) from Proposition~\ref{prop:affine293}.
Put \(e=v(c)\ge1\) and \(h=s-1\). After taking out the
first hyperbolic pair, write the remaining graph lattice as
\[
  M=P_h\oplus W\oplus O t_0\oplus S_h .
\]
Here \(P_h\) consists of the \(h\) primal isotropic lines
and \(S_h\) of their opposite lines. In the central graph
coordinates \(W=\{(w,0):w\in O^k\}\) and \(t_0=(0,1)\).
The quotient of the central block by its terminal line
is identified with \(\Lambda=O^k\). For the first unipotent
blocks \((l,H,z)\) put \(X=cH,\ Y=cl,\ Z=cz\). Denote the
\(W\)-coordinate of \(X\) by \(X_W\) and the restriction
of \(Y\) to \(W\) by \(Y_W\). Write
\(t=\min(e,v(X_W))\), \(u'=\min(e,v(Y_W))\), taking the
minimum over the coordinates and \(v(0)=+\infty\).

\begin{lemma}[A pair of opposite elementary divisors]\leavevmode
\label{lem:graphdepth293}
For every \(nU_L\in\Omega'\), the elementary divisors of
\(A\) relative to \(\Lambda\) have the form
\((-d,0,\ldots,0,d)\) with
\[
                \max(0,e-t-u')\le d\le e .
\]
\end{lemma}
\begin{proof}
Write \(g=n^{-1}\gamma_c=k_0 b'\). Both first endpoints
are fixed by \(B'\), giving \(X,Y,Z\) integral as in
Proposition~\ref{prop:marked291}. In the order
\((u_1,M,v_1)\) the matrix is
\[
   g^{-1}=\begin{pmatrix}
       1&Y/c&Z/c\\ 0&I&X/c\\ -c&-Y&1-Z
   \end{pmatrix}.
\]
Let \(\lambda_1,\ldots,\lambda_h\) be the remaining
isotropic dual covectors in their flag order, with
coordinate sections \(q_j\) in \(S_h\). Write \(X_S\)
for the \(S_h\)-component of \(X\). The first \(j+1\)
dual rows of \(g^{-1}\) in this ordering are
\[
 (-c,-Y,1-Z),\qquad (0,\lambda_i,\lambda_i X/c)
                       \quad (1\le i\le j).
\]
Their exterior product is integral. Indeed the upper
block radical preserves each initial isotropic list by
unitriangular changes, its central block \(D\) fixes the
isotropic graded lines, and \(k_0^{-1}\) is integral.
Taking the minor on a \(W\)-coordinate, \(q_1,\ldots,q_{j-1}\)
and \(v_1\) gives \(\pm Y_{W,i}\lambda_j X/c\) at
coordinate \(i\). Taking the full minor on
\(q_1,\ldots,q_h,v_1\) gives \(\pm D_1\), where
\(D_1=1-Z+Y X_S/c\). Therefore
\(\pi^{u'}X_S/c\) is integral and \(D_1\in O\). If
\(u'=e\) the first assertion follows directly from
integrality of \(X_S\).

Project \(g^{-1}\) of the row lattice to the column as
follows. Intersect with the rational subspace where the
last coordinate and the \(S_h\)-coordinates vanish, then
quotient by the span of \(u_1,P_h,t_0\). At the reference
lattice this gives \(\Lambda\). The flag group \(B'\)
preserves both subspaces and acts on the quotient by its
\(p\)-component, since \(D\) fixes \(t_0\) and \(N\) is
trivial on the quotient. As \(g^{-1}=b'^{-1}k_0^{-1}\),
the output lattice is \(A\Lambda\).

For an integral input \((a_0,m,\alpha)\), write
\(m=m_P+w+bt_0+m_S\). The intersection conditions read
\[
 m_S=-X_S\alpha/c,\qquad
 Y(m_P+bt_0)+Y_Ww\equiv D_1\alpha\pmod {cO}.
\]
The congruence expresses existence of an integral \(a_0\)
making the last coordinate zero. The projected vector is
\(w+X_W\alpha/c\). Modulo \(\Lambda\) all projected vectors
therefore lie in one cyclic module of length at most \(e\).
There is at most one negative exponent, of absolute value
at most \(e\). If \(e-t-u'>0\), take \(\alpha=\pi^{u'}\)
and \(m_P=b=0\). The first condition is integral, and
\(Y_W\Lambda=\pi^{u'}O\) permits the stated congruence.
The projected vector has exact denominator length
\(e-t-u'\).

In the split dual lattice reverse the endpoints. The
contragredient of the central block \(D\) preserves the
annihilator of \(t_0\), whose reference lattice has basis
dual to \(W\). On that annihilator it acts by \(p^{-{\mathsf t}}\).
Impose this extra vanishing along with the isotropic
intersection conditions and project past the first primal
lines on the reversed side. The projected dual lattice is
\(A^{-{\mathsf t}}\Lambda^\vee\). Indeed on the dual
row the transformation of \(g^{-1}=b'^{-1}k_0^{-1}\)
is \(g^{\mathsf t}=b'^{\mathsf t}k_0^{\mathsf t}\),
and \(b'^{\mathsf t}\) induces \(A^{-{\mathsf t}}\)
on the dual subquotient.
For clarity, the dual transformation of \(g^{-1}\) is
\(g^{\mathsf t}\). In reversed block order it has middle
block \(I\) and a single column above the last diagonal
entry, equal up to sign to \(Y^{\mathsf t}/c\). Its
projection on the \(W^\vee\)-coordinates of an integral
input is consequently an integral vector plus an integral
multiple of this one projected column. The imposed
vanishing conditions only restrict the allowable inputs.
Thus the dual projection modulo \(\Lambda^\vee\) is
likewise contained in one cyclic module of length at
most \(e\). There is at most one positive exponent of
\(A\). Finally \(|\det A|=1\) by
Proposition~\ref{prop:affine293}, so the two absolute
values coincide. This proves the claim.
\end{proof}

\subsection{The two split valuations and the common scalar}\label{nav:IV-51}
At a split finite place use the first component, with
uniformizer \(\pi\), residue cardinality \(P\), and
valuation \(v\). Write \(c_+,c_-\) for the two components
of \(c=Q\xi\), \(e_\pm=v(c_\pm)\), and \(L=v(\lambda_+)\).
Put \(\epsilon=1\) above \(Q\), zero elsewhere. Write
\(a_\pm=e_\pm-\epsilon\ge0\), \(h=s-1\), and
\(n=k+2s-1\). At \(T\) use the graph basis and the enlarged
support \(KB'\) of Proposition~\ref{prop:affine293}. At
other split places use the original support \(KB\). For
\(g=n_1^{-1}w_{c,\lambda}\) in that support let \(A\) be
the inverse column action of the period factor, or the
inverse quotient action of its central block at \(T\).
Put \(A_0=\lambda_+ A\).

\begin{lemma}[Elementary divisors after a central translation]\leavevmode
\label{lem:scalardepth294}
Write the unipotent blocks of \(n_1\) as \((l,H,z)\).
Set \(X=c_+ H,\ Y=c_- l,\ Z=c_- z\).
On the column and dual column use the truncated depths
\(t=\min(a_+,v(X_W/\pi^\epsilon))\) and
\(u'=\min(a_-,v(Y_W/\pi^\epsilon))\). These are nonnegative
on the indicated support. The only possible nonzero
exponents of \(A_0\) are \(-d_-\), \(d_+\), with
\[
  0\le d_-\le e_+,\quad 0\le d_+\le e_- ,\qquad
  d_+-d_-=e_--e_+-(2h+1)L .
\]
If \(S=t+u'+2\epsilon\), then away from \(T\) one has
\(d_-+d_+\ge e_++e_--2S+|L|\).
At \(T\) the same bound holds for \(L<0\), and for
\(L\ge0\) it holds with \(-L\) in place of \(|L|\).
\end{lemma}
\begin{proof}
In block order \(u,M,v_1\) the inverse is
\[
  g^{-1}=
  \begin{pmatrix}
    0&0&1/c_+\\
    0&\lambda_+^{-1} I&\lambda_+^{-1}X/c_+\\
    c_-&Y&Z
  \end{pmatrix}.
\]
The middle part of \(gu\) is \(-X\), and its first entry
is \(c_+(lH-z)\). Thus \(X,Y\in\pi^\epsilon O\) in
their blocks. Both \(Z\) and \(c_+(lH-z)\) are integral
and are \(1\bmod\pi\) above \(Q\). These are the first
primal and dual endpoint tests, which use the actual
level above \(Q\).

First use the primal projection of
Lemma~\ref{lem:graphdepth293} by quotienting the terminal
line. At other places use the integral splitting of the
unit terminal line with its column. Write \(q_1,\ldots,q_h\)
for the integral sections of the isotropic dual covectors
\(\eta_i\) on \(M\), and \(X_S\) for this part of \(X\).
The first dual row is \((c_-,Y,Z)\). Each subsequent
row in the isotropic list is
\(\lambda_+^{-1}(0,\eta_i,\eta_i X/c_+)\).
Every initial wedge of these rows is integral, by the
flag and the integral row factor. For the first \(j+1\)
rows, take columns \(q_1,\ldots,q_{j-1},v_1\) and one
additional column. Taking a \(W\)-coordinate or the
initial \(u\)-coordinate gives, up to sign,
\(\lambda_+^{-j}a\,\eta_j X/c_+\), respectively for
\(a=Y_{W,i}\) or \(a=c_-\). The whole list on
\(q_1,\ldots,q_h,v_1\) gives, up to sign,
\(\lambda_+^{-h}D_1\), with \(D_1=Z-YX_S/c_+\).
Set \(u_c=\min(e_-,v(Y_W))=u'+\epsilon\). Since the
tested \(a\)'s generate \(\pi^{u_c}O\), the valuations of
\(\pi^{u_c}\eta_j X/c_+\) and \(D_1\) are at least
\(jL\) and \(hL\), respectively.

Multiply the output column lattice by \(\lambda_+\).
It becomes \(A_0\Lambda\). A projected vector has form
\(w+X_W\alpha/c_+\) for integral \(w,\alpha\), so its
denominators give at most one negative exponent of length
at most \(e_+\). Take \(\alpha=\pi^{u_c+h\max(0,-L)}\).
The required opposite coordinates \(-X_S\alpha/c_+\)
are integral. With the two discarded primal components
zero, the last equation is
\(Y_W w\equiv-D_1\alpha\bmod c_-O\). It is soluble
because its right side is in \(\pi^{u_c}O\). This also
covers \(u_c=e_-\). The middle vector now gives
\[
              d_-\ge e_+-S-h\max(0,-L),
\]
since \(t+\epsilon\) is the column valuation of \(X_W\)
truncated at \(e_+\). A nonpositive lower bound needs
no choice of vector.

Reverse the endpoints in the dual split component. Its
two \(c\)'s are interchanged, the central scalar is
\(\lambda_+^{-1}\), and the new vector and row are
\(-Y^{\mathsf t},-X^{\mathsf t}\). This follows by taking
the inverse transpose of the exchange and of the original
upper unipotent. Away from \(T\) quotient the dual terminal
line, again a direct integral summand. At \(T\) instead
impose the additional equation of annihilating the primal
terminal vector \(t_0\), as in Lemma~\ref{lem:graphdepth293}.
The coordinate section for this last equation is
\(t_0^\vee\) in the graph dual basis, so the \(W^\vee\)
coordinates are precisely the duals of the preceding
\(W\)-coordinates. The augmented dual list has one more
row. Its initial wedges are integral also for \(KB'\),
since the group \(D\) fixes \(t_0\) and the upper radical may
add earlier isotropic vectors to it. The same minors
and congruence argument, with \(h'=h+1\) at \(T\) and
\(h'=h\) elsewhere, give
\[
      d_+\le e_-,\qquad
      d_+\ge e_--S-h'\max(0,L).
\]
The dual output after scaling is \(A_0^{-{\mathsf t}}\Lambda^\vee\).
Its negative denominators are cyclic as in the primal
projection, so \(A_0\) has at most one positive exponent.

The exchange has determinant \(-c_+/c_-\) on the first
plane and \(\lambda_+^n\) on \(M\). The row integral
factor has unit determinant and the unipotents have
determinant one. Hence
\(v(\det A_0)=e_--e_+-(n-k)L\), where \(n-k=2h+1\).
For \(L<0\), add this determinant valuation to twice
the lower bound for \(d_-\). For \(L\ge0\), subtract it
from twice the lower bound for \(d_+\). This gives the
stated lower bounds for the sum.
\end{proof}

\subsection{The support of the finite test on the first plane}\label{nav:VI-48}
Use \(u,v,M,J,\gamma=\gamma_c\) of Proposition~\ref{prop:plane281}
at a finite place of \(F\). Write \(O\) for its integers,
\(\pi\) for a uniformizer, \(O_E\) for the local quadratic
integer algebra, and \(L=(O_Eu+O_Ev)\perp L_M\)
for the lattice with its first hyperbolic plane split off.
At \(T\) use the actual
graph lattice in \(M\). Put \(K=K_v^0\), \(J_K=J(F_v)\cap K\)
and
\[
 X=(B(F_v)K\cap\gamma B(F_v)K)/K,\qquad
 m_v(\gamma)=\sum_{tK\in J(F_v)\backslash X}
  \frac{\operatorname{vol}(J_K)}
       {\operatorname{vol}(J(F_v)\cap tKt^{-1})}.
\]
This quotient of volumes is independent of the normalization
on \(J(F_v)\).

\begin{theorem}[A first-vector test for finite cells]\leavevmode
\label{thm:finite281}
Let \(U\) be the unipotent radical of the stabilizer of \(Eu\)
and \(U_L=U(F_v)\cap\operatorname{Stab}(L)\).
With \(U_L\) of volume one, \(m_v(\gamma)\) is the local
unsigned finite mass divided by the identity mass, in
quotient measures. It equals the number of right \(U_L\)-cosets
in
\[
                 \Omega=\{n\in U(F_v):n^{-1}\gamma\in K B(F_v)\}.
\]
Write \(a=1\) above \(Q\) and \(a=0\) elsewhere. If
\(cO_E=\pi^a O_E\), there is one such coset. Both decorations
then have the same finite column coordinate on every cell
at this place.
\end{theorem}
\begin{proof}
In order \(u,M,v\) write the elements of \(U\) as
\[
 n(w,z)=
 \begin{pmatrix}1&-w^*&z\\0&1&w\\0&0&1\end{pmatrix},
 \qquad z+\bar z=-h(w,w),\qquad w^*=h(w,\cdot).
\]
This follows by imposing preservation of the hyperbolic
pairing and of the form on \(M\). Inversion sends \((w,z)\)
to \((-w,\bar z)\). Integrality is equivalent to
\(w\in L_M,\ z\in O_E\), since \(L_M\) is self-dual.
The first-plane parabolic decomposition gives \(B=J\ltimes U\).
At full lattice level, a product belongs to \(K\) precisely
when both factors do, by its diagonal block on \(M\). This
applies also to the graph at \(T\). At a place above \(Q\)
the integral elements of \(U\) have, on the first split
factor, only upper off-diagonal entries in the first row and
the last column. The prohibited entries of
Proposition~\ref{prop:index90} are inside the middle block.
Replacing the original opposite first vector by \(v\) adds a
multiple of \(u\) integral at these odd places and preserves
these conditions. Thus \(U_L\subset K\) also there.
The middle diagonal block again gives
\(B(F_v)\cap K=J_K\ltimes U_L\). Conjugation by \(J(F_v)\)
preserves Haar measure on \(U(F_v)\). It acts unitarily on
the \(E_v\)-space of \(w\), with \(F_v\)-determinant of absolute
value one, and fixes the center.

Every element of \(X\) has the form \(jnK\), \(j\in J(F_v)\),
\(n\in\Omega\), since \(j\) commutes with \(\gamma\).
If \(jnK=n'K\), the middle diagonal block of
\((n')^{-1}jn\) gives \(j\in J_K\).
The remaining condition is \(jnj^{-1}U_L=n'U_L\).
The stabilizer of \(nU_L\) for this \(J_K\)-action is precisely
\(J(F_v)\cap nKn^{-1}\). Each summand in \(m_v\) is consequently
the size of its orbit of right \(U_L\)-cosets. Multiplying
its unnormalized stabilizer reciprocal by \(\operatorname{vol}K\)
gives its local mass in \(J(F_v)\backslash G(F_v)\).
For the identity orbit that mass is
\(\operatorname{vol}K/\operatorname{vol}(B(F_v)\cap K)\).
The decomposition of this last compact group and the
normalization of \(U_L\) give the stated ratio.

Every \(b\in B(F_v)\) fixes \(u\). Every \(k\in K\) sends \(u\)
to an integral vector and, above \(Q\), fixes it modulo \(\pi\)
in both factors. Indeed upper unitriangular reduction fixes
the first column and, in the dual action, the last covector.
Hence membership in \(\Omega\) implies
\[
 n(w,z)^{-1}\gamma u-u=c\bar z\,u-cw+cv\ \in\ \pi^a L.
\]
This bounds both \(w,z\) in fractional lattices, so the
coset count is finite. At minimal valuation it forces
\(n\in U_L\). Conversely \(U_L\subset\Omega\), since
\(\gamma\in K\). Every representative is then \(jK\) with
\(\gamma^{-1}jK=jK\), proving the last claim.
\end{proof}

\subsection{The two column lattices of a first-plane orbit}\label{nav:VI-49}
Work at the split place of Theorem~\ref{thm:split281} in its
coordinates. Write \(\Lambda=O^k\) for the definite column
lattice on the first split factor. Two lattices are neighbors
here if their intersection has index \(b\) in both. Denote
by \(x,y\) the \(W,W^\vee\) components of the residues
\(X,Y\) of that theorem.

\begin{proposition}[The relative positions of the two columns]\leavevmode
\label{prop:lattice282}
In the finite coset count of Theorem~\ref{thm:split281} the
first column lattice is \(\Lambda\). The second is a neighbor
exactly when \(i=j=\infty\), \(x\ne0\), and \(y\ne0\).
Every neighbor occurs \(u_b=b^{2s-2}(b-1)^2\) times. All
remaining terms give \(\Lambda\), with multiplicity
\[
       a_b=b^{2s-2}(b-1)
          \left[(2+2s(b-1))b^k-1\right].
\]
\end{proposition}
\begin{proof}
For a lattice of the form \(b' L\), \(b'\in B(F_v)\), its
column lattice is the image in \(W_0\) of its intersection
with the perpendicular of the isotropic flag space and of the
terminal vector, under projection modulo the flag space.
Indeed the unipotent factor of \(b'\) acts trivially on the
middle quotient, and the Levi factor acts on the column by
its \(H\)-block. The same description works on the first split
factor by imposing annihilation by the opposite flag covectors.
For \(n(l,H,z)L\) the column is consequently \(\Lambda\).
For \(g^{-1}L=\gamma_c^{-1}n(l,H,z)L\) use the inverse matrix
in Theorem~\ref{thm:split281}. An integral input
\((t,m,\beta)\) maps to a vector with middle component
\(m+X_0\beta/c\) and last component
\(-ct-Y_0m+(1-Z_0)\beta\). Thus the latter vanishes for an
integral \(t\) exactly when
\(Y_0m\equiv(1-Z_0)\beta\bmod c\).
The middle component must also be annihilated by
\(\lambda_1,\ldots,\lambda_s\), including the terminal covector.
Project its remaining coordinates to \(W_0\).

If \(i=j=\infty\) these conditions put \(m\) in the direct
sum of the first \(s-1\) flag coordinates and \(O^k\).
The second column is precisely
\[
 \Lambda'=\{w+x_0\beta/c:\ w\in\Lambda,\ \beta\in O,\
             y_0w\equiv(1-Z_0)\beta\bmod c\},
\]
where \(x_0,y_0\) are the component lifts. When \(y=0\)
the successful cosets have \(1-Z_0\) a unit, forcing
\(\beta\in cO\). This gives \(\Lambda'=\Lambda\). When
\(x=0\) the same unit permits any \(w\), again giving
equality. For \(x,y\ne0\), the intersection with \(\Lambda\)
is the inverse image of \(\ker y\) modulo \(c\). Indeed
integrality requires \(\beta\in cO\) and \(yx=0\); conversely
one may take \(\beta=0\) on that hyperplane. The image
modulo \(\Lambda\) is the full line of \(x_0/c\), since
\(y\ne0\) permits every \(\beta\). This proves the neighbor
condition.

If \(i<\infty\), \(X_0\) has only entries in \(P_{i-1}\)
and \(1-Z_0\) is a unit. For any desired \(w\in\Lambda\)
take \(m=w\) and choose \(\beta\) by the congruence.
There are no fractional column coordinates. If \(j<\infty\),
the equation for \(\lambda_j\) forces \(\beta\in cO\).
Here \(Y_0\) is supported in \(D_{j-1}\), so every integral
column is obtained with \(m=w,\ \beta=t=0\).

There are \(b^{2s-2}b(b^k-1)(b^{k-1}-1)\) cosets in
the neighbor case, by the count of nonzero orthogonal pairs
and the unrestricted \(Z\). The integral group
\(\mathrm{GL}_k(O)\subset J_K\) acts on the cosets by
conjugation and on the resulting column lattices by its
usual action. It is transitive on neighbors by elementary
divisors, whose exponents are \(1,-1\) and \(k-2\) zeros.
There are \(d_b=(b^k-1)/(b-1)\) choices for the intersection
as an index-\(b\) sublattice of \(\Lambda\). Each has
\(d_b-1\) index-\(b\) superlattices besides \(\Lambda\),
all producing distinct neighbors with that intersection.
Dividing the count by \(d_b(d_b-1)\) gives \(u_b\)
for \(k\ge2\), as in the original ranks. Subtracting the
neighbor count from Theorem~\ref{thm:split281} gives \(a_b\).
\end{proof}

\subsection{The last additive coordinate at a split auxiliary prime}\label{sec:signedfibers331}
Fix a point \(x\) of an interval in Section~\ref{sec:missing}, its field
\(F=E_x^+\), and an admissible tuple \(\xi\). Let \(v\in Q(\xi)\) be
one of the places where the twist is nontrivial on units. The conditions
on the tuple say that \(v\) splits in \(E_x\), that its ideal norm is a
prime \(p\), and that \(p\notin v^2\). The residue degree and the
ramification index at \(p\) are therefore both one. In particular
\(F_v=\QQ_p\) over \(\QQ_p\), and we use \(O=\ZZ_p\) and uniformizer
\(p\). The excluded set contains 2. The character on units kills
\(1+pO\), so its inverse gives a nontrivial character \(\chi\) of
\(\mathbb F_p^\times\). These conditions are recorded in
\srcloc{Formalism/LocalToGlobal/Divisibility/Tuples}{63--67,110--128}.

Take \(s=m-2\ge1\), \(a_j=k_0+j\), \(n_j=a_j+jt\),
\(r=a_s+1\) and \(d=n_{s+1}\), with the marked planes and ordered
bases of Theorem~\ref{thm:splitpath322}. Here \(A_j\) is the marked
\(a_j\)-plane of the residue core at stage \(j\), and \(U=A_{s+1}\).
Fix a row \(l\in O^{1\times d}\) with \(\bar l|_U\ne0\), and put
\(U_0=\ker(\bar l|_U)\) and \(C_s=O/p^sO\). For \(\nu\in C_s^\times\)
choose a lift \(\widetilde\nu\) and use the frequency ball
\[
       {\cal B}_\nu=p^{-s-1}\widetilde\nu l+O^{1\times d}.
\]
The Fourier convention is \(\psi(-zv')\) for column \(v'\), with
\(\psi\) of conductor \(O\). The effective row \(z\) includes the scalar
in the Fourier kernel. In the top test it is the last matrix row,
multiplied by that scalar. In particular identifying a global Bessel
datum specifies a scalar but does not specify the last matrix row.
The two arguments in the top test appear in
\srcloc{Development/AbelianVarietyTheory/Definitions/TwistedInnerProducts/StandardOrderExtendedBasisApplyLowBlock}{796,919--920}.

Write \(z_0=0\), and for \(1\le j\le s\) let \(z_j\) be the marked
core basis vector in the representative \(h_j\) of
Theorem~\ref{thm:splitpath322}. Write \(\mathcal P_l\) for the following
set of integral core matrices \(X_0,\ldots,X_s\) and integral top
matrices \(B\).
For \(0\le j<s\) require \(X_jz_j-z_{j+1}\in p^jO^{n_{j+1}}\).
The \(a_j\) marked columns of \(X_j\) followed by
\(v_j=p^{-j}(X_jz_j-z_{j+1})\) must reduce to a basis of \(A_{j+1}\).
At stage \(s\) the marked columns of \(X_s\) must reduce to a basis of
\(U_0\). The marked columns of \(B\) must reduce to an isomorphism
of its two marked \(r\)-planes. Unmarked columns are arbitrary
integral columns.

For \(j<s\) let \(u_j\) be the residue determinant of the indicated
basis in its original orientation, including the fixed ordering sign.
Let \(u_+\) be the marked determinant of \(B\).
Appending a variable last column \(v'\in U\) to the marked columns
of \(X_s\) gives a determinant \(L(v')=u_s\bar l(v')\) for a unique
unit \(u_s\in\mathbb F_p^\times\). Define
\[
    D=u_+\prod_{j=0}^s u_j\in\mathbb F_p^\times,\qquad
    Y=p^{-1}lX_sz_s\bmod p^s\in C_s .
\]
The second coordinate exists because \(\bar X_sz_s\in U_0\).
Use product additive Haar with integral lattices of volume one.

\begin{theorem}[The signed integral on every core position]\leavevmode
\label{thm:signedfibers331}
On the ball \({\cal B}_\nu\) the local path, with the compact
weights and initial endpoint of Theorem~\ref{thm:splitmass323},
has support precisely \(\mathcal P_l\) in its position variables.
Each fiber of \((D,Y)\) in this support has the same measure
\({\cal V}\) as the position set in that theorem. On the fiber
\((D,Y)=(b,y)\) the path has the constant value
\[
    c_\chi\,\chi(\bar\nu)^{-1}\chi(b)\,\psi(p^{-s}\widetilde\nu\widetilde y),
    \qquad
    c_\chi=\eta_\chi\,p^{\,r-1-(s+1)d}G(\chi)^h,\qquad
    h=1+s(s+1)/2 ,
\]
where \(\widetilde y\) is any lift,
\(G(\chi)=\sum_{a\ne0}\chi(a)\psi(-p^{-1}\widetilde a)\)
on \(\mathbb F_p^\times\), and \(\eta_\chi\) is the actual product
of unramified splitting scalars in Theorem~\ref{thm:splitpath322}.
As the positions vary, the full frequency support consists of the
rows \(z\) for which \(p^{s+1}z\) is integral and its reduction
is nonzero on \(U\).

For any integrable weight on the positions and this frequency ball,
let \(H_\nu(b,y)\) be its mean on the indicated fiber times the
whole ball, with respect to Haar probability on that product. The
weighted integral is
\begin{equation}\label{eq:signedfibers331}
 {\cal V} c_\chi\,\chi(\bar\nu)^{-1}
  \sum_{b\in\mathbb F_p^\times}\sum_{y\in C_s}
  H_\nu(b,y)\chi(b)\psi(p^{-s}\widetilde\nu\widetilde y).
\end{equation}
It vanishes whenever \(H_\nu(b,y)\) is invariant under adding
\(p^{s-1}C_s\) to \(y\) for each \(b\).
\end{theorem}
\begin{proof}
Theorem~\ref{thm:ramifiedcoset324} leaves a single chain of
intermediate double cosets. In its representatives, integrality at
an interior edge requires integrality of \(X_j\) and \(v_j\).
On the marked residue space the first \(a_j+1\) columns reduce to
the indicated core columns, all with zero lower entries. They must
span the whole marked target core, since the remaining \(s-j\)
columns must already span the flag quotient of that dimension.
The conditions are sufficient by the triangular lower block
calculation of Theorem~\ref{thm:ramifiedflag324}. Each of its
\(s-j\) free diagonal residues contributes a factor \(G(\chi)\);
its marked core determinant contributes \(\chi(u_j)\).

At the last edge use the variable change \(v'=-X_sz_s+p^s v\).
The core must again be integral. The integral in \(v\) is constant
on residue classes except for its Fourier character. Translation
by \(pO^d\) forces \(p^{s+1}z\) to be integral. Integration on
the marked plane then forces the residual frequency to annihilate
the image of the first \(r-1\) marked columns. These columns must be
independent and lie in \(U\), by the determinant test. If the
residual frequency vanishes on \(U\), the last residue sum is a
multiple of \(\sum_a\chi(a)=0\). On \({\cal B}_\nu\) the residual
row is \(\bar\nu\bar l\). Hence the image is exactly \(U_0\).
In the last determinant coordinate \(L(v)\) the residue sum now
has parameter \(\bar\nu/u_s\). It contributes
\(\chi(u_s)\chi(\bar\nu)^{-1}G(\chi)\) by changing variables in
the sum. This also proves sufficiency, since \(|G(\chi)|=\sqrt p\)
by Theorem~\ref{thm:lastcolumn321}. Every row satisfying the stated
frequency condition is covered by taking \(l=p^{s+1}z\) and \(\nu=1\).

The translation of \(v'\) has the remaining phase
\(\psi(zX_sz_s)=\psi(p^{-s}\widetilde\nu\widetilde Y)\).
Adding an integral row to the center has no effect. Changing a lift
of \(\nu\) by \(p^s\) times an integer can change the frequency ball,
but gives the same phase on these positions. The top test supplies
\(\chi(u_+)\). The interior compact integrals cancel their reciprocal
weights by Theorem~\ref{thm:splitcompact323}. At the last orbit all
contributing \(v'\) reduce to \(-\bar X_sz_s\), a nonzero vector in
\(U_0\), since \(s\ge1\) and \(z_s\) is a marked basis vector. Thus
the same primitive orbit and its entire stabilizer weight occur.
Combining the Jacobians and splitting factors of
Theorem~\ref{thm:splitpath322} gives \(c_\chi\). The preceding
calculation accounts for every factor whose value changes with
the positions.

The units acting on the first marked column of \(B\) permute the
values of \(D\) transitively and preserve \(Y\) and measure.
Choose an integral column \(w\) with \(lw=1\), and an integral row
\(e\) with \(ez_s=1\). For \(u\in O\), translation
\(X_s\mapsto X_s+puwe\) fixes all the residue determinants and
adds \(u\) to \(Y\). It preserves measure and all the other positions.
Hence all \((p-1)p^s\) fibers have equal measure. The fiber
\((1,0)\) is the one counted in Theorem~\ref{thm:splitmass323}
for a marked coordinate row. For general \(l\), counting ordered
bases of \(U_0\) gives the same number, since all hyperplanes of
\(U\) have dimension \(r-1\). The lifting condition \(Y=0\) selects
one of \(p^s\) equal fibers by the displayed translation. The other
factors in the count are independent of \(l\). This proves the
assertion about \({\cal V}\). The ball has measure one, so summing
the integral over the fibers gives \eqref{eq:signedfibers331}.
The sum of the additive phase on \(p^{s-1}C_s\) is zero, as \(\nu\)
is a unit and the character there has order \(p\).
\end{proof}

For example a weight supported on a single \(b\), constant in \(y\),
gives zero on every primitive frequency. The weight
\(\overline{\chi(b)\psi(p^{-s}\widetilde\nu\widetilde y)}\)
on all the fibers gives \((p-1)p^s{\cal V}c_\chi\chi(\bar\nu)^{-1}\).
Both calculations use the full position domain. For a discrete
absolutely summable class expansion, sum its original weights in
each fiber and omit the Haar factor \({\cal V}\) in
\eqref{eq:signedfibers331}. This follows by partitioning the
original sum. At several places the same partition may be taken
jointly, including when the class weights have global dependencies.

\begin{theorem}[Elementary divisors at the residue characteristic]\leavevmode
\label{thm:localsmith331}
Use one common array of weights on \(\mathbb F_p^\times\times C_s\)
in \eqref{eq:signedfibers331}. Take all primitive \(\nu\) and any
set of distinct nontrivial \(\chi\). Remove from each output its
factor \({\cal V}\eta_\chi p^{r-1-(s+1)d}\). Work at any place
above \(p\) of a cyclotomic splitting field, with \(v_p(p)=1\).
For \(\chi=\omega^{-a}\) in this completion, where \(\omega\) is
the Teichm\"uller character and \(1\le a\le p-2\), the nonzero
elementary divisors contributed by \(\chi\) have valuations
\[
 \frac{ha}{p-1}+
 \frac{j+(p-1)\displaystyle\sum_{c=1}^{s-1}
      p^{c-1}\lfloor j/\varphi(p^c)\rfloor}{\varphi(p^s)},
 \qquad 0\le j<\varphi(p^s).
\]
The full multiset is the union over the chosen \(\chi\).
\end{theorem}
\begin{proof}
Let \(\zeta=\psi(p^{-s})\). It is a primitive \(p^s\)-th root.
On the \(b\)-coordinate the character matrix splits integrally
at \(p\), by character orthogonality and the invertibility of
\(p-1\). Each retained character gives one copy of evaluation of
\(1,T,\ldots,T^{p^s-1}\) at the primitive roots
\(\zeta^{u_0},\ldots,\zeta^{u_{n-1}}\), where \(n=\varphi(p^s)\)
and \(u_0<\cdots<u_{n-1}\) are the positive units below \(p^s\).
There is also the scalar \(G(\chi)^h\) and the row unit
\(\chi(\bar\nu)^{-1}\). Divide the powers \(T^j\) for \(j\ge n\)
by the monic \(\Phi_{p^s}\) to remove redundant columns by integral
column operations.

For \(0\le j<n\) use the monic basis
\(P_j(T)=\prod_{i<j}(T-\zeta^{u_i})\), an integral change of
basis with determinant one. Its evaluation matrix is lower triangular.
If \(i\ne k\), the valuation of \(\zeta^{u_i}-\zeta^{u_k}\) is
\(p^{v_p(u_i-u_k)}/n\). Indeed the quotient of the roots has
order \(p^{s-v_p(u_i-u_k)}\), and
\(\Phi_{p^b}(1)=p\) gives valuation \(1/\varphi(p^b)\) for the
difference from one of a primitive \(p^b\)-th root. All its
conjugate differences have the same valuation, since exponents
prime to \(p\) change them by units.

For each \(1\le c<s\) the list of the \(u_i\) modulo \(p^c\)
cycles through all \(\varphi(p^c)\) unit residues in increasing
order in each successive block. Among the first \(j\) terms each
residue occurs at least \(\lfloor j/\varphi(p^c)\rfloor\) times,
and the residue of \(u_j\) occurs exactly that often. Expanding
\(p^b=1+(p-1)\sum_{c=1}^{b}p^{c-1}\) in the sum of the \(j\)
pairwise valuations shows that \(v_p(P_j(\zeta^{u_j}))\) is
the second summand in the formula. All entries below it have
at least this valuation. After factoring the diagonal entry from
each column the remaining lower triangular matrix is integral with
diagonal one. Its inverse is integral. The diagonal valuations
increase with \(j\), which gives the claimed elementary divisors.
Finally Lemma~\ref{lem:gauss89} gives \(v_p(G(\chi))=a/(p-1)\).
Changing the primitive additive character in that lemma multiplies
the sum by a root of unity. This accounts for the first summand.
\end{proof}


\subsection{Indices of the successive lattices in a split Bessel flag}\label{sec:exchangeindex336}
Use the split coordinates and character of Lemma~\ref{lem:exchangeroot336}.
In the chosen basis let \(\Lambda\) be the full coordinate lattice,
with \(A\)-part \(O^a\). For \(d\ge1\) let \(K(d)\) be the group of
automorphisms congruent to one modulo \(\pi^d\) on \(\Lambda\), and
let \(K(0)=\GL(\Lambda)\). Let \(c\) be characterized by
\[
       \{t\in k:\psi_i(tO)=1\}=\pi^cO,\qquad D=\max(0,d-c).
\]
The simple characters \(\psi_i\) are all equal in this frame.
For \(x_i\in k^a\) define \(M_0=O^a\), \(M_i=M_{i-1}+Ox_i\),
and let \(r_i\) be the least nonnegative integer with
\(\pi^{r_i}x_i\in M_{i-1}\). Put \(r_0=0\).

\begin{theorem}[A bound for every contributing tuple]\label{thm:flagindex336}
Suppose the row \(F\) is \(K(d)\)-fixed. If the value
\(L(R(P)F)\), with \(P=\prod_{i=1}^s X_i(x_i)\), is nonzero, then
\[
 x_i\in\pi^{c-d-r_{i-1}}M_{i-1},\qquad
 r_i\le iD,\qquad [M_i:M_0]\le q^{D i(i+1)/2}.
\]
Consequently the joint exchange integral is a finite sum of values
of this same functional, with at most
\[
       q^{\,asd+aD\binom{s+1}{2}+D\binom{s+1}{3}}
\]
terms. Each term carries the measure
\(\prod_i\mu_i(\pi^dO^a)\). The estimates hold with arbitrary fixed
data at the other places, including right translates there.
\end{theorem}
\begin{proof}
Set \(x_0=0\). One has \(P^{-1}e_l=e_l-x_{l-1}\). For each \(i\)
we can find \(w\in\sum_{l\le i}Oe_l\) whose \(e_i\)-coefficient
is \(\pi^{r_{i-1}}\) and for which \(P^{-1}w\in\Lambda\).
For \(i=1\) take \(e_1\). For \(i\ge2\), write
\(\pi^{r_{i-1}}x_{i-1}=m+\sum_{j<i-1}t_jx_j\), with
\(m\in O^a\), \(t_j\in O\). Take
\(w=\pi^{r_{i-1}}e_i-\sum_{j<i-1}t_je_{j+1}\).
Its displacement under \(P^{-1}\) is \(-m\).

Let \(y\in\pi^d M_{i-1}^\vee\), where \(M^\vee\) consists of
covectors integral on \(M\). Define
\[
 b=y-\sum_{j\ge i}y(x_j)d_j^*,\qquad n=1+w b .
\]
The covector \(b\) annihilates the span of \(e_1,\ldots,e_i\).
The matrix \(n\) belongs to \(N(k)\) by the flag description.
Conjugating it by \(P\) on the right gives
\[
 P^{-1}nP=1+(P^{-1}w)(bP),\qquad
 bP=y+\sum_{j<i} y(x_j)d_j^* .
\]
All the entries of \(bP\) belong to \(\pi^dO\), so this last
matrix fixes \(F\). This also holds for \(d=0\), since the inverse
is obtained by subtracting the same square-zero matrix.
It follows from \(N\)-equivariance that a nonzero value at \(P\)
forces \(\chi(n)=1\).

The only simple coordinate of \(wb\) that can contribute is its
\((e_i,d_i)\) entry. The earlier \(d_j\) are annihilated by \(b\)
and \(w\) has no later row entries. On the opposite half \(n\)
fixes the columns of the simple character, by the same
inverse-adjoint calculation as in Lemma~\ref{lem:exchangeroot336}.
Thus \(\psi_i(-\pi^{r_{i-1}}y(x_i))=1\) for every such \(y\).
The annihilator of \(M_{i-1}^\vee\) is \(\pi^cM_{i-1}\), as
one sees in a basis of \(M_{i-1}\). This proves the first
inclusion. It implies \(r_i\le r_{i-1}+D\). The quotient
\(M_i/M_{i-1}\) is cyclic over \(O\) with annihilator
\(\pi^{r_i}O\), so its order is \(q^{r_i}\). Multiplication of
the successive indices proves the bound for \(M_i\).

Each \(X_i(\pi^d O^a)\) fixes \(F\), so the joint integrand
is constant on products of these cosets. The earlier choices
modulo \(\pi^dO^a\) determine \(M_{i-1}\). The allowed \(x_i\)
lie in \(\pi^{-iD}M_{i-1}\), which contains \(\pi^dO^a\).
The number of its cosets is at most
\(q^{a(d+iD)+D\binom{i}{2}}\). Multiplying over \(i\) gives the
assertion and additivity of Haar measure gives the weight.
\end{proof}

For the integral on the original word \(u\), take the product of
these bounds over \(v\in S_u\) with its own \(q_v,d_v,c_v\).
If the prescribed row is \(R(g)F_0\), with \(F_0\) fixed by
\(K(d_0)\), one may take \(d=d_0+t_{\max}-t_{\min}\),
where \(t_{\max},t_{\min}\) are the largest and smallest
elementary-divisor valuations of \(g\) on \(\Lambda\).
Indeed the two integral basis changes in a Smith decomposition
preserve every \(K(j)\), and conjugation by its diagonal part
lowers any entry valuation by at most \(t_{\max}-t_{\min}\).
Checking the inverse gives the lattice-stabilizer assertion also
when \(d_0=0\). This evaluates the joint integral on the
prescribed translate throughout the same interval.


\subsection{The exchange integral for a lattice in the middle block}\label{sec:middlecollapse337}
Retain the flag, the roots \(X_i(x)=1+xd_i^*\) and the ordinary
Bessel functional \(L\) of Lemma~\ref{lem:exchangeroot336}.
At a terminal split place the annihilator of \(O\) under the simple
additive character is \(O\). Take \(s\ge1\).
Write the full lattice on the chosen half as
\[
 \Lambda=\bigoplus_{i=1}^s Oe_i\ \oplus M\ \oplus
                  \bigoplus_{i=1}^s Of_i ,
 \qquad M\subset A_k\oplus kz .
\]
Here \(M\) is any full lattice in the indicated space. Put
\(J=M\cap A_k\), and define \(d\in\mathbb Z\) by
\(\{b:w+bz\in M\text{ for some }w\in A_k\}=\pi^{-d}O\).
Let \(F\) be fixed by the stabilizer of \(\Lambda\), allowing
a character of its determinant. All other local data are fixed.
\begin{theorem}[The surviving cosets in the middle block]\label{thm:middlecollapse337}
If \(L(R(\prod_i X_i(x_i))F)\ne0\), then
\[
       x_1,\ldots,x_{s-1}\in J,\qquad z-x_s\in M .
\]
For \(d<0\) the exchange integral is zero. For \(d\ge0\), choose
\(x_*\in A_k\) with \(z-x_*\in M\). Its value is exactly
\[
 \left(\prod_{i<s}\mu_i(J)\right)\mu_s(\pi^dJ)
 \sum_{x\in(x_*+J)/\pi^dJ} L(R(X_s(x))F).
\]
In this notation the sum uses one representative of each coset.
There are \(q^{ad}\) such cosets, where \(a=\dim_k A_k\).
\end{theorem}
\begin{proof}
Write \(M^\vee\) for the covectors integral on \(M\), and likewise
for \(J^\vee\). Every \(y\in J^\vee\) extends to some
\(b_0=y+\alpha z^*\in M^\vee\). In fact projection identifies
\(M/J\) with the free rank-one module \(\pi^{-d}O\). Lift a generator,
prescribe an integral value there, then extend the resulting
functional on \(M\) by \(k\)-linearity.

Proceed from \(i=1\) to \(s\). Once the preceding \(x_j\) lie in
\(J\), their shears have determinant one and preserve \(\Lambda\).
They commute with all the other shears and may be removed from the
argument of \(F\). Write \(P=\prod_{j=i}^s X_j(x_j)\), which fixes
\(e_i\), and take any \(b_0=y+\alpha z^*\in M^\vee\). Extend it
by zero on the flags. Then
\[
 b=b_0P^{-1}
   =y-\sum_{j=i}^{s-1}y(x_j)d_j^*+
                                  (\alpha-y(x_s))z^* .
\]
The element \(n=1+e_i b\) is in \(N(k)\). Each flag vector moves
by an earlier one, the action on the middle quotient is trivial,
and the inverse adjoint on the other half moves only \(f_i\).
Moreover \(P^{-1}nP=1+e_i b_0\) has determinant one and preserves
\(\Lambda\), as does its inverse. Evaluating \(L(R(nP)F)\) first
by this stabilizer and then by \(N\)-equivariance shows that
\(\chi(n)=1\) is necessary. The simple coordinate at \(i\) is
the only contributing one. For \(i<s\) its value is \(-y(x_i)\).
As \(y\) ranges through \(J^\vee\), triviality of its character
forces \(x_i\in J\). Indeed, in a basis of \(J\) each coordinate
is in the annihilator of \(O\). For \(i=s\) the simple coordinate
is \(\alpha-y(x_s)=b_0(z-x_s)\). Applying the same annihilator
calculation to a basis of \(M\) gives \(z-x_s\in M\).

This last condition has a solution exactly when \(1\in\pi^{-d}O\),
and its solution set is then \(x_*+J\). The first \(s-1\) integrations
now give their respective volumes. For \(t\in\pi^dJ\) the final
shear \(X_s(t)\) adds \(bt\in J\) to every \(w+bz\in M\).
It fixes the row, so the remaining integrand is constant on each
indicated coset. Their number is \([J:\pi^dJ]=q^{ad}\).
\end{proof}

For the transported terminal graph of Proposition~\ref{prop:graphexchange336},
one has \(z\in M\) and \(d=1\). The formula uses exactly one sum
over \(J/\pi J\) at each terminal place, with its original Haar
weight. At all \(f\) such places this requires at most \(T^{af}\)
ordinary Bessel evaluations. Some of these evaluations can vanish.


\subsection{The exchange at the two auxiliary split primes}\label{sec:auxiliaryflag339}
Fix a place over \(g_1\) or \(g_2\) of Theorem~\ref{thm:levels88}.
Use the half marked for the first position in Section~\ref{sec:exchange336}.
Let \(k\) be this local field, \(O\) its valuation ring, \(\pi\) a
uniformizer, and \(q\) its residue cardinality. Use the split flag
of Lemma~\ref{lem:exchangeroot336}, with chosen half ordered as
\[
       e_1,\ldots,e_s,\ A,\ z,\ f_s,\ldots,f_1 .
\]
The auxiliary primes are odd and split completely, and \(\varpi\)
in \(z=\varpi e\) is a unit there. The two isotropic strings are
integrally dual and the column lattice is unimodular at these
places. Let \(\Gamma\) be the preimage of the signed symmetric row
level of Proposition~\ref{prop:index90} and \(\lambda_\Gamma\)
its character. In the half defining this compact type the residue
group is upper unitriangular and the entries from the marked middle
line to \(A\) are zero. Its character is the chosen nontrivial
additive residue character on the sum of the adjacent entries other
than those at the two specified junctions, together with the
entry from \(z\) to \(e_s\). This description is
\srcloc{Formalism/LocalToGlobal/Columns/GenericUnipotentCharacter}{605}.
In the conjugate half take the inverse adjoint of these matrices.

Let \(b_1,\ldots,b_{2s+1}\) be the ordered string
\(e_1,\ldots,e_s,z,f_s,\ldots,f_1\), omitting \(A\), and let
\(Z_j(t)=1+b_jt b_{j+1}^*\). In this section \(b_j^*\) is the
chosen coordinate functional. Write \(\psi_j(t)=\chi(Z_j(t))\)
and \(c_j\) for the integer determined by
\(\{t:\psi_j(tO)=1\}=\pi^{c_j}O\).

\begin{lemma}[The residue row character on this flag]\label{lem:auxiliaryroots339}
For \(1\le i\le s\), the subgroups \(X_i(\pi O^a)\) and
\(Y_i(O^{a*})\) preserve the row and act trivially in its compact
character. On \(N\cap\Gamma\), the character \(\lambda_\Gamma\)
is the product of its values on the \(2s\) string entries:
each \(Z_j(O)\) belongs to \(\Gamma\), and its character is
nontrivial on \(O/\pi O\). The global Bessel character on \(N(k)\)
is the product of the nontrivial additive characters \(\psi_j\)
on those same entries.
\end{lemma}
\begin{proof}
All \(X_i(\pi O^a)\) reduce to the identity, in either half.
In the half defining the compact type, \(Y_i\) adds an \(A\)-coordinate
to \(e_i\). This is an allowed upper entry. The only possible
adjacent entry of that form is at the junction between \(e_s\)
and the first column coordinate, which the character omits. In
the conjugate half the inverse adjoint has an entry from \(f_i\)
to \(A\) in the defining half. The intervening marked line \(z\)
separates every column coordinate from \(f_i\), so these entries
are allowed and have trivial character as well.

On \(N\) the entire middle block is the identity. In the defining
half all surviving terms of the compact character are therefore
the successive entries of the string, including the extra entry
from \(z\) to \(e_s\). Each has a nonzero residue multiplier. On
passing to the opposite half the string reverses and dualizes.
On these entries inversion is negation: each joins consecutive
graded blocks of the parabolic flag, including the middle block
at its two ends. A product through an intermediate strictly
upper block cannot contribute. The inverse adjoint therefore
permutes the entries, multiplies each by a unit and introduces
a sign. This proves the compact assertions in either convention.
For the Bessel character, the chosen half gives the first \(s\)
entries directly by its defining trace, as in
Lemma~\ref{lem:exchangeroot336}. The inverse adjoint supplies the
last \(s\) by the same calculation. The factors in the split trace,
the Hermitian marked-line pairing and its chosen multiplier are
nonzero, proving the last assertion.
\end{proof}

\begin{theorem}[A common auxiliary translate for both positions]\label{thm:auxiliaryflag339}
There is a diagonal \(t\in G(k)\), acting identically on \(A\),
with the following properties. Its eigenvalues \(l_j\) on \(b_j\)
satisfy \(l_1=1\) and
\[
 v(l_{j+1})-v(l_j)=1-c_j,\qquad
       \psi_j((l_j/l_{j+1})w)=\lambda_\Gamma(Z_j(w))
                           \quad(w\in O).
\]
On any row \(F\) of the specified compact type and any column,
the complete exchange at this place on \(R(t)F\) equals its
ordinary period on \(R(t)F\) multiplied by
\[
             C_v=\prod_{i=1}^s\mu_i(\pi l_{i+1}^{-1} O^a).
\]
The character of the translated compact row agrees with \(\chi\)
on \(N\cap t\Gamma t^{-1}\). The local compact type admits a
smooth compactly supported test with nonzero ordinary local
pairing for the matching column character.
\end{theorem}
\begin{proof}
The pairing \((a,w)\mapsto\psi_j(aw)\) identifies
\(\pi^{c_j-1}O/\pi^{c_j}O\) with the characters of \(O/\pi O\).
Indeed its left kernel is zero by the conductor definition, and
the groups have the same order. Choose a representative \(a_j\)
of the nontrivial character in the lemma. Set \(l_{j+1}=l_j/a_j\).
These prescriptions define the diagonal in the chosen split
half and the inverse adjoint in the other. This torus normalizes
\(N\). For \(n=t n_0t^{-1}\) with \(n_0\in N\cap\Gamma\), its
\(j\)-th string entry is \((l_j/l_{j+1})\) times that of \(n_0\).
Multiplying the displayed character identities proves agreement
on the full intersection.

The translated row is fixed by \(Y_i(l_i O^{a*})\) and
\(X_i(\pi l_{i+1}^{-1}O^a)\), by the lemma and conjugation.
Under \(\psi_i\) the annihilator of the first lattice is
\(\pi^{c_i}l_i^{-1} O^a\), equal to the second lattice by the
valuation relation. The finite Fourier sum in
Theorem~\ref{thm:exchangefourier336} has just the zero coset.
It thus removes this exchange for the volume stated. Apply this
successively to the tail functionals, starting at \(i=1\).
Determinant twists of the row do not affect these shears.

For the last assertion write \(K_t=t\Gamma t^{-1}\) with
transported character \(\lambda_t\). Extend \(\lambda_t\) by
zero from \(K_t\) to \(G(k)\), and extend \(\lambda_t|_{H\cap K_t}\)
by zero from \(H\cap K_t\) to \(H(k)\). These are smooth tests of
compact support and the row has the required right character.
If the column is nonzero at \(h\), the row at \(nh\) is supported
exactly on \(n\in N\cap K_t\). Character agreement cancels the
unipotent character and the conjugate column cancels the value
at \(h\). The double integral is the product of the two positive
compact subgroup volumes.
\end{proof}

In the absolute-trace normalization \(\theta=-1/\sqrt{-d_0}\)
up to sign, all the string multipliers are units at these split
primes. Thus \(c_j=0\), \(v(l_j)=j-1\) and
\[
                 C_v=q^{a\binom{s}{2}}\prod_i\mu_i(O^a).
\]
The proof of the exchange identity itself uses just the two
trivial shear actions. It also applies at these levels with a
trivial row character and the indicated valuations; the positive
local test in the theorem uses the stated signed character.


\subsection{The auxiliary translate with least loss on the theta tests}\label{sec:centeredflag341}
Use an auxiliary half over $g_1$ or $g_2$ of
Theorem~\ref{thm:auxiliaryflag339} with its absolute-trace
normalization. Write $q$ for its residue size and $\pi$ for a
uniformizer. The ordered string of lines has length $d_0=2s+1$.
The rest of the first row is the column space $A$. Take $t$ equal
to the identity on $A$ and with string eigenvalues
$\pi^{-s}l_1,\ldots,\pi^{-s}l_{d_0}$, where the $l_i$ are those
of Theorem~\ref{thm:auxiliaryflag339}. The resulting string
valuations are $-s,-s+1,\ldots,s$.

\begin{proposition}[Exchange with the centered auxiliary translation]\leavevmode\label{prop:centeredflag341}
At this translate the character on the compact intersection still
agrees with the Bessel character. The local scalar removing the
first-position exchanges is
\[
                   q^{-a s(s+1)/2}\prod_{i=1}^{s}\mu_i(O^a).
\]
The determinant of $t$ is a unit. Replacing the auxiliary translates
in Theorem~\ref{thm:markedglobal339} by these elements at both
auxiliary primes gives its global comparison with the displayed
local scalars, the same $C_T$ and the same $D_T^s$.
\end{proposition}
\begin{proof}
The characters on consecutive string entries depend only on ratios
of consecutive eigenvalues. All ratios are unchanged by the common
factor. In the proof of Theorem~\ref{thm:auxiliaryflag339} the row
is fixed by the two shear lattices with scales $l_i$ and
$\pi l_{i+1}^{-1}$ respectively. Replacing the eigenvalues multiplies
these two lattices by $\pi^{-s}$ and $\pi^s$. They are still
annihilators of each other, so the only surviving exchange coset is
zero. Its measure at step $i$ is the measure of
$\pi^{s+1-i}O^a$ since $v(l_{i+1})=i$. The product is as stated.
The sum of the centered valuations is zero, giving the determinant
claim. The local changes take place at disjoint factors, so the
successive integrations of Theorem~\ref{thm:markedglobal339} apply.
\end{proof}

In the growing tower of ranks $n_i=r+it_{\rm r}$ use the matrix
indicator on each auxiliary edge from Proposition~\ref{prop:split173}.
Its support is $Y_i+\pi^{c_i}\mathcal M_i$, where
$Y_i=\binom{I_{n_i}}{0}$ and $\mathcal M_i=M_{n_{i+1},n_i}(O)$.
Extend $t$ at every vertex by the identity in the additional
coordinates. Write $g_i$ for these extensions, and put
\[
   L_i=\frac{d_0(d_0^2-1)}{6}
             +\left(\frac{n_i+n_{i+1}}2-d_0\right)s(s+1).
\]

\begin{theorem}[The optimal overlap on each auxiliary theta edge]\leavevmode\label{thm:auxoverlap341}
The normalized absolute coefficient of the original and translated
indicators at the auxiliary place is $q^{-L_i}$. On their common
subdivision each indicator has exactly $q^{L_i}$ cells. The values
of the translated indicator on those cells are algebraic units.

This coefficient is largest among simultaneous translations at the
vertices which preserve all the centers $Y_i$ and induce the
prescribed character ratios on the original diagonal string.
Here the initial translation preserves that string and acts diagonally
on its lines, with any common shift of the valuations allowed. The
extension on the complementary coordinates at each vertex may be
arbitrary. One choice attains these bounds on every edge at once.
\end{theorem}
\begin{proof}
In the displayed extension the reference lattice splits into the
string and the other coordinate lines at every vertex. The
exponents on the Hom lattice are all target exponents minus source
exponents, as in Theorem~\ref{thm:thetacells340}. Among string pairs
half the sum of their absolute differences is
$\sum_{k=1}^{d_0-1}k(d_0-k)=d_0(d_0^2-1)/6$. For each zero exponent
in the source or target complement the other string contributes
$\tfrac12\sum_{k=-s}^s|k|=s(s+1)/2$. The remaining pairs contribute
zero. This gives $L_i$. The sum of the signed differences is zero,
so both intersection indices are $q^{L_i}$. The centers agree.
Dividing the measure of the intersection by the geometric mean of
the two support measures gives the coefficient. In the matrix Weil
action the half density is one. Its splitting character at the pair
is a root of unity. Indeed a smooth character of each general linear
group factors through determinant by elementary matrices, and its
restriction to the compact unit group has finite image. Thus the
coefficient of each individual cell in the indicator is a unit.

For the upper bound consider any allowed sequence $b_i$. The
identity $b_{i+1}Y_i=Y_i b_i$ identifies a stable copy of the same
string at every vertex. Its lattice is a direct summand in each
reference frame. We first record the effect of passing to a stable
subspace. Let $W$ be a lattice in a vector space with automorphism
$B$ preserving a subspace $V_1$. Write $W_1=W\cap V_1$, $W_2$ for
the quotient lattice, and $B_1,B_2$ for the induced maps. The
intersection $J=W\cap BW$ meets $V_1$ in $W_1\cap B_1W_1$ and
its image in the quotient is contained in $W_2\cap B_2W_2$.
Multiplying the subspace and quotient indices shows that
$[W:J][BW:J]$ is at least the product of the corresponding two
index products on the subspace and quotient. This follows by
applying the index formula to $J\subset W$ and then to $J\subset BW$.

Apply this observation to $Y\mapsto b_{i+1}Yb_i^{-1}$ on the matrix
lattice. First use the subspace of maps with image in the string.
On that subspace and on its quotient then use the maps zero on the
source string. The four resulting lattices are Hom lattices between
the source and target strings and their complements, since the
reference string inclusions split integrally. Thus the negative
base-$q$ logarithm of the normalized overlap is at least the sum
of the four graded costs. On the string the exponents are
$m,m+1,\ldots,m+2s$ for some $m\in\mathbb Z$ fixed by the first
vertex. The string-to-string cost stays $d_0(d_0^2-1)/6$.
For each Smith exponent $v$ on a complementary quotient, the mixed
cost includes $\tfrac12\sum_{k=0}^{2s}|m+k-v|$. This is minimized
at the median $v=m+s$, with value $s(s+1)/2$. This can be checked
by comparing the sum at consecutive integers. The cost between the
two complementary quotients is nonnegative. These bounds give
$L_i$ on every edge. The chosen centered diagonal has $m=-s$ and
all complementary exponents zero, so all inequalities are attained.
\end{proof}


\subsection{The integral cost of isolating a vector marking}\label{sec:pointlattice353}

At one unrefined ordinary initial class \(a\) take its actual
rank \(k\) lattice \(M\), the prescribed coefficient weights
\(\lambda_\tau\), and the product of vector spaces \(\Omega_0\)
of Section~\ref{sec:vectorchain353}. Put \(f=|S_T|\) and
\(Q=T^{fk}\). Enlarge the coefficient field \(J\) to contain a
primitive \(T\)-th root \(\zeta\). Set \(u=\zeta-1\), and put
\[
 z_{\min}=\prod_\tau\tau(\bar\pi)^{\lambda_{\tau,k}},
 \qquad z_{\max}=\prod_\tau\tau(\bar\pi)^{\lambda_{\tau,1}} .
\]
Write \(D_a\) for the column Weyl lattice in the unrefined
coefficient frame of \(M\) over \(\mathcal O_J\). At a refinement
\(\xi\), denote the actual column lattice by \(D_{a,\xi}\).
It depends only on the tuple of lines spanned by the \(\xi_P\).
Let \(S_a^\circ\) consist of the separately zero-sum tables
with value at \(\xi\) in \(D_{a,\xi}\). Independence of the
refined initial coordinates at this class is proved in
Section~\ref{sec:pathnorm352}.

\begin{theorem}[Evaluation and prescribed support in the transformed lattice]\leavevmode
\label{thm:pointlattice353}
Fix \(v=(v_P)\in\Omega_0\) with every \(v_P\ne0\). For each
\(P\) choose \(d_P\in k_P^\times\) different from one.
Use the transform \(\mathcal F_0\) of
Theorem~\ref{thm:vectorchain353}. Then
\[
 \begin{aligned}
   \{(\mathcal F_0 H)(v):H\in S_a^\circ\}
         &=u^f z_{\min} D_a,\\
   \left\{y:
      \Big(\bigotimes_P(\delta_{v_P}-\delta_{d_Pv_P})\Big)y
                         \in\mathcal F_0 S_a^\circ\right\}
         &=Q\,u^{-f} z_{\max}D_a .
 \end{aligned}
\]
Both are identities of fractional lattices on every sufficiently
distant period interval.
\end{theorem}
\begin{proof}
For one product of covector lines choose representatives \(c_P\).
The coefficient lattice \(D_{a,c}\) is constant on that product
of scalar orbits. The separate augmentation kernels there
are generated by tensors of differences of point masses, as in
Theorem~\ref{thm:pathnorm352}. If \(c_P(v_P)=0\) at some place,
each such generator transforms to zero at \(v\). Otherwise
each nontrivial difference gives there a factor which is a
difference of distinct \(T\)-th roots of unity. Its ratio to
\(u\) is an algebraic unit: \(\zeta^j-1\) is divisible by \(u\)
and conversely \(u\) by \(\zeta^j-1\) when \(T\nmid j\), by
the geometric sum in each direction. Independence of the
orbits shows that the evaluation lattice is \(u^f\sum_cD_{a,c}\),
where the sum is over tuples with \(c_P(v_P)\ne0\) at every \(P\).

For the second assertion put \(D(v)=\bigotimes_P
(\delta_{v_P}-\delta_{d_Pv_P})\). Fourier inversion gives
the unique preimage of \(D(v)y\) at \(\xi\) as
\[
   Q^{-1}\prod_P
       \big(\psi_P(-\xi_P(v_P))-\psi_P(-d_P\xi_P(v_P))\big)y.
\]
It has separate sums zero since \(D(v)y\) does. For a tuple
of allowed lines its coefficient is \(Q^{-1}u^f\) times a
unit. For every other tuple its coefficient vanishes.
The condition on \(y\) is therefore \(y\in Q u^{-f}\bigcap_cD_{a,c}\),
over the same allowed tuples.

We calculate the sum and intersection in the actual coefficient
frames. Let \(M(c)\) be the sublattice of \(M\) defined by the
covector kernels at \(P\mid\pi\). The lattice used for column
coefficients is the ordinary linear dual of the Hermitian dual of
\(M(c)\), in each coefficient embedding. This follows from the
definition of the column on the Hermitian dual of the marked
kernel and the inverse transpose action on its coefficients, as
in Proposition~\ref{prop:normalframe349}. Equivalently see
\srcloc{Formalism/LocalToGlobal/Columns/StdLatUnimodular}{117},
\srcloc{Formalism/LocalToGlobal/Columns/CoefCol}{89} and
\srcloc{Formalism/LocalToGlobal/Divisibility/CoefMod}{326}.
At a split pair \(P,\bar P\) the Hermitian form pairs the two
local lattices perfectly. Taking the Hermitian dual exchanges
the two places and takes the ordinary dual there. Taking the
second, ordinary dual recovers the original sublattice from
the opposite place. Thus the coefficient lattice at \(\bar P\)
is, in the identified unrefined frame, \(W_{\lambda_\tau}(L_{c_P})\)
of Lemma~\ref{lem:openflags353}. At \(P\) it uses the entire
lattice. This also proves that all refinements agree away from \(T\).

At a fixed coefficient prime over \(T\), distinct embeddings
\(\tau\) induce distinct primes of \(E\): each completion
\(E_P=\mathbb Q_T\) has only one embedding over \(\mathbb Q_T\)
into an algebraic closure. The affected factors correspond to
\(\tau\) inducing \(\bar P\mid\bar\pi\); each then depends
on its own \(c_P\). At such a factor apply
Lemma~\ref{lem:openflags353} with \(\mathcal U=\{v_P\}\).
Its cutoff holds since all initial spreads are \(O(xb+x^C)\) by
\eqref{eq:letters111} and \(\log T\gg m x^\eps\) on this family.
Sums over the independently chosen \(c_P\) factor as tensors
of sums. For intersections take ordinary lattice duals, tensor
the sums of the duals and dualize again. The unrefined factors
stay in their original frames. Normalizing the valuation by
\(v(T)=1\), each \(\tau(\bar\pi)\) has valuation one in an
affected factor and zero otherwise. Hence locally the sum is
\(z_{\min}D_a\) and the intersection is \(z_{\max}D_a\).
At every other prime both scalars are units and every lattice
is \(D_a\). These local identities prove the two global ones.
\end{proof}

\begin{corollary}[The two weight extremes on the common fields]\leavevmode
\label{cor:isolatedcost353}
For the coded ordinary columns, the normalized logarithmic finite
cost of evaluation at a single \(v\) relative to \(D_a\) is
\(o(b/x^A)\) for every fixed \(A>0\). If there is a high
embedding, at some prime above the same \(T\) the valuation
difference between the specified-support and evaluation lattices
in the theorem, multiplied by \(\log T\), is at least
\((\lceil xb\rceil+\lceil xb/\exp(x^{\eps/3})\rceil)\log T\).
Both assertions concern the same word and coefficient module.
\end{corollary}
\begin{proof}
In rank \(k\) at the initial vertex the algebraic weight is obtained
by subtracting \(\rho_k(i)=(k+1)/2-i\) from \(u_i\). Thus
\(\lambda_k=-b_*\) at a high embedding and
\(\lambda_k=(k-2)/2\) elsewhere by \eqref{eq:letters111}, where
\(b_*=\lceil xb/\exp(x^{\eps/3})\rceil\). At the high embedding
\(\lambda_1=\lceil xb\rceil\). The two extremes are independent
of the word for a fixed translate. One has \(v_T(u)=1/(T-1)\),
by evaluating the \(T\)-th cyclotomic polynomial at one and
using the unit ratios above. The only finite cost for evaluation
is therefore bounded in absolute value by
\(f(b_*+k+1)\log T\), with supremum over all primes above \(T\)
in the normalization \(v(T)=1\). Both \(f,k\) and \(\log T\)
are polynomially bounded in \(x\). This proves the first estimate.
Choose a prime dividing \(\tau(\bar\pi)\) for one high embedding.
At this prime the difference in valuations is
\[
    fk-\frac{2f}{T-1}
        +\sum_{\tau':\,v(\tau'(\bar\pi))=1}
                      (\lambda_{\tau',1}-\lambda_{\tau',k}) .
\]
The first term after combining the two Fourier contributions is
nonnegative, as are all spreads, and the chosen spread is
\(\lceil xb\rceil+b_*\). This proves the lower bound.
\end{proof}


\subsection{Simultaneous interpolation of vector markings}\label{sec:grid354}

Retain the actual coefficient lattice at the ordinary class \(a\)
of Theorem~\ref{thm:pointlattice353}. At each \(P\mid\pi\) choose
a set \(G_P\) of \(s_P\ge1\) distinct nonzero vectors in \(M/PM\).
Put \(G=\prod_P G_P\), \(s=\sum_P s_P\), and
\(D_a^G=\bigoplus_{v\in G}D_a\). Let \(E_G\) take the values of
a vector table on \(G\). Values outside \(G\) are free. We use
the same unnormalized Fourier transform, root \(u=\zeta-1\), and
weight scalar \(z_{\min}\) as in that theorem.

\begin{theorem}[Integral interpolation on a product of residue sets]\leavevmode
\label{thm:grid354}
Let \(h=\max_\tau(\lambda_{\tau,1}-\lambda_{\tau,k})\).
Suppose \(T>h+k\) and, at each \(P\),
\[
       s_P+1\le T-1,\qquad T-\binom{s_P+1}{2}>h .
\]
Then the lattice of simultaneous evaluations satisfies
\[
     u^s z_{\min} D_a^G
       \ \subseteq\ E_G\mathcal F_0 S_a^\circ
       \ \subseteq\ u^f z_{\min}D_a^G .
\]
For \(s_P\le\lfloor\sqrt T\rfloor\) on every sufficiently distant
common interval, these hypotheses hold for all selected words.
In that range one integer \(B_x=T^{1+f(b_*+k)}\) satisfies
\[
        B_xD_a^G\ \subseteq\ E_G\mathcal F_0 S_a^\circ
                     \ \subseteq\ B_x^{-1}D_a^G ,
       \qquad \log B_x=o(b/x^A)\quad(A>0\ \text{fixed}).
\]
The integer is common to all the stated grids and words with
this terminal prime.
\end{theorem}
\begin{proof}
The upper inclusion is Theorem~\ref{thm:pointlattice353} at
each coordinate. For the lower inclusion enumerate the vectors
at one place as \(v_1,\ldots,v_{s_P}\). Use the list consisting
of all \(v_i\) and all \(v_i-v_j\) for \(i<j\). It has at most
\(\binom{s_P+1}{2}\) members, each nonzero. A covector \(c_P\)
nonzero on this list takes distinct, nonzero values \(t_i\) on
the \(v_i\). Fix such a covector and an index \(j\), and write
\(z_i=\psi_P(t_i)\). The \(z_i\) are distinct nonidentity
\(T\)-th roots because the residue field has prime order. Use
the coefficients of the polynomial
\[
       p_j(Z)=Z(Z-1)\prod_{i\ne j}(Z-z_i)
                  =\sum_{d=1}^{s_P+1}p_{j,d} Z^d
\]
as the coefficients of point masses at \(d c_P\) on the
covector side. All \(d\) are distinct nonzero residues. Their
coefficients have sum zero since \(p_j(1)=0\). The transform
of this array at \(v_i\) is \(p_j(z_i)\). It vanishes for
\(i\ne j\). At \(v_j\) it is \(u^{s_P}\) times a unit of
\(\mathcal O_J\), by the unit calculation for differences of
roots in Theorem~\ref{thm:pointlattice353}.

Take the tensor product of these arrays for one chosen slot of
\(G\) and multiply it by a vector in \(D_{a,c}\). All its
nonzero entries have the same product of covector lines, hence
belong to their prescribed coefficient lattice. Its transform
on \(G\) is supported on the chosen slot and takes every value
in \(u^s D_{a,c}\) there, by multiplication by the inverse
of the unit. Vary the good covector independently at each place.
At any prime over \(T\) the lattice computation in the proof of
Theorem~\ref{thm:pointlattice353} applies with the lists just
chosen. Lemma~\ref{lem:openflags353} gives
\(\sum_c D_{a,c}=z_{\min}D_a\) after localization. Elsewhere
the same equality holds since the sets of choices are nonempty
and the lattices coincide. Equality follows globally by
localization. Summing these constructions on all the slots proves
the first assertion.

On the given family \(h/T\) tends to zero uniformly, by
\eqref{eq:letters111} and Theorem~\ref{thm:common111}. For
\(s_P\le\sqrt T\) the number of excluded hyperplanes is at
most \((T+\sqrt T)/2\), which gives the two cutoff inequalities.
In the normalization \(v(T)=1\) one has
\(s\,v(u)\le f\sqrt T/(T-1)<1\) eventually, since \(f\) is
polynomial in \(x\) and \(\log T\gg m x^\eps\). At any prime
above \(T\), \(|v(z_{\min})|\le f(b_*+k)\) by
Corollary~\ref{cor:isolatedcost353}. Outside \(T\) all the
scalars are units. Comparing valuations gives the stated two
integer inclusions. Finally \(b_*=\lceil xb/\exp(x^{\eps/3})\rceil\)
while \(k,f,\log T\) are polynomially bounded in \(x\). The
claimed rate for \(\log B_x\) follows.
\end{proof}

\section{Local factors for a reversed block}\label{nav:g-IX-2}

\subsection{The split factors for two unrestricted local valuations}\label{nav:II-65}
Use \(a_\pm,L,\epsilon\) of Lemma~\ref{lem:scalardepth294}
and put \(m_0=a_++a_-\), \(C_k=4(k+1)^4\). Write
\(F_v(c,\lambda)\) for the orbital operator on the
original local columns, relative to the identity mass.

\begin{proposition}[A uniform split bound for one moving line]\leavevmode
\label{prop:splitall294}
At a split place on the support with \(m_0>0\) one has
\[
 \|F_v(c,\lambda)\|\le
   C_k^{m_0} P^{(k+4s)m_0/2+k\epsilon}
   \begin{cases}
      1,&v\nmid T,\\
      P^{m_0/2+s|L|},&v\mid T .
   \end{cases}
\]
When \(m_0=0\) the norm is at most one. If the original
support is empty, its operator is zero.
\end{proposition}
\begin{proof}
The first-plane decomposition \(B=J\ltimes U\) and the
commutation in Proposition~\ref{prop:representative294}
give the counting measure and relative translations of
Theorem~\ref{thm:finite281}. Indeed the stabilizer ratio
of a finite cell is the size of its orbit of right
\(U_L\)-cosets under \(J\cap K\), by the same diagonal
block argument. At a full column level its coefficient
is a translation compressed to \(C=\mathrm{GL}_k(O)\)
invariants. At \(v\mid Q\), the support is \(K_1B\)
and \(C\) acts on its cosets as in
Proposition~\ref{prop:markedpowers292}. The new isometry
still commutes with \(C\). The compact conjugation average
therefore uses the same principal-congruence ambiguity
and the same scalar character phases. At these places
each orbit average is bounded by \(\Xi(A_0)\) by
Lemmas~\ref{lem:conjugation292} and \ref{lem:scalardepth294}.
Here translating by a central scalar preserves the
spherical majorant, since the inducing character for
\(\Xi\) is trivial on the center.

At \(T\) sum first over the enlarged compact orbits of
Proposition~\ref{prop:affine293}, inserting a zero at
every unsuccessful coset. The decomposition, the success
test \(y\beta^{-1}\in S_0\) and the density bound
\(|\phi(y)|\le1\) there use only commutation with \(D_0\),
which still holds. The affine change \(y\mapsto y\beta^{-1}\)
now has linear part \(A\) of arbitrary determinant.
The weighted composition has norm at most
\(|\det A|^{-1/2}\) by additive change of variables.
Lemma~\ref{lem:tensor293}, after division by that scalar,
gives the local orbit majorant
\[
      P^{(e_--e_+-n L)/2}\,\Xi(A_0).
\]
Indeed \(v(\det A)=e_--e_+-nL\). The two exponents of
\(A_0\), and this determinant, are constant on each compact
orbit. Assign its majorant to each member to sum the
orbit sizes. The depths may vary among the members.

We count using \(X,Y,Z\) from Lemma~\ref{lem:scalardepth294}.
Right multiplication by an integral first unipotent
chooses representatives of \(H,l\) modulo the integral
vector and row lattices. At depths \(t,u'\) there are at
most \(P^{k(m_0-t-u')+(n-k)m_0}\) such choices.
With \(H,l\) fixed, the two endpoint tests for \(z\)
each prescribe a coset of a fractional integral lattice,
of respective denominator lengths \(a_-,a_+\).
Their intersection contains at most \(P^{\min(a_-,a_+)}\)
classes modulo \(O\). This continues to hold after the
choice of lifts, because right integral translations
preserve support. For \(m_0=0\) this gives at most one
coset and a single unitary translation has norm at most
one.

For \(m_0>0\) the spherical bound is
\(k^{d_-+d_+}P^{-(k-1)(d_-+d_+)/2}\).
Use \(e_++e_--2S=m_0-2(t+u')-2\epsilon\)
in Lemma~\ref{lem:scalardepth294}. Apart from the
\(L\)-terms, the contribution of the \(k\) coordinates
and the negative spherical exponent is at most
\[
       (k+1)m_0/2-(t+u')+(k-1)\epsilon
\]
in the exponent of \(P\). The other coordinates and
the center cost at most \((n-k+1/2)m_0\). Their sum is at
most \((k+4s)m_0/2+k\epsilon\). Away from \(T\) the
remaining \(L\)-term is nonpositive. At \(T\) combine it
with the determinant contribution. For \(L\ge0\) the sum
of the \(L\)-terms is \(((k-1)-n)L/2\le0\). For \(L<0\)
it is \((n-k+1)|L|/2=s|L|\). The remaining determinant
term \((e_--e_+)/2\) is at most \(m_0/2\) there.
Finally there are at most
\((a_++1)(a_-+1)\le2^{m_0}\) depth pairs and
\(d_-+d_+\le m_0+2\epsilon\le3m_0\).
These factors fit in \(C_k^{m_0}\), proving the bound.
\end{proof}

\subsection{Finite mass for a single reversed block}\label{nav:II-66}
\begin{proposition}[A bound from the two prefix determinants]\leavevmode
\label{prop:blockfinite295}
Let \(w\) be the supported representative in
Proposition~\ref{prop:block295} with cross block
\(C=cS_j\), and put \(A_\xi=|N_{E/\mathbb Q}\xi|^{1/2}\).
The product of its finite unsigned masses relative to
the identity is at most
\[
                  (Q^f A_\xi)^{\,2rj}.
\]
The same expression bounds its ordinary finite column
operator in norm on the original level.
\end{proposition}
\begin{proof}
At a split place let \(c_\pm\) be the two components,
\(e_\pm=v(c_\pm)\), and \(P=Nv\). Work in the actual
integral first-component basis, including the graph at
\(T\). Write the upper blocks of \(n_0\in U\) as
\(l,H,Z\) and the row lattice as \(\Lambda\). Put
\(C_+=c_+S_j\). The intersection of \(g^{-1}\Lambda\)
with the first \(j\)-plane for \(g=n_0^{-1}w\) is
\(C_+^{-1}L_+\), where
\[
 L_+=\{v\in O^j: Hv\in O^n,\ (lH-Z)v\in O^j\}.
\]
Indeed applying \(g^{-1}=w^{-1}n_0\) to an integral
input with last coordinate \(v\) gives zero middle and
last coordinates precisely when its middle coordinate
is \(-Hv\) and its first coordinate is \((lH-Z)v\).
If \(g\in KB\), its inverse lattice is a period
translate of \(\Lambda\). The first \(j\)-plane is stable
under \(B\), and \(B\) acts on it by upper unitriangular
matrices. Its intersection lattice therefore has the
same covolume as \(O^j\). Consequently
\([O^j:L_+]=P^{m_+}\), where \(m_+=je_+\).
Right multiplication of \(n_0\) by \(U_L\) preserves
\(g^{-1}\Lambda\), so this intersection and \(L_+\) are
well defined for a coset.

In the reversed dual component the three upper blocks
are \(-H^{\mathsf t},-l^{\mathsf t},(lH-Z)^{\mathsf t}\).
Their rectangular column and their product minus the
central block are \(-l^{\mathsf t}\) and \(Z^{\mathsf t}\).
The same intersection calculation gives a lattice
\(L_-\subset O^j\) of index \(P^{m_-}\), \(m_-=je_-\),
on which these two matrices are integral. The sign
permutation of the exchange is integral on both sides.
All arguments apply to the original row level because
it stabilizes the chosen row and dual lattices.

There are at most \((2P^{j-1})^m\) sublattices of
index \(P^m\) in \(O^j\). Each is reached by a chain of
\(m\) steps of index \(P\), and from any free rank-\(j\)
lattice there are \((P^j-1)/(P-1)\le2P^{j-1}\) such
steps. For fixed \(L_+\), each row of \(H\), modulo
integral rows, belongs to \(L_+^\vee/(O^j)^\vee\)
of order \(P^{m_+}\), by elementary divisors.
There are \(n\) such rows. For fixed \(L_-\), the \(n\)
rows of \(l^{\mathsf t}\) and the \(j\) rows of
\(Z^{\mathsf t}\) allow at most \(P^{(n+j)m_-}\)
residue choices. To use this count, choose the exact
lifts of \(H,l\) first by right integral multiplication.
The two kernel lattices persist under this operation.
Thereafter \(Z\) varies only modulo the integral center
in a given right coset. Thus the local mass is at most
\[
  (2P^{j-1})^{m_++m_-}P^{n m_++(n+j)m_-}
       \le P^{rj(e_++e_-)} ,
\]
using \(2\le P\) and \(n+2j=r\).

At an inert place use the quadratic integer ring with
residue cardinality \(P'=P^2\), and put \(e=v_w(c)\).
The same lattice \(L_+\) has index \((P')^{je}\).
Now the other rectangular block is the negative adjoint
of \(H\) for the middle form and \(lH-Z=Z^*\), with
star on the square center denoting conjugate transpose.
Thus \(L_+\) alone gives at most
\((P')^{(n+j)je}\) row residue choices for \(H,Z^*\).
To lift an integral shift of \(H\) to \(U_L\), choose one
central entry from each off-diagonal pair and solve the
diagonal trace equations by integral trace surjectivity
in the unramified quadratic extension. For a fixed
rectangular lift, two admissible central blocks with
the same residue modulo integral matrices differ by an
integral skew matrix and are equivalent. Including
the choices of \(L_+\) gives at most
\((P')^{rje}=P^{2rje}\).

The factors of \(|N_{E/\mathbb Q}c|^{1/2}=Q^f A_\xi\)
are \(P^{(e_++e_-)/2}\) and \(P^e\), respectively.
Multiplication proves the mass estimate by
Proposition~\ref{prop:block295}. Each term in its
coset count acts by a relative unitary column translation,
with compact averaging if needed, and a scalar character
of modulus one. This gives the same bound for the operator.
\end{proof}

\subsection{The two prefix lattices at the marked primes}\label{nav:II-67}
\begin{proposition}[The finite mass with the row congruences imposed]\leavevmode
\label{prop:blockmarks296}
Write \(A_\xi=|N_{E/\mathbb Q}\xi|^{1/2}\) for the nonzero
integer in \(c=Q\xi\) from Proposition~\ref{prop:block295}.
The full finite unsigned mass relative to the identity is at most
\[
                    Q^{j^2 f}A_\xi^{2rj}.
\]
The same bound holds for the finite operator norm on the given
column.
\end{proposition}
\begin{proof}
Use the notation of Proposition~\ref{prop:blockfinite295} at
a place over \(Q=g_1g_2\). Both primes split completely. Put
\(P=Nv,\ \pi O=\mathfrak p_v,\ a_\pm=e_\pm-1\), and let
\(\Lambda\) be the row lattice in the first component.
For a supported \(g=u_0^{-1}w=k_0b_0\), with \(k_0\) in the
actual row level and \(b_0\in B\), its first-plane image
satisfies
\[
     gP_j\cap\Lambda=k_0(P_j\cap\Lambda).
\]
The reduction of the original row level preserves the first
\(j\)-plane in this lattice. This follows from its upper block
form in Lemma~\ref{lem:flags211}. Therefore both the middle
and last projections of every vector in this intersection
vanish modulo \(\pi\). Its vectors have last projection
\(v\in L_+\) and middle projection \(-Hv\), and every
\(v\in L_+\) occurs, by the explicit description in
Proposition~\ref{prop:blockfinite295}. Hence
\[
           L_+\subset\pi O^j,\qquad H L_+\subset\pi O^n.
\]
In the reversed dual component the reduction also preserves the
first \(j\)-plane. It is the annihilator of the initial
subspace of codimension \(j\) in the original upper block
flag. The identical argument
gives \(L_-\subset\pi O^j\) and \(l^{\mathsf t} L_-\subset
\pi O^n\).

Now \(L'_\pm=\pi^{-1}L_\pm\) have indices \(P^{ja_\pm}\).
Each row of \(H\) is integral on \(L'_+\) and each row of
\(l^{\mathsf t}\) on \(L'_-\). The \(j\) central rows of
\(Z^{\mathsf t}\) are integral on \(\pi L'_-\), of index
\(P^{j(a_-+1)}\). Counting lattices and row residues as in
Proposition~\ref{prop:blockfinite295} gives, with
\(m'_\pm=ja_\pm\), a local bound
\[
 (2P^{j-1})^{m'_++m'_-}
 P^{n m'_+ +(n+j)m'_-+j^2}
                 \le P^{rj(a_++a_-)+j^2}.
\]
Choosing fixed row lifts by right integral multiplication
preserves support and the kernel lattices, as in that proof.
At places away from \(Q\) use its unmarked bounds directly.
Their factors and \(P^{rj(a_++a_-)}\) above \(Q\) multiply to
\(A_\xi^{2rj}\). There are \(f\) places of \(F\) over each
marked prime, giving the claimed prefactor. Taking norms of
unitary translations, also after compression, proves the
operator assertion.
\end{proof}

\subsection{The finite operator for all blocks in a composition}\label{nav:II-68}
\begin{proposition}[Finite cosets over the internal unipotent]\leavevmode
\label{prop:compositionfinite298}
Put \(\delta_k=100(s+1+\log_2(2(k+1)))\). For any
supported composition of Proposition~\ref{prop:composition298}
with \(2\le j\le s\), the finite signed operator has norm at most
\[
     |N_{E/\mathbb Q}G_j|^{(k+\delta_k)/2}
                          \prod_{p=1}^j |N_{E/\mathbb Q}G_p|^{u_p}.
\]
\end{proposition}
\begin{proof}
At a finite place fix \(z\) in the prefix unipotent. To count
the accompanying \(u_0\) use \(g=u_0^{-1}w_{C(z),\lambda}\).
On the support write \(g=k_0 b\) with \(k_0\) in the original
row level and \(b\) in the period subgroup, enlarged to \(B'\)
when counting the affine orbits over \(T\). On the \(j\)-plane
let \(b_1\) be the matrix of \(b^{-1}\), which is upper
unitriangular. The first \(j\) columns of \(gb^{-1}=k_0\)
are integral in the row basis. Their bottom block is
\(C(z)b_1=z^*C b_1\), because the last row block of \(u_0^{-1}\)
is \((0,0,I_j)\). In the split case this argument applies
in both the primal and reversed dual components.

Apply Lemma~\ref{lem:pivots298}. All its minors are integral,
so \(v(G_p)\ge0\) and its \(u_p\) entries in column \(p\)
lie in \(G_p^{-1}O\). At an inert place their coordinate
volume costs at most \(\prod_p(P^2)^{u_p v(G_p)}\), with
\(P=Nv\). At a split place conjugation interchanges the
two components of \(z\). Counting them both gives at most
\(\prod_p P^{u_p(v(G_{p,+})+v(G_{p,-}))}\).
These volume bounds also bound the numbers of possible
right \({\cal V}_L\)-cosets. Indeed if an element \(z\)
admits a supported \(u_0\), the same holds for each \(zv_0\)
with \(v_0\in{\cal V}_L\), by replacing \(u_0\) with
\(\widehat v_0^{-1}u_0\widehat v_0\). This changes the
original unipotent by right integral multiplication.
Each such coset has coordinate volume one. Multiplication
of the local bounds gives the displayed product over \(p\).

For the remaining estimate at fixed \(z\), the prefix
lengths \(m_\pm\) of Proposition~\ref{prop:blocksplit297}
are now \(v(\det C(z)_\pm)=v(G_{j,\pm})\). To see that
its proof applies, the intersection of the inverse lattice
with the \(j\)-plane is still \(C(z)_+^{-1}L_+\),
and \(L_+\) is still the kernel defined by \(H,lH-Z\)
in Proposition~\ref{prop:blockfinite295}. The determinant
of the triangular period block on that plane is one. The
same statements in the reversed dual give \(m_-=v(G_{j,-})\)
and the kernel defined by \(l^{\mathsf t},Z^{\mathsf t}\).
The middle projection in Lemma~\ref{lem:blocklattice297}
still has the formula \(m+Hv\) with \(lm+Zv\) integral.
It only uses that the outer exchange blocks are invertible.
Thus both projected-length estimates and the determinant
relation for the column hold with these lengths.
The compact groups used in the split proof commute with
\(\widehat z\) and with \(w\). At \(T\) their action on
the middle block and its success test are the original affine
ones. Hence the same compact-orbit estimate applies at each
fixed \(z\), allowing the actual scalar phases of its terms.
The residue count for the rectangular and central coordinates
used only the lengths of the two kernels. We obtain the split
bound in Proposition~\ref{prop:blocksplit297} with
\(M=v(G_{j,+})+v(G_{j,-})\).

At an inert place the adjoint kernel and the unimodular
column calculation in Proposition~\ref{prop:blockfinite297}
apply with \(m=v(G_j)\). Indeed \(lH-Z=Z^\dagger\) and
the middle projection are unchanged, and \(C(z)\) enters
the index calculation only through its determinant.
Taking the product of these bounds for the fixed-\(z\)
fibers costs at most \(|N_{E/\mathbb Q}G_j|^{(k+\delta_k)/2}\)
by the final comparison in Proposition~\ref{prop:blockfinite297}.
This bound is independent of \(z\). Multiplying by the
number of possible internal cosets proves the claim.
\end{proof}

\subsection{Combining the real and finite transverse factors}\label{nav:II-69}
For a supported transverse label with \(\mathcal G>0\), use the
arithmetic endpoints \(a_i\) of Proposition~\ref{prop:phase301}.
Write \(j_i=J_i,\ h_i=j_i-j_{i-1}\), and put
\[
 H_i=\prod_\tau\max(1,|\tau a_i|),\qquad
 A_-=\max\!\left(1,\ Q^{-df}\prod_\tau
                     (|\tau D_d|-T|\tau D_s|)\right).
\]
The terminal height \(H_t\) is calculated from
\(\Delta_+=D_d+TD_s\) at every place. In particular
\(A_-\le H_t\) by the triangle inequality.
Let \(B_\gamma\) denote the sum of absolute normalized
cell traces in Proposition~\ref{prop:tracemass301}.

\begin{theorem}[A bound for a transverse label at all places]\leavevmode
\label{thm:product301}
Take any retained fraction \(\vartheta\ge1/2\) as in that
proposition. For arbitrary \(\theta_i\ge0\) of total weight
\(L=\sum_i j_i\theta_i\le r^2/2\), one has
\[
 B_\gamma \le
 \frac{\exp(Cfkd+CfdL\log b)}{(k!)^{df}}\,
          Q^{-dfk} A_-^k\prod_{i=1}^t H_i^{-3-\theta_i}.
\]
All original finite places, the terminal covector level and
the restriction to the specified terminal character are
included.
\end{theorem}
\begin{proof}
Set \(G=|N_{F/\QQ}\mathcal G|\) and \(S=d(d-1)\).
For the original block scalars \(c_i\) put
\(b_i=2h_i(d-j_i)+\binom{h_i}{2}\).
The finite formula and the free-entry volume bound in
Theorem~\ref{thm:allfinite215} give
\[
 \frac{\mathcal M_\gamma^{=}}{\mathfrak m_{k-1}}
 \le \exp(O(f2^{-(k-2d)}))\,
     G^{k-1} Q^{f(2a-S-2d(k-1))}
                          \prod_i |N_{E/\QQ}c_i|^{b_i}.
\]
Here the \(Q\)-power includes the coincidence condition on
the two reduced flags. In multiplying the finite and real
bounds we keep this saving.

Put \(B_i^+=\prod_\tau|Q^{j_i}\tau a_i|\), with \(B_0^+=1\).
These are at least one since the \(Q^{j_i}a_i\) are nonzero
algebraic integers. At each block with \(b_i>0\) one has
\(|N_{E/\QQ}c_i|=(B_i^+/B_{i-1}^+)^{2/h_i}\).
Indeed this is the ratio of successive endpoint determinants.
Only a final singleton can have its pivot changed by forming
\(\Delta_+\), and then \(b_t=0\).
Telescoping gives the exact finite pivot factor
\[
        \prod_{i<t}(B_i^+)^{h_i+3h_{i+1}}\,
                       (B_t^+)^{h_t-1}.
\]
The exponents are nonnegative and sum to at most \(4d\).
Each \(B_i^+\le Q^{j_i f}H_i\). Also, for every \(p\)
with \(u_p>0\), the product of the absolute arithmetic pivots
up to \(p\), taken over all real places, is at least one.
Inside block \(i\) it is the geometric interpolation of
\(B_{i-1}^+\) and \(B_i^+\), because the pivots there have
equal moduli. Consequently the inverse between-block factors
in Proposition~\ref{prop:phase301} cost at most
\((2T)^{2f\sum_p p u_p}\). This argument also covers a
nonintegral individual block scalar.

In the product of the remaining real powers the factors
\(G^{-\kappa}\) cancel all but \(G^s\) in the finite mass,
since \(k-1-\kappa=s\). The power of \(\prod_\tau L_\tau\)
is \(2\kappa-E_0=\kappa-2ds-s^2\), between zero and \(k\)
on our intervals. Hence that power contributes at most
\(Q^{df(2\kappa-E_0)} A_-^k\).
To bound the residual \(G^s\), expand
\(\mathcal G=|\Delta_+|^2-2T\operatorname{Re}
(\Delta_+\overline{D_s})\). At a real place this is at most
\((1+2T)\max(1,|\tau\Delta_+|)^2\max(1,|\tau D_s|)\).
Thus \(G^s\) costs at most
\(\exp(O(fd^2\log b))H_t^{2s}\) times \(H_{t-1}^s\) when
\(D_s\ne0\). In that case \(j_{t-1}=s\).
When \(D_s=0\) the last factor can be omitted.

Combining these powers with the finite pivot factor leaves
positive exponents \(v_i\) on the \(H_i\) whose sum is at
most \(4d+3s\). Use \(\eta_i=v_i+3+\theta_i\) in
Proposition~\ref{prop:phase301}. Then
\(\sum_i j_i\eta_i\le O(d^2)+L\le r^2\). The cost of the
fixed \(O(d^2)\) in these tests fits in \(\exp(Cfkd)\).

The terminal class count contributes \((T^k-1)^{-f}\)
by Proposition~\ref{prop:tracemass301}. After including it,
the remaining exponent of \(T\), divided by \(f\), is
\[
       dE_0+2\sum_p p u_p-s(k+d)-k=O(d^3).
\]
Here \((1-T^{-k})^{-f}\le\exp(Cf)\). Since
\(\log T=O(d\log b)\) and \(d^3\log b=o(k)\),
the remaining \(T\)-cost is at most \(\exp(Cfkd)\).
The powers of two have the same bound. Adding the finite
pivot bounds, the \(Q\)-exponent divided by \(f\) is at most
\[
          -2d(k-1)+d(2\kappa-E_0)+O(d^2)
                     \le -dk+O(d^2).
\]
Its \(O(d^2)\) term also fits the stated exponential.
Finally Theorem~\ref{thm:complement209} gives
\(\mathfrak m_{k-1}/\mathfrak m_k\le e^{Cfk}(k!)^{-f}\),
since \(|\Delta_F|,\zeta_F(k)\ge1\).
Together with the \(s\) factorials per place in the real
integral this gives \(d\) factorials and proves the theorem.
\end{proof}

\subsection{The stabilizer volume of a single terminal flag}\label{nav:II-70}
Take the original fine row level \(K^0\). For a supported
transverse cell let \(\Gamma_D\) be its rational common
unipotent stabilizer in the monomial frame coordinates of
Proposition~\ref{prop:tracemass301}. Put
\(a_0=\sum_i\binom{h_i}{2}\) and
\(c_E=2^{-f}\sqrt{|\Delta_E|}\).

\begin{proposition}[A uniform covolume for counting individual cells]\leavevmode
\label{prop:covol304}
If \(\mathcal G>0\) and \(H_*<b^{Mfd}\), the original
unipotent covolume at every fine cell satisfies
\[
    \operatorname{vol}(\Gamma_D\backslash D_\infty)
                         \le c_E^{a_0}\exp(Cfd^6\log b).
\]
\end{proposition}
\begin{proof}
Choose primitive integral representatives of the two
decorations in the row \(\Lambda\), with cross matrix
\(Z=L M_0 V\) as in Lemma~\ref{lem:classes304}. Put
\(u=a\bar a\). The changes from the monomial frames to
these integral frames can be taken to be \(L^*,V\).
Indeed their cross matrix is \(Z\). Any initially given
pair of upper changes yielding the same \(Z\) differs
from this choice by a simultaneous stabilizer of \(M_0\).
It has inverse-adjoint paired actions on the two frame
spans and extends by the identity on their perpendicular.
It belongs to the common unipotent group in
Lemma~\ref{lem:stab205}. Inner changes of coordinates in
this group preserve its Haar measure. A transport of the
whole framed row simply transports the measure with it.
Thus covolume can be computed in these frame coordinates.

In monomial coordinates an element of \(D\) acts on the
first frame by a block diagonal matrix of upper unitriangular
blocks \(A_i\), on the second by the conjugate inverse
transpose transported through \(M_0\), and on \(C\) by the
identity. Inside each block the scalar in \(M_0\) cancels.
For a nonzero integral element \(l\), if \(A_i\equiv1\bmod l\)
entrywise and \(l\OO_E\) is conjugation invariant, then
the second matrix is also the identity modulo \(l\).
Both \(L^*,V\) have entries in \(u^{-1}\OO_E\) and their
inverses in \(u^{-(d-1)}\OO_E\). Conjugating the two actions
into the integral bases therefore gives the identity
modulo \(l u^{-d}\OO_E\).

Take \(l=Q\mathcal G\,u^d\), which is a nonzero integer
of \(F\). Put \(U_0=\Lambda\cap Y_0,\ U_1=\Lambda\cap Y_1\)
and \(U'=U_0\oplus U_1\). In the integral frames its Gram
determinant is \(\pm\mathcal G\). Thus \(\mathcal G(U')^\vee
\subset U'\) by the adjugate identity. The projection of
\(\Lambda\) onto \(Y_0+Y_1\) lies in \((U')^\vee\) by
pairing against \(U'\). For the elements just constructed,
the difference from the identity maps \(U'\) into
\(Q\mathcal G U'\) and hence maps \(\Lambda\) into
\(Q U'\subset Q\Lambda\). Its inverse has the same
property. Such an element preserves the lattice at
all places, including its modification above \(T\), and
is the identity on the residue row above \(Q\). It
therefore belongs to the original fine stabilizer.

The product of the upper groups with all off-diagonal
entries in \(l\OO_E\) has covolume
\((c_E|N_{E/\QQ}l|)^{a_0}\) in \(D_\infty\). In entry
coordinates Haar measure is product area. Successively
quotienting by the entries on each upper diagonal gives
the product of their additive lattice covolumes, since
the cross terms in multiplication involve earlier
diagonals. This subgroup is contained in \(\Gamma_D\).
The bound \(\log G=O(fd^2\log b)\) follows from the height
calculation in Theorem~\ref{thm:rationalrank302}.
Lemma~\ref{lem:classes304} gives
\(\log |N_{E/\QQ}l|=O(fd^4\log b)\), including
\(|N_{E/\QQ}\mathcal G|=G^2\).
Since \(a_0\le d^2\), the proposition follows.
\end{proof}

\subsection{A reversed block of arbitrary length}\label{nav:IV-52}
Fix \(2\le j\le s\). In the first \(j\) signed planes
choose integral Witt pairs as in Proposition~\ref{prop:plane281}.
Write \(P_j,P_j^*\) for their primal and opposite subspaces
and \(M_j\) for the perpendicular of these planes, of
dimension \(n=r-2j\). Let \(S_j\) denote the signed
permutation matrix whose antidiagonal reads
\(1,-1,\ldots,(-1)^{j-1}\). Put \(C=cS_j\). In block
order \(P_j,M_j,P_j^*\) consider the isometry
\[
  w=\begin{pmatrix}
     0&0&(C^*)^{-1}\\
     0&\lambda I&0\\
     C&0&0
  \end{pmatrix},\qquad \lambda\bar\lambda=1,
\]
where star on \(C\) is conjugate transpose. Use the
one-word orbital operator \(O_w\) in identity normalization.

\begin{proposition}[The common stabilizer and its unipotent quotient]\leavevmode
\label{prop:block295}
This representative covers the supported relevant labels
with a single reversed prefix block of length \(j<s\)
in Theorem~\ref{thm:suffix237}. It also covers the labels
with one reversed isotropic block in the full common
column case of Proposition~\ref{prop:class236} at \(j=s\).
In either case \(c=Q\xi\), \(0\ne\xi\in\mathcal O_E\),
and \(T\xi^j\lambda\) is integral. Put \(J=B\cap wBw^{-1}\).
Then \(B=J\ltimes U\), where \(U\) is the unipotent
radical for the whole \(j\)-plane \(P_j\).
\end{proposition}
\begin{proof}
The first \(j\) vectors pair across the two flags by \(C\),
and \(w\) multiplies the other vectors and the terminal
vector by \(\lambda\). The two cited classifications give
the rational double coset from these data. The last prefix
determinant divided by \(Q^j\) is a sign times \((c/Q)^j\).
Its integrality implies integrality of \(c/Q\) by taking
every prime valuation. For \(j<s\),
Proposition~\ref{prop:finite237} gives integrality even
of \(\xi^j\lambda\). For \(j=s\) the terminal congruence
in Proposition~\ref{prop:class236} gives the assertion
with \(T\).

An element of \(J\) preserves \(P_j\) and \(P_j^*\).
It preserves their nondegenerate sum and its perpendicular,
so it is block diagonal. Its primal block ranges over
upper unitriangular \(j\)-square matrices, with the
inverse adjoint on the opposite space. The signed
reversal conjugates lower unitriangular matrices into
upper unitriangular ones, so each such upper block
preserves the second prefix as well. On \(M_j\) the
two residual Bessel groups coincide because the second
frame is scalar times the first. Thus \(J\) is the
product of this prefix unitriangular group and the
residual Bessel group. The block diagonal projection of
\(B\) has image exactly \(J\), with kernel \(U\).
The exchange normalizes \(J\). It centralizes the middle
group. On the prefix group it takes inverse adjoints
and reverses the indices with unit signs. This preserves
Haar measure at every place, as can also be seen by
filtering upper entries by superdiagonal. Conjugation
by \(J\) preserves Haar on \(U\) because the prefix
determinants are one, the middle action is unitary and
the center transforms by congruence on the prefix.

Let \(U_L=U\cap\operatorname{Stab}(\Lambda_{\rm row})\)
at a finite place. It is contained in the actual row
level. Above \(Q\) its only nonidentity blocks go from
a later block to an earlier one across the first
\(j\)-plane or its opposite. These entries belong to
the row unipotent radical of Lemma~\ref{lem:flags211}.
At the graph and the other places the row level is the
full stabilizer. Passing to the diagonal blocks gives
\(B\cap K=(J\cap K)\ltimes U_L\). Hence the ratio of
finite unsigned mass to the identity mass is the count
of right \(U_L\)-cosets satisfying \(u_0^{-1}w\in KB\).
Indeed every supported first representative is \(b u_0K\)
with \(b\in J\), and the stabilizer ratio counts the
\(J\cap K\)-orbit of \(u_0U_L\), exactly as in
Theorem~\ref{thm:finite281}. Normalization of \(J\) by
\(w\) is sufficient for this argument.
\end{proof}

\subsection{The projected column lattice of a whole block}\label{nav:IV-53}
Use \(w\) from Proposition~\ref{prop:block295} at a finite
place. In either split component use the local integer ring
\(O\), and at an inert place use the integer ring of the
quadratic extension. Write \(P_0=\#(O/\pi)\) and use lengths
over \(O\). For a rectangular matrix \(M\), write \(t(M)\)
for the length of the image of \(z\mapsto Mz\bmod O^{u}\)
on the standard integral domain, where \(u\) is the number
of rows. This length is the sum of the absolute values of the
negative elementary divisors and is unchanged by transpose.
Let \(H,l,Z\) be as in Proposition~\ref{prop:blockfinite295}.

In the middle space of dimension \(n=k+2(s-j)+1\), write a
vector in the order \(m=(p,b,s_0)\), where \(b\) consists of
the \(k\) column coordinates. The coordinates \(s_0\) are
set to zero on taking the column subquotient, and the
coordinates \(p\) are then discarded. Put \(q=\dim s_0\).
In the first split component \(q=h=s-j\), since one sets the
residual opposite entries to zero and quotients by the
residual primal vectors and the terminal line. Use the same
description in the reversed dual component away from \(T\).
At \(T\) use the integral graph basis and the enlarged group
\(B'\) of Proposition~\ref{prop:affine293}. In the reversed
dual component also impose annihilation of the primal terminal
line, giving \(q=h+1\). These are the integral subquotients
in Lemma~\ref{lem:graphdepth293}. In particular the two
column coordinate blocks are dual to one another. At an
inert place use the decomposition with \(q=h\) and the
unimodular column form.

\begin{lemma}[An index calculation for the projected lattice]\leavevmode
\label{lem:blocklattice297}
Suppose \(g=u_0^{-1}w\in KB\), or \(g\in KB'\) for the
enlarged counting set at \(T\). Let \(A\) be the inverse
column action of the corresponding period factor and put
\(A_0=\lambda A\) in the chosen component. Set
\[
     X=t(H_W),\qquad Y=t(l_W),\qquad L=v(\lambda),
 \qquad m_-=t\bigl([\,l\ \ Z\,]\bigr),
\]
where \(W\) denotes the \(k\) column coordinates.
If \(d_-\) is the sum of the negative denominator lengths of
\(A_0 O^k\) relative to \(O^k\), then
\[
         X+Y-m_-+qL\ \le\ d_-\ \le X.
\]
There are at most \(j\) negative elementary divisors.
In the split case \(m_-=jv(c_-)\) in the first component.
In the inert case \(m_-=jv(c)\).
\end{lemma}
\begin{proof}
The inverse lattice is \(g^{-1}\Lambda\), with \(\Lambda\)
the row lattice in the chosen coordinates. Intersect it with
the space where the last \(j\) entries are zero and quotient
by the first \(j\) entries. Multiply the resulting middle
lattice by \(\lambda\). Inverting the exchange and applying
it after \(u_0\) as in Proposition~\ref{prop:blockfinite295}
shows that this scaled middle lattice consists of
\[
 m+Hv,\quad (m,v)\in{\cal T},\qquad
 {\cal T}=\{(m,v)\in O^n\oplus O^j:lm+Zv\in O^j\}.
\]
Indeed the omitted first input is forced to be \(-lm-Zv\)
to make the last output zero. The first output is unrestricted.
The volume of \({\cal T}\), with unit integral cubes, is
\(P_0^{-m_-}\). In the split case \(m_-\) is the length
of \(O^j/L_-\) in Proposition~\ref{prop:blockfinite295},
by transposing \([\,l\ \ Z\,]\) and taking its kernel modulo
the lattice. That proof uses only the two preserved prefix
planes and unit determinants and therefore applies to \(B'\)
as well. In the inert case the adjoint test is equivalent
to integrality of \(Hv,Z^\dagger v\) on the primal kernel,
since the middle form is unimodular and
\(lH-Z=Z^\dagger\). This gives the asserted length there.

Map \({\cal T}\) onto the \(s_0\)-projection of the scaled
middle lattice by \((m,v)\mapsto s_0+H_Sv\), with \(S\)
denoting those rows. Its image has Haar volume \(P_0^{-qL}\).
In fact
\(g^{-1}\Lambda=b_0^{-1}\Lambda\) for an element \(b_0\)
of the period group in use. After taking the first middle
subquotient, its action on the \(s_0\)-quotient is triangular
with all diagonal entries one. At \(T\) the extra dual
coordinate is a quotient on which the contragredient affine
group acts by one. Multiplication by \(\lambda\) contributes
the factor \(|\lambda|^q=P_0^{-qL}\), using the absolute
value normalized by \(P_0\).

Let \(V_0\subset O^j\) consist of \(v\) for which the zero
fiber of this map contains a point above \(v\). In that fiber
one must take \(s_0=-H_Sv\). At a fixed admissible \(v\),
the allowed \((p,b)\) form a coset of the kernel modulo \(O^j\)
of the matrix \([\,l_p\ \ l_W\,]\). Set
\(b_* = t([\,l_p\ \ l_W\,])\ge Y\).
The fiber volume in \((p,b,v)\)-coordinates is therefore
\(P_0^{-b_*-[O^j:V_0]_{\rm len}}\), where the bracket
denotes the length of the quotient. The change of additive
coordinates from \(s_0\) to \(s_0+H_Sv\) has determinant
one. Each nonempty fiber of the map on \({\cal T}\) is a
translate of its kernel, with this same volume. Fubini's
theorem now gives
\[
           [O^j:V_0]_{\rm len}=m_- -qL-b_*.
\]
The projected column lattice is \(A_0 O^k\), by the two
invariant subspaces defining the subquotient. Its image modulo
\(O^k\) is exactly \(H_W V_0\bmod O^k\), because the other
summand \(b\) is integral. The full image \(H_W O^j\bmod O^k\)
has length \(X\) and the loss on restriction to \(V_0\)
is at most \([O^j:V_0]_{\rm len}\). This proves both bounds.
The image is generated by at most \(j\) elements since
\(V_0\) is a lattice in \(O^j\).
\end{proof}

\subsection{The unipotent quotient for a composition of the prefix}\label{nav:IV-54}
Suppose \(j=j_1+\cdots+j_t\), where every \(j_i\ge1\).
Let \(J_i=\sum_{a\le i}j_a\). On \(1\le p\le j\) put
\(i(p)=\min\{i:p\le J_i\}\) and define the products
\[
       G_0=1,\qquad G_p=\prod_{a=1}^p c_{i(a)},\qquad
       u_p=j-J_{i(p)}.
\]
At an endpoint \(G_{J_i}\) agrees up to sign with its
leading cross determinant. These products are defined also
at intermediate positions, where a leading cross minor can vanish.
Use the ordinary single-word operators \(O_w\). The level and
the terminal frequency are those of Section~\ref{sec:missing}.

\begin{proposition}[The stabilizer and its quotient for several blocks]\leavevmode
\label{prop:composition298}
Every relevant supported pair with \(0<j=s-\dim(X_0\cap X_1)<s\),
and every transverse pair with full common column (\(j=s\)),
has a representative of the block form in
Proposition~\ref{prop:block295}, with
\[
           C=\bigoplus_{i=1}^t c_i S_{j_i}.
\]
Here \(a_i=G_{J_i}/Q^{J_i}\) are nonzero integers of \(E\),
\(Ta_t\lambda\) is integral, and \(\lambda\bar\lambda=1\).
Let \({\cal V}\) be the group of block upper unitriangular
\(j\)-square matrices over \(E\) with identity diagonal blocks
of sizes \(j_i\). Write \(\widehat z=\operatorname{diag}
(z,I,z^{-*})\). If \(U\) is the radical of the whole
\(j\)-plane, the unipotent quotient in \(J\backslash B\)
has unique coordinates \(\widehat z\,u_0\) with
\(z\in{\cal V},u_0\in U\). Its measure is product Haar.
\end{proposition}
\begin{proof}
Theorem~\ref{thm:suffix237} and Proposition~\ref{prop:class236}
classify the indicated pairs by the blocks, the cross
coefficients and the common scalar on the remaining frame.
The displayed isometry realizes these data on the fixed
rational Witt pairs. The supported endpoint minors divided
by the indicated powers of \(Q\) are integral by
Proposition~\ref{prop:shift203}. The terminal assertion comes
from Proposition~\ref{prop:finite237} for \(j<s\) and from
Proposition~\ref{prop:class236} for \(j=s\).

Set \(J=B\cap wBw^{-1}\). Its elements stabilize both
opposite \(j\)-planes and are therefore block diagonal in
the decomposition by these planes and the middle. On the
primal prefix a matrix in \(J\) is upper unitriangular.
Its inverse adjoint, transported by the signed reversals,
must also be upper unitriangular for the second prefix.
The order of the blocks is the same in the two lists,
whereas each list within a block is reversed. Consequently
the permissible upper entries are exactly those inside a
single diagonal block. Equivalently the two triangular
conditions on the block indices impose both upper and lower
triangularity, and impose just upper triangularity inside
each diagonal block. Thus \(J\) is the product of these
diagonal unitriangular groups and the residual middle
Bessel group. Conjugation by \(w\) preserves \(J\) and
its Haar measures, by inverse adjoint and reversal inside
each block and the identity on the residual group.

Every upper unitriangular prefix factors uniquely into its
block diagonal part and an element \(z\) of \({\cal V}\).
Splitting off \(U\) from \(B\) now proves the quotient
description. These coordinates have product measure since
the action of the prefix unipotents on \(U\) has determinant
one. All entries of \({\cal V}\) strictly above the block
diagonal have independent additive Haar measure, as follows
from triangular multiplication. At a finite place write
\({\cal V}_L\) for their integral group and \(U_L\) for the
integral radical. Both are in the actual row level,
including above \(Q\) by the block description of
Lemma~\ref{lem:flags211}. Taking diagonal blocks of integral
elements gives the level decomposition used in
Proposition~\ref{prop:block295}. In particular the finite
mass relative to the identity counts the supported right
\((\widehat{\cal V}_L U_L)\)-cosets of \(\widehat z u_0\).
For fixed representatives of \(z{\cal V}_L\), count
right \(U_L\)-cosets in the fiber. The compact column
groups used for averaging fix \(z\), since they act in
the middle block.

In these coordinates
\[
 \widehat z^{-1}w=w_{C(z),\lambda},\qquad C(z)=z^*C.
\]
The same identity holds after the common Euclidean Witt
normalization at infinity, which replaces \(C\) by \(C/2\).
Changing the rational opposite vectors is an upper central
change of frame as in Lemma~\ref{lem:blockphysical295}.
It acts trivially on the section character, and conjugation
has determinant one on the full upper unipotent. On the
additional coordinates \(z\) the change from rational
to unit primal vectors is uniform scaling of the entire
prefix and leaves its matrix unchanged. Normalizing each
of the \(\sum_pu_p\) extra additive coordinates to adelic
quotient volume one contributes the reciprocal of
\(2^{-f}\sqrt{|\Delta_E|}\). These factors cost at most
\(2^{f\sum_pu_p}\), since \(|\Delta_E|\ge1\).
\end{proof}

\subsection{The signed factor at an added split prime}\label{nav:VI-52}
Let \(c_j=jQ\sqrt{-q}\), with \(j\) in the set
\(\mathcal T\) of Theorem~\ref{thm:tail283}. At a place
\(v\) over a prime \(p\mid j\) we have \(F_v=\QQ_p\),
the row and column split, and the prescribed levels are the
full lattice stabilizers. Put \(K_H=\mathrm{GL}_k(\ZZ_p)\).
Normalize the finite factor by the identity orbital factor
as in Theorem~\ref{thm:phase282}. The additive character of
the ordinary seed is retained in both occurrences.

\begin{proposition}[A bound for the signed extra factors]\leavevmode
\label{prop:split287}
At \(v\mid p\mid j\) the finite factor between
\(K_H\)-invariant columns is of the form
\[
 F_v(c_j)=a'_v I+u'_v A_v,\qquad
 |a'_v|\le a_p,\qquad |u'_v|\le u_p ,
\]
where \(a_p,u_p\) are the unsigned multiplicities of
Proposition~\ref{prop:lattice282} with residue size \(p\).
In particular, writing \(R=r-1=k+2s\) and
\(H_x=4(k+s+1)^2\), its operator norm on the prescribed
word is at most \(H_xp^R\). If \(m_f(c_0)\) denotes
the unsigned finite mass for \(c_0=Q\sqrt{-q}\), the
full finite operator for \(c_j\) has norm at most
\[
           m_f(c_0)\,H_x^{\,f\omega(j)}j^{Rf},
\]
where \(\omega(j)\) counts distinct rational prime factors
and \(\omega(1)=0\).
\end{proposition}
\begin{proof}
The first split component of \(c_j\) has valuation one.
Thus Theorem~\ref{thm:split281} and
Proposition~\ref{prop:lattice282} apply to its actual unit
part. In the unipotent fiber the quotient mass counts each
successful integral right coset once. On any one such coset
the two column arguments have relative position equal to
the identity double coset or to the neighbor double coset.
This follows by projecting the two integral lattices onto
the middle column as in that proposition. The character
multiplier has modulus one. Changing an integral representative
leaves the correlation on spherical vectors unchanged because
the column compact factors act trivially there and the
unipotent lattice character is trivial on integral entries.
These entries are integral for the physical unit \(\varpi\)
at \(p\). Equivalently one can average over choices of
representatives, in which case the multiplier has modulus
at most one.

For a unitary translation action \(R_H\), let \(e_H\) denote
the orthogonal average over \(K_H\). The operators
\(e_H R_H(h)e_H\) agree for all \(h\) in one double coset
by Haar invariance on its two sides. In the neighbor case
this operator equals \(A_v/[d_p(d_p-1)]\) on spherical
vectors, by summing the \(d_p(d_p-1)\) right cosets.
Each distinct neighbor has unsigned multiplicity \(u_p\), and
the identity has total multiplicity \(a_p\). Summing the
possible phases now gives the two scalar inequalities
without any condition on their signs. Spherical convolution
at other places commutes with these operations.

The explicit counts are
\[
 a_p=p^{2s-2}(p-1)\{(2+2s(p-1))p^k-1\},\qquad
 u_p=p^{2s-2}(p-1)^2 .
\]
They give \(a_p\le(2+2s)p^{k+2s}\) and \(u_p\le p^{2s}\).
Use the neighbor estimate of Proposition~\ref{prop:satake287}
and the triangle inequality to obtain \(H_x p^R\).

At every place not dividing \(j\), the unchanged valuations
give the same unsigned upper bound as at \(c_0\), by
Theorems~\ref{thm:finite281} and \ref{thm:product281}.
Weight one there at a minimal valuation and the full
unsigned mass at the places over \(q\) bound the operators
on the entire finite column. More explicitly, the finite
pairing disintegrates into the unipotent choices and a
common integration over the column quotient, as in
Theorem~\ref{thm:phase282}. Each tuple of local choices
gives the commuting product of the corresponding local
translations and phases. Averaging on spherical vectors at
the new places therefore inserts exactly the local factors
just bounded. There are \(f\) places of \(F\) of residue
size \(p\) above each \(p\mid j\), since \(p\) splits
completely in \(E\). Multiplying gives the global bound.
Compression at the end can only decrease the operator norm.
\end{proof}

\subsection{A local bound at every power of the terminal prime}\label{nav:VI-53}
\begin{proposition}[The first-plane factor on the original terminal graph]\leavevmode
\label{prop:terminal293}
At each \(v\mid T\) with \(e=v(c)=v(\xi)\ge1\), the signed
first-plane operator between the \(C_\ell\)-fixed columns
satisfies
\[
        \|F_v(Q\xi)\|\le [4(k+1)^2]^e (Nv)^{(k+4s)e}.
\]
At \(e=0\) the factor is the identity. The bound is
uniform in the unit part of \(\xi\) and the primitive
residue covector specifying the graph.
\end{proposition}
\begin{proof}
Use the counting normalization of
Theorem~\ref{thm:finite281}. The original summands belong
to the finite set \(\Omega'\) of
Proposition~\ref{prop:affine293}. Insert zeros at all the
other cosets in this set. The same proposition bounds the
average on each orbit by \(\Xi(A)\), with no compact index.
The elementary divisors of \(A\) are constant on this
orbit, since conjugating \(\beta\) by \(D_0\) and changing
the compact representative can multiply \(A\) only by
elements of \(C\) on either side. By
Lemmas~\ref{lem:tensor293} and \ref{lem:graphdepth293},
the bound for an orbit of size \(m_0\) is at most
\(m_0 k^{2e}P^{-(k-1)d}\), where \(P=Nv\). Equivalently
assign the same weight \(k^{2e}P^{-(k-1)d}\) to each
member when summing this majorant. The depths \(t,u'\)
in the last lemma may be read separately on each member.

For a fixed pair of depths there are at most
\[
       P^{k(2e-t-u')+(2(n-k)+1)e}
\]
members, with \(n=r-2=k+2s-1\). Indeed \(X,Y\) modulo
\(c\) are a column and row over \(O/cO\) in the graph
basis. Their \(k\) column coordinates on each side give
the first exponent. Their other coordinates give
\(2(n-k)e\). After fixing representatives for those two
blocks, the central variable allows at most \(P^e\) choices
by the right integral multiplication used in
Proposition~\ref{prop:marked291}. Both first endpoints
are integral in the enlarged set, which suffices for
this count.

Let \(j=e-t-u'\). The remaining part of the exponent is
\[
      k(2e-t-u')-(k-1)d
       \le k(e+j)-(k-1)\max(0,j)\le (k+1)e .
\]
Since \(2(n-k)+1=4s-1\), this gives the stated power of
\(P\). The number of depth pairs is at most
\((e+1)^2\le4^e\). At \(e=0\),
Theorem~\ref{thm:finite281} gives the unique integral
coset, and the representative \(1\) has
\(\gamma_c\in K\) and trivial relative character. The
argument applies on each prescribed local constituent
and hence on their automorphic sum.
\end{proof}

\subsection{High powers of the physical block parameter}\label{nav:VI-54}
\begin{lemma}[Two physical-frequency estimates for one block]\leavevmode
\label{lem:blockphysical295}
For \(j,n\) as in Proposition~\ref{prop:block295}, let
\(R=r-1\), choose a real \(1\le\ell\le j\), and take
\(A=3\ell R\) or \(A=6\ell R\). At a real
place \(\tau\), the positive local orbital integral divided
by its diagonal norm, with ordinary coordinate measure on
the block radical, is at most
\[
     \exp(C Rj\ell s\log b)\,|\tau c|^{-jA}.
\]
The field factor for the product in rational quotient
measures is at most one.
\end{lemma}
\begin{proof}
First use unit terminal frequency. The upper central
unipotent that replaces each integral opposite vector
by its Euclidean orthogonal Witt partner fixes the
isotropic flag. It is trivial on the Bessel character,
as it has only entries from opposite to primal isotropic
vectors. It commutes with the middle column group.
The Euclidean Witt pairs can be used from the outset
because their opposite vectors differ from the
rational ones by multiples of the corresponding primal
vectors. Changing the representatives by this upper
central element in both flags transports their stabilizer
by conjugation and preserves the quotient integrals.
In the unit pairs the matrix in
Theorem~\ref{thm:matrixfiber295} is \(C_0=\tau C/2\).

On \(J_\infty\backslash G_\infty\) use Cauchy--Schwarz
with squared-integral weight
\((l_j(g)/l_j(w^{-1}g))^A\) and its reciprocal, as in
Lemma~\ref{lem:power291}. The first squared section norm
descends to \(B_\infty\backslash G_\infty\). Its weight
disintegrates over \(U\) by Proposition~\ref{prop:block295}
and Theorem~\ref{thm:matrixfiber295}. The remaining
diagonal moment has power \(2A-2(n+j)\) on \(l_j(g)\).
For the second squared norm, change variables by \(w\).
The exchange normalizes \(J\) with modulus one, and its
inverse has cross block \(C_0^*\) of the same determinant
modulus. Thus both squared-integral bounds agree.
These positive disintegrations are justified by Tonelli.
Since \(B\) fixes the exterior product of the first \(j\)
unit isotropic vectors,
\(l_j(g)\le\nu(g)^j\), with \(\nu\) the minimum norm of a
coset representative. For these powers
\[
       0<j(2A-2(n+j))\le 12j\ell R\le12j^2R\le4r^2
\]
for sufficiently large \(x\), since \(j\le s\) and
\(s^2=o(r)\) by Theorem~\ref{thm:rangenew50}. The
local moment in Theorem~\ref{thm:moment286} costs at
most \(\exp(CRj\ell s\log b)\). For a nonintegral power
use the next integer, since \(\nu\ge1\).
The integrand defining
\(K_{n,j,A}\) decreases with \(A\), and \(A\ge2n+3j\)
on these intervals. At \(\alpha=n+3j/2\) the gamma
quotients in Theorem~\ref{thm:matrixfiber295} are at
most one for these \(n\). Thus
\(K_{n,j,A}\le\pi^{nj+j^2/2}\), also within this cost.

At terminal frequency \(\varpi_\tau\), the scaling that
multiplies the primal flag vectors by \(\varpi_\tau\)
changes \(C_0\) into \(C_0/T\). Its modulus on \(U\)
is \(T^{j(n+j)}\). Indeed the rectangular coordinates
receive one complex scalar of modulus \(\sqrt T\) per
entry in addition to a middle unitary matrix of determinant
of modulus one, and the skew center receives scalar
\(T\) in \(j^2\) real coordinates. In the ratio of the
\(J\)- and \(B\)-quotient densities this contributes
\(T^{-j(n+j)}\), by the change of variables in
Proposition~\ref{prop:scale284}. The full physical
factor contributed by the exchange is consequently
\[
             (2/|\tau c|)^{jA} T^{j(A-n-j)} .
\]
Use \(\log T=O(s\log b)\) from Theorem~\ref{thm:prime210}.

Finally choose rational compatible measures by giving
\(U_L\) volume one at all finite places and \([U]\)
adelic volume one. The rectangular additive lattice gives
\(nj\) copies of complex covolume
\(2^{-f}\sqrt{|\Delta_E|}\), as in
Proposition~\ref{prop:measure283}. Graph changes on a
split pair have opposite determinant valuations and the
total covolume is unchanged. The center contributes
\(j(j-1)/2\) further complex coordinates and \(j\) skew
scalar coordinates of covolume \(\sqrt{\Delta_F}\) each.
The central extension has the product covolume by
integration first on the center, since the rectangular
projection of every integral radical is onto. Passing
to Euclidean unit Witt pairs scales the rectangular
entries by \(\sqrt2\) and the central entries by \(2\),
with real Jacobian \(2^{nj+j^2}\) at each place. Thus
the field factor multiplying the ordinary coordinate
densities is the reciprocal of
\[
 (2^{-f}\sqrt{|\Delta_E|})^{nj+j(j-1)/2}
       (\sqrt{\Delta_F})^j 2^{(nj+j^2)f}.
\]
This denominator is at least one. The normalization gives
\([B]\) and \([J]\) the same volume, since \(B=J\ltimes U\).
The common compact real column contracts the two factors
by Schur orthogonality as in Theorem~\ref{thm:real282}.
Its dimension factor cancels against the diagonal. On the
other part of \(J\) the two actions are scalar characters
and coincide for relevant labels. This establishes the
stated product normalization for the positive majorant.
\end{proof}

\subsection{The split operator for a whole reversed block}\label{nav:VI-55}
At a split place let \(P=Nv,\ m_\pm=jv(c_\pm)\),
\(M=m_++m_-\), \(h=s-j\), \(q_0=n-k=2h+1\), and
\(L=v(\lambda_+)\). On the original local column let \(F_v\)
be the signed orbital operator relative to the identity.

\begin{proposition}[All split places, including the terminal graph]\leavevmode
\label{prop:blocksplit297}
Uniformly for \(2\le j\le s\),
\[
       \|F_v\|\le [2(k+1)]^M P^{(k/2+32(s+1))M}.
\]
\end{proposition}
\begin{proof}
Use \(A_0=\lambda_+ A\) and \(X,Y\) of
Lemma~\ref{lem:blocklattice297} in the first component.
Write \(d_+,d_-\) for the sums of the positive exponents
and of the absolute negative exponents of \(A_0\), respectively.
In the reversed dual, the rectangular blocks become
\(-H^{\mathsf t},-l^{\mathsf t}\) in exchanged positions,
the two prefix lengths interchange, and the scaled output
lattice has matrix \(A_0^{-{\mathsf t}}\).
Applying that lemma in both components gives
\[
 \begin{aligned}
 X+Y-m_-+hL &\le d_-\le X\le m_+,\\
 X+Y-m_+-h'L &\le d_+\le Y\le m_- ,
 \end{aligned}
\]
where \(h'=h+1\) over \(T\) and \(h'=h\) elsewhere.
The upper bounds for \(X,Y\) follow from the kernel lattices
\(L_+,L_-\) of Proposition~\ref{prop:blockfinite295}.
The determinant of \(g=u_0^{-1}w\) has valuation
\(m_+-m_-+nL\). All period graded pieces outside the column
have determinant one. Scaling \(A\) on the column gives
\[
 d_+-d_-=m_--m_+-q_0 L,\qquad q_0|L|\le2M .
\]
The second inequality uses \(d_-+d_+\le X+Y\le M\).
Put \(D=d_-+d_+\). For \(L<0\) add the determinant valuation
to twice the lower bound for \(d_-\). For \(L\ge0\) subtract
it from twice the bound for \(d_+\). Away from \(T\) this gives
\[
        D\ge 2(X+Y)-M+|L|.
\]
At \(T\) it gives \(D\ge 2(X+Y)-M-L\) for both signs.

We specify how the phases and original levels enter. The
finite mass is the count of right \(U_L\)-cosets by
Proposition~\ref{prop:block295}. At a full column level their
translations are compressed to compact invariants. Above
\(Q\) write \(KB=K_1B\) as in Proposition~\ref{prop:markedpowers292}.
Conjugation by \(C=\mathrm{GL}_k(O)\) stabilizes \(K_1,B,U_L\)
and commutes with \(w\), since the middle block of \(w\) is
scalar. Two decompositions of the same matrix modulo \(K_1\)
change its column action by the principal congruence subgroup,
which fixes each original vector. Thus averaging coset terms
on each \(C\)-orbit averages the conjugates of \(A\) with
phases of modulus at most one. Temperedness at these places
is Proposition~\ref{prop:tempered292}. The regular representation
and compact averaging proof of Lemma~\ref{lem:conjugation292}
bounds this average by \(\Xi(A)\) before using any exponents
of \(A\). This also bounds every full-level compression on a
tempered constituent. The spherical coefficient is unchanged
on multiplication by a central scalar, because its inducing
character is trivial on the center. Hence \(\Xi(A)=\Xi(A_0)\).

At \(T\) use instead the \(D_0\)-orbits of the enlarged
cosets \(u_0^{-1}w\in KB'\) in Proposition~\ref{prop:affine293}.
This set is finite by the kernel count of
Proposition~\ref{prop:blockfinite295}, which still applies
since \(B'\) preserves both prefix determinant tests.
The scalar middle block again commutes with \(D_0\), which
preserves the integral radical. The original success test and
the relative column translation are obtained by factoring the
central block \(\beta\) into the affine and stabilizer factors
in that proposition. Its change-of-covector formula thus applies
to every such orbit with the actual phases, extending vectors by
zero outside the original shell. The composition by the affine
map with linear part \(A\) has operator norm at most
\(|\det A|^{-1/2}\) by change of variables. The compact
tensor estimate in Lemma~\ref{lem:tensor293}, whose proof
gives \(\Xi(A)\) for every \(A\), bounds the orbit average by
\[
              P^{(m_- -m_+-nL)/2}\Xi(A_0).
\]
The determinant and \(\Xi(A_0)\) are constant on each
averaging orbit. Its size can therefore be counted by giving
this bound to each of its members.

For a fixed \(X,Y\) count the contributing cosets. The two
prefix kernels have at most \(P^{jM}\) choices by
Proposition~\ref{prop:blockfinite295}. A rectangular \(k\)-by-\(j\)
matrix with image length \(X\) modulo the lattice has at most
\(P^{(k+j)X}\) possible residues. Its kernel in \(O^j\) has
index \(P^X\) and at most \(P^{jX}\) choices, and each row
on that kernel allows at most \(P^X\) residues. Apply this
to \(H_W\) and to \(l_W^{\mathsf t}\). The other rectangular
rows cost at most \(P^{q_0M}\), and the central rows cost
at most \(P^{j m_-}\). As in the prefix count, choose fixed
rectangular lifts by right integral multiplication before
counting the center. Altogether the count is at most
\[
                       P^{k(X+Y)+(3j+q_0)M}.
\]
There are at most \(2^M\) length pairs, since
\((m_++1)(m_-+1)\le2^M\).

At most \(j\) elementary divisors of each sign occur by
Lemma~\ref{lem:blocklattice297}. Put \(v_0=k-4j+3>0\).
The spherical factor is at most
\(k^M P^{-v_0D/2}\) by Lemma~\ref{lem:blocksphere297}.
Away from \(T\) its combination with \(P^{k(X+Y)}\) has
exponent at most \(v_0M/2+(k-v_0)M\le(k/2+4j)M\).
Above \(T\) the extra exponent from inserting the determinant
and using the local lower bound for \(D\) is at most
\[
                  M/2+(v_0-n)L/2 .
\]
For \(L\ge0\) the second term is nonpositive. For \(L<0\)
it is at most \((q_0+4j)|L|/2\le(1+4j)M\).
These estimates and \(q_0=2(s-j)+1\) give the stated
bound, including \(M=0\). An empty original support has
zero operator.
\end{proof}

\subsection{Summation over the rational flag parabolic}\label{nav:VI-16}
\label{sec:parabolic129}

Let \(P\) stabilize the isotropic flag used to define \(N\).
We show that enlarging the summation from \(B(F_+)\) to
\(P(F_+)\) preserves the full norm of every seed of
Section~\ref{sec:identity127}. This evaluates an infinite family
of rational cosets. Write \(\mathbb A=\mathbb A_{F_+}\).
Throughout, either of the two finite levels in that section can
be used. In particular it is the original level at the auxiliary
primes \(g_1,g_2\), with unipotent reduction at both.

\begin{lemma}[Character transport compatible with finite support]
\label{lem:transport129}
Let \(\lambda\in P(F_+)\) and
\(\beta,\beta'\in\{1,\varpi\}\). Suppose that
\[
   \chi_\beta(\lambda n\lambda^{-1})=\chi_{\beta'}(n)
                 \quad(n\in N(\mathbb A)),\qquad
                     \lambda_f\in B_f K B_f .
\]
Then \(\beta=\beta'\) and \(\lambda\in B(F_+)\).
\end{lemma}
\begin{proof}
Use the Witt splitting from Section~\ref{sec:char94}. On the
one-dimensional flag quotients the Levi component acts by
\(a_1,\ldots,a_s\in E^\times\). Its action on the paired quotients
is by \(\bar a_i^{-1}\), and its middle component is \(M\in U(W)(F_+)\)
for \(W=W_0\perp Ee\). This decomposition follows by reading the
diagonal blocks of a flag-preserving isometry. Removing them leaves
an element of \(N(F_+)\), whose conjugation fixes every character
of \(N(\mathbb A)\).

The first \(s-1\) simple coordinates acquire factors \(a_i/a_{i+1}\).
Each can be tested on the elementary change
\(u_{i+1}\mapsto u_{i+1}+t u_i,\ v_i\mapsto v_i-\bar t v_{i+1}\).
Varying \(t\) in one complex completion and using the nondegenerate
real trace pairing forces \(a_i/a_{i+1}=1\). Write \(a\) for the
common value, also when \(s=1\).
To test the last simple block, take any linear form \(\xi\) on
\(W\) at a complex place. Write \(\xi(w)=h(z,w)\). There is a
unipotent isometry given on the last Witt plane and \(W\) by
\[
 w\longmapsto w+\xi(w)u_s,\qquad
 v_s\longmapsto v_s-z-\tfrac12 h(z,z)u_s,\qquad u_s\longmapsto u_s.
\]
The middle pairings are preserved, the cross terms with \(v_s\)
cancel, and the last image has norm zero. This element has the
specified last block and all preceding simple coordinates zero.
Conjugating the block by the Levi replaces \(\xi\) by
\(a\,\xi\circ M^{-1}\). The character equality for all \(\xi\)
therefore gives \(\beta a M^{-1}e=\beta'e\). These are equalities
of rational vectors, as can already be checked in a complex
embedding. Since \(h(e,e)=-1\) and \(M\) is unitary, we get
\begin{equation}
            a\bar a=\frac{\beta'\bar\beta'}{\beta\bar\beta}.
            \label{eq:norm129}
\end{equation}

In a finite decomposition \(\lambda_f=b_1 k b_2\), the middle
element \(k\) belongs to \(K\cap P_f\). The two elements of \(B_f\)
act trivially on every one-dimensional quotient of the isotropic
flag. Thus the image of \(k\) on each such quotient is \(a\).
At every finite place the image of the compact subgroup
\(K_v\cap P(F_{+,v})\) in the multiplicative group of a field has
valuation zero. This holds in both factors at split places.
It follows that \(a\) is a global unit. The two different terminal
frequencies have norms \(1,T\), so \eqref{eq:norm129} forces
\(\beta=\beta'\).

Now \(a\bar a=1\), and every conjugate of the algebraic integer
\(a\) has modulus one. Thus \(a\) is a root of unity. Indeed the
monic polynomials with roots the conjugates of \(a^j\), counted
over the embeddings of \(E\), have integral coefficients bounded
independently of \(j\). Their roots comprise a finite set.
At each auxiliary prime any element of the row level has
unipotent reduction. Its characteristic polynomial in the
standard lattice reduces to \((X-1)^r\). Applying this to the
above \(k\) shows that \(a\) reduces to one at places above both
\(g_1\) and \(g_2\). The order of a root of unity reducing to one
in residue characteristic \(p\) is a power of \(p\). Indeed its
prime-to-\(p\) part \(z\) still reduces to one, and, if it had
order \(d>1\), the equality \(1+z+\cdots+z^{d-1}=0\) would reduce
to \(d=0\) in the residue field. The two distinct auxiliary primes
therefore give \(a=1\). We now have \(Me=e\), so \(M\) acts on
\(W_0=e^\perp\cap W\) by an element of \(H(F_+)\). The Levi and
unipotent components together put \(\lambda\) in \(B(F_+)\).
\end{proof}

\begin{theorem}[The full parabolic sum preserves the seed norm]
\label{thm:parabolic129}
For the seeds of \eqref{eq:seed127} define the parabolic sum by
\begin{equation}
      (\mathcal P^{P}\Phi)(g)
       =\sum_{\gamma\in B(F_+)\backslash P(F_+)}\Phi(\gamma g)
            \quad\hbox{in }L^2\hbox{ along the unipotent quotient}.
       \label{eq:sum129}
\end{equation}
It defines an element of \(L^2(P(F_+)\backslash G(\mathbb A))\).
For any finite combination of the chosen types, frequencies and
columns its squared norm there equals the squared norm of the
input on \(B(F_+)\backslash G(\mathbb A)\). In particular its
Gram operator on the refined zero-sum columns is precisely the
scalar in \eqref{eq:main127}. Distinct words and distinct terminal
frequencies have orthogonal images.
\end{theorem}
\begin{proof}
Fix \(g\) and representatives for the countable rational coset set.
As a function of \(n\in N(F_+)\backslash N(\mathbb A)\), the term
with index \(\gamma\) and frequency \(\beta\) is
\[
      \Phi(\gamma n g)=
        \chi_\beta(\gamma n\gamma^{-1})\Phi(\gamma g).
\]
Here the transformation law holds on the entire group, since left
translation by \(N_f\) preserves the finite support. If two terms
are both nonzero, their finite supports imply
\((\gamma\gamma'^{-1})_f\in B_f K B_f\).
If their characters on \(N(\mathbb A)\) agree, conjugating the
variable by \(\gamma'\) and using Lemma~\ref{lem:transport129}
gives equal frequencies and equal rational cosets.
The distinct characters are orthogonal on the compact unipotent
quotient. Right translation by an element on which their ratio
differs from one proves this directly.

For a finite sum of inputs, first collect all types and columns
at the same frequency into \(\Phi_\beta\). Put
\[
    D(g)=\sum_{\gamma\in B(F_+)\backslash P(F_+)}
                  \sum_{\beta} |\Phi_\beta(\gamma g)|^2 .
\]
The nonnegative function \(D\) is invariant on the left by
\(P(F_+)\) and by \(N(\mathbb A)\). Unfolding by a fundamental
region for the discrete group \(P(F_+)\) and using Tonelli gives
\[
 \int_{P(F_+)\backslash G(\mathbb A)}D(g)\,dg
   =\sum_\beta\int_{B(F_+)\backslash G(\mathbb A)}
                              |\Phi_\beta(g)|^2\,dg<\infty
\]
by Theorem~\ref{thm:identity127}. Thus almost every unipotent
fibre has square-summable coefficients. Orthogonality proves
convergence of the series there in \(L^2\), independently of the
ordering. Its squared norm on that fibre with Haar probability
is \(D(g)\).

Here the base of the fibres is
\(P(F_+)N(\mathbb A)\backslash G(\mathbb A)\). The subgroup in
this quotient is closed because it is the inverse image in
\(P(\mathbb A)\) of the closed discrete rational Levi under the
map to the adelic Levi. It is unimodular, since conjugation by the rational
Levi on \(N(\mathbb A)\) has modulus one by the product formula
for the determinant of the adjoint action on the unipotent Lie
algebra. Haar disintegration therefore uses the invariant measure
on \(N(F_+)\backslash N(\mathbb A)\) in each fibre, with the base
measure scaled to allow Haar probability on that fibre. On a
measurable section of the quotient the partial sums are measurable,
and their fibrewise \(L^2\) limits give a measurable \(L^2\)
function. Equivalently, the integrated squared tails tend to zero
by dominated convergence with dominating function \(D\) on that
section. Changing a representative translates the fibre or
permutes the rational indices and conjugates the unipotent
variable by a rational Levi element. These operations preserve
the fibre measure and the sum. This defines \eqref{eq:sum129} on
the stated quotient.

Integrating its squared norm gives the displayed unfolded
integral. The different \(\Phi_\beta\) are orthogonal on the
input side by Theorem~\ref{thm:identity127}, which also gives
the mixed inner products of the types within one frequency.
Polarization now proves all the Gram assertions.
\end{proof}

The intermediate sum \(\mathcal P^P\) includes every rational
element preserving the isotropic flag. To obtain automorphic row
forms on \(G(F_+)\backslash G(\mathbb A)\) one must also sum over
\(P(F_+)\backslash G(F_+)\). The terms from different such cosets
couple different flags, whose unipotent radicals need not agree.
The orthogonality proof above uses their common radical at the
parabolic stage. The estimates for pairs of different flags and
the comparison on the integral theta lattice remain open here.

\subsection{One normalization for the identity term on the common interval}\label{nav:VI-27}
Keep the shortened field \(E\), the interval \([V,V^{3/2}]\), and
the common prime \(T\) and word family of Theorem~\ref{thm:common111}.
Put \(f=[E^+:\QQ]\), \(s=m/2-1\), \(r=2s+k+1\). We use the
probability measure on the compact column quotient \([H]\).
Here an identity term is the contribution to the Gram pairing
of two Poincar\'e sums from equal rational cosets in the two
summation indices. Its computation at the original and translated
row levels is given in
Section \ref{nav:VI-15}.

\begin{theorem}[Common normalization in the probability norm]
\label{thm:commonnorm128}
At each of the two row levels \(K\) the identity term is
\(C_K\) times the probability inner product on the multiplicity
space of each selected column type at its induced column level.
The same positive \(C_K\) serves all the words and both terminal
frequencies. At the translated level it also serves the zero-sum
refinement. The mixed identity term between two different
word-frequency pairs vanishes. Writing the translated level
as \(K_\sigma'\), one has
\begin{equation}
 \frac{C_{K_\sigma'}}{C_{K_\sigma}}=a_{x,T}^2,\qquad
 a_{x,T}=T^{\,fs(r-s)/2}(T^k-1)^{f/2}\in\ZZ_{>0}.
 \label{eq:ratio128}
\end{equation}
Thus the positive square root needed to compare these two
identity-term norms belongs to \(\QQ\) on the original field family.
\end{theorem}
\begin{proof}
Write \(M_K\) for the number of finite column classes and
\(d_\rho\) for the dimension of a column representation \(\rho\).
Fix one unit vector \(w\) in it. In the class of \(h_i\), the scalar
function corresponding to the class covector \(\ell_i\) has
value \(\ell_i(\rho(h_\infty)w)\). The classes have equal mass
\(1/M_K\) and trivial stabilizers. Averaging the projection onto
\(\CC w\) over \(H_\infty\) gives the identity divided by \(d_\rho\)
with probability Haar measure, by Schur's lemma and the trace.
The squared norm of the column is therefore
\[
        (M_K d_\rho)^{-1}\sum_i\|\ell_i\|^2 .
\]
The same formula for inner products follows by polarization.
Varying the covectors gives one copy of the full multiplicity
space, since evaluation at a fixed nonzero vector on the irreducible
type is injective on the space of intertwining maps and its
dimension is \(M_K d_\rho\).

With the Haar measures of Theorem \ref{thm:identity127}, the
coefficient multiplying the probability inner product is
\[
       C_K=M_K\,\frac{\operatorname{vol}_f(K)\,v_N\,
                  \operatorname{vol}(H_\infty)}
                 {\operatorname{vol}_f(N_f\cap K)} .
\]
In particular the factor \(d_\rho\) has cancelled. The stated
mixed terms vanish by that theorem, and restriction to the
zero-sum subspace preserves the formula. Each unrefined class
has \((T^k-1)^f\) refinements by Corollary~\ref{cor:rank111}.
Proposition \ref{prop:vol127} gives the unipotent index
\(T^{fs(r-s)}\), and the row volumes agree. Taking the ratio
gives \eqref{eq:ratio128}. Finally \(E^+\) contains
\(C_{l'}=\QQ(\zeta_{2^{l'+2}})^+\) of degree \(2^{l'}\).
Our choice has \(2^{l'}\ge x^{2+8\eps}/\operatorname{ord}(\psi)\),
so \(l'\ge1\) for all sufficiently large \(V\).
The degree \(f\) is even. Both exponents in \(a_{x,T}\) are
integers.
\end{proof}

\begin{corollary}[Integral cost and normalized trace on the same fields]
\label{cor:cost128}
Let \(D_{\rm aux}\) be the denominator of the auxiliary finite
Whittaker pairing on its integral residue bases. For every fixed
\(A>0\), uniformly at all the points of the selected intervals,
\begin{equation}
 \log(D_{\rm aux}a_{x,T})=o(b/x^A),\qquad b=\exp(x^\eps).
 \label{eq:height128}
\end{equation}
If \(J^{\rm id}_{c,\beta}\) denotes the Gram operator of the
identity term for word \(c\), restricted to the refined zero-sum
multiplicity space with the probability norm, then
\begin{equation}
 \frac{\Gamma_c^2}{D\,C_{K_\sigma'}}
           \operatorname{tr}(J^{\rm id}_{c,\beta})
    =1+O\!\left(\exp(-x^{1/4})\right)
           \qquad(\beta=1,\varpi).
 \label{eq:tr128}
\end{equation}
Here \(D=\mathcal V(G)I_\sigma\prod_\tau W_r(\lambda_\tau)\)
is the row quantity of Theorem~\ref{thm:common111}.
\end{corollary}
\begin{proof}
Corollary \ref{cor:content125} gives
\[
 D_{\rm aux}=\prod_{q\in\{g_1,g_2\}}
                  \left(\prod_{j=2}^{k}(q^j-1)\right)^f .
\]
Its logarithm is at most \(fk(k+1)\log(g_1g_2)/2\).
Equation~\eqref{eq:ratio128} gives
\(\log a_{x,T}\le f(s(r-s)+k)\log T/2\).
The ranks and \(f\) are bounded by powers of \(x\),
\(g_i\le2(kb)^8\), and \(\log T\) is bounded by a power of \(x\).
For the last bound apply Lemma~\ref{lem:size97}. The lists of
Theorem~\ref{thm:common111} satisfy its parameter bound, and all
the chosen primes lie in its terminal range. Thus the sum of the
two logarithms is bounded by \(x^{C_{\psi,\eps}}\) for large \(x\).
For any fixed \(A,C\), \(x^{C+A}e^{-x^\eps}\) tends to zero,
which proves the asserted uniform bound.

Theorem~\ref{thm:commonnorm128} gives
\(J^{\rm id}_{c,\beta}=C_{K_\sigma'}\mathrm{Id}\).
The dimension of the space is \(R_c^\circ\).
Substitution of its estimate in Corollary~\ref{cor:rank111}
proves \eqref{eq:tr128}.
\end{proof}

This computes the finite denominator and the change in the
identity-term norm on the same fields, with their costs in
\eqref{eq:height128}. A vector may have a positive identity
contribution and zero projected Poincar\'e sum, because the
remaining rational cosets can cancel it. Controlling those terms
and identifying the resulting value on the integral theta lattice
are still required for the period estimate in this section.

\section{Split and inert Hecke operators}\label{nav:g-IX-3}

\subsection{Counting the two decorated flags at a split prime}\label{nav:VIII-48}
We compute the next valuation in Theorem~\ref{thm:finite281}.
Let \(v\) be a split place outside \(QT\), with full row
lattice stabilizer and residue field \(\kappa\) of size \(b\).
Use \(k=\dim_E W_0\), \(h=s-1\), and \(n=k+2h+1\), the
dimension of \(M\) on the first split factor.

\begin{theorem}[Exact local mass when the pairing has valuation one]\leavevmode
\label{thm:split281}
If \(cO_E=\pi O_E\), the unsigned finite mass relative to the
diagonal is
\[
 m_v(\gamma_c)=b^{2h}
   \left[b^{2k}+\bigl(b-3+2s(b-1)^2\bigr)b^k+1\right].
\]
\end{theorem}
\begin{proof}
Use the same letter \(c\) for its first component, a
uniformizer. The first split row group is \(\mathrm{GL}_{n+2}\).
Write an element of \(U\) by its three upper blocks
\((l,H,z)\), a row of length \(n\), a column of length \(n\)
and a scalar. In this proof \(H\) denotes that column.
Set \(g=n(l,H,z)^{-1}\gamma_c\), using block order \(1,n,1\).
Direct multiplication gives
\[
 g=
 \begin{pmatrix}
 1+c(lH-z)&-l&lH-z\\-cH&1&-H\\c&0&1
 \end{pmatrix},
 \qquad
 g^{-1}=
 \begin{pmatrix}
 1&l&z\\0&1&H\\-c&-cl&1-cz
 \end{pmatrix}.
\]
The first column of \(g\) and the last row of \(g^{-1}\) must
be integral. Thus \(cH,cl,cz\) are integral and
\((cl)(cH)\equiv0\bmod c\).
Right multiplication by the integral group adds integral
vectors to \(l,H\), and adds \(z_0+lH_0\) to \(z\) when its
last column entries are \(H_0,z_0\). Fix representatives of
residue scalars, using zero for zero. Every right integral
coset satisfying these conditions therefore has a unique
representative \(H=X_0/c,\ l=Y_0/c,\ z=Z_0/c\), where
\(X_0,Y_0,Z_0\) are the chosen lifts, coordinate by coordinate,
of \(X\in\kappa^n\), \(Y\in(\kappa^n)^\vee\),
\(Z\in\kappa\) with \(YX=0\). Reducing \(l,H\) first keeps
\(cz\) integral. Once both lifts are fixed the remaining
choice is \(z\) modulo the integers. This also proves the
converse for the two integrality conditions.

On the middle space take coordinate vectors \(p_1,\ldots,p_s\)
for the remaining \(s-1\) primal flag vectors followed by
the terminal line. Write \(\lambda_1,\ldots,\lambda_s\) for
the corresponding dual flag in \(M^\vee\). The two sets of
coordinates are disjoint except that \(p_s,\lambda_s\) pair
to one. Multiplying the terminal vector or covector by the
fixed units of norm product \(-T\) does not affect these
tests. Put \(P_i=\langle p_1,\ldots,p_i\rangle\) and
\(D_i=\langle\lambda_1,\ldots,\lambda_i\rangle\) over \(\kappa\),
with \(P_0=D_0=0\). The condition \(g\in KB\) is precisely
that on each side the transformed list can be made integral
and primitive by adding earlier entries to later ones. Indeed
these are the two factors of the column changes of
Lemma~\ref{lem:flag215}. After such changes the two lists
still have their original pairings. The integral isometry
in that lemma carries them to the reference decoration.

Let \(i\) be the first index with \(Yp_i\ne0\), and \(j\)
the first with \(\lambda_j X\ne0\). Use \(\infty\) when the
corresponding list vanishes. Put
\(A=1+Y_0X_0/c-Z_0\), \(D=1-Z_0\), and write \(\bar A\)
for the residue of \(A\).
The primal first vector is \((A,-X_0,c)\).
The following vectors are \(p_a-(Y_0p_a/c)u\).
When \(i=\infty\) they are all the original coordinate vectors.
The list then passes unless \(X\in P_s\) and \(Z=1\).
In fact its reduction is independent exactly when
\(\bar A\ne0\) or \(X\notin P_s\); if \(X\in P_s\) the
coordinate lifts have \(Y_0X_0=0\). If a prefix has dependent
reductions, its exterior product is divisible by \(c\) and
every allowed change preserves that product.

If \(i<\infty\), the first \(i-1\) vectors after the
first are unchanged. For a successful list this prefix
must already span a primitive lattice. Removing the pole
at \(p_i\) requires \(u\bmod c\) in the span of the prefix
reductions. To see this, a linear combination of primitive
independent columns lying in \(c^{-1}L\) has all coefficients
in \(c^{-1}O\), by an integral left inverse. Multiply by
\(c\) and reduce. Consequently \(X\in P_{i-1}\) and
\(\bar A\ne0\). In this case \(Y_0X_0=0\), so \(Z\ne1\).
Conversely assume \(X\in P_{i-1}\) and \(Z\ne1\).
That prefix spans the vector \(Au+c\,v\) together with the
earlier middle coordinates. Here \(v\) is the last coordinate
vector. After the prefix, all unwanted multiples of \(u/c\)
can be canceled with multiples of \(Au+cv\), leaving integral
vectors that project to \(p_a\) in the middle coordinates.
All prefix reductions are independent, as the first vector
has a unit \(u\)-entry. This proves sufficiency.

On the dual side the first row is \((-c,-Y_0,D)\)
and the subsequent rows are
\(\lambda_a+(\lambda_aX_0/c)v^\vee\), with \(v^\vee\) the
last coordinate covector. The same argument using these rows
gives the following test. For \(j=\infty\) it passes unless
\(Y\in D_s\) and \(Z=1\). For \(j<\infty\) it passes
exactly when \(Y\in D_{j-1}\) and \(Z\ne1\).

When \(i<\infty\) and the primal test passes, \(X\in P_{i-1}\)
implies \(j=\infty\). For each \(i\) there are \(b^{i-1}\)
choices of \(X\), \((b-1)b^{n-i}\) choices of \(Y\), and
\(b-1\) choices of \(Z\). The case \(j<\infty\) has the same
count upon exchanging vectors and covectors. These disjoint
cases total \(2s(b-1)^2b^{n-1}\).

If \(i=j=\infty\), the two middle residues belong to
\(P_h\oplus W\) and \(D_h\oplus W^\vee\), respectively,
where \(W\) is the \(k\)-dimensional residue column.
The flag components give \(b^{2h}\) choices.
Write \(x,y\) for the two remaining components.
The condition \(yx=0\) has \(b^k+(b^k-1)b^{k-1}\) solutions
by separating \(x=0\). Each allows \(b\) choices of \(Z\),
except that \(Z=1\) is excluded for the \(2b^k-1\) pairs
with \(x=0\) or \(y=0\). Adding the previous two cases and
using \(n-1=k+2h\) proves the formula.
\end{proof}

\subsection{The good Hecke operators on the prescribed column}\label{nav:VIII-52}
Let \(H=U(W_0)\) be the definite rank-\(k\) column of
Section~\ref{sec:char94}. Use any original word of
\eqref{eq:letters111}, its column level, and its real type
or its dual. At an unmarked split place \(v\) put \(P=Nv\),
\(K_v=\mathrm{GL}_k(\mathcal O_v)\) and give \(K_v\) volume
one. Write \(D_v\) for the unnormalized sum over the
index-\(P\) sublattices, and \(A_v\) for the sum over the
neighbors whose intersection has index \(P\) in each lattice.

\begin{proposition}[Spherical neighbor bounds on the word space]\leavevmode
\label{prop:satake287}
On the \(K_v\)-fixed automorphic column space of each
prescribed real type,
\[
       \|D_v\|\le kP^{k/2},\qquad
       \|A_v\|\le (k^2+2)P^k .
\]
The bounds hold for every unmarked split \(v\) on the common
fields and intervals, with the displayed constants independent
of the field and the word.
\end{proposition}
\begin{proof}
We check the base-change input on this column. The group
is compact at all real places and has no proper rational
parabolic, since its Hermitian form is positive definite
at a real place. Its automorphic constituents are discrete,
unitary and cuspidal. The field \(F=E^+\) has even degree and at least two
real places, since it contains
\(\QQ(\sqrt{q d_0})\). The two imaginary quadratic
subfields of coprime discriminants ensure that \(E/F\) is
unramified at every finite place, as proved in
Theorem~\ref{thm:norm91}. At a nonsplit finite place \(-1\)
is a norm, so each pair of positive diagonal entries splits
as a hyperbolic plane there. Thus \(H\) is quasi-split at
all finite places and is the definite unitary group in
Labesse's base-change theorem.

At each real place the column Harish--Chandra entries are
the \(u_i\). Consecutive interior entries differ by
\(4\Delta+\Delta(c_{i+1}-c_i)\ge3\Delta\).
The first gap is at least \(4k\Delta\) by the endpoint
formula. At an ordinary place
\(u_{k-1}-u_k=4\Delta-\Delta c_{k-1}\ge3\Delta\),
and the high-place choice only increases this gap.
Duality reverses and negates the list and preserves the gaps.
Thus all gaps exceed one for sufficiently large \(x\).
The compact types are cohomological. For a definite unitary
group the archimedean discrete-series and pseudo-coefficient
hypotheses of Labesse's Corollary 5.3 hold because a compact-group
character projects onto its specified irreducible by ordinary
Schur orthogonality. This is also the definite case stated
in Theorem 2.2 and Remark (1) of the paper of
Guerberoff on modularity lifting cited below.

Base change consequently gives an isobaric sum of discrete
automorphic representations of general linear groups over
\(E\). Its local parameter at a split place agrees with that
of the column there, and its holomorphic infinitesimal
character at a complex place is the column list, up to the
dual convention. The discrete-spectrum classification of
M\oe glin and Waldspurger expresses a noncuspidal discrete
summand as a single string of length at least two over a
unitary cuspidal representation. The twists in such a string
differ successively by the idelic norm. At a complex place
\(|z|_{\CC}=z\bar z\), so this would give two holomorphic
entries differing by one. There is no such pair on this
column. Every discrete summand is therefore unitary cuspidal.
The unramified parameter is their multiset union, with each
constituent unramified, since the direct sum has trivial
inertia and zero monodromy. As in the proof of
Proposition~\ref{prop:basechange258}, the unramified bound
of Luo, Rudnick and Sarnak, also applied to the dual,
places every Satake modulus between \(P^{-1/2}\) and
\(P^{1/2}\). Unitary characters have modulus one.
The inputs here are
\href{https://webusers.imj-prg.fr/~michael.harris/CBCM.pdf}{Labesse, Corollary 5.3},
\href{https://doi.org/10.1112/S0010437X10005154}{the definite-unitary base-change statement},
\href{https://doi.org/10.24033/asens.1595}{M\oe glin and Waldspurger on the discrete spectrum}
and \href{https://doi.org/10.1090/pspum/066.2/1703764}{Luo, Rudnick and Sarnak on unramified parameters}.

For precision, the eigenvalue of \(D_v\) is
\(P^{(k-1)/2}\sum_{i=1}^k\alpha_i\) for one of the two dual
Satake conventions. In the upper triangular normal form
for an index-\(P\) sublattice, a diagonal uniformizer in
position \(i\) allows \(P^{k-i}\) choices of entries to
its right, taken modulo that uniformizer. The normalized
spherical vector in the unramified principal series has
value \(P^{-(k+1-2i)/2}\alpha_i\) on these representatives
by the Borel modulus. This proves the eigenvalue formula,
which passes to the spherical constituent.
Thus its absolute value is at most \(kP^{k/2}\).
Summation on a compact arithmetic quotient decomposes into
these unitary automorphic constituents, so the bound holds
on the whole word space. There are
\(d_P=(P^k-1)/(P-1)\le 2P^{k-1}\) index-\(P\) sublattices.
Going down to one and then up either returns to the original
lattice or passes through the unique intersection with a
neighbor. Thus \(A_v=D_v^*D_v-d_P I\), with the same
identity in the opposite order under the dual convention.
Taking norms gives the asserted estimate.
\end{proof}

\subsection{The eigenvalue of the inert neighbor operator}\label{nav:VIII-53}
Retain the hyperspecial column at the inert place of
Proposition~\ref{prop:neighbors289}. Its rank in Section~\ref{sec:missing}
is even, \(k=2l\). Write \(X_1,\ldots,X_l\) for the normalized
unramified inducing characters evaluated at \(\pi\) on the
diagonal torus of \(U(k)\) in a hyperbolic basis.

\begin{proposition}[An inert spherical bound on the original column]\leavevmode
\label{prop:sphere289}
The eigenvalue of the neighbor operator is
\[
  P^{k-1}\left(\sum_{i=1}^l(X_i+X_i^{-1})
                    +\frac{1-P^{-k}}{1+P^{-1}}\right).
\]
On the column of each prescribed word one has
\(\|A_v\|\le(k+1)P^k\). The signed finite factor of a first
plane of valuation one consequently has norm at most
\(B_xP^{r-1}\), where \(B_x=4(k+s+1)^2\).
\end{proposition}
\begin{proof}
Use the upper triangular Borel with unipotent radical \(U_+\)
and diagonal \(t_\nu\) of exponents
\((\nu_1,\ldots,\nu_l,-\nu_l,\ldots,-\nu_1)\).
In the normalized Satake transform of a double-coset indicator
\(\phi\) the coefficient at \(\nu\) is
\(\delta(t_\nu)^{1/2}\int_{U_+}\phi(t_\nu n)\,dn\).
The integral compact subgroup has volume one. This transform
is invariant under permutations and inversions of the torus
variables and evaluates to the Hecke eigenvalue at the
inducing characters.
\href{https://math.umd.edu/~tjh/satake24.pdf}{Haines and
Rostami, equation (9.2.2) and Section 10.1}.

Let \(t=\operatorname{diag}(\pi^{-1},1,\ldots,1,\pi)\).
Use \(\phi=1_{K_HtK_H}\) by Proposition~\ref{prop:neighbors289}.
Every minor of an element of this double coset has valuation
at least \(-1\), by exterior powers of its two integral factors
and of \(t\). In an upper triangular matrix the principal minor
formed from all negative diagonal exponents therefore gives
\(\sum_i|\nu_i|\le1\). For \(t\) itself the condition
\(tn\in K_HtK_H\) is equivalent to \(n\) integral. Indeed
the first row of \(tn\) bounds the first row of \(n\).
For \(1<i<k\), the minor in rows \(1,i\) and columns \(1,j\),
\(j\ge i\), equals \(\pi^{-1}n_{ij}\), which bounds all
remaining entries. Conversely integral \(n\) belongs to \(K_H\).
Conjugation by \(t\) dilates the \(k-2\) entries over \(D\)
of the first-line unipotent radical by \(\pi^{-1}\) and its
one skew central coordinate by \(\pi^{-2}\). It fixes the
middle unipotent quotient. Thus \(\delta(t)^{1/2}=P^{k-1}\).
Weyl invariance gives this same coefficient for each of the
\(2l\) nonconstant monomials.

To determine the constant, evaluate on the trivial
representation, whose eigenvalue is the number \(d\) of
neighbors. Let \(Z_j\) count all isotropic vectors, including
zero, in a nonsingular residue Hermitian space of rank \(j\).
After diagonalizing, the norm of one coordinate takes each
nonzero value \(P+1\) times and zero once. Hence \(Z_0=1\)
and \(Z_{j+1}=Z_j+(P+1)(P^{2j}-Z_j)\). Induction gives
\(M_k=Z_k-1=(P^k-1)(P^{k-1}+1)\) for even \(k\).
By Proposition~\ref{prop:neighbors289},
\[
      \frac d{P^{k-1}}=\sum_{i=1}^l P^{2i-1}
                             +\sum_{j=0}^{l-1}P^{-2j}.
\]
On the trivial representation the normalized parameters,
up to the Weyl action, are \(P^{k+1-2i}\). Indeed in the
upper triangular Lie algebra the root spaces for \(e_i-e_j\)
and \(e_i+e_j\), \(i<j\), each have dimension two over
\(C\), being paired matrix entries over \(D\). Those for
\(2e_i\) have one skew entry. The half sum with multiplicities
therefore has coordinates \(k+1-2i\), giving the stated
values for the inducing character \(\delta^{-1/2}\).
Subtracting the \(2l\) monomials from \(d/P^{k-1}\) leaves
\((1-P^{-1})\sum_{j=0}^{l-1}P^{-2j}
=(1-P^{-k})/(1+P^{-1})\).

For the prescribed definite column the proof of
Proposition~\ref{prop:satake287} gives an isobaric sum of
unitary cuspidal representations over \(E\), using the real
gaps of the very same word. Labesse's transfer agrees with
local base change at every unramified place. At this inert
place its parameters are \(X_i,X_i^{-1}\) for \(1\le i\le l\).
The unramified bound used in Proposition~\ref{prop:satake287}
now gives \(|X_i|,|X_i^{-1}|\le(Nw)^{1/2}=P\).
\href{https://webusers.imj-prg.fr/~michael.harris/CBCM.pdf}
{Labesse, Corollary 5.3} and
\href{https://www.imj-prg.fr/fa/bpFiles/Minguez.pdf}
{M\'inguez, Theorem 4.1}.
The constant inside parentheses in the eigenvalue is between
zero and one. The asserted bound follows on each unitary
constituent and hence on their orthogonal sum. Finally
Proposition~\ref{prop:neighbors289} bounds the signed operator
by \(P^{2h}(P-1)+P^{2h}(P^2-1)(k+1)P^k\).
Here \(2h+2+k=r-1\), giving the claimed common bound.
\end{proof}

\subsection{Temperedness of the column at split places}\label{nav:VIII-54}
\begin{proposition}[The split local components of the prescribed column]\leavevmode
\label{prop:tempered292}
Let \(F=E^+\). Every automorphic constituent \(\sigma\)
of the definite column \(H=U(W_0)\) with one of the real
types in Proposition~\ref{prop:satake287}, or its dual,
has tempered local component at every place of \(F\) that
splits in \(E\). This includes ramified representations
at \(g_1,g_2\) and \(T\).
\end{proposition}
\begin{proof}
Fix such a constituent. The proof of
Proposition~\ref{prop:satake287} verifies the hypotheses
of definite unitary base change for these fields. The
transfer is an isobaric sum \(\boxplus_i\Pi_i\).
The individual discrete summands are conjugate self-dual,
and the local identification holds at every split place.
Each summand is unitary cuspidal. Indeed a residual
string of length at least two would contribute
two entries differing by one to the holomorphic
infinitesimal character at infinity. All the gaps of the
ordered column list exceed one.
\href{https://webusers.imj-prg.fr/~michael.harris/CBCM.pdf}
{Labesse, Corollary 5.3} and
\href{https://doi.org/10.1112/S0010437X10005154}
{Guerberoff, Theorem 2.2}.

Write \(d_i\) for the rank of \(\Pi_i\). At each complex
place, each of its two infinitesimal-character lists is a
sublist of the corresponding list for the transfer. The
entries of each full list are distinct and in
\((k-1)/2+\mathbb Z\), by the highest weight formula for
the compact column and its conjugate. The rank \(d_i\)
requires entries in \((d_i-1)/2+\mathbb Z\) for an
algebraic infinitesimal character. We describe explicitly
a unitary conjugate self-dual twist to correct the parity.

Take \(K=\mathbb Q(\sqrt{-q})\subset E\), with its quadratic
idele class character \(\eta\) on \(\mathbb Q\). On the
image of \(\mathbb C^\times\) in the idele class group
\(C_K\), prescribe \(z\mapsto z/|z|\), using one complex
embedding. On the image of \(C_{\mathbb Q}\) prescribe
\(\eta\). To see that the prescriptions join continuously,
write rational idele classes using a positive archimedean
scalar and finite units \(\widehat{\mathbb Z}^{\times}\).
The positive scalars are already in the first subgroup.
If a rational finite unit idele represents a complex scalar
in \(C_K\), their quotient is a principal idele. Its finite
components are rational and fixed by conjugation. The
principal element thus belongs to \(\mathbb Q^\times\).
Since it is a finite unit everywhere, it is \(1\) or \(-1\).
The prescriptions agree on these elements because
\(\eta_\infty\) is the sign character and \(\eta\) is
trivial on principal ideles. The product subgroup is closed.
Indeed its norm-one part is the image of the compact group
\(S^1\times\widehat{\mathbb Z}^{\times}\), and positive
real scalars split off by the idele norm. This also proves
continuity of the joined character. The character
extension theorem for locally compact abelian groups
extends it to a unitary character \(\mu\) of \(C_K\).
For \(x\in C_K\) we have
\(\mu(x)\mu(\bar x)=\eta(N_{K/\mathbb Q}x)=1\) by
quadratic reciprocity in idele class form.

Pull back \(\mu\) by \(N_{E/K}\). The resulting character
\(\nu\) is unitary and conjugate self-dual. At each complex
place its two infinitesimal entries are \(1/2,-1/2\), in
one order. If \(k-d_i\) is odd twist \(\Pi_i\) by
\(\nu\circ\det\), and otherwise use \(\Pi_i\) itself.
The two lists now have the required parity. After sorting
each and subtracting the rank-\(d_i\) half sum of positive
roots, both yield integral dominant highest weights, since
consecutive entries differ by at least one. This proves
regular algebraicity of the twisted cuspidal representation.
The cohomological characterization of regular algebraic
cuspidal representations of general linear groups gives
the archimedean cohomology required in Caraiani's theorem.
The twist has preserved conjugate self-duality. For
\(d_i\ge2\) that theorem gives temperedness at every
finite place of \(E\). For \(d_i=1\) this follows directly
from unitarity. Removing a unitary character twist
preserves temperedness.
\href{https://homepage.univie.ac.at/harald.grobner/papers/GLn(D).pdf}
{Grobner and Raghuram, Theorem 6.3},
\href{https://doi.org/10.1112/S0010437X10005154}
{Guerberoff, proof of Theorem 2.3} and
\href{https://arxiv.org/abs/1010.2188}
{Caraiani, Theorem 1.2 and Corollary 5.9}.

At a split place the local Langlands parameter of
\(\sigma_v\) in either general linear factor is the
direct sum of those of the \(\Pi_{i,w}\), up to the dual
for the other factor. A finite direct sum of tempered
parameters is tempered, since its restriction to the
Weil group in the Weil-group-times-\(\mathrm{SL}_2\)
convention is bounded. The tempered local correspondence
for \(\mathrm{GL}_k\) proves the assertion.
\end{proof}

\subsection{The inert coefficient and the global finite factor}\label{nav:III-59}
\begin{proposition}[The finite operator for every rank-one label]\leavevmode
\label{prop:finiteall294}
Let \(c=Q\xi\) and \(\lambda\) represent a supported label
of Proposition~\ref{prop:representative294}. Put
\[
  A_\xi=|N_{E/\mathbb Q}\xi|^{1/2},\qquad
  D_Q(\xi)=\prod_{\substack{v\mid Q\\a_+(v)+a_-(v)>0}}Nv .
\]
With \(C_k=4(k+1)^4\), the finite operator has norm at most
\[
       D_Q(\xi)^k\,
       A_\xi^{\,k+6s+1+2\log C_k/\log2}.
\]
\end{proposition}
\begin{proof}
At a nonsplit place \(E_w/F_v\) is unramified quadratic
and both levels are full. The norm-one scalar \(\lambda\)
is a unit. Put \(e=v_w(c)=v_w(\xi)\), \(P=Nv\). In the
unitary upper radical \(l=-h_M(H,\cdot)\) in the usual
Hermitian notation. Use \(X=cH,\ Y=\bar c l,\ Z=\bar c z\).
The first endpoint conditions in the exchange matrix
make \(X,Y,Z\) integral. Here \(Y=-h_M(X,\cdot)\), so
the two truncated column depths coincide, say \(t\).

The column denominators have at most one negative exponent
of length \(d\le e\). Moreover \(d\ge e-2t\).
For detail, take the isotropic dual rows of the matrix
in Lemma~\ref{lem:scalardepth294} over \(\mathcal O_{E_w}\)
with \(c_+=c,\ c_-=\bar c\). Use the quotient by the
terminal line, an integral orthogonal summand here. The
minors on one column coordinate or the initial coordinate
give \(\pi^t X_S/c\) integral, since \(\lambda\) is a
unit. The full minor makes \(Z-YX_S/c\) integral. The
input with last coordinate \(\pi^t\) admits an integral
column \(w\) solving the last congruence modulo \(\bar c\).
It gives denominator \(e-2t\) if this number is positive.
The upper bound and cyclicity follow from the single
fractional column \(X_W/c\) in the projection. Self-duality
of the physical column under its unitary Levi gives
positive exponent \(d\) as well. The unit \(\lambda\)
preserves this column lattice.

The spherical compression at \(d\) has norm at most
\((k+1)^dP^{-(k-2)d}\) by Proposition~\ref{prop:cartan290}.
The integral right unipotents permit fixed lifts for \(X\)
modulo \(c\), since the trace map on the unramified
quadratic integer ring is onto, including in residual
characteristic two. In fact an arbitrary integral shift
of \(H\) can be realized by completing it with a central
coordinate of the required trace. For fixed \(H\), the
possible \(z\)'s differ by trace-zero elements in
\(\pi^{-e}\mathcal O_{E_w}\), taken modulo integral
trace-zero elements. This gives at most \(P^e\) choices,
since the trace-zero lattice has rank one over
\(\mathcal O_{F_v}\). Thus at depth \(t\) the count is at
most \(P^{2k(e-t)+(2(n-k)+1)e}\). Writing \(j=e-2t\),
we have \(2k(e-t)-(k-2)\max(0,j)\le(k+2)e\).
Together with \(n-k=2s-1\) this gives the local bound
\(C_k^{2e}P^{(k+4s+1)e}\), including the at most
\(e+1\) depths and the factor \((k+1)^e\). At valuation
zero the integral exchange and unit central scalar give
one coset.

At a split place, integrality of \(\xi\lambda\) in both
components gives \(-a_+\le L\le a_-\). Thus
\(|L|\le a_++a_-=m_0\). Proposition~\ref{prop:splitall294}
is bounded by \(C_k^{m_0}P^{(k+6s+1)m_0/2}\) with the
extra \(P^k\) at active places above \(Q\). The local
factors of \(A_\xi\) are \(P^{m_0/2}\) at split places
and \(P^e\) at inert places. Since \(P\ge2\), the sum of
all the exponents of \(C_k\) just used is at most
\(2\log A_\xi/\log 2\). The common-column unfolding of
Proposition~\ref{prop:representative294} multiplies the
finite factors. This proves the assertion.
\end{proof}

\subsection{Bounds on the finite operator for every prescribed word}\label{nav:V-46}
Let \(F_q^{\mathrm{col}}=\prod_{v\mid q}F_v\) be the operator
of Theorem~\ref{thm:phase282}, and set
\[
 d_b=\frac{b^k-1}{b-1},\qquad
 L_b=b^{\,2s-2+k}(b-1)(1+2s(b-1)),\qquad
 U_b=a_b+u_b d_b(d_b-1)\frac{3b-1}{b(b+1)} .
\]
The \(q\) in the subscript denotes the rational field prime.
The column forms are on \(H=U(W_0)\) at the finite level of
Theorem~\ref{thm:midlevel131}, with the given irreducible
real word type. Use its probability inner product.

\begin{theorem}[Positive lower and upper bounds after terminal projection]\leavevmode
\label{thm:bounds282}
On each prescribed word and every one of its terminal
character spaces or zero-sum spaces, the finite operator satisfies
\[
       \left(\prod_{v\mid q} L_{b_v}\right)1
       \ \le\ F_q^{\mathrm{col}}\ \le\
       \left(\prod_{v\mid q} U_{b_v}\right)1,
       \qquad b_v=\mathrm Nv .
\]
The same inequalities hold as quadratic forms on the
coordinates retained in Corollary~\ref{cor:trace235}.
On the common intervals its operator norm, divided by the
unsigned mass \(m_f(\gamma_c)\), is at most
\(4^{\,\#\{v\mid q\}}q^{-f/2}\) for all sufficiently large
\(x\).
\end{theorem}
\begin{proof}
At one place let \(D_v\) sum the translations to the
index-\(b\) sublattices of the column lattice. On the
spherical level its adjoint sums over index-\(b\)
superlattices, by inversion of the compact double coset in
the unimodular local group. A two-step return has \(d_b\)
paths to the starting lattice, one through each intermediate
lattice. A path to a distinct endpoint has its intersection
with the starting lattice as the unique intermediate lattice.
Thus \(D_v^*D_v=d_b\,1+A_v\). Reversing the two steps gives
the same identity, with the sum of the lattices as intermediate
lattice. It follows that
\[
                 F_v=(a_b-u_b d_b)1+u_b D_v^*D_v,
                 \qquad a_b-u_b d_b=L_b>0 .
\]
This proves the lower bound without a spectral assumption.

There is no vector of this word type fixed by
\(\mathrm{SL}_k(F_v)\). Indeed the simply connected group
\(\mathrm{SU}(W_0)\) is split and noncompact at \(v\).
Strong approximation away from \(v\) makes its rational
points dense in its adeles away from \(v\), including the
real coordinates. For a smooth local-invariant function
on \([H]\), approximate \(gag^{-1}\) away from \(v\), with
\(a\in\mathrm{SU}(W_0)(\mathbb A_F)\). At \(v\) the
difference lies in the local special linear group, which
is normal. Rational left invariance and local right
invariance give equality at \(g\) and \(ga\) by continuity.
Such a function is therefore trivial under all real special
unitary factors. The original word has nonconstant highest
weight there, as the consecutive entries in its cloud are
separated by the gaps of Theorem~\ref{thm:common111}.
The same conclusion holds for Hilbert vectors by smoothing:
the invariant subspace is stable under translations and
the real compact type projector. This is the application of
strong approximation used in Theorem~\ref{thm:mix277}.

The operator \(A_v/[d_b(d_b-1)]\) is probability convolution
on the double coset with exponents \(1,-1\) and zeros.
The local matrix-coefficient estimate in the proof of
Lemma~\ref{lem:coefficient199} gives norm at most
\((3b-1)/(b(b+1))\) on this subspace. Its local argument
uses only three coordinates, the size of the residue field
and the absence of special-linear invariant vectors.
It applies since \(k\ge3\) on the original intervals.
For the general local method see
\href{https://gauss.math.yale.edu/~ho2/doc/mat.pdf}
{H.~Oh, Uniform pointwise bounds for matrix coefficients}.
This proves \(F_v\le U_b\,1\). All the \(F_v\) commute,
are positive, and commute with the terminal translations,
which act at other places. Multiplication gives the bounds,
and then compression preserves the quadratic-form inequalities.

By Theorem~\ref{thm:split281},
\(m_v=a_b+u_b d_b(d_b-1)\) and \(a_b/m_v=O(s b^{2-k})\)
uniformly. For the original ranks \(s b^{3-k}\) tends to
zero uniformly for \(b\ge q\). Hence \(U_b/m_v\le4/b\)
eventually. The product of all \(b_v\) over \(q\) is
\(q^{f/2}\) by the ramification computation in
Theorem~\ref{thm:product281}.
\end{proof}

\subsection{The inert operator and the product over the finite places}\label{nav:V-53}
\begin{proposition}[The finite product for an arbitrary single block]\leavevmode
\label{prop:blockfinite297}
At an inert place with \(P=Nv\) and \(m=jv(c)\), the signed
local operator has norm at most
\[
              [2(k+1)]^m P^{(k+32(s+1))m}.
\]
For the full product over the finite places one has
\[
  \|F(c,\lambda)\|\le (Q^f A_\xi)^{kj+\Delta_j},\qquad
  \Delta_j=100j\bigl(s+1+\log_2(2(k+1))\bigr).
\]
\end{proposition}
\begin{proof}
Use the unramified quadratic ring with residue cardinality
\(P^2\). Since the terminal prime splits, the terminal line
at this place is an integral orthogonal summand. The column
is unimodular. Hence the column part of \(l\) is the negative
adjoint of \(H_W\) with an integral unimodular factor, and
\(X=Y=t\) in Lemma~\ref{lem:blocklattice297}. Also
\(\lambda\bar\lambda=1\) gives \(v(\lambda)=0\).
The scaled column action is unitary and its exponents are
paired with their negatives. Multiplication by \(\lambda\)
is an integral central action and does not change its
hyperspecial double coset. The lemma bounds the common
denominator length by \(2t-m\le d_-\le t\le m\), with
at most \(j\) nonzero pairs.

The primal kernel has at most \((P^2)^{jm}\) choices.
The column rows at length \(t\) have at most
\((P^2)^{(k+j)t}\) residues, by the row-kernel count in
Proposition~\ref{prop:blocksplit297}. The \(q_0=n-k\)
remaining rectangular rows and the \(j\) central rows cost
at most \((P^2)^{(q_0+j)m}\). Integral rectangular changes
lift in the unramified radical, including in residual
characteristic two, by the trace argument of
Proposition~\ref{prop:blockfinite295}. Thus the count at
\(t\) is at most \(P^{2kt+2(3j+q_0)m}\). The hyperspecial
compression is at most \((k+1)^m P^{-(k-4j+2)d_-}\)
by Lemma~\ref{lem:blocksphere297}. Insert \(d_-\ge2t-m\);
the exponent from this last power and \(P^{2kt}\) is at
most \((k+8j)m\), since \(t\le m\). There are at most
\(2^m\) choices of \(t\). As \(j\le s\) and \(q_0=2(s-j)+1\),
the claimed local bound follows.

Multiply this estimate and Proposition~\ref{prop:blocksplit297}.
The local factors of \(|N_{E/\mathbb Q}c|^{1/2}=Q^f A_\xi\)
are \(P^{m/j}\) at an inert place and \(P^{M/(2j)}\) at a
split place. Replace \(2(k+1)\) by
\(P^{\log_2(2(k+1))}\), valid also for \(P=2\). The
exponents in the two local estimates are then bounded by
the corresponding exponents of the asserted product.
The original constituent spaces and their tensor products
obey the same bounds, by decomposing the unitary
automorphic column into its local representations and
taking the operator norm on each constituent.
\end{proof}

\section{Unramified and transverse factors}\label{nav:g-IX-4}

\subsection{The unmarked root integrals}\label{nav:III-45}
Use Proposition~\ref{prop:positions264} at a finite place \(v\).
Write \(E_v=E\otimes_F F_v\), with integer ring \(O\), also
when \(E_v\) is split. Let \(U_\ell\) denote the upper
unitriangular group over \(E_v\), and put
\(\chi(u)=\psi_E(\sum_{j<\ell}u_{j,j+1})\). The character
\(\psi_E\) uses the absolute trace to \(\QQ_p\), \(p\) below
\(v\), and the standard additive character. Put
\(X_*=(d b_1,\ldots,d^\ell b_\ell)\). Write a test in the first
position as a function of the \emph{scaled} columns, so the
indicators in that proposition are applied after dividing
column \(j\) by \(d^j\). Keep the second position \(p\)
in the unscaled row coordinates. For a local test \(\Phi\) set
\[
   \mathcal E_v(g)=
      \int_{U_\ell}\overline{\chi(u)}
         [\omega(g)\Phi](X_*u,p_*)\,du,\qquad
   R_v=\{g\in U(S)(F_v):gX_*=X_*u\text{ for some }u\in U_\ell\}.
\]
This integral takes values in the Schwartz space of the
residual oscillator. It has compact integration support
because \(X_*\) has independent columns and admits a linear
left inverse. Haar measure on \(U_\ell\) is product additive
measure in its upper entries, by triangularity of multiplication.

\begin{proposition}[Support at an unmarked finite place]\leavevmode
\label{prop:flagunmarked264}
For \(v\) unmarked, take the test of
Proposition~\ref{prop:positions264} with the residual lattice
vacuum. Let \(K_v\) be the stabilizer of \(\Lambda_{S,v}\).
Then \(\mathcal E_v(g)\ne0\) holds exactly on \(R_vK_v\).
\end{proposition}
\begin{proof}
Write \(P\) for the stabilizer of the flag spanned successively
by \(b_1,\ldots,b_r\), and \(N_P\) for its unipotent radical.
We record how the evaluation changes under this parabolic.
In the first polarization the action of \(g\in U(S)\) is
substitution \(x\mapsto g^{-1}x\) followed by a fixed invertible
operator on the residual variable, including the splitting
scalar. Indeed the oscillator splits across the orthogonal
hyperbolic part of \(T\) and \(T_0\); on the first part the
substitution intertwines both translations and dual modulations.
Uniqueness of the Heisenberg representation leaves only a
scalar. In particular the residual operator is independent of
the upper entries of \(u\).

For \(g\in P\), the induced action on the second position is
the action on the quotient of \(S\) by the perpendicular of
\(\langle b_1,\ldots,b_r\rangle\). Evaluation at \(p_*\)
after this action equals an invertible operator on the core
applied to evaluation at \(g^{-1}p_*\). To see this directly,
displacements into \(\langle b_1,\ldots,b_r\rangle\) act by
all modulations of the position \(p\). Conjugation by \(g\)
preserves these displacements, with the dual change of their
characters. The kernel of evaluation at \(p_0\) is spanned by
differences of a modulation and its value at \(p_0\) acting
on tests. Indeed a locally constant test vanishing at \(p_0\)
has compact position support away from \(p_0\); cover it by
finitely many compact open pieces, on each of which a
modulation separates that piece from \(p_0\). Subdivision
into disjoint pieces gives the stated span. Applying \(g\)
and then \(g^{-1}\) gives the invertible map on the evaluation
quotients. This is the Heisenberg construction in the local
Schr\"odinger model.
\href{https://www.math.toronto.edu/~skudla/castle.pdf}{S.~Kudla,
\emph{Notes on the local theta correspondence}, Chapter I}.

Every element of \(R_v\) is in \(P\); on the isotropic columns
its matrix is upper unitriangular in the scaled frame.
On the opposite rows the matrix is consequently lower
unitriangular. It fixes \(p_*\), since that position is zero
or is supported on the last row. Substituting \(u_0u\) for
\(u\) in the integral under a left element of \(R_v\) changes
the integral by an invertible core operator and a scalar of
absolute value one from \(\chi\) (the splitting is already in
the operator). This proves left invariance of its nonvanishing
set. In particular \(N_P\subset R_v\), since adding earlier
isotropic columns to the last, possibly nonisotropic, column
is allowed. Right invariance under \(K_v\) follows from the
integral vacuum in Proposition~\ref{prop:rowfinite262}.

We use \(U(S)(F_v)=PK_v\). For completeness, at an
unramified nonsplit place intersect any rational isotropic
flag of the given dimensions with the unimodular lattice.
These saturated submodules have a primitive adapted basis.
Unimodularity supplies integral vectors pairing dually to
its isotropic vectors. Subtract one half of their Gram matrix
times the isotropic frame to make the dual vectors isotropic.
Here the residue characteristic is odd. The hyperbolic pairs
split off; the remaining unimodular complements of the same
rank are isometric by the orthonormal-basis argument of
Lemma~\ref{lem:orbvol263}. This gives transitivity of \(K_v\)
on such flags. In the split case the two isotropic components
amount to a flag in a free space and the annihilators of a
dual flag. Together they are nested because the original
flag is isotropic. Intersecting with the lattice and taking a
single adapted basis proves the same transitivity in all
residue characteristics. Inverting the orbit decomposition
gives \(PK_v\).

Modulo the left and right operations already treated it
suffices to evaluate \(g=tg_0\), where \(t b_j=t_j b_j\)
for \(j\le r\), \(t\) acts by \(\bar t_j^{-1}\) on the opposite
columns, and \(g_0\in U(S_0)(F_v)\). The diagonal position
conditions on \(g^{-1}X_*u\) give \(t_j^{-1}\in O\) for all
\(j\le r\). If \(c=2r+1\) there is also the condition
\[
  g_0^{-1}b_\ell\in
       (I_\ell^{(c)})^{-1}\Lambda_{S_0,v}.
\]
Indeed division of column \(j\) by \(d^j\) leaves its diagonal
entry unchanged. For \(k<j\) its \(b_k\)-component is
\(t_k^{-1}d^{k-j}u_{kj}\). The original lattice is the direct
sum of the stated lines and the core. Thus the upper entries
range independently over
\[
       u_{kj}\in t_k\,\bigl(dAB\bigr)^{j-k}O\quad(k<j),
\]
where products denote localized fractional ideals.
This uses \(I_k^{(c)}/I_j^{(c)}=(AB)^{j-k}\).
Proposition~\ref{prop:tracelat253} gives
\((d)AB=\mathfrak D_{E/\QQ}^{-1}\) locally at every unmarked
place. The largest fractional ideal on which \(\psi_E\) is
trivial is the inverse different. This holds for the split
product as well, by varying one component at a time.
Orthogonality in the entry \(u_{k,k+1}\) therefore forces
\(t_k\in O\) for every \(k<\ell\), since its integral is zero
unless the additive character is trivial on the entire ideal.
All other entries contribute ordinary volumes. The row and
core operators in this computation are independent of \(u\).

For \(c=2r\) the last remaining entry is \(t_r\).
On the opposite rows \(g^{-1}p_*\) has last row
\(\bar t_r p_{*,r}\). Its value on the middle column of \(T\)
must belong to \(O\), by the row-ideal computation of
Proposition~\ref{prop:positions264}. The value of \(p_{*,r}\)
there is a unit. Thus \(t_r\in O\) as well.
For \(c=2r+1\) the simple roots have already treated all
\(t_1,\ldots,t_r\). Moreover \(I_\ell^{(c)}\bar I_\ell^{(c)}=1\),
so the lattice in the display for \(g_0^{-1}b_\ell\) is
unimodular, contains \(b_\ell\), and has the original lattice
stabilizer on \(S_0\). The vector has norm \(-1\); the integral
norm-orbit argument of Lemma~\ref{lem:orbvol263} places it
in the compact orbit of \(b_\ell\). Hence \(g_0\) is a product
of an element fixing \(b_\ell\) and a lattice stabilizer.
For \(c=2r\) all of \(U(S_0)\) already fixes the frame.
All \(t_j\) and \(t_j^{-1}\) are now integral, so \(t\) itself
is in \(K_v\). This proves necessity in both cases.

At \(g=1\) the row and column tests take value one on the
product just computed, and every simple-root character is
trivial there. Integration gives its positive volume times
the nonzero residual vector. The left and right invariances
prove sufficiency.
\end{proof}

\subsection{Spherical bounds with several nonzero elementary divisors}\label{nav:III-61}
\begin{lemma}[A spherical bound uniform in the number of moving directions]\leavevmode
\label{lem:blocksphere297}
Let \(P=Nv\) and suppose at most \(j\) positive and at most
\(j\) negative elementary divisors occur in a column translate
\(a\), with \(2\le j\le s\). Denote their total absolute length
by \(D\). At a split place the spherical coefficient \(\Xi(a)\)
of Lemma~\ref{lem:conjugation292} satisfies
\[
       \Xi(a)\le k^D P^{-(k-4j+3)D/2}.
\]
At an inert place write \(D=2d_-\). The norm of the
hyperspecial compression on the given automorphic column
is at most \((k+1)^{d_-} P^{-(k-4j+2)d_-}\).
\end{lemma}
\begin{proof}
Normalize the lattice stabilizer to have mass one and let
\(\phi_a\) be the indicator of its double coset. Order the
diagonal exponents increasingly in the split group. As in
Proposition~\ref{prop:cartan290}, the normalized Satake transform
of \(\phi_a\) is bounded coefficient by coefficient by the
product of the transforms of \(D\) translates of one coordinate
of valuation \(1\) or \(-1\). Indeed the convolution of two
double-coset indicators is at least the indicator of the
product double coset, by integrating over the first right
coset at the product representative. The Satake integral
preserves positivity. The convolution transforms multiply.
At the trivial torus parameter defining \(\Xi\), each
one-coordinate transform evaluates to \(kP^{(k-1)/2}\)
by the split calculation of Proposition~\ref{prop:satake287}.

The degree of \(\phi_a\) is at least the upper Borel modulus
\(\delta_B(a)\). The normalized Satake coefficient at \(a\)
is at least \(\delta_B(a)^{1/2}\) by integration over the
integral unipotents. Its opposite Weyl coefficient is equal
to it. Evaluation of this opposite monomial at the trivial
representation multiplies it by \(\delta_B(a)^{1/2}\),
and all contributions are nonnegative. For ordered exponents
\(v_1\le\cdots\le v_k\) the logarithm to base \(P\) of this
modulus is \(\sum_i(2i-k-1)v_i\). Nonzero entries occur in
the first or last \(j\) positions, so this sum is at least
\((k-2j+1)D\). Dividing the numerator by the degree gives
the first bound.

At an inert place the hyperspecial Cartan decomposition in
Proposition~\ref{prop:cartan290} gives \(j\) or fewer
nonnegative exponents \(a_1\ge a_2\ge\cdots\) on hyperbolic
pairs, with \(d_-=\sum_i a_i\). The representative is the
product of \(d_-\) conjugates of the one-pair translate.
At the absolute values of the Satake parameters of our
column, the transform of this translate is bounded by
\((k+1)P^k\) by Proposition~\ref{prop:sphere289}. The same
positivity argument bounds the numerator by its \(d_-\)-th
power. The Borel modulus is
\[
                  P^{\,2\sum_i (k-2i+1)a_i}.
\]
For each pair of indices \(i<i'\) the two upper root
coordinates are quadratic-field additive lines with weights
\(a_i-a_{i'}\) and \(a_i+a_{i'}\). The long root is one
ground-field line with weight \(2a_i\). Multiplying these
factors gives the formula, also in residue characteristic
two. The degree argument applies to this group as well.
Since \(a_i=0\) for \(i>j\), division by the degree proves
the second bound.
\end{proof}

\subsection{Making both terminal vectors isotropic}\label{nav:IV-55}
Let \(d=s+1\) and let \(Y,Y'\) be a relevant supported pair with
\(\mathcal G>0\) as in Lemma~\ref{lem:forms203}. Absolute squares
over \(E\) mean \(x\bar x\). Adjoin to the Hermitian space \(V\)
an orthogonal vector \(t\) of norm one, and let the original group
act on \(\widetilde V=V\perp Et\) by fixing \(t\). Add
\(\varpi t\) to the final vector of each frame. An upper
unitriangular change of either old frame makes exactly the same
change of its extended frame. Indeed the last row of such a
change is the last row of the identity. An element fixing \(t\)
preserves the extended decorated flag exactly when it preserves
the original one.

\begin{proposition}[Opposite maximal isotropic spaces and the fixed vector]\leavevmode
\label{prop:extend300}
Both extended frames are totally isotropic. Their cross matrix is
\(Z+T E_{dd}\), of determinant \(\Delta=D_d+TD_s\ne0\).
Their span \(S\) is nondegenerate. Its perpendicular is positive
definite of rank \(k\) at every real place. If \(t_\perp\) is the
component of \(t\) in this perpendicular, then
\[
              h(t_\perp,t_\perp)=\mathcal G/|\Delta|^2 .
\]
The extended cross matrix can be put in monomial form with the
original reversed blocks. For a final block of length at least two
the original pivots are unchanged. For a final singleton its pivot
increases by \(T\).
\end{proposition}
\begin{proof}
The added summand cancels the norm \(-T\) of each terminal vector
and leaves its pairings against the preceding vectors zero.
It contributes \(T\) to the terminal cross entry, whence the
assertion about \(Z\) and the determinant by cofactor expansion.
The strict inequality \(|\tau D_d|>T|\tau D_s|\) of
Proposition~\ref{prop:shift203} implies \(\Delta\ne0\).
In the ordered basis of both extended frames the Gram matrix of
\(S\) has zero diagonal blocks and off-diagonal blocks \(M\)
and \(M^*\), with \(M=Z+T E_{dd}\). This matrix is nonsingular
and represents \(d\) hyperbolic planes. The extended space has
signature \((k+d,d)\), giving the assertion about the perpendicular.

The column of pairings of the two frames against \(t\) consists
of two copies of \(\bar\varpi\,{\bf e}_d\), where \({\bf e}_d\)
is the last coordinate vector. The inverse Gram matrix has
off-diagonal blocks \(M^{-*}\) and \(M^{-1}\). Thus the norm of
the projection of \(t\) to \(S\) is
\(2T\operatorname{Re}(D_s/\Delta)\), since
\((M^{-1})_{dd}=D_s/\Delta\). Subtracting this quantity from
one gives the displayed formula by expanding \(\Delta\bar\Delta\).

Take the representatives giving the monomial form for the original
pair. If its final block begins at \(q<d\), column \(q\) has just
one nonzero entry, in row \(d\), say \(b\).
Subtract \((T/b)\) times column \(q\) from column \(d\) of
the extended cross matrix. This is a permitted change of the
second frame and restores its original monomial matrix. In the
singleton case the extended matrix is already monomial and its
last pivot is nonzero by \(\Delta\ne0\). In particular
\(\Delta\in Q^d\mathcal O_E\) for a finitely supported pair by
Proposition~\ref{prop:shift203}.
\end{proof}

At a real place take the original Euclidean unit Witt vectors
\(p_i,q_i\) for \(i\le s\), with \(u_i=\sqrt2\,p_i\), and set
\(p_d=(t+e)/\sqrt2,\ q_d=(t-e)/\sqrt2\).
The middle block in this basis is the positive \(k\)-space.
Write \(P_0,P_-\) for the column matrices of the \(p_i,q_i\).
Conjugate the relative isometry by the frequency change of
Proposition~\ref{prop:scale284}, which multiplies \(p_i\) by
\(\tau\varpi\) for \(i\le s\) and fixes \(e,t\).
It sends the original first extended frame to
\(\sqrt2\,\tau\varpi\,P_0\). Consequently the monomial
cross matrix \(C\) in the unit-frequency coordinates is the
one in the proposition evaluated at \(\tau\) and divided by
\(2T\) in every entry. Write \(\delta_p\) for the product of
the moduli of its first \(p\) row pivots. In particular
\(\delta_d=|\tau\Delta|/(2T)^d\).

\subsection{A common rational subspace of the two columns}\label{nav:IV-56}
In a matching cell of Theorem~\ref{thm:ranktrace302}, let
\(\Lambda\) be its row lattice, \(Y_i=X_i+Ez_i\) its two
framed spaces and \(P_i=Y_i^\perp/X_i\) their positive
columns. Write \(L_i\) for the image of \(\Lambda\cap Y_i^\perp\)
in \(P_i\), and let \(L_C=\Lambda\cap C\) in the common
positive perpendicular \(C=(Y_0+Y_1)^\perp\). The quotient
maps induce isometric injections \(p_i:C\to P_i\).
Indeed the Hermitian form on \(Y_i^\perp\) descends modulo
its radical \(X_i\), and its restriction to \(C\) is positive
definite. Here \(\dim_E C=k-1\).

\begin{theorem}[The rank of the difference of the column maps]\leavevmode
\label{thm:rationalrank302}
The comparison gives a rational isometry \(\beta:P_0\to P_1\)
with \(\beta L_0=L_1\). Set \(D=p_1-\beta p_0\),
\(h=\operatorname{rank}_E D\), \(K=L_C\cap\ker D\) and
\(L'=(DC)\cap L_1\). Then
\[
 \delta(K)\delta(L')\prod_{v\mid Q}q_v^{\,2(h-\ell_v(c))_+}
       \le 2^{2fh}\delta(L_C),\qquad
 2fh\log(Q/2)\le\log\!\frac{G}{Q^{2df}I}+2S(c).
\]
Here \(y_+=\max(y,0)\), \(G=|N_{F/\QQ}\mathcal G|\),
and \(I\) is the saturation index of Theorem~\ref{thm:saturation220}.
On the complementary range in Theorem~\ref{thm:ranktrace302}
one has \(h=O(d^2)\) uniformly. Thus \(p_1\) and \(\beta p_0\)
agree on a rational subspace of dimension at least
\(k-1-O(d^2)\) in \(C\). On the columns retained in
Theorem~\ref{thm:columns235} one can also insert the lower
bounds \(\delta(K)\ge Y_{k-1-h}\), \(\delta(L')\ge Y_h\),
where \(Y_0=1\).
\end{theorem}
\begin{proof}
Put \(\gamma_0=1,\gamma_1=\gamma\). The rational maps
\(\gamma_i^{-1}\) identify \(P_i\) with the fixed model
column. Transport the comparison \(\alpha\) of
Proposition~\ref{prop:loops302} between these two coordinates
to obtain \(\beta\). In the finite decompositions
\(\gamma_i^{-1}t=b_i h_i\), \(h_i\in K\), the action of
\(b_i\) on the model column is its column component \(y_i\).
Thus the column lattice on side \(i\) is \(y_i\) times the
model column lattice in these coordinates. At \(T\) the
model lattice in the first split factor is the inverse
image of the kernel of the terminal covector and in the
second factor it is the full dual parent lattice, by
Lemma~\ref{lem:termparent221}. Thus \(K_{\rm line}^*\)
preserves this lattice as well as all other local column
lattices. The condition \(y_1^{-1}\alpha y_0\in K_{\rm line}^*\)
now gives \(\beta L_0=L_1\). In particular \(D L_C\subset L_1\).

At \(v\mid Q\) both isotropic frames reduce to the same
submodule \(U\) in the row. Reduction of the two middle
quotients gives \(U^\perp/U\). Indeed restriction of the
self-dual row pairing onto the dual of either primitive
isotropic module is onto, and its kernel stays exact on
reduction, as in Theorem~\ref{thm:unique221}. The terminal
vector has unit norm here. Its residue line in the middle
quotient agrees for the two frames because the supporting
frame changes reduce into \(\bar N\rtimes U_k\) of
Lemma~\ref{lem:flags211}. Orthogonal projection off this
line is integral and commutes with reduction. Both column
reductions therefore identify with its common perpendicular.
The two projections of each integral vector \(c_0\in L_C\)
are represented in this quotient by the same residue
row vector.

To identify the action of \(\beta\), write a row vector in
the frame of \(t\). In the model frame of the \(i\)-th flag its
image has matrix \(h_i\) after removal of \(b_i\). On the
residual column the radical of the flag acts trivially, so
this change has matrix \(u_{i,v}\). Going from row coordinates
on side zero to row coordinates on side one by \(\beta\)
thus has matrix
\(u_{1,v}^{-1}\overline{(y_1^{-1}\alpha y_0)_v}u_{0,v}
=A_{\alpha,v}\) in the first split factor. In the other
factor the matrix is its inverse transpose in dual bases,
by the perfect pairing of the two residual columns. In both
factors the difference from the identity has rank \(\ell_v(c)\).
It follows that the reduction of \(D:L_C\to L_1\) has rank
at most \(\ell_v(c)\) in each factor. This implication
only uses the images of integral row vectors and applies
even if \(p_i(L_C)\) has smaller rank on reduction.

The lattice \(L'\) is saturated in \(L_1\), so reduction of
\(L'\) injects locally into reduction of \(L_1\). At either
of the two primes of \(E\) above \(v\mid Q\), at least
\((h-\ell_v(c))_+\) of the \(h\) elementary divisors of
\(D L_C\subset L'\) are divisible by the local uniformizer.
The residue field has size \(q_v\) in each factor. Their
indices multiply, giving
\[
                [L':D L_C]\ge
                     \prod_{v\mid Q} q_v^{\,2(h-\ell_v(c))_+}.
\]

For the archimedean inequality give \(L_C/K\) the orthogonal
quotient metric. For any positive Hermitian lattice \(M\)
of rank \(a\), use its Euclidean covolume in the \(f\) complex
embeddings divided by \(c_E^a\), with
\(c_E=2^{-f}\sqrt{|\Delta_E|}\).
This agrees with \(\delta(M)=[M^\vee:M]^{1/2}\) when \(M\)
is integral. For a free lattice the Gram determinant at each
embedding is the real volume factor and their product is
the square root of its norm from \(E\). Passage to a free
sublattice of finite index gives the same identity for a
projective lattice by multiplicativity of indices and duality.
A basis of \(K\) as a \(\ZZ\)-lattice extended by lifts of a
basis of \(L_C/K\) shows
\(\delta(L_C)=\delta(K)\delta(L_C/K)\) by orthogonal
projection of the lifts. The normalizing powers of \(c_E\)
also multiply. The map induced by \(D\) on the real
orthogonal quotient has operator norm at most two at each
place, since it is the difference of two isometries on \(C\).
Its volume expansion in real dimension \(2h\) at each of
the \(f\) places is at most \(2^{2h}\). Consequently
\[
 [L':D L_C]\delta(L')=\delta(DL_C)
       \le 2^{2fh}\delta(L_C)/\delta(K).
\]
This proves the first inequality, with rank-zero determinants
equal to one. Both \(K\) and \(L'\) are integral, so their
determinants are at least one. Use
\(\sum_{v\mid Q}\log q_v=f\log Q\),
\((h-\ell_v)_+\ge h-\ell_v\), and the identity
\(\delta(L_C)=G/(Q^{2df}I)\) of Theorem~\ref{thm:saturation220}
to obtain the second inequality.

In the remaining height range each \(H_i<b^{Mfd^2}\).
Writing \(\Delta_+=D_d+TD_s=Q^d a_t\), the identity
\(\mathcal G=|\Delta_+|^2-2T\operatorname{Re}
(\Delta_+\overline{D_s})\) gives
\[
 G\le (1+2T)^f Q^{(2d+s)f} H_t^2
                 \max(1,H_{t-1})
\]
when \(D_s\ne0\), because then \(D_s=Q^s a_{t-1}\).
For \(D_s=0\) simply use \(G\le Q^{2df}H_t^2\).
Since \(\log Q\asymp\log b\) and \(\log T=O(d\log b)\),
we obtain \(\log G=O(fd^2\log b)\). The bound on \(S(c)\)
in Theorem~\ref{thm:ranktrace302} and \(I\ge1\) now imply
\(h=O(d^2)\).

Finally \(p_1K\subset L_1\) is an isometric copy of \(K\).
Its saturation has the same rank and determinant at most
\(\delta(K)\), and \(L'\) is already saturated. A matching
retained pair has both columns in the retained set, which
is a union of scalar orbits. Applying the primitive
determinant bounds on \(L_1\) gives the last assertions.
\end{proof}

\subsection{The weighted local estimate on the last row}\label{nav:IX-33}
Use precisely the period field, row and normalizations in
Theorem~\ref{thm:last244}. Choose the representatives and the
common orthonormal basis of Lemma~\ref{lem:basis63} for the
finite set \(S\) of primes \(p\ge5\) inert in \(K\), unramified
in \(E\), and prime to the level parameter \(Q\) of the finite
CM test. Use the marked vacuum and rational torsion translates of
order dividing \(Q\), with coefficients integral above \(S\), as
in Theorem~\ref{thm:one61}. The fixing-level representatives may
be taken to move only at good completely split primes. With the
common basis these are precisely finite sums of the class germs
to which Theorem~\ref{thm:direct76} applies. This also allows
several such tests with a common \(Q\) and integral coefficients.
Let \(A_n,A_s\) be common multiples of the respective stabilizer
orders. View the right side of Theorem~\ref{thm:last244} as a
bilinear form \(\mathcal B_{\rm row}\) on the arrays of exterior tensors.
The coordinate minor formula extends it to the ambient tensor
spaces. Thus its matrix can be studied before it is paired with
a particular incoming automorphic form.

\begin{theorem}[Divisibility of the matrix for the prescribed last row]
\label{thm:frob244}
For \(p\in S\), choose one half \(H_C\) in each even cycle of
\(\pi_p\) on the CM embeddings, as in Theorem~\ref{thm:direct76},
using the same choice for all matrix entries. Denote by
\(d_M(u,h)=\max(0,\lceil(u-h(M-1))/M\rceil)\). Put
\[
 g_p=\sum_{C\ {\rm even}}\sum_{\nu\in H_C}
          \sum_{\substack{j\ge1\\j\ {\rm odd}}}
 \left[
   \max\{d_{p^j}(L_\nu,ns),\
           n\,d_{p^j}(\lambda_{\nu,n},s)\}
                   -L_{\pi_p^j\nu}
 \right]_+ .
\]
Every matrix entry of \(A_nA_s\mathcal B_{\rm row}\) has valuation at least
\(g_p\) at every place above \(p\), where \(v_p(p)=1\). In
particular it takes the same bound on integral arrays in the
actual Cartan lattices. The product of the compact dimensions in
the period normalization satisfies \(\log\mathfrak d=O_{\psi,\eps}(x^C)\)
for a fixed \(C\).

Let \(k\) and \(B_\nu\) be the rank parameter and weight offsets
of Lemma~\ref{lem:degree181}, and \(b=\exp(x^\eps)\).
The guaranteed local bounds on the same interval satisfy
\[
 \sum_{p\in S}g_p\log p\ \ge\
  (k-1)\sum_{p\in S}(\log p)\sum_{C\ {\rm even}}\sum_{\nu\in H_C}
       \sum_{\substack{j\ge1\\j\ {\rm odd}}}
       [B_\nu/p^j-B_{\pi_p^j\nu}]_+
       -O_{\psi,\eps}\!\left(x^{C'}e^{x^{\eps/3}}\right)
\]
for a fixed \(C'\). The error is \(o(b^\eta)\) for every fixed
\(\eta>0\). The choice of half cycles may depend on \(p\).
\end{theorem}
\begin{proof}
Fix a pair of class representatives and a pair of coordinate
basis tensors. Their polynomial is a product of minors. At
embedding \(\nu\) every monomial has total degree \(L_\nu\).
Moreover exactly \(\lambda_{\nu,n}\) factors have size \(n\)
and each uses all \(n\) columns. All other factors have
nonnegative degree in every column. Thus the degree in each
individual column is at least \(\lambda_{\nu,n}\), regardless
of the smaller subsets chosen in that tensor.

For one monomial denote by \(\alpha_{ab,\nu}\) its entry
exponents. With \(M=p^j\), the sum of the floors of the \(ns\) entry exponents
at \(\nu\) divided by \(M\) is at least \(d_M(L_\nu,ns)\).
Indeed the sum of the remainders is at most \(ns(M-1)\).
Applying the same argument in each column gives the second
lower bound \(n\,d_M(\lambda_{\nu,n},s)\). The sum of the
entry exponents at the opposite embedding is exactly
\(L_{\pi_p^j\nu}\). In \eqref{eq:direct76}, the sum of the
individual positive parts is at least the positive part of the
sum. This proves the stated contribution for every \(\nu,j\)
of this single monomial. The same halves work throughout.
Summing these contributions and taking integral linear combinations
of the ordinary derivatives proves the bound for each minor
product. The integers \(A_n,A_s\) clear its class denominators.
This proves the matrix assertion uniformly over all class pairs,
including any signed combination with integral coefficients.

The dimension formula for each of the two Schur representations
has factors \((\lambda_i-\lambda_j+j-i)/(j-i)\), after padding
by zeros. Each factor is at most \(1+\lambda_1\).
The ranks and the number of embeddings are bounded by powers
of \(x\), and \(\log(1+\max_\tau L_\tau)=O_{\psi,\eps}
(x^\eps+\log(2x))\) from \eqref{eq:lastdegree181} and the
prescribed offsets. This gives the asserted dimension bound.

In the same degree formula write \(L_\nu=(k-1)B_\nu+J_x\).
The other terms of \eqref{eq:lastdegree181} do not depend
on \(\nu\). The \(v_i\), the spacings, and the ranks in
that formula are polynomially bounded in \(x\), so
\(|J_x|+ns\le x^{C_1}\) on sufficiently distant intervals.
Since \(d_M(L_\nu,ns)\ge L_\nu/M-ns\), the contribution to
\(g_p\) of a positive \(B_\nu/M-B_{\pi_p^j\nu}\) is at least
\((k-1)(B_\nu/M-B_{\pi_p^j\nu})-x^{C_1}\). A nonpositive
difference needs no error term.

For completeness the actual offsets are given following the
weight lists in Proposition~\ref{prop:short107} and the
construction of the terminal prime in Section~\ref{nav:VI-26}. Write
\(X'=\exp(x^{\eps/3})\) and let \(n_0\) be the degree parameter in
\(w''=b/(n_0 X')\) there. It is bounded by a fixed power of \(x\). The three initial offsets are
\(\lceil xb/X'\rceil+\lfloor b/X'\rfloor\),
\(\lceil b\rceil\), and \(\lceil b/(n_0X')\rceil\).
A parity correction can add one. In particular \(B_\nu>0\)
and \(R=\max B_\nu/\min B_\nu\le x^{C_2}X'\).
Only powers \(p^j\le R\) can have a positive difference.
The number of available \(\nu\)'s is at most \(f\). Each
integer at most \(R\) is a power of at most one prime, hence
\(\sum_{p^j\le R}\log p\le R\log(2R)\). The accumulated
error is at most \(fx^{C_1}R\log(2R)\). This has the stated
order and is \(o(\exp(\eta x^\eps))\) for each \(\eta>0\).
\end{proof}

For an algebraic incoming and outgoing array one can subtract
their denominator costs explicitly. If their least coordinate
valuations at a place are \(a_v,b_v\), the normalized value
of \(A_n A_s\mathcal B_{\rm row}\) on them has valuation at least
\(g_p+a_v+b_v\), by bilinearity. The matrix estimate itself
does not require either vector to survive the preceding lifts.
Using its value in a product formula for the full period requires
the arithmetic size of those actual vectors and a nonzero
contraction with the signed terminal projection.

\subsection{Congruences for paired frames}\label{nav:II-52}
At a marked place \(v\), the split unitary group of \(S\) is
\(\mathrm{GL}_s(F_v)\), \(s=\dim_E S\). Write \(K\) for the
stabilizer of the lattice, \(K(a)\) for its principal congruence
subgroup, \(Q=Nv\), and \(\pi\) for a uniformizer. Use the
unscaled columns \(b_1,\ldots,b_\ell\) of
Proposition~\ref{prop:positions264}. Their Gram matrix is zero
when \(c=2r\). When \(c=2r+1\) its only nonzero entry is the
last diagonal entry, equal to \(-1\). Denote these ranks by
\(h=0,1\), respectively.

\begin{lemma}[The marked orbit of an ordered frame]\leavevmode
\label{lem:flagmarked264}
Let \(b'_1,\ldots,b'_\ell\) have the same Gram matrix as the
given frame. Then
\[
 b'_j-b_j\in\pi^a L_{j,v}\quad(1\le j\le\ell)
 \quad\Longleftrightarrow\quad
    (b'_1,\ldots,b'_\ell)\in K(a)(b_1,\ldots,b_\ell).
\]
With the \(K\)-orbit assigned invariant probability, the
measure of this congruence class for \(a\ge1\) equals
\[
 \left[
 Q^{(a-1)(2s\ell-\ell^2)}\!
 \prod_{j=0}^{\ell-1}(Q^s-Q^j)\,
 Q^{h(s-\ell)}
 \prod_{j=0}^{\ell-h-1}(Q^{s-\ell}-Q^j)
 \right]^{-1}.
\]
An empty product is one.
\end{lemma}
\begin{proof}
Identify the split halves of the self-dual lattice with a free
module \(M\) over \(\mathcal O_v\) and its dual. Locally choose
generators of \(I_j^{(c)}\) in each half. Multiply the \(w\)-part
of column \(j\) by its first generator and the covector from
its \(\bar w\)-part by the second. The resulting columns form
a matrix \(V\) in \(M^\ell\), and the rows form a matrix \(Y\)
with values in \(\mathcal O_v^\ell\). The columns extend to a
basis, and the rows to a dual basis, since the staircase lines
are individual integral direct summands. Their product \(YV=N\)
is prescribed by the Gram matrix and the chosen generators.
In the central case the two generators have unit product since
the middle ideal has product one with its conjugate. Hence the
reduction of \(N\) has rank \(h\). The conditions on \(b'_j\)
are exactly \(V'\equiv V,\ Y'\equiv Y\bmod\pi^a\), with
\(Y'V'=N\). The split group sends a pair to
\((gV,Yg^{-1})\).

Prescribe \(g_0V=V'\), extending \(g_0\) by the identity on
the other elements of a basis extending the columns of \(V\).
Then \(g_0\equiv1\bmod\pi^a\). Put \(D=Y-Y'g_0\).
It vanishes on \(V\) and is divisible by \(\pi^a\).
The rows of \(Y'\) have a unit maximal minor and therefore
admit an integral right inverse \(J\). Now \(g=g_0+JD\) satisfies
\(gV=V'\), \(Y'g=Y\) and \(g\equiv1\bmod\pi^a\). In particular
it is invertible. Conversely any such congruence transformation
preserves the indicated residue classes. This proves the
orbit identity.

Over the residue field there are
\(\prod_{j=0}^{\ell-1}(Q^s-Q^j)\) independent ordered columns.
After fixing them, choose a complement of dimension \(s-\ell\).
The row map on the columns is fixed by \(N\). Its map on the
complement must surject onto the quotient of the target by
\(\operatorname{im}N\), of dimension \(\ell-h\). There are
\(\prod_{j=0}^{\ell-h-1}(Q^{s-\ell}-Q^j)\) such surjections,
by choosing independent rows. Each has \(Q^{h(s-\ell)}\)
lifts to the target. All these pairs lie in one orbit. Indeed
the prescribed isomorphism between the two column spaces
intertwines their row maps. Extend it on the intersections with
the row kernels to an isomorphism of the two full kernels.
On the sum of the column space and the kernel this gives a
well-defined isomorphism. Complete it using lifts on each side
of a basis of the target modulo \(\operatorname{im}N\).

In lifting a solution from depth \(n\) to \(n+1\), the linear
condition on the corrections over the residue field is
\(\Delta Y\,V+Y\,\Delta V=D_0\), where \(D_0\) is prescribed.
The first term alone attains every \(\ell\)-square matrix
because \(V\) has independent columns. There are therefore
\(Q^{2s\ell-\ell^2}\) lifts. They are in one orbit over
\(\mathcal O_v/\pi^{n+1}\) above the residue-field orbit.
Lifting the residue-field transformation and applying the preceding
right-inverse construction over this ring maps either pair to
the other. Every invertible matrix over the quotient lifts to \(K\)
by lifting its entries. Thus the count is exactly the number
of distinct reductions of the compact orbit. The orbit
identity identifies each fiber with a translate of the
congruence orbit. Their invariant probabilities are equal,
which proves the formula.
\end{proof}

\subsection{The marked integrals on the original core chain}\label{nav:IX-36}
Use the finite level from Theorem~\ref{thm:first250}. Write \(S_0\)
for its marked places, \(w(v)\) for the chosen half over \(v\), and
\(\mathfrak m=\prod_{v\in S_0}w(v)^{a_v}\) for the original ideal.
At step \(i<m\) put \(u=n_i\) and \(N=n_i+d\).
In the map from the larger lattice to the smaller one use the
\(w(v)\)-position space at a marked place. Choose bases of the two
lattices there, put \(M_v=\operatorname{Mat}_{u,N}(\mathcal O_v)\),
and let \(M_v^\vee\) be its full additive dual on the other half.
The depth equations are those of Proposition~\ref{prop:depth245}.
Normalize the additive measure by \(\operatorname{vol}(M_v)=1\)
and group measure by
\(\operatorname{vol}(\mathrm{GL}_u(\mathcal O_v))=1\).

\begin{theorem}[All marked factors at the prescribed depth]\leavevmode
\label{thm:markedcore253}
For every original field and every step, one integral finite test
at the marked places gives a scalar local doubling pairing on the
entire subspace fixed by \(K_u(a_v)=1+\varpi_v^{a_v}
\operatorname{Mat}_u(\mathcal O_v)\).
It is fixed on the larger side by \(K_N(a_v)\). Its scalar at
\(s=d/2\) is
\[
 D_{i,v}=(Nv)^{-a_v(uN+u^2)}
                   \prod_{j=1}^{u}\big(1-(Nv)^{-j}\big)^{-1}.
\]
There is a common integral residue class giving all the marked
tests for a given pair of lattices. Over the entire chain one has
\[
 (N\mathfrak m)^{-\sum_{i=0}^{m-1}(2n_i^2+d n_i)}
           \ \le\ \prod_{i=0}^{m-1}\prod_{v\in S_0}D_{i,v}
           \ \le\ 1 .
\]
\end{theorem}
\begin{proof}
Fix \(v\in S_0\) and abbreviate \(a=a_v,\ q=Nv\).
Proposition~\ref{prop:tracelat253} makes the two halves of the
original lattice \(M_v\) and \(M_v^\vee\). The congruence on the
chosen half replaces its displacement lattice by \(\varpi_v^a M_v\),
leaving the frequency lattice \(M_v^\vee\). Take
\(\phi_v=1_C\), where \(C=(I_u\ \ 0)+\varpi_v^a M_v\).
Its values are integers. Translation by \(\varpi_v^a M_v\)
preserves \(C\), and modulation by \(M_v^\vee\) is trivial on
its support. For a mixed displacement with components \(x,\ell\)
in these two lattices, the phase in the symmetric Weyl operator
for \(2(x,\ell)\), besides the pure translations and modulations,
is the character evaluated at \(\pm2\langle x,\ell\rangle\).
It is one by duality, including in residue characteristic two.
This verifies both depth equations for the actual test.

The determinant splitting characters at the auxiliary place have
conductor exponent one, and those at every other marked place
have exponent two, the stipulated depth there. Since the
auxiliary depth is positive, both characters are trivial on
\(1+\varpi_v^a\mathcal O_v\) everywhere. In the split
Schr\"odinger action the two groups act on the two sides of
the rectangle by change of basis, multiplied by the determinant
characters and real powers of absolute determinants. Both sorts
of scalar are one on the indicated congruence groups.
Consequently these subgroups fix \(\phi_v\).
For the exponents see
\srcloc{Development/AbelianVarietyTheory/Definitions/Divisors/BasePlaceNotMemChiGridBadPlaces}{70--82,125--139}.

The first \(u\) columns of any element of \(C\) make an
invertible square matrix congruent to \(I_u\). If \(hC\)
meets \(C\) for \(h\in\mathrm{GL}_u(F_v)\), two such
square matrices have quotient \(h\), forcing \(h\in K_u(a)\).
Conversely each such \(h\) fixes \(\phi_v\). Its oscillator
matrix coefficient on the smaller group is therefore
\(q^{-auN}1_{K_u(a)}\), also for the conjugated oscillator.
At the theta point the doubling section restricts to this
coefficient. Integrating it against the source matrix coefficient
of \(b,b'\) fixed by \(K_u(a)\) gives
\(q^{-auN}\operatorname{vol}(K_u(a))\langle b,b'\rangle\),
with the conjugate pairing if the conjugate convention is used.
Reduction onto \(\mathrm{GL}_u(\mathcal O_v/\varpi_v^a)\) is
surjective since invertibility is detected modulo \(\varpi_v\).
Counting independent columns modulo \(\varpi_v\) and then
the \(q^{(a-1)u^2}\) lifts gives the displayed scalar.
This identifies the indicator of Proposition~\ref{prop:split173}
inside the original depth space in the orientation of this core.

To realize the residues simultaneously, set
\(P=\Lambda_{\rm big},\ Q=\Lambda_{\rm small}\).
Finite projectivity of \(P\) gives
\(\operatorname{Hom}(P,Q)=P^\vee\otimes Q\). Reduction modulo
\(\mathfrak m\) identifies the quotient with the full module of
maps of the two reduced lattices. By the Chinese remainder
theorem this is the product of the local matrix modules at
the \(w(v)\). The chosen residues therefore lift to one
element of \(\operatorname{Hom}_{\mathcal O_E}(P,Q)\).
Finally \(\sum_{v\in S_0}a_v\log Nv=\log N\mathfrak m\) because
each place is split. Dropping the inverse Euler factors, all
at least one, gives the lower bound. The upper bound follows
from \(\operatorname{vol}(C)\le1\) and
\(\operatorname{vol}(K_u(a))\le1\).
\end{proof}

\subsection{The opposite line and its finite Fourier coefficients}\label{nav:IX-37}
Let \(C_0,C_1,C_2\) be the first three even cores on a field of
Theorem~\ref{thm:firstsquare255}. Write \(Y\) for the bottom
row space, \(S=C_1\), and \(s=\dim_E S\). The first two horizontal
targets are \(V_1=S\) and \(Z=V_2\). In the given frames
\[
 C_0=Y\perp Ez,\qquad Z=C_2\perp Ez',\qquad
 h(z,z)=1,\quad h(z',z')=-1 .
\]
The isometry \(V_1=S\) preserves the two lattices since the added
summand there has dimension zero. Write \(\mathfrak A\) for the
ideal of Proposition~\ref{prop:tracelat253} and \(\mathfrak m\)
for the original marking. The lattice of \(Ez\) is \(\mathcal O_Ez\).

\begin{proposition}[The exact transfer of the marked line test]\leavevmode
\label{prop:oppline256}
The lattice on \(Ez'\) is \(\mathfrak l z'\), where
\[
      \mathfrak l=\frac{\mathfrak A\mathfrak m}
                             {\bar{\mathfrak A}\bar{\mathfrak m}} .
\]
The Hermitian adjoint, identifying the two underlying lines by
\(z\mapsto z'\) and accounting for the opposite sign of their
forms, carries the whole original depth displacement lattice in
\(\operatorname{Hom}(S,Ez)\) onto that in
\(\operatorname{Hom}(Ez',S)\). Put \(D_L=(N\mathfrak m)^s\).
In the marked positions the corresponding Fourier transform of
the line indicator in Theorem~\ref{thm:firstsquare255} is
\(D_L^{-1}\) times a sum of exactly \(D_L\) distinct admissible
coset indicators, with root-of-unity coefficients. This is in
the normalization giving the original unshifted position lattice
volume one at the marked places.
\end{proposition}
\begin{proof}
For a one-dimensional added block, the grid assigns the ratio
of the two modified oscillator ideals to its generator. Those
ideals are \(\mathfrak A\bar{\mathfrak m}^{-1}\) and
\(\bar{\mathfrak A}\mathfrak m^{-1}\), giving the displayed
formula. This also proves \(\mathfrak l\bar{\mathfrak l}=1\),
as required for the negative line lattice to be self-dual.
The rank-one specializations of the defining formulas are in
\srcloc{Development/AbelianVarietyTheory/Definitions/Divisors/CmHostImaginaryGeneratorNeZero}{132--134}
and \srcloc{Development/AbelianVarietyTheory/Definitions/Divisors/BasePlaceNotMemChiGridBadPlaces}{80--88,178--189}.

Use \(h_S\) and the positive scalar form on the coordinate line
in defining \(x^\dagger\). The old alternating trace form is
the trace of \(-\xi_0 h(x,y)\) with the ordinary Hermitian map
pairing \(h\), since the source is \(-S\). The new source is
\(-Z\), whose restriction to \(Ez'\) is positive. The new map
pairing on \(x^\dagger,y^\dagger\) is \(\overline{h(x,y)}\)
by cyclicity of the matrix trace. Since \(\bar\xi_0=-\xi_0\),
the two alternating trace forms agree. This is an \(F\)-rational
symplectic map. It intertwines right action by \(g^{-1}\) with
left action by \(g\), since \((g^{-1})^\dagger=g\) on \(U(S)\).

At unmarked places the displacement coefficient ideal on a
Hom lattice is \(\mathfrak A\). At the selected component of
a marked place it is multiplied by \(\mathfrak m\), because
\(\mathfrak A\) is a unit there. Thus the old line displacement
lattice is the finite completion of
\(\mathfrak A\mathfrak m\operatorname{Hom}(\Lambda_S,\mathcal O_E)\).
Perfectness of the Hermitian form on \(\Lambda_S\) carries its
adjoint onto \(\bar{\mathfrak A}\bar{\mathfrak m}\Lambda_S\).
The new line displacement lattice is the completion of
\(\mathfrak A\mathfrak m\,\mathfrak l^{-1}\Lambda_S\).
The formula for \(\mathfrak l\) identifies them at every place.

At a marked \(v\), set \(a=a_v,\ q_v=Nv\). Write \(M_v\) for
the integral old row lattice of positions and \(M_v^\vee\) for
its additive dual, a lattice of columns. The adjoint exchanges
the two split components. Returning to the chosen position
component therefore Fourier transforms the old function
\(1_{x_0+\varpi_v^a M_v}\), where \(x_0\) is the indicated
coordinate row. Its transform is
\[
   q_v^{-as}\,\psi_v(\pm\langle x_0,y\rangle)\,
                        1_{\varpi_v^{-a}M_v^\vee}(y).
\]
Indeed integrating a character on the translated lattice gives
zero unless its restriction to the lattice is trivial, and in
that case gives its value at the translate times the volume.
The phase is constant modulo \(M_v^\vee\), and there are
\(q_v^{as}\) cosets in the displayed support. In new line
coordinates this support is the full integral Hom lattice since
\(\mathfrak l\) has valuation \(a\) on the chosen half. Its
depth sublattice is precisely \(M_v^\vee\). Each of these
cosets is therefore an allowed integral test. Multiplying the
formulas over \(v\in S_0\) proves the claim. At an unmarked
split place the unshifted indicator transforms into its dual
lattice indicator with the reference volume as scalar. At an
inert place transport the polarization of the hyperbolic
factor \(S\) by the rational adjoint. The unshifted vacuum is
then its transported vacuum. These common reference factors
are independent of the marked residues.
\end{proof}

\clearpage
\section{References}\label{nav:bibliography-XIX}
\begingroup\def\section*#1{}\begin{thebibliography}{99}\small
\bibitem{XIX-XIV-IV-molev106} A. I. Molev, \emph{Gelfand--Tsetlin bases for classical Lie algebras}, Handbook of Algebra, vol.~4 (2006), 109--170. \href{https://www.maths.usyd.edu.au/u/pubs/publist/preprints/2002/molev-18.pdf}{Theorem 2.1}.
\bibitem{XIX-XIV-IV-serre211} J.-P. Serre, \emph{Bounds for the orders of the finite subgroups of \(G(k)\)}, in \emph{Group Representation Theory}, EPFL Press, 2007. \href{https://arxiv.org/pdf/1011.0346}{Lecture I, Section 1.1, Theorem 1}.
\end{thebibliography}\endgroup
\end{document}
